\documentclass[11pt]{article}
\usepackage[T1]{fontenc}
\usepackage[margin=1in]{geometry}
\usepackage{amsmath,amssymb,amsthm,mathtools,mathrsfs}
\usepackage{booktabs,longtable,array}
\usepackage[colorlinks=true,linkcolor=blue,citecolor=blue,urlcolor=blue,hypertexnames=false]{hyperref}
\hypersetup{pdftitle={A conditional construction of a hereditarily atomic domain without ACCP},pdfauthor={}}
\theoremstyle{plain}
\newtheorem{introtheorem}{Theorem}
\renewcommand{\theintrotheorem}{\Alph{introtheorem}}
\newtheorem{theorem}{Theorem}[section]
\newtheorem{proposition}[theorem]{Proposition}
\newtheorem{lemma}[theorem]{Lemma}
\newtheorem{corollary}[theorem]{Corollary}
\theoremstyle{definition}
\newtheorem{definition}[theorem]{Definition}
\newtheorem{hypothesis}[theorem]{Hypothesis}
\theoremstyle{remark}
\newtheorem{remark}[theorem]{Remark}
\DeclareMathOperator{\Frac}{Frac}
\DeclareMathOperator{\Pic}{Pic}
\DeclareMathOperator{\Cl}{Cl}
\DeclareMathOperator{\Spec}{Spec}
\DeclareMathOperator{\Max}{Max}
\DeclareMathOperator{\Gal}{Gal}
\DeclareMathOperator{\ord}{ord}
\DeclareMathOperator{\im}{im}
\DeclareMathOperator{\divi}{div}
\newcommand{\Z}{\mathbb Z}
\newcommand{\Q}{\mathbb Q}
\newcommand{\F}{\mathbb F}
\newcommand{\OO}{\mathcal O}
\newcommand{\EA}{\mathfrak m_A}
\newcommand{\EB}{\mathfrak m_B}
\setlength{\emergencystretch}{2em}
\title{A conditional construction of a hereditarily atomic domain without ACCP}
\author{}
\date{October 22, 2026}


\begin{document}
\maketitle
\begin{abstract}
We construct a tower $\Q(\sqrt{65})=F_0\subset F_1\subset\cdots$ of totally
real number fields, each quadratic over the previous one, whose ring of
integers $D$ is an atomic almost Dedekind domain that does not satisfy the
ascending chain condition on principal ideals (ACCP).  The domain $D$ has
exactly two non-invertible maximal ideals.  A conjecture of Coykendall,
Gotti and Hasenauer is equivalent to the statement that every hereditarily
atomic domain, that is, every domain all of whose subrings are atomic,
satisfies ACCP.  We prove that the ring of integers of a tower built in
this way is hereditarily atomic, and therefore contradicts the conjecture,
provided every sufficiently late step $F_n/F_{n-1}$ satisfies two
arithmetic conditions.  The first, a Selmer condition (S), is the
injectivity of an explicit $\F_2$-linear map defined by Hilbert symbols at
the two dyadic primes.  The second, an odd unit condition (O), asks, for
each odd prime $p$ ramified in $F_n$, that one explicit vector lie in the
$\Z_p$-span of the images of certain relative units in a product of local
unit quotients at the ramified places above $p$.  The proof controls factorizations in an
arbitrary subring through a copy of $\Z[1/2]/\Z$ in the Picard group of
$D$, through the residue fields of $D$, and through approximation by a
single global unit.  Each condition concerns one quadratic step.  We give a
finite sufficient test for (O) and a reformulation of (S) through an
auxiliary quadratic extension.  Whether some tower satisfies both
conditions at every late step is left open.
\end{abstract}
\noindent\textbf{AI authorship disclosure.}
The research reported here and this manuscript, including its exposition
and proofs, were produced by the AI models GPT-6 Astra and Claude Opus 5.5.
Separate instances of these models reviewed the individual sections.  No
human authorship or human verification is asserted; AI review does not
constitute human verification.

\tableofcontents




\section{Introduction}\label{sec:introduction}

An integral domain is \emph{atomic} if every nonzero nonunit is a finite
product of irreducible elements.  It satisfies the \emph{ascending chain
condition on principal ideals} (ACCP) if every ascending chain of principal
ideals is eventually constant.  Every domain that satisfies ACCP is atomic
\cite[p.~2]{AAZ}, and Grams showed that the converse fails
\cite{Grams}.  Throughout, subrings contain the identity.

Coykendall, Gotti and Hasenauer call a domain \emph{hereditarily atomic} if
every subring is atomic \cite[Definition~2.1]{CGH}, and \emph{hereditarily
ACCP} if every subring satisfies ACCP \cite[Section~6]{CGH}.  They
conjecture that the two properties coincide \cite[Conjecture~6.1]{CGH}.
The conjecture is equivalent to the statement that every domain without
ACCP contains a subring that is not atomic \cite[Proposition~6.2]{CGH}.
Equivalently, every hereditarily atomic domain satisfies ACCP.  For
commutative cancellative monoids the corresponding statement is a theorem
of Gotti and Li: every hereditarily atomic monoid satisfies ACCP
\cite[Theorem~4.2]{GL}.  For domains, \cite[Theorem~6.4]{CGH} proves that
the two properties coincide under each of seven additional conditions.  In
characteristic zero, each of these conditions either requires an element
that is transcendental over $\Q$ or forces the domain itself to be
non-atomic.

This paper looks for a counterexample among rings of algebraic integers,
where no element is transcendental.  Let $K$ be an algebraic extension of
$\Q$ and let $D$ be its ring of integers.  Two known facts shape the
search.  First, the field $K$ is hereditarily atomic exactly when $D$ is a
Dedekind domain \cite[Theorem~4.4(1)]{CGH}.  The question therefore
concerns $D$, not $K$.  Second, every element of $D$ is integral over
$\Z$, so every subring $S\subseteq D$ satisfies $S^\times=S\cap D^\times$
(the proof of Lemma~\ref{lem:units} applies).  By a result of Grams (see \cite[Proposition~5.1]{CGH}),
if $D$ satisfied ACCP then so would every subring of $D$.  A counterexample
of this form must therefore fail ACCP itself; in particular $D$ cannot be
a Dedekind domain.

When $K$ is the union of a tower of number fields in which the
ramification index stays bounded along every sequence of primes $V_n$ of
$F_n$ with $V_{n+1}\mid V_n$, $D$ is an
\emph{almost Dedekind domain}: its localization at every maximal ideal is
a discrete valuation ring \cite[\S36]{Gilmer} (see
Section~\ref{sec:nt-background} and Lemma~\ref{lem:geometry}).  Divisibility in $D$ is then
read off from valuations, as in a Dedekind domain.  Two features
distinguish $D$ from a Dedekind domain.  A maximal ideal of $D$ need not be
finitely generated.  The Picard group $\Pic(D)$, the group of invertible
ideals modulo principal ideals, is the direct limit of the ideal class
groups of the fields in the tower, and it can be large.  An almost Dedekind
domain need not be atomic (see the discussion in \cite[Section~2]{CGH}).
Our construction uses both features to make $D$ atomic without ACCP.  Since
such a $D$ is not Dedekind, some overring of $D$ is not atomic
\cite[Theorem~3.6]{CGH}.  This overring contains $D$ properly, since $D$
itself is atomic, so it is not a subring of $D$, and this does not settle
the question.

\subsection{Main results}\label{sec:intro-results}

Our first result is unconditional.

\begin{introtheorem}\label{thm:A}
There is a tower
$F_0=\Q(\sqrt{65})\subset F_1\subset F_2\subset\cdots$ of totally real
number fields with $[F_n:F_{n-1}]=2$ for every $n\ge1$ such that the ring
of integers $D$ of $K=\bigcup_nF_n$ has the following properties.
\begin{enumerate}
\item[\textup{(a)}] $D$ is an almost Dedekind domain.
\item[\textup{(b)}] $D$ is atomic.
\item[\textup{(c)}] $D$ does not satisfy ACCP.
\item[\textup{(d)}] Exactly two maximal ideals of $D$ are not invertible.
Their residue fields are $\F_{47}$ and $\F_7$.
\end{enumerate}
\end{introtheorem}

Section~\ref{sec:construction} constructs the tower
(Theorem~\ref{thm:construction}).  Properties (a)--(d) follow from
Lemma~\ref{lem:geometry}, Theorem~\ref{thm:ambient} and the residue fields
recorded in (T1).  In $F_0$ the primes $47$ and $7$ split.  At
every later step one prime above $47$ and one prime above $7$ split again,
and along every other sequence of primes $V_n$ of $F_n$ with
$V_{n+1}\mid V_n$ the primes are eventually inert in each step.  Following the primes that keep splitting gives two chains of
primes of residue degree one: primes $A_n$ of $F_n$ above $47$ and primes
$B_n$ of $F_n$ above $7$, with
\[
 A_{n-1}\OO_{F_n}=A_nP_n,\qquad B_{n-1}\OO_{F_n}=B_nQ_n .
\]
The two non-invertible maximal ideals $\EA$ and $\EB$ of $D$ are the limits
of the chains $(A_n)$ and $(B_n)$.  The primes $P_n$ and $Q_n$ split off
from the chains at level $n$ give invertible maximal ideals of $D$.  The
failure of ACCP is explicit: there are elements $y_n\in\OO_{F_n}$ with
$(y_n)=A_nB_n$, starting with $y_0=-16+3\sqrt{65}$, such that
$y_0D\subsetneq y_1D\subsetneq y_2D\subsetneq\cdots$.

Each step of this tower is a \emph{scaled step}
(Definition~\ref{def:scaled-step}).  Informally, $F_n=F_{n-1}(y_n)$, where
$y_n$ is a root of a polynomial $X^2-tX+y_{n-1}$, so that $y_n$ has
relative norm $y_{n-1}$.  The coefficient $t\in\OO_{F_{n-1}}$ is given by
an explicit formula \eqref{eq:scaled-family} in terms of $y_{n-1}$, an
auxiliary element $c$ with $(c)=A_{n-1}R$ for an auxiliary prime $R$, an
integer $H\ge5$ with the scale $h=2^H$, and a power $u=\epsilon_H^k$, with
$k$ odd, of a unit $\epsilon_H$ of $F_0$ chosen so that $u-1$ has $2$-adic
valuation exactly $H$ at both dyadic primes.  The name refers to the
scale $h$.  A second element $\beta_n=c(u-1)/h+uy_n$ also has relative norm
$y_{n-1}$, and $(\beta_n)=A_nQ_n$.  Consequently the element
\[
 z_n=\frac{y_n\beta_n}{y_{n-1}}
\]
has relative norm one and $(z_n)=A_nP_n^{-1}$.  Section~\ref{sec:factorization}
isolates six properties \textup{(T1)--(T6)} of a tower with such elements,
and the tower of Theorem~\ref{thm:A} has all of them.  Our second result is
conditional.

\begin{introtheorem}\label{thm:B}\label{thm:main}
Let $F_0\subset F_1\subset\cdots$, with elements $y_n,\beta_n\in\OO_{F_n}$,
be a tower that satisfies the hypotheses \textup{(T1)--(T6)} of
Section~\ref{sec:factorization} and in which every step
$F_{n-1}\subset F_n$ is a scaled step in the sense of
Definition~\ref{def:scaled-step}.  Let $D$ be the ring of integers of
$K=\bigcup_nF_n$.  Suppose that, for every sufficiently large $n$, the step
$F_n/F_{n-1}$ satisfies both the Selmer condition \textup{(S)} of
Section~\ref{sec:dyadic} and the odd unit condition \textup{(O)} of
Section~\ref{sec:odd}.  Then $D$ is hereditarily atomic and does not
satisfy ACCP.
\end{introtheorem}

Theorem~\ref{thm:B} is proved at the end of Section~\ref{sec:odd}.  Both conditions concern a single step.  Write
$F=F_{n-1}$ and $L=F_n=F(\sqrt\delta)$, where $\delta=t^2-4y_{n-1}$, and
let $T$ be the set of all finite primes of $F$ that ramify in $L$.  The
field $F$ has exactly two \emph{dyadic} primes, that is, primes above $2$.

\emph{The Selmer condition.}  Let $\Delta(F,\delta)$ be the group of square
classes $[x]\in F^\times/F^{\times2}$ of totally positive elements $x$ such
that $v_{\mathfrak p}(x)$ is even at every prime $\mathfrak p$ outside $T$
and the two dyadic primes, and such that at each $\mathfrak p\in T$ the
element $x$ is a square or $\delta$ times a square in the completion
$F_{\mathfrak p}$.  It is a finite $\F_2$-vector space that contains
$[\delta]$.  The Hilbert symbols $(x,-\delta)_v$ at the two dyadic places
$v$ (Section~\ref{sec:nt-background}) define an $\F_2$-linear map
$\lambda_\delta$ from $\mathrm{Sel}_\delta=\Delta(F,\delta)/\langle[\delta]\rangle$
to $\F_2^2$.  Condition (S) is that $\lambda_\delta$ is injective.  Hilbert
reciprocity places the image of $\lambda_\delta$ in the line of vectors
with coordinate sum zero, and $\mathrm{Sel}_\delta$ is never zero
\eqref{eq:selmer-dimension}.  Thus (S) says that $\mathrm{Sel}_\delta$ has
exactly two elements and that its nonzero class $[x]$ satisfies
$(x,-\delta)_v=-1$ at both dyadic places.  The formal definition is
\eqref{eq:S}.

\emph{The odd unit condition.}  Let $p$ be an odd prime that ramifies in
$F_n$; by (T1), $p\ge5$.  At each place $v$ of $F_n$ above $p$ with ramification index
$e(v/p)=2$, let $W_v$, the \emph{coefficient ring} of $(F_n)_v$, be the
ring of integers of the largest unramified subfield of the completion
$(F_n)_v$.  The quotient
$C_v=\OO_{(F_n)_v}^\times/W_v^\times$ is a free $\Z_p$-module of finite
rank (Lemma~\ref{lem:odd-layer}).  Let $M_{n,p}$ be the product of the
groups $C_v$ over all such places $v$.  The \emph{relative units}, the
units of $F_n$ whose norm to $F_{n-1}$ is $1$, map to $M_{n,p}$, and so
does $z_n$, which is a unit at these places.  Condition (O) at $p$ asks
that the image of $z_n$ lie in the $\Z_p$-span of the images of those
relative units whose residues at the two dyadic primes of $F_n$ have
trivial $p$-primary component.  Two features of (O) matter.  It is a
condition in the whole product $M_{n,p}$: the image of $z_n$ must be
approximated at all places above $p$ simultaneously, not at each place
separately.  And it asks for membership in a $\Z_p$-span: knowing that this
span has full rank, that is, finite index in the submodule of $M_{n,p}$
on which the nontrivial automorphism of $F_n/F_{n-1}$ acts by $-1$, does
not suffice.  The formal definition is \eqref{eq:O},
which imposes (O) at every odd prime ramified in $F_n$.

\subsection{How the proof works}\label{sec:intro-strategy}

The proof of Theorem~\ref{thm:B} controls factorizations in an arbitrary
subring $S\subseteq D$ with three tools: the Picard group of $D$, the
residue fields of $D$, and approximation of the elements $z_n$ by global
units.

\emph{The Picard group: bounded factorizations and the failure of ACCP.}
Divisibility in $D$ is decided by valuations
(Lemma~\ref{lem:geometry}).  Since all maximal ideals of $D$ other than
$\EA$ and $\EB$ are invertible, $\Pic(D)$ decides which products of them
are principal.  The classes of the split-off primes generate a copy of
$\Z[1/2]/\Z$ in $\Pic(D)$, in which $[P_jD]$ corresponds to $2^{-j-1}$ and
$[Q_jD]$ to $-2^{-j-1}$ (Lemma~\ref{lem:classes}).  This uses (T3): no
nontrivial ideal class of $F_{n-1}$ becomes principal in $F_n$.  Let
$\kappa$ and $n_1$ be integers with $\kappa\sum_{j>n_1}2^{-j-1}<1$.  A
proper nonzero ideal $\prod_{j>n_1}P_j^{e_j}$ with $0\le e_j\le\kappa$ has
class represented by a number strictly between $0$ and $1$, so it is not
principal.  Hence an element $x$
with $v_{\EA}(x)=0$ or $v_{\EB}(x)=0$, called \emph{one-sided}, has bounded
factorization lengths in every subring of $D$
(Proposition~\ref{prop:onesided}).  Mixing the two chains cancels the
classes: $y_{n-1}/y_n$ generates $P_nQ_n$, which gives the strict chain
$y_0D\subsetneq y_1D\subsetneq\cdots$.  An element with both valuations
positive splits in $D$ into two one-sided factors, so $D$ is atomic
(Theorem~\ref{thm:ambient}).  These arguments occupy
Section~\ref{sec:factorization}.

\emph{The residue fields: which orders a subring can contain.}  Each step
has a prime that ramifies in it and is the only ramified prime above its
rational prime (an \emph{isolated ramified prime},
Definition~\ref{def:isolated} and hypothesis (T5)).  This forces every proper subfield of $K$ to be a
number field, so a subring $S$ with fraction field different from $K$ is a
bounded factorization domain (Proposition~\ref{prop:subfields}).  Suppose
instead that $S$ has an element with unbounded factorization lengths.  The
long factorizations produce elements of $S$ supported on late primes
$P_j,Q_j$.  At a maximal ideal $V$ of odd residue characteristic different
from $\EA,\EB$, the two-primary components of the residues of $y_n$ have
order $2^{n+o_V}$ for large $n$, with residue offsets $o_V$ that distinguish the maximal ideals
over the same rational prime (T6).  These orders force $S$ to map onto
every odd residue field of $D$, and to separate the maximal ideals of $D$
over each odd prime (Proposition~\ref{prop:residues}).  Teichm\"uller
lifting, which lifts each nonzero residue to a root of unity of order
prime to $p$ in the completion (Section~\ref{sec:nt-background}), then determines the
completions of the orders $S\cap F_n$: at odd primes unramified throughout
the tower they are maximal (Proposition~\ref{prop:unramified}); above a
ramified prime, for all sufficiently large $n$, they are maximal at the
places of ramification index one, and at a place $V$ of ramification index
two they are fixed by one nonnegative integer $s_V$ independent of $n$
(Proposition~\ref{prop:orders}).  Section~\ref{sec:orders} carries this out.

\emph{Unit approximation: moving a factorization into the subring.}  Let
$x\in S$ have unbounded factorization lengths and positive valuation at
both $\EA$ and $\EB$.  We split $x$ in $D$ into two one-sided factors.  A
factor with a given divisor is determined only up to a unit of $D$, and it
lies in $S$ only if it lies in the order $S\cap F_N$ for some $N$.  An order
differs from the maximal order at only finitely many primes, so this is a
finite list of local conditions.  By the previous paragraph, for large $N$ there are none at the odd
places of ramification index one, so the conditions sit at the dyadic places and at the
odd places of ramification index two.  The factor differs from an element of
$S$ by a product of powers of elements $z_n$ (Lemma~\ref{lem:split}).
So it suffices that every late $z_n$ be approximable by global units
simultaneously at all these places.  Condition \eqref{eq:JT} of
Section~\ref{sec:joint} states this approximation in a compact group
$\mathcal U_N$ that is the product of the unit group of
$\OO_{F_N}\otimes\Z_2$ and the quotients $C_v$ at the odd places of
ramification index two.  Under this condition one global unit moves both
factors into $S\cap F_N$ (Lemma~\ref{lem:return}), and each factor is
bounded in $S$ because it is one-sided.  Hence $S$ is atomic
(Theorem~\ref{thm:joint}).

\emph{Class field theory: from approximation to (S) and (O).}  At the
dyadic places, class field theory identifies the obstruction to
approximating an element by global units with the Artin symbol of its
ideal, that is, the image of the ideal under the Artin map of class field
theory (Section~\ref{sec:nt-background}).  This symbol lies in the Galois
group of the maximal abelian pro-$2$ extension unramified outside $2$ and
split at the real places (Lemma~\ref{lem:ray-closure}).  Under (S), the
relevant quotient of this group has exponent two
(Proposition~\ref{prop:ray-exponent-two}), and Kummer theory, which
describes quadratic extensions by square classes
(Section~\ref{sec:nt-background}), identifies its quadratic characters
with the classes of $\mathrm{Sel}_\delta$ (Lemma~\ref{lem:ray-selmer}).  The Artin symbol of
$(z_n)=A_nP_n^{-1}$ then vanishes, because every class of
$\mathrm{Sel}_\delta$ is a square locally at $A_{n-1}$; this gives
Theorem~\ref{thm:dyadic} in Section~\ref{sec:dyadic}.  At the odd places of
ramification index two, the quotients $C_v$ are pro-$p$, and
approximation becomes exactly condition (O)
(Lemma~\ref{lem:odd-membership}); for $p\ge5$ the $p$-adic logarithm
identifies each $C_v$ with the coefficient ring $W_v$, which gives ranks
and a finite test.  Proposition~\ref{prop:assembly}
combines the two, and Section~\ref{sec:odd} completes the proof of
Theorem~\ref{thm:B}.  Let $U^0_{n,p}$ be the $\Z_p$-span in the
definition of (O), and $M^-_{n,p}$ the $(-1)$-eigenspace of $M_{n,p}$
described in Section~\ref{sec:intro-results}.
Equality $U^0_{n,p}=M^-_{n,p}$ implies (O) at $p$, and Nakayama's lemma for
compact $\Z_p$-modules reduces it to a rank computation modulo $p$.  This
gives a finite sufficient test (Proposition~\ref{prop:odd-certificate}).

\subsection{What is proved and what remains}\label{sec:intro-status}

Theorem~\ref{thm:A} is unconditional.  Theorem~\ref{thm:B} requires (S)
and (O) at every sufficiently late step of one tower.  Each condition is a
statement about a single quadratic step $F_n/F_{n-1}$ and its unit group.
Condition (S) is a question of linear algebra over $\F_2$ on a finite group
of square classes.  For (O), Proposition~\ref{prop:odd-certificate} gives
a sufficient test at each prime $p$ by a finite computation with finitely
many relative units.  Appendix~\ref{app:comparison} reformulates (S)
through the auxiliary quadratic extension $F(\sqrt{\delta_+})$, where
$\delta_+$ is an explicit factor of the discriminant $\delta$, and
Appendix~\ref{app:precision} shows that the construction can prescribe, at
finitely many odd primes in each step, how closely $z_n$ is approximated by
elements of the coefficient rings (Proposition~\ref{prop:finite-precision}).
We do not construct an infinite tower in which every late step satisfies
both (S) and (O).  Section~\ref{sec:selection} states this remaining
problem as an existence statement for one-step extensions of finite towers
(Hypothesis~\ref{hyp:selection}), shows that it yields a hereditarily
atomic domain without ACCP (Proposition~\ref{prop:selection}), and explains
why prescribed local precision does not by itself give (O).

\subsection{Organization}\label{sec:intro-organization}

Section~\ref{sec:preliminaries} fixes notation and collects the background
from factorization theory and algebraic number theory, with references.
Section~\ref{sec:factorization} states the hypotheses (T1)--(T6) and
proves the factorization results that use only them.
Section~\ref{sec:construction} constructs the tower of
Theorem~\ref{thm:A}.  Section~\ref{sec:orders} determines the orders that a
subring with unbounded factorization lengths must contain, and
Section~\ref{sec:joint} proves hereditary atomicity from unit
approximation.  Sections~\ref{sec:dyadic} and~\ref{sec:odd} reduce unit
approximation to (S) and (O) and prove Theorem~\ref{thm:B}.
Section~\ref{sec:selection} formulates the remaining existence problem.
Appendix~\ref{app:precision} proves the local precision results,
Appendix~\ref{app:comparison} reformulates (S), and
Appendix~\ref{app:dependencies} records the dependencies among the main
results.

\section{Preliminaries}\label{sec:preliminaries}

\subsection{Notation and conventions}\label{sec:notation}

All rings are commutative domains with identity, and subrings contain that
identity.  Throughout, $p$ and $\ell$ denote rational primes, and $p$ is
odd unless stated otherwise.  The letter $\mathfrak m$ always denotes a
maximal ideal, and generic primes of a number field are written $\mathfrak
p$ and $\mathfrak q$.  A generic number field is written $\mathcal E$.
Places are written $v$ and $w$.  For a place $v$ of a number field
$\mathcal E$ we write $\mathcal E_v$ for the completion, $\OO_{\mathcal
E_v}$ for its ring of integers, $k_v$ for its residue field, $f_v$ for the residue
degree and $e(v/p)$ for the ramification index over $p$.  At a finite
place, $v$ also denotes the normalized valuation, and we write
$v_{\mathfrak p}$ when the place is given by a prime $\mathfrak p$; on
$\Q_2$ and its unramified extensions, $v_2$ is normalized by $v_2(2)=1$.
For a finite abelian group, or an element of one, the subscript $(2)$
denotes the two-primary part.  Ideal class groups are written
$\Cl(\mathcal E)$, ray class groups $\Cl_{\mathfrak f}(\mathcal E)$ for a
modulus $\mathfrak f$, and $\Pic$ denotes the Picard group.  For
$x\in\mathcal E^\times$ we write $[x]$ for its class in
$\mathcal E^\times/\mathcal E^{\times2}$, called its \emph{square class},
and $[x]_{\mathfrak p}$ for its class in $\mathcal E_{\mathfrak
p}^\times/\mathcal E_{\mathfrak p}^{\times2}$.

\emph{The tower.}  The fields are $F_0=\Q(\sqrt{65})\subset
F_1\subset\cdots$ with $[F_n:F_{n-1}]=2$.  We write $\sigma_n$ for the
nontrivial automorphism of $F_n/F_{n-1}$, $E_n=\OO_{F_n}^\times$ for the
unit group, and $E_n^1$ for the group of \emph{relative units}, the kernel
of the norm map $E_n\to E_{n-1}$.  We put $K=\bigcup_nF_n$ and $D=\OO_K$.
The letter $n$ always denotes a level of the tower; $m$, $N$ and $N'$
denote further levels.  The elements $y_n$, $\beta_n$, $z_n$ and $z_n'$
are those of hypothesis (T2), and the residue offsets $o_V$ are those of
(T6) (Section~\ref{sec:factorization}).

\emph{Primes.}  Hypothesis (T1) fixes two chains of primes $(A_n)$ and
$(B_n)$: $A_n$ and $B_n$ are primes of $F_n$ with residue fields $\F_{47}$
and $\F_7$, and
\[
 A_{n-1}\OO_{F_n}=A_nP_n,\qquad B_{n-1}\OO_{F_n}=B_nQ_n .
\]
Thus $A_{n-1}$ and $B_{n-1}$ split in $F_n$, and $P_n$ and $Q_n$ are the
primes split off from the chains $(A_n)$ and $(B_n)$ at level $n$.  A
prime $P_j$ or $Q_j$ of $F_j$, its unique extension to each later field,
and its extension to $D$ are denoted by the same letter.  The maximal ideals
of $D$ determined by the two chains are $\EA$ and $\EB$; they are the two
non-invertible maximal ideals of $D$ (Theorem~\ref{thm:ambient}).  A
maximal ideal of $D$ is usually denoted $V$, with $V_n=V\cap\OO_{F_n}$,
normalized valuation $v_V$ and residue field $k(V)=D/V$; likewise
$k(\mathfrak p)$ denotes the residue field of a prime $\mathfrak p$ of a
number field, for instance $k(V_n)=\OO_{F_n}/V_n$.  The tower primes $A_n$, $B_n$, $P_j$, $Q_j$,
the auxiliary prime $R$ and the detecting primes $G_i$ are written in
italic capitals.

\emph{One quadratic step.}  When we study a single step we write $F=F_{n-1}$,
$L=F_n$, $\sigma=\sigma_n$, $A=A_{n-1}$, $B=B_{n-1}$, $y=y_{n-1}$, and
$A\OO_L=A'P$, $B\OO_L=B'Q$.  The parameters of a scaled step are an integer
$H\ge5$, the scale $h=2^H$, a unit $\epsilon=\epsilon_H$ of $F_0$, an
exponent $k$, the unit $u=\epsilon^k$ and an element $c$ of $F$.  They
determine the coefficient $t\in\OO_F$ of \eqref{eq:scaled-family}, and
$y_n$ is a root of $X^2-tX+y$.  The discriminant of the step is
\[
 \delta=t^2-4y=\frac{\delta_-\delta_+}{(chu)^2},\qquad
 \delta_-=c^2-h^2y,\qquad
 \delta_+=c^2(u-1)^2-h^2y(u+1)^2,
\]
so that $L=F(\sqrt\delta)$.  The auxiliary prime $R$ satisfies
$(c)=AR$, and the detecting primes $G_i$ (Definition~\ref{def:detecting})
are primes of $F$, ramified in $L$, whose residue characters separate the
square classes of virtual units (Section~\ref{sec:nt-background}).  The set of dyadic primes of $F$ is
$\mathcal D$, and $T$ is the set of all finite primes of $F$ ramified in
$L$.

\emph{Groups and maps.}  The letter $U$ is used for unit groups:
$U_L=\prod_{v\mid2}(1+2\OO_{L_v})$ is the group of dyadic principal units
of $L$, $\mathcal U_N$ is the compact group of \eqref{eq:TN} in which unit
approximation is measured, and $U^0_{n,p}\subseteq U^{\mathrm{all}}_{n,p}$
are the $\Z_p$-modules of \eqref{eq:odd-images}.  The group of square
classes of virtual units of a number field $\mathcal E$ is $\mathcal
V(\mathcal E)$ (Section~\ref{sec:nt-background}).  The Selmer group $\Delta(F,\delta)$,
its quotient $\mathrm{Sel}_\delta$ and the map $\lambda_\delta$ are
defined in Section~\ref{sec:dyadic}; the modules $M_{n,p}$ and the maps
$\mu_{n,p}$ and $\operatorname{res}_{n,p}$ are defined in
Section~\ref{sec:odd}.  Hilbert symbols $(x,b)_v$ take values in
$\{\pm1\}$, and $[x,b]_v\in\F_2$ is defined by
$(-1)^{[x,b]_v}=(x,b)_v$.  At a place $v$ of $F_n$ of residue
characteristic $p$, $W_{n,v}$ is the coefficient ring of $(F_n)_v$
(Section~\ref{sec:nt-background}), always written with the level first;
$W_v$ abbreviates it when the level is fixed.

Letters such as $x$, $b$, $e$, $g$, $m$, $r$, $s$, $\xi$, $\gamma$, $\psi$, $I$,
the indeterminate $X$, and the letter $\chi$ for a quadratic character,
are also used as local variables; each such use is declared where it
occurs.  Apart from these, every symbol keeps one meaning throughout the
paper.

\subsection{Factorization background}\label{sec:factorization-background}

Let $S$ be a domain.  An \emph{irreducible element}, or \emph{atom}, of $S$
is a nonzero nonunit that is not a product of two nonunits.  The domain
$S$ is \emph{atomic} if every nonzero nonunit is a finite product of atoms.
It satisfies the \emph{ascending chain condition on principal ideals}
(ACCP) if every chain $x_1S\subseteq x_2S\subseteq\cdots$ of principal
ideals is eventually constant.  It is a \emph{bounded factorization
domain} (BFD) if it is atomic and, for each nonzero nonunit $x$, the
number of atoms in a factorization of $x$ into atoms is bounded
\cite[p.~1]{AAZ}.  By \cite[Theorem~2.4]{AAZ}, $S$ is a BFD if and only if,
for each nonzero nonunit $x$, the lengths $r$ of the expressions
$x=x_1\cdots x_r$ with all $x_i$ nonunits of $S$ are bounded.  This is the
form used in Section~\ref{sec:factorization}, where an
element with such a bound is called \emph{bounded} in $S$.  There are
implications
\[
 \text{BFD}\ \Longrightarrow\ \text{ACCP}\ \Longrightarrow\ \text{atomic}
\]
\cite[p.~2 and Section~2]{AAZ}.  Neither implication reverses: a monoid
domain of Anderson, Anderson and Zafrullah satisfies ACCP but is not a BFD
\cite[Example~2.1]{AAZ}, and Grams constructed an atomic domain without
ACCP \cite{Grams}.  For a systematic treatment of factorization theory we
refer to \cite{GHK}.

Hereditary atomicity is defined in Section~\ref{sec:introduction}.  We use
one descent property of ACCP.  Let $S$ be a subring of a domain that
satisfies ACCP, and suppose that every element of $S$ that is a unit of the
larger domain is a unit of $S$.  Then $S$ satisfies ACCP, by a result of
Grams (see \cite[Proposition~5.1]{CGH}).  If the larger domain is the ring of integers
of an algebraic extension of $\Q$, this unit condition holds automatically,
because every element is integral over $\Z$ (the proof of
Lemma~\ref{lem:units} applies).  For a subring $S$ of $D$ it reads
$S^\times=S\cap D^\times$, and it shows that a nonunit of $S$ is a nonunit
of $D$, so an element that is bounded in $D$
is bounded in every subring of $D$ that contains it, with the same bound.

An \emph{almost Dedekind domain} is a domain whose localization at every
maximal ideal is a discrete valuation ring \cite[\S36]{Gilmer},
\cite[Section~2.2]{FHL}.  An almost Dedekind domain has dimension at most
one and is a Prüfer domain \cite[p.~434]{Gilmer}, that is, every nonzero
finitely generated ideal is invertible \cite[Section~2]{CGH}.  It is a
Dedekind domain exactly when it is Noetherian \cite[Section~2]{CGH}.  An
invertible ideal is finitely generated.  Hence a maximal ideal of an
almost Dedekind domain is invertible if and only if it is finitely
generated.  As in Section~\ref{sec:notation}, $v_V$ is the normalized
valuation at a maximal ideal $V$; divisibility is then decided by these
valuations (Lemma~\ref{lem:geometry}).  For $x\ne0$ in an almost Dedekind
domain $S$, the \emph{divisor} of $x$ is the family $(v_V(x))_V$ indexed
by the maximal ideals $V$ of $S$, and its \emph{support} is the set of
those $V$ with $v_V(x)>0$.  The Picard group
$\Pic(S)$ of a domain $S$ is the group of invertible fractional ideals
modulo nonzero principal fractional ideals.  One-sided elements
(Definition~\ref{def:onesided}), which carry the boundedness argument, are
defined in Section~\ref{sec:factorization}.

\subsection{Number-theoretic background}\label{sec:nt-background}

This subsection collects the facts from algebraic number theory that the
paper uses.  We assume familiarity with Dedekind domains, the splitting of
primes in extensions, units and ideal class groups.  Each item below says
which factorization question it serves and gives a precise reference; we
use Milne's notes \cite{MilneANT,MilneCFT} throughout.

\paragraph{Towers of number fields.}  The facts in this paragraph hold for
every tower of number fields; we state them for the tower
$F_0\subset F_1\subset\cdots$, with $K=\bigcup_nF_n$ and $D=\OO_K$.  Then
$D=\bigcup_n\OO_{F_n}$, and the maximal ideals $V$ of $D$ correspond
bijectively to the sequences $(V_n)_{n\ge0}$ of primes $V_n$ of $F_n$ with
$V_{n+1}\mid V_n$, via $V_n=V\cap\OO_{F_n}$.  In a step $F_n/F_{n-1}$ of
degree two, a prime $\mathfrak p$ of $F_{n-1}$ with ramification indices
$e_i$ and residue degrees $f_i$ in $F_n$ satisfies $\sum_ie_if_i=2$
\cite[Theorem~3.34]{MilneANT}.  So $\mathfrak p$ either splits into two
primes, or is inert, or ramifies.  If $F_n/F_{n-1}$ is Galois and $w$ is a
place of $F_n$ above a place $v$ of $F_{n-1}$, then $(F_n)_w/(F_{n-1})_v$ is
Galois with group the decomposition group of $w$
\cite[Proposition~8.10]{MilneANT}; in particular $(F_n)_w=(F_{n-1})_v$ when
$v$ splits.  The localization $D_V$ is the union of the localizations of
the rings $\OO_{F_n}$ at $V_n$.  If the ramification index $e(V_n/p)$ is
the same for all large $n$, a uniformizer at one such level remains a
uniformizer at all later levels, and $D_V$ is a discrete valuation ring
(Lemma~\ref{lem:geometry}).

Every invertible ideal of $D$ is finitely generated, hence is extended
from an ideal of some $\OO_{F_n}$.  An ideal of $\OO_{F_n}$ that becomes
principal in $D$ already becomes principal in some $\OO_{F_m}$, since the
generator and the coefficients witnessing the equality lie in a finite
level (proof of Lemma~\ref{lem:classes}).  Hence $\Pic(D)$ is the direct
limit of the class groups $\Cl(F_n)$ under the maps induced by extension
of ideals.  These maps need not be injective; an ideal class that becomes
principal in a larger field is said to \emph{capitulate}.  Hypothesis (T3)
excludes capitulation in every step.  For factorization this is the
decisive point: the valuations of an element determine its divisor, and
$\Pic(D)$ decides which divisors are principal.

\paragraph{Units and square classes.}  The unit group
$\OO_{\mathcal E}^\times$ of a number field $\mathcal E$ with $r_1$ real
and $r_2$ complex places is finitely generated of rank $r_1+r_2-1$
\cite[Theorem~5.1]{MilneANT}.  If $\mathcal E$ is totally real, its torsion
units are $\pm1$, so $\OO_{\mathcal E}^\times/\OO_{\mathcal E}^{\times2}$
is an $\F_2$-vector space of dimension $[\mathcal E:\Q]$.  An element
$x\in\mathcal E^\times$ is a \emph{virtual unit} if $v_{\mathfrak p}(x)$ is
even for every prime $\mathfrak p$, that is, if $x\OO_{\mathcal E}=I^2$ for
a fractional ideal $I$.  We write
\[
 \mathcal V(\mathcal E)=\{[x]\in\mathcal E^\times/\mathcal E^{\times2}:
   v_{\mathfrak p}(x)\equiv0\pmod2\text{ for every prime }\mathfrak p\}
\]
for the group of square classes of virtual units.  Sending $[x]$ to the
class of $I$ gives an exact sequence
\[
 0\longrightarrow\OO_{\mathcal E}^\times/\OO_{\mathcal E}^{\times2}
 \longrightarrow\mathcal V(\mathcal E)
 \longrightarrow\Cl(\mathcal E)[2]\longrightarrow0,
\]
so $\dim_{\F_2}\mathcal V(\mathcal E)=[\mathcal E:\Q]+\dim_{\F_2}\Cl(\mathcal
E)[2]$ when $\mathcal E$ is totally real.  Units are the only freedom in changing a factor of an
element without changing its divisor, so the images of global units in
local unit groups are the decisive objects in Sections~\ref{sec:joint}
to~\ref{sec:odd}.  The detecting primes of Section~\ref{sec:construction}
are chosen so that every nonzero class of $\mathcal V(F_{n-1})$ has a
nonsquare residue at one of them, and Proposition~\ref{prop:scaled-step}
uses them to obtain (T3) and (T4).

\paragraph{Local fields at an odd prime.}  Let $\mathcal K$ be a finite
extension of $\Q_p$ with ring of integers $\OO$, maximal ideal $\mathfrak
m$ and residue field $k_{\mathcal K}$ of order $q$.  The ring $\OO$ contains
exactly $q$ roots of $X^q-X$, and reduction maps them bijectively and
multiplicatively onto $k_{\mathcal K}$; they are the \emph{Teichmüller representatives}
\cite[Remark~7.36]{MilneANT}.  This holds for every prime $p$, and
Section~\ref{sec:dyadic} uses it at $p=2$.  The field $\mathcal K$ contains a largest
unramified extension $\mathcal K_0$ of $\Q_p$, obtained by adjoining the
roots of unity in $\mathcal K$ of order prime to $p$
\cite[Proposition~7.50 and Corollary~7.51]{MilneANT}.  Its ring of
integers $W$ is the \emph{coefficient ring} of $\mathcal K$; it contains
the Teichmüller representatives.  Suppose that $\mathcal K/\mathcal K_0$
has degree two, so that it is totally ramified.  Then
$\OO=W\oplus\pi W$ for every uniformizer $\pi$
\cite[Proposition~7.55 and Remark~7.56]{MilneANT}.  Since $p$ is odd,
completing the square in the minimal polynomial of a uniformizer gives a
uniformizer $\pi$ with $\pi^2\in pW^\times$; we call such a $\pi$ a
\emph{square uniformizer}.  Newton's lemma lifts approximate roots of
polynomials to roots \cite[Theorem~7.32]{MilneANT}.

The $p$-adic exponential and logarithm are controlled by the following
bound.  Let $v$ be the normalized valuation of $\mathcal K$ and $e=v(p)$
its absolute ramification index, and put
$\mathfrak a=\{x\in\mathcal K:v(x)>e/(p-1)\}$.  For $i\ge1$ one has
$v_p(i!)\le(i-1)/(p-1)$ \cite[Example~7.29]{MilneANT}, and
$v_p(i)\le(i-1)/(p-1)$ because $p^j\ge1+j(p-1)$.  Hence, for
$x\in\mathfrak a$, the terms $x^i/i!$ and $x^i/i$ have valuation at least
$v(x)+(i-1)\bigl(v(x)-e/(p-1)\bigr)$, which exceeds $v(x)$ for $i\ge2$
and tends to infinity.  So the series $\exp(x)$ and $\log(1+x)$ converge
on $\mathfrak a$, and $v(\exp(x)-1)=v(\log(1+x))=v(x)$.  The same bounds
show that, for $x,x'\in\mathfrak a$ and $\mu=\min\{v(x),v(x')\}$, each
term of total degree $d$ in the expansions of the composed series
$\exp(\log(1+X))$, $\log(\exp(X))$ and $\log((1+X)(1+Y))$ at $(x,x')$ has
valuation at least $\mu+(d-1)\bigl(\mu-e/(p-1)\bigr)$.  Therefore the formal
identities $\exp(\log(1+X))=1+X$, $\log(\exp(X))=X$ and
$\log((1+X)(1+Y))=\log(1+X)+\log(1+Y)$ may be evaluated on $\mathfrak a$,
and $\log$ is an isomorphism from the multiplicative group $1+\mathfrak a$
onto the additive group $\mathfrak a$, with inverse $\exp$ (compare the
proof of \cite[Chapter~III, Lemma~2.4]{MilneCFT}).  In particular, if
$e<p-1$, then $\mathfrak a=\mathfrak m$, and $\log$ maps $1+\mathfrak m^r$
onto $\mathfrak m^r$ for every $r\ge1$.

These facts describe the orders a subring can contain.  If $x\in\OO$ has
residue $\bar x\in k_{\mathcal K}$, the powers $x^{q^r}$ converge to the
Teichmüller representative of $\bar x$.  So a subring whose image in a
residue field is the whole field contains, after completion, all
Teichmüller representatives and hence the whole coefficient ring.  This is
how Section~\ref{sec:orders} reduces membership in an order to one
nonnegative integer $s_V$ at each ramified place
(Proposition~\ref{prop:orders}).  For $p\ge5$ and $[\mathcal K:\mathcal K_0]\le2$, so that $e\le2<p-1$, the
logarithm makes the quotients $\OO^\times/W^\times$ into free $\Z_p$-modules, on which
condition (O) is linear (Lemma~\ref{lem:odd-layer}).  By (T1), only
primes $p\ge5$ occur in (O); Section~\ref{sec:odd} explains why $p=3$ is
excluded.

\paragraph{Hilbert symbols.}  Let $v$ be a place of a number field
$\mathcal E$.  For $x,b\in\mathcal E_v^\times$ the Hilbert symbol
$(x,b)_v\in\{\pm1\}$ equals $1$ exactly when $b$ is a norm from
$\mathcal E_v(\sqrt x)$, and it is bimultiplicative.  At a finite place
this is \cite[Chapter~III, Theorem~4.4(a),(d)]{MilneCFT}.  At a real
place it is elementary (compare \cite[Chapter~III, Remark~4.8]{MilneCFT}):
$b$ is a norm from $\mathbb R(\sqrt x)$ exactly when $x>0$ or $b>0$, so
$(x,b)_v=-1$ exactly when $x<0$ and $b<0$.  In particular, at a real place
$(x,b)_v=1$ whenever $x>0$ or $b>0$ there.

Let $v$ be a finite place of odd residue characteristic and $\pi$ a
uniformizer at $v$.  For $x\in\mathcal E_v^\times$ with $v(x)$ even, write
$\chi_v(x)\in\F_2$ for $0$ if the residue of the unit $x\pi^{-v(x)}$ is a
square and $1$ otherwise; we call $\chi_v$ the \emph{residue character}
at $v$.  Another uniformizer changes $\pi^{-v(x)}$ by an even power of a
unit, so $\chi_v(x)$ does not depend on $\pi$.  The nonzero squares form a
subgroup of index two in $k_v^\times$, so $\chi_v$ is a homomorphism on
the group of elements of even valuation, and it vanishes on squares.  By
Newton's lemma \cite[Theorem~7.32]{MilneANT}, $\chi_v(x)=0$ exactly when
$x$ is a square in $\mathcal E_v$.  For $x,b\in\mathcal E_v^\times$ with
$v(x)$ even, we have $[x,b]_v=v(b)\chi_v(x)$, that is,
$(x,b)_v=(-1)^{v(b)\chi_v(x)}$.  For $x,b\in\mathcal E^\times$ with $x$ a
unit at $v$ this is \cite[Chapter~VIII, 5.8]{MilneCFT}; it extends to
local $x$ and $b$, since both sides depend only on the square classes of
$x$ and $b$ in $\mathcal E_v^\times$, which contain elements of
$\mathcal E^\times$.  In general write $x=\pi^{2k}w$ with $w$ a unit;
bimultiplicativity gives $(x,b)_v=(\pi^k,b)_v^2(w,b)_v=(w,b)_v$, and
$\chi_v(x)=\chi_v(w)$.  In particular $(x,b)_v=1$ when $v(x)$ and $v(b)$
are both even, for instance when $x$ and $b$ are both units at $v$.
For $x,b\in\mathcal E^\times$ the product formula $\prod_v(x,b)_v=1$ holds
over all places of $\mathcal E$ \cite[Chapter~V, Example~5.4]{MilneCFT}.
In a quadratic step $L=F(\sqrt \delta)$, the symbols decide whether an
element of $F$ is locally a norm from $L$.  The product formula is the
global relation among them that we use.
It places the values of the map $\lambda_\delta$ of
Section~\ref{sec:dyadic} on the line of vectors with coordinate sum zero,
and Section~\ref{sec:dyadic} also uses it to show that a local condition at
$A_{n-1}$ holds automatically.

\paragraph{Ray class groups and primes in ray classes.}  A modulus
$\mathfrak f$ of a number field $\mathcal E$ is a nonzero ideal $\mathfrak
f_0$ of $\OO_{\mathcal E}$ together with a set of real places.  The ray
class group $\Cl_{\mathfrak f}(\mathcal E)$ is the group of fractional
ideals prime to $\mathfrak f_0$ modulo the principal ideals $(x)$ with
$x\equiv1\pmod{\mathfrak f_0}$ and $x$ positive at the real places of
$\mathfrak f$.  It is finite, and it sits in an exact sequence
\[
 0\longrightarrow\OO_{\mathcal E}^\times/\OO_{\mathcal E,\mathfrak
 f,1}^\times\longrightarrow
 \prod_{v\mid\mathfrak f\text{ real}}\{\pm1\}\times
   (\OO_{\mathcal E}/\mathfrak f_0)^\times\longrightarrow\Cl_{\mathfrak
   f}(\mathcal E)
 \longrightarrow\Cl(\mathcal E)\longrightarrow0,
\]
where $\OO_{\mathcal E,\mathfrak f,1}^\times$ is the group of units
congruent to $1$ modulo $\mathfrak f_0$ and positive at the real places of
$\mathfrak f$ \cite[Chapter~V, Theorem~1.7]{MilneCFT}.  By class field
theory there is a finite abelian extension of $\mathcal E$, the ray class
field, whose Galois group is identified with $\Cl_{\mathfrak f}(\mathcal
E)$ by the Artin map, which sends a
prime $\mathfrak p$ not dividing $\mathfrak f_0$ to its Frobenius element
\cite[Chapter~V, Theorems~3.5 and~3.6, Corollary~3.7 and
Example~3.9]{MilneCFT}.  The Chebotarev density theorem
\cite[Chapter~VIII, Theorem~7.4]{MilneCFT}, together with the fact that the
primes of absolute residue degree at least two have Dirichlet density zero
\cite[Chapter~VI, Proposition~4.5 and Corollary~4.6]{MilneCFT}, shows that every class of
$\Cl_{\mathfrak f}(\mathcal E)$ contains infinitely many primes of absolute residue
degree one.  In Section~\ref{sec:construction}, an element $c$ with
$(c)=A_{n-1}R$ satisfying finitely many congruence and sign conditions
requires the auxiliary prime $R$ to lie in a prescribed class of a ray class
group, and Chebotarev's theorem supplies such primes.  In
Section~\ref{sec:dyadic} the moduli $2^m\OO_F$, for all $m$, give the
Galois group in which the obstruction to dyadic unit approximation lies
(Lemma~\ref{lem:ray-closure}).

\paragraph{Quadratic Kummer theory.}  Let $\mathcal E$ be a number field,
and fix an algebraic closure.  Sending a finite
abelian extension of $\mathcal E$ of exponent two inside this closure to
the group of elements of $\mathcal E^\times$ that are squares in the
extension is a bijection onto the subgroups of $\mathcal E^\times$ that
contain $\mathcal E^{\times2}$ with finite index.  The extension
corresponding to such a subgroup is generated by the square roots of its
elements, and its degree equals that index \cite[Chapter~VII, Appendix~A,
Theorem~A.3]{MilneCFT}.  In particular, the quadratic extensions of
$\mathcal E$ correspond to the nonzero square classes.  If
$x\in\mathcal E^\times$ is a unit at a finite place $v$
of odd residue characteristic, then $\mathcal E(\sqrt x)$ is unramified at
$v$ \cite[Chapter~VII, Appendix~A, Proposition~A.5]{MilneCFT}.  If $v(x)$
is odd, then $\mathcal E(\sqrt x)$ is ramified at $v$.  Hence
$\mathcal E(\sqrt x)$ is ramified at such $v$ exactly when $v(x)$ is odd:
if $v(x)$ is even, choose $a\in\mathcal E^\times$ with $v(a)=v(x)/2$;
then $\mathcal E(\sqrt x)=\mathcal E(\sqrt{x/a^2})$ and $x/a^2$ is a unit
at $v$.  Locally, with $a$ a power of a uniformizer, the same argument and
Newton's lemma show that $\mathcal E_v(\sqrt x)$ is the unramified
quadratic extension of $\mathcal E_v$ exactly when $x$ is a nonsquare unit
times an even power of a uniformizer, that is, when $v(x)$ is even and
$\chi_v(x)=1$.  At a real place, $\mathcal E(\sqrt x)$ is split exactly
when $x$ is positive there.  In
Section~\ref{sec:dyadic}, the group that obstructs dyadic unit
approximation is a Galois group, and Kummer theory identifies the
quadratic characters of the relevant quotient of this group with the
classes of $\mathrm{Sel}_\delta$ (Lemma~\ref{lem:ray-selmer}).  This turns an
approximation problem in a profinite group into linear algebra over $\F_2$.

\paragraph{Curves over finite fields.}  For a smooth projective
geometrically connected curve of genus $g$ over $\F_q$, the number of its
$\F_q$-rational points differs from $q+1$ by at most $2g\sqrt q$
\cite[Section~1]{MilneWeil}.  In Section~\ref{sec:construction} this bound
supplies values of $c$ that keep a given odd prime inert in the next step
(Lemma~\ref{lem:scaled-holding}) and, once a residue field is large,
values that give the residue of $y_n$ a prescribed two-primary component
(Lemma~\ref{lem:scaled-residues}).  This is what allows the construction
to prescribe the orders of the two-primary components of residues, and hence the residue offsets of (T6).

\section{Almost Dedekind towers and factorization}\label{sec:factorization}

This section isolates the properties of the tower that the factorization
arguments use, and it proves everything about the ring $D$ itself that
Theorem~\ref{thm:A} asserts.  Section~\ref{sec:hypotheses} states six
hypotheses (T1)--(T6) on a tower of quadratic extensions and says where
each one is used.  Only (T1)--(T3) are needed in the present section.
Hypothesis (T1) describes how primes split in the tower.  It makes $D$ an
almost Dedekind domain in which divisibility is decided by valuations, as
in a Dedekind domain (Lemma~\ref{lem:geometry}).  Hypotheses (T2) and (T3)
supply principal ideals and ideal classes that embed a copy of
$\Z[1/2]/\Z$ in $\Pic(D)$ (Lemma~\ref{lem:classes}).  Comparing classes in
this subgroup shows that a nonzero nonunit whose valuation vanishes at one
of the two maximal ideals $\EA,\EB$ has bounded factorization lengths in
every subring of $D$ that contains it (Proposition~\ref{prop:onesided}).
Section~\ref{sec:joint} applies this to the two factors into which an
element of a subring is split.  Finally, Theorem~\ref{thm:ambient} shows
that $D$ is atomic, does not satisfy ACCP, and has exactly two
non-invertible maximal ideals.  Hypotheses (T4)--(T6) are used only from
Section~\ref{sec:orders} on, to control the subrings of $D$.

\subsection{The structural hypotheses}\label{sec:hypotheses}

Let $F_0=\Q(\sqrt{65})\subset F_1\subset\cdots$ be a tower of totally real
number fields with $[F_n:F_{n-1}]=2$, and let $\sigma_n$, $E_n$, $K$ and
$D=\OO_K$ be as in Section~\ref{sec:notation}.  We write $\Max D$ for the
set of maximal ideals of $D$.  A \emph{branch} is a sequence
$(V_n)_{n\geq m}$ of primes $V_n$ of $F_n$ with $V_{n+1}\mid V_n$.  The
branches with $m=0$ correspond bijectively to the maximal ideals $V$ of
$D$, via $V_n=V\cap\OO_{F_n}$ (Section~\ref{sec:nt-background}), and we
use the same letter for a maximal ideal and its branch.  Primes $P_j$ and
$Q_j$ follow the convention of Section~\ref{sec:notation}.

Hypothesis (T5) uses the following notion.  If a step $F_n/F_{n-1}$ has
an isolated ramified prime, defined below, then every subfield of $F_n$ contains $F_{n-1}$ or is
contained in it (Lemma~\ref{lem:witness}).  This is what forces every
proper subfield of $K$ to be a number field.

\begin{definition}\label{def:isolated}
Let $L/F$ be a quadratic extension of number fields.  A prime $G$ of $F$
is an \emph{isolated ramified prime} of $L/F$ if its residue
characteristic $\ell$ is unramified in $F$, $F_G=\Q_\ell$, $G$ ramifies in
$L$, and every other prime of $F$ over $\ell$ is unramified in $L/F$.
\end{definition}

We impose the following hypotheses.  Each concerns the fields of the
tower and their rings of integers, not the subrings of $D$.
\begin{enumerate}
\item[\textup{(T1)}] There are primes $A_n,B_n$ of $F_n$ with
$(F_n)_{A_n}=\Q_{47}$ and $(F_n)_{B_n}=\Q_7$ such that
\[
 A_{n-1}\OO_{F_n}=A_nP_n,\qquad
 B_{n-1}\OO_{F_n}=B_nQ_n,
\]
where $A_0=(47,\sqrt{65}-21)$ and $B_0=(7,\sqrt{65}-3)$.  Thus $A_{n-1}$
and $B_{n-1}$ split in $F_n$.  We call $(A_n)$ and $(B_n)$ the two
\emph{chains}, and we call $P_n$ and $Q_n$, the second primes of $F_n$
over $A_{n-1}$ and $B_{n-1}$, the primes \emph{split off} from the chains.
We write $\EA,\EB$ for the maximal ideals of $D$ determined by the chains
$(A_n)$ and $(B_n)$.  They have residue fields $\F_{47}$ and $\F_7$.
Moreover:
\begin{itemize}
\item every $P_n$ and every $Q_n$ is unramified and inert in each
subsequent step;
\item the two dyadic primes of $F_0$ are unramified and inert in every
step;
\item for each odd rational prime $p\notin\{47,7\}$ there is a level after
which all primes over $p$ are unramified and inert in every step;
\item every branch over $47$ or $7$ other than $(A_n)$ and $(B_n)$ is
eventually unramified and inert, that is, unramified and inert in every
step from some level on;
\item absolute ramification indices at odd primes are at most two, and
$3$, $7$ and $47$ are unramified in every $F_n$.
\end{itemize}

\item[\textup{(T2)}] There are $y_0=-16+3\sqrt{65}$ and
$y_n,\beta_n\in\OO_{F_n}$ for $n\geq1$ satisfying
\[
 (y_n)=A_nB_n,\quad (\beta_n)=A_nQ_n,\quad
 N_{F_n/F_{n-1}}y_n=N_{F_n/F_{n-1}}\beta_n=y_{n-1}.
\]
Define
\begin{equation}\label{eq:siblings}
 z_n=\frac{y_n\beta_n}{y_{n-1}},\qquad
 z_n'=\frac{y_n}{\beta_n}.
\end{equation}
Then $(z_n)=A_nP_n^{-1}$, $(z_n')=B_nQ_n^{-1}$, and both relative norms
are one.

\item[\textup{(T3)}] For each $n\geq1$, extension of ideals induces an
injective homomorphism $\Cl(F_{n-1})\to\Cl(F_n)$, and $[A_0]$ has order
two in $\Cl(F_0)$.

\item[\textup{(T4)}] The unit norms satisfy
$N_{F_n/F_{n-1}}(E_n)\subseteq E_{n-1}^{2}$ for every $n$.

\item[\textup{(T5)}] Every step $F_n/F_{n-1}$ has an isolated ramified
prime $G$ (Definition~\ref{def:isolated}).

\item[\textup{(T6)}] If $V\in\Max D\setminus\{\EA,\EB\}$ has odd residue
characteristic, then there is an integer $o_V$, the \emph{residue offset}
of $V$, such that, for all sufficiently large $n$, the two-primary
component of the residue of $y_n$ at $V$ has order $2^{n+o_V}$.  Distinct
such $V$ over the same rational prime have distinct residue offsets.
\end{enumerate}

The assertions of (T2) about $z_n$ and $z_n'$ follow from the displayed
identities.  Indeed $(y_{n-1})=A_{n-1}B_{n-1}$ (for $n=1$ because
$y_0\in A_0\cap B_0$ and $N_{F_0/\Q}(y_0)=-329=-7\cdot47$), so $y_{n-1}\OO_{F_n}=A_nP_nB_nQ_n$ by (T1), which gives
the two principal ideals, and the relative norm of $z_n$ is
$y_{n-1}\cdot y_{n-1}/y_{n-1}^2=1$, and that of $z_n'$ is
$y_{n-1}/y_{n-1}=1$.

In (T6), the residue field $D/V$ is an algebraic extension of a finite
field, so every element of its multiplicative group has finite order.
The \emph{two-primary component} of such an element is its component of
two-power order, which is therefore well defined.  Hypothesis (T6) is a
statement about the residues of $y_n$ at the finite levels, required for
large $n$ separately for each $V$; no starting level uniform in $V$ is
required.

\paragraph{The role of each hypothesis.}
\begin{itemize}
\item (T1) determines the maximal ideals of $D$ and makes $D$ almost
Dedekind (Lemma~\ref{lem:geometry}).  Its local clauses (two unramified
dyadic primes, ramification index at most two at odd primes, $3$
unramified) restrict the local fields that occur in
Sections~\ref{sec:orders}--\ref{sec:odd}.
\item (T2) supplies the elements used throughout.  The $y_n$ give the
strictly increasing chain $y_0D\subsetneq y_1D\subsetneq\cdots$, and the
conjugates $\sigma_n(\beta_n)$ split an element into two one-sided factors
(Theorem~\ref{thm:ambient}).  The $z_n$ and $z_n'$ give the class
identities of Lemma~\ref{lem:classes}, and they are the elements that
Section~\ref{sec:joint} approximates by global units.
\item (T3) excludes capitulation of ideal classes.  It makes the map
$\iota$ of Lemma~\ref{lem:classes} injective, and this injectivity drives
Proposition~\ref{prop:onesided}.
\item (T4) keeps the two-primary components of the residues of global
units at a fixed odd maximal ideal $V\ne\EA,\EB$ inside a finite group
(Lemma~\ref{lem:unitresidues}).  This is used in
Proposition~\ref{prop:residues}.
\item (T5) forces every proper subfield of $K$ to be a number field
(Lemma~\ref{lem:witness} and Proposition~\ref{prop:subfields}).
\item (T6) lets a subring with unbounded factorization lengths recover
every odd residue field of $D$ and separate the maximal ideals of $D$ over
each odd prime (Proposition~\ref{prop:residues}).
\end{itemize}
Section~\ref{sec:construction} constructs a tower that satisfies all six
hypotheses (Theorem~\ref{thm:construction}).

\subsection{Valuations}\label{sec:valuations}

The first lemma shows that, for divisibility, $D$ behaves like a Dedekind
domain: each element has a valuation at every maximal ideal, and these
valuations decide divisibility and units.  The difference from the
Dedekind case is recorded in \eqref{eq:tails}: the valuation of an element
at $\EA$ reappears at infinitely many of the maximal ideals $P_j$, and
similarly for $\EB$ and the $Q_j$.  The facts on splitting and completions
used in the proof are recalled in Section~\ref{sec:nt-background}.

\begin{lemma}\label{lem:geometry}
Under \textup{(T1)}, $D$ is almost Dedekind.  Denote the normalized
valuation at $V\in\Max D$ by $v_V$.  For every $0\ne x\in D$,
\begin{equation}\label{eq:tails}
 v_{P_j}(x)=v_{\EA}(x),\qquad v_{Q_j}(x)=v_{\EB}(x)
 \quad(j\gg0),
\end{equation}
and $x$ has positive valuation at only finitely many maximal ideals
outside $\{\EA,\EB,P_j,Q_j:j\geq1\}$.  Each $P_j,Q_j$ is invertible.  An
element of $D$ is a unit of $D$ exactly when all its valuations are zero,
and $x'\mid_Dx$ exactly when $v_V(x')\leq v_V(x)$ at every $V\in\Max D$.
More generally, an element $x\in K^\times$ lies in $D$ exactly when
$v_V(x)\geq0$ at every $V\in\Max D$, and lies in $D^\times$ exactly when
all its valuations are zero.
\end{lemma}

\begin{proof}
We show that each localization is a discrete valuation ring, then
compare valuations along the two chains, then count the remaining
maximal ideals.

\emph{Step 1: $D$ is almost Dedekind.}  Let $V\in\Max D$ have residue
characteristic $p$.  The local ring $D_V$ is the increasing union of the
local rings of the $\OO_{F_n}$ at $V_n$.  The ramification index
$e(V_n/p)$ is nondecreasing in $n$.  By (T1) it is at most two if $p$
is odd, and it is constant if $p=2$, because the dyadic primes of $F_0$ are
unramified in every step.  Hence it is constant from some level on.  From
that level on, a uniformizer at one level remains a uniformizer at every
later level.  Every nonzero element of $D_V$ lies in one of these later
local rings, where it is a power of this uniformizer times a unit.  Thus
$D_V$ is a valuation domain with value group $\Z$, that is, a discrete
valuation ring.

\emph{Step 2: the chains.}  Place $x$ in $F_m$.  By (T1) and
Section~\ref{sec:nt-background}, the completions along the chains $(A_n)$
and $(B_n)$ are all $\Q_{47}$ and $\Q_7$, so each step along a chain is
unramified of local degree one.  The prime $P_j$ with $j>m$ lies over
$A_{j-1}$, has completion $\Q_{47}$, and is unramified in every later
step.  Therefore $v_{P_j}(x)=v_{A_{j-1}}(x)=v_{A_m}(x)=v_{\EA}(x)$ for every
$j>m$.  The same argument for $(B_n)$ gives $v_{Q_j}(x)=v_{\EB}(x)$ for
$j>m$.  This proves \eqref{eq:tails}.

\emph{Step 3: the remaining maximal ideals.}  If $x\in V$ and $V$ lies over
the rational prime $p$, then $N_{F_m/\Q}(x)\in V\cap\Z=p\Z$.  So only the
finitely many rational primes dividing $N_{F_m/\Q}(x)$ can occur.  For a
rational prime other than $47$ and $7$, (T1) supplies a common level after
which no prime over it splits, so finitely many maximal ideals of $D$ lie
over it.  Over $47$ and $7$, consider the primes of $F_0$ conjugate to
$A_0$ and $B_0$.  The set of infinite branches through one of them is the
inverse limit of the finite sets of primes above it at the successive
levels, and is therefore compact.  Each such branch is eventually inert by
(T1), so some finite-level prime on it has exactly one infinite
continuation.  The set of branches through that prime is then an open
singleton.  A compact discrete space is finite, so these branches are
finitely many.  The branches through $A_0$ are $(A_n)$ and, for each
$j\geq1$, the branch of $P_j$, because $P_j$ is inert in every later step;
the same holds for $B_0$ with the $Q_j$.  Hence all remaining maximal
ideals over $47$ and $7$ are $\EA$, $\EB$ and the $P_j,Q_j$.  This proves
the finiteness assertion.

\emph{Step 4: invertibility of $P_j$ and $Q_j$.}  The ideal $P_j\OO_{F_n}$
remains prime for every $n\geq j$, so $P_jD$ is maximal.  It is invertible
because extension of ideals preserves the inverse of the invertible ideal
$P_j\subset\OO_{F_j}$.  The same argument applies to $Q_j$.

\emph{Step 5: units and divisibility.}  Inside $K$ we have
$D=\bigcap_{V\in\Max D}D_V$.  Hence $x/x'\in D$ exactly when
$v_V(x)-v_V(x')\geq0$ for every $V$, and $x^{-1}\in D$ exactly when
$v_V(x)\leq0$, that is $v_V(x)=0$, for every $V$.  The same equality
$D=\bigcap_VD_V$ gives the criteria for elements $x\in K^\times$: $x\in D$
exactly when every $v_V(x)\geq0$, and $x,x^{-1}\in D$ exactly when every
$v_V(x)=0$.
\end{proof}

Support and divisor of elements of $D$ are as in
Section~\ref{sec:factorization-background}.

\subsection{The Picard group}\label{sec:picard}

Since $D$ is almost Dedekind, the valuations of an element determine its
principal ideal, and the question of which ideals are principal is
answered by $\Pic(D)$.  By Section~\ref{sec:nt-background}, $\Pic(D)$ is the
direct limit of the class groups $\Cl(F_n)$, and (T3) makes every map in
this direct system injective.  The next lemma locates the classes of the
primes of the two chains in a copy of $\Z[1/2]/\Z$.

The mechanism is visible already at the first steps.  The ideal $A_0$ has
class of order two, which corresponds to $\tfrac12$.  In $F_1$ it splits
as $A_1P_1$, and the element $z_1$ of (T2) shows that $A_1$ and $P_1$ have
the same class.  So the class of $A_0$ is twice the class of $A_1$, which
therefore corresponds to $\tfrac14$.  Continuing, $[A_nD]$ corresponds to
$2^{-n-1}$.  This matches the divisors.  For $m>n$, iterating (T1) gives
$A_n\OO_{F_m}=A_mP_{n+1}\cdots P_m$, and correspondingly
$2^{-n-1}=2^{-m-1}+\sum_{j=n+1}^{m}2^{-j-1}$.  The primes of the chain
$(B_n)$ and the $Q_j$ carry the negative classes, because $y_n$ and $z_n'$
are principal.

\begin{lemma}[The classes of $A_n$, $B_n$, $P_j$, $Q_j$ in $\Pic(D)$]\label{lem:classes}
Under \textup{(T1)--(T3)}, there is an injective homomorphism
$\iota:\Z[1/2]/\Z\longrightarrow\Pic(D)$ such that
\begin{equation}\label{eq:classes}
\begin{aligned}
 [A_nD]&=\iota(2^{-n-1}+\Z),&[B_nD]&=-[A_nD],\\
 [P_jD]&=\iota(2^{-j-1}+\Z),&[Q_jD]&=-[P_jD].
\end{aligned}
\end{equation}
\end{lemma}

The lemma says nothing about the rest of $\Pic(D)$, and nothing more is
needed.

\begin{proof}
We compute the classes at each finite level, show that no nonzero class
of a finite level vanishes in $D$, and then pass to the limit.

\emph{Step 1: classes at level $n$.}  The principal ideal
$(z_n)=A_nP_n^{-1}$ gives $[P_n]=[A_n]$ in $\Cl(F_n)$.  Consequently
extension sends $[A_{n-1}]$ to $[A_nP_n]=2[A_n]$.  By (T3), $[A_0]$ has
order two, and if $[A_{n-1}]$ has order $2^n$, then injectivity of
extension shows that $2[A_n]$ has order $2^n$, so $[A_n]$ has order
$2^{n+1}$.  By induction, $[A_n]$ has order $2^{n+1}$ for every $n$.  The
principal ideal $(y_n)=A_nB_n$ gives $[B_n]=-[A_n]$, and
$(z_n')=B_nQ_n^{-1}$ gives $[Q_n]=[B_n]=-[P_n]$.

\emph{Step 2: no capitulation in $D$.}  A class of a finite level that
becomes principal in $D$ already becomes principal in a finite level.
Indeed, an equality $ID=uD$ involves finitely many generators of $I$ and
of $I^{-1}$, the element $u$, and finitely many coefficients witnessing
both inclusions $ID\subseteq uD$ and $uD\subseteq ID$.  All these data lie
in a common finite level.  Thus (T3) prevents a nonzero class of a finite
level from vanishing in $\Pic(D)$, and the class of $A_nD$ in $\Pic(D)$
has order $2^{n+1}$.

\emph{Step 3: the limit.}  By Step 1, $2[A_nD]=[A_{n-1}D]$ in $\Pic(D)$.
So the cyclic subgroups generated by the classes $[A_nD]$ are compatible,
and $\iota(a2^{-n-1}+\Z)=a[A_nD]$ is a well-defined homomorphism on
$\Z[1/2]/\Z$, the union of its cyclic subgroups generated by
$2^{-n-1}+\Z$.  It is injective because $[A_nD]$ and $2^{-n-1}+\Z$ have the
same order $2^{n+1}$.  The identities \eqref{eq:classes} follow from
Step 1.
\end{proof}

\subsection{Bounded factorizations of one-sided elements}\label{sec:onesided}

Let $S$ be a domain and $x$ a nonzero nonunit of $S$.  We consider
expressions $x=x_1\cdots x_r$ with arbitrary nonunits $x_i$ of $S$, not
assumed irreducible.  We say that $x$ is \emph{bounded} in $S$, or has
\emph{bounded factorization lengths} in $S$, if the lengths $r$ of these
expressions are bounded.  A longest such expression is a product of
atoms, since a factor that is not an atom could be split into two
nonunits to give a longer expression.  So a bounded element has a
factorization into atoms, and a domain in which every nonzero nonunit is
bounded is a BFD (Section~\ref{sec:factorization-background}).  This
convention lets us prove boundedness before atomicity is known.

The next lemma lets us compare factorizations in a subring $S$ with
valuations in $D$: a nonunit of $S$ is a nonunit of $D$.

\begin{lemma}[Units in subrings]\label{lem:units}
For every subring $S\subseteq D$, one has
$S^\times=S\cap D^\times$.
\end{lemma}

\begin{proof}
Every element of $D$ is integral over $\Z\subseteq S$.  Let $s\in S$ have
an inverse in $D$, and let $s^{-k}+a_{k-1}s^{-(k-1)}+\cdots+a_0=0$ with
$a_i\in\Z$.  Multiplying by $s^{k-1}$ expresses $s^{-1}$ as
$-(a_{k-1}+a_{k-2}s+\cdots+a_0s^{k-1})\in S$.
\end{proof}

\begin{definition}\label{def:onesided}
A nonzero $x\in D$ is \emph{one-sided} if $v_{\EA}(x)=0$ or
$v_{\EB}(x)=0$, that is, if $x\notin\EA\cap\EB$.
\end{definition}

The proof below shows that, in a factorization of a one-sided element
$x$, every factor has positive valuation at $\EA$ or $\EB$ or at one of
finitely many maximal ideals determined by $x$.  An element of
$\EA\cap\EB$ behaves differently: the classes of $P_n$ and $Q_n$ cancel,
and $y_{n-1}/y_n$ has divisor $P_nQ_n$.  So
$y_0=y_n\prod_{k=1}^n(y_{k-1}/y_k)$ is a product of $n+1$ nonunits of $D$
for every $n$, and $y_0$ is not bounded in $D$.

\begin{proposition}\label{prop:onesided}
Assume \textup{(T1)--(T3)}.  Let $S\subseteq D$ be a subring.  Each
one-sided nonzero nonunit of $S$ has bounded factorization lengths in
$S$, with a bound that can be chosen independently of $S$.
\end{proposition}

\begin{proof}
We bound the number of factors that can meet $\EA$, and show that every
other factor must have positive valuation at one of finitely many
maximal ideals fixed in advance.

\emph{Step 1: the cutoff.}  Suppose first that
$(v_{\EA}(x),v_{\EB}(x))=(\kappa_A,0)$ with $\kappa_A>0$.  By
\eqref{eq:tails}, we may choose $n_1$ such that $v_{P_j}(x)=\kappa_A$ and
$v_{Q_j}(x)=0$ for $j>n_1$, and such that
$\kappa_A\sum_{j>n_1}2^{-j-1}<1$.  Let $\mathcal H$ be the set of maximal
ideals at which $x$ has positive valuation, other than $\EA$ and the
$P_j$ with $j>n_1$.  It is finite by Lemma~\ref{lem:geometry}.

\emph{Step 2: the factors.}  Consider a factorization $x=x_1\cdots x_r$ in
nonunits of $S$.  Each factor is a nonunit of $D$ by
Lemma~\ref{lem:units}, and divides $x$ in $D$.  Since valuations are
additive, at most $\kappa_A$ factors have positive valuation at $\EA$.  A
remaining factor has valuation zero at $\EA$ and at $\EB$.  By
Lemma~\ref{lem:geometry} it therefore has finite support, and this
support is nonempty because the factor is a nonunit of $D$.  Suppose that
this support were contained in $\{P_j:j>n_1\}$.  Then the principal ideal
of the factor would be $\prod_{j>n_1}P_j^{e_j}D$, with
$0\leq e_j\leq\kappa_A$ because the factor divides $x$, and with not all
$e_j$ zero.  By Lemma~\ref{lem:classes} its class would be
\[
 \iota\left(\sum_{j>n_1}e_j2^{-j-1}+\Z\right),
\]
where the finite sum lies strictly between zero and one.  This class is
nonzero because $\iota$ is injective, which contradicts principality.
Hence every remaining factor has positive valuation at a member of
$\mathcal H$.

\emph{Step 3: the bound.}  Again by additivity of valuations,
\[
 r\leq \kappa_A+\sum_{V\in\mathcal H}v_V(x).
\]
The case $(v_{\EA}(x),v_{\EB}(x))=(0,\kappa_B)$ with $\kappa_B>0$ is
symmetric.  If $v_{\EA}(x)=v_{\EB}(x)=0$, the whole support of $x$ is
finite, each factor has positive valuation somewhere in it, and
$r\leq\sum_Vv_V(x)$.  These bounds involve only $x$ and $D$, not $S$.
\end{proof}

\subsection{The ring of integers of the tower}\label{sec:ambient}

We now prove the properties of $D$ stated in Theorem~\ref{thm:A}.  The
conjugates $\sigma_n(\beta_n)$ of the elements $\beta_n$ of (T2) split an element with both valuations positive
into two one-sided factors, the elements $y_n$ give an ascending chain of
principal ideals, and finite generation distinguishes $\EA$ and $\EB$ from
the other maximal ideals.

\begin{theorem}\label{thm:ambient}
Assume \textup{(T1)--(T3)}.  The almost Dedekind domain $D$ is atomic,
does not satisfy ACCP, and has exactly two non-invertible maximal ideals,
namely $\EA$ and $\EB$.
\end{theorem}

\begin{proof}
We prove the three assertions in turn.

\emph{Step 1: atomicity.}  Let $x$ be a nonzero nonunit of $D$.  If $x$ is
one-sided, Proposition~\ref{prop:onesided} with $S=D$ shows that $x$ is
bounded in $D$, so a longest expression of $x$ is a factorization into
atoms.  It remains to treat $x$ with $\kappa_A=v_{\EA}(x)$ and
$\kappa_B=v_{\EB}(x)$ both positive.  By \eqref{eq:tails}, choose $n_1$
such that $v_{P_j}(x)=\kappa_A$ and $v_{Q_j}(x)=\kappa_B$ for $j>n_1$.
The automorphism $\sigma_n$ interchanges the two primes $A_n,P_n$ over
$A_{n-1}$ and the two primes $B_n,Q_n$ over $B_{n-1}$, so for each $n$
the element $\sigma_n(\beta_n)$ has ideal $P_nB_n$.  Choose $\kappa_B$
distinct indices $i_1,\ldots,i_{\kappa_B}>n_1$, and put
$x'=\prod_k\sigma_{i_k}(\beta_{i_k})$.  Then
$(v_{\EA}(x'),v_{\EB}(x'))=(0,\kappa_B)$, the valuations of $x'$ at the
$P_j$ are at most one, those at the $Q_j$ are at most $\kappa_B$, and
all its support outside $\EA,\EB$ consists of primes $P_j,Q_j$ with
$j>n_1$.  By Lemma~\ref{lem:geometry}, $x'\mid_Dx$.  Both $x'$ and $x/x'$
are one-sided nonunits of $D$, since $v_{\EB}(x')=\kappa_B>0$ and
$v_{\EA}(x/x')=\kappa_A>0$.  Proposition~\ref{prop:onesided} gives atomic
factorizations of both, and their product is an atomic factorization of
$x$.

\emph{Step 2: failure of ACCP.}  Since $(y_{n-1})=A_{n-1}B_{n-1}$, (T1) gives
$y_{n-1}\OO_{F_n}=A_nP_nB_nQ_n$, so the quotient $y_{n-1}/y_n$ has ideal
$P_nQ_n$ in $\OO_{F_n}$.  It is therefore an integral nonunit of $D$, and
\[
 y_0D\subsetneq y_1D\subsetneq y_2D\subsetneq\cdots
\]
is a strictly increasing chain of principal ideals.

\emph{Step 3: the non-invertible maximal ideals.}  Let
$V\notin\{\EA,\EB\}$.  By prime avoidance, choose $x\in
V\setminus(\EA\cup\EB)$.  Its support is finite by
Lemma~\ref{lem:geometry}.  Adjoin to $x$ an element of $D$ that is a
uniformizer at $V$ (a numerator of a uniformizer of $D_V$) and, for every
other maximal ideal in the support of $x$, an element of $V$ outside that
maximal ideal.  These finitely many elements generate $V$, since the
ideal they generate is contained in $V$ and agrees with $V$ after
localization at every maximal ideal.  A finitely generated nonzero ideal
of an almost Dedekind domain is invertible
(Section~\ref{sec:factorization-background}), so $V$ is invertible.  In
contrast, any finite set of elements of $\EA$ lies in every sufficiently
late $P_j$, by \eqref{eq:tails}, and so cannot generate $\EA$.  Likewise
$\EB$ is not finitely generated.  Since every invertible ideal is finitely
generated, $\EA$ and $\EB$ are precisely the non-invertible maximal ideals.
\end{proof}

Lemma~\ref{lem:geometry} and Theorem~\ref{thm:ambient}, together with the residue fields recorded in (T1), prove parts
(a)--(d) of Theorem~\ref{thm:A} for every tower satisfying (T1)--(T3);
part (a) uses only (T1).  Theorem~\ref{thm:A} then follows from the
existence of such a tower, which Theorem~\ref{thm:construction} provides.

\section{Construction of the tower}\label{sec:construction}

This section constructs the tower of Theorem~\ref{thm:A}.
Theorem~\ref{thm:construction} produces an infinite tower
$F_0=\Q(\sqrt{65})\subset F_1\subset\cdots$ that satisfies the hypotheses
\textup{(T1)--(T6)} of Section~\ref{sec:factorization} and in which every
step is a scaled step (Definition~\ref{def:scaled-step}).  With
Lemma~\ref{lem:geometry} and Theorem~\ref{thm:ambient}, this proves
Theorem~\ref{thm:A}.  Scaled steps are also the steps to which
Theorem~\ref{thm:B} applies.

\emph{What one step must achieve.}  The tower is built one quadratic step
at a time.  Suppose that $F=F_{n-1}$ has been constructed, together with
primes $A=A_{n-1}$ and $B=B_{n-1}$ and an element $y=y_{n-1}$ with
$(y)=AB$.  The next field is $L=F(\alpha)$, where $\alpha$ is a root of a
polynomial $X^2-tX+y$ with $t\in\OO_F$, and $\alpha$ becomes $y_n$.  Then
$N_{L/F}\alpha=y$ automatically, and the construction consists in choosing
$t$.  The field $L$ must be totally real, since the assumptions at the
next level require it; the step arranges this by making the discriminant
of $X^2-tX+y$ totally positive.  Beyond this, the step must achieve five
things.
\begin{enumerate}
\item[(1)] \emph{The chains continue} (hypotheses (T1) and (T2)).  The
primes $A$ and $B$ split in $L$; the ideal $(\alpha)$ is the product of one
prime $A'$ above $A$ and one prime $B'$ above $B$; and a second element
$\beta$ of norm $y$ generates $A'Q$, where $Q$ is the other prime of $L$
above $B$.  Together with the injectivity in (2), these identities give the
class of $A'$ order $2^{n+1}$ in $\Cl(L)$.
\item[(2)] \emph{No capitulation and square unit norms} (hypotheses (T3)
and (T4)).  No nontrivial ideal class of $F$ becomes principal in $L$, and
the norm of every unit of $L$ is a square in $F$.  Both follow from making
finitely many chosen primes of $F$, the detecting primes of
Definition~\ref{def:detecting}, ramify in $L$.  The residue characters at
these primes detect every nonzero class of the group $\mathcal V(F)$ of
square classes of virtual units (Section~\ref{sec:nt-background}).
\item[(3)] \emph{All other branches are eventually inert} (hypothesis
(T1)).  At each step, finitely many primes of $F$ carry an inert
prescription: $L$ is required to be unramified and inert there.  New
ramification is allowed only over rational primes that are unramified in
$F$, which keeps the absolute ramification indices at odd primes at most
two.
\item[(4)] \emph{Residue offsets} (hypothesis (T6)).  At finitely many odd
primes, the step prescribes whether the prime splits or is inert, and it
prescribes the two-primary component of the residue of $\alpha$.  Along a
branch that stays inert, the relation
$N_{F_n/F_{n-1}}y_n=y_{n-1}$ doubles the order of this component at each
step, once the residue fields have order congruent to $1$ modulo $4$ and
the component is nontrivial.  Prescribing the component at finitely many
steps therefore fixes the residue offsets $o_V$ and makes them distinct
over each rational prime.
\item[(5)] \emph{An isolated ramified prime} (hypothesis (T5)).  One of
the detecting primes is ramified in $L$, and every other prime of $F$ over
the same rational prime is unramified in $L$.
\end{enumerate}

\emph{Why the scaled family.}  Apart from the existence of $\beta$, each
requirement is a local condition on $L$ at finitely many primes of $F$.  The
scaled family \eqref{eq:scaled-family} of Section~\ref{sec:scaled-family}
meets the requirement on $\beta$ by an identity, and it turns every local
condition into a congruence condition on a single element $c\in F$.  For
$\beta=a+u\alpha$ with $a\in F$ and $u\in F^\times$ one has
$N_{L/F}\beta=a^2+aut+u^2y$, so $N_{L/F}\beta=y$ exactly when
$t=(y(1-u^2)-a^2)/(au)$.  The family takes $a=c(u-1)/h$, where $h=2^H$ is
the scale and $u$ is a power of a fixed unit of $F_0$ with $v_2(u-1)=H$ at
both dyadic primes \eqref{eq:dyadic-depth}.  This makes $a$ a unit at the
dyadic primes, which allows them to stay inert.  The discriminant of
$X^2-tX+y$ is then $\delta_-\delta_+/(chu)^2$, where $\delta_-$ and
$\delta_+$ are explicit polynomials in $c$, $u$ and $h^2y$.  At a prime $V$
of $F$ at which this discriminant is nonzero, the splitting of $V$ in $L$
and the residues of $\alpha$ therefore depend only on the classes of $c$
and $u$ modulo a sufficiently high power of $V$.  Finally, $c$ is chosen
with $(c)=AR$ for an auxiliary prime $R$.  The factor $A$ places $\alpha$
and $\beta$ at the same prime of $L$ above $A$, and $R$ is chosen with
$u\equiv-1\pmod R$, which makes $t$ integral.  Such a prime $R$, compatible
with all congruence conditions on $c$, exists by the Chebotarev density
theorem applied to a ray class field (Section~\ref{sec:nt-background}).

\emph{Organization.}  Section~\ref{sec:base-field} computes the arithmetic
of the base field $F_0=\Q(\sqrt{65})$ (Lemma~\ref{lem:base-field}).
Section~\ref{sec:scaled-family} states the assumptions under which one
step is made and defines the scaled family.
Sections~\ref{sec:inert-prescriptions} and~\ref{sec:residue-prescriptions}
solve the local problems at a single odd prime:
Lemma~\ref{lem:scaled-holding} realizes an inert prescription, and
Lemma~\ref{lem:scaled-residues} realizes a prescribed two-primary
component of the residue of $\alpha$ once the residue field is large.
Both lemmas count points on curves over finite fields.  The rest of the
section assembles the step and the tower.
Lemma~\ref{lem:order-avoidance} finds primes with prescribed Frobenius in
a finite abelian extension at which a fixed unit has residue order prime
to given odd primes;
Definition~\ref{def:detecting} defines detecting primes;
Proposition~\ref{prop:scaled-step} combines all local conditions into one
quadratic step; and Definition~\ref{def:scaled-step} defines scaled steps
by the properties of that step used in later sections.  A finite procedure
(Section~\ref{sec:offsets}) then separates the residue offsets, and Definition~\ref{def:schedule} and
Theorem~\ref{thm:construction} assemble the infinite tower.

\subsection{The base field}\label{sec:base-field}

The tower starts from the real quadratic field $F_0=\Q(\sqrt{65})$.  Its
class group has order two and is generated by the class of $A_0$, which
starts the chain of classes $[A_n]$ of order $2^{n+1}$.  The unit
$\epsilon_0=8+\sqrt{65}$ supplies the unit $\epsilon_H$ of the scaled
family.  The following lemma collects the facts about $F_0$ that the
construction uses.  Each can be checked by hand.  Put
\[
 \omega=\frac{1+\sqrt{65}}2,\qquad \epsilon_0=8+\sqrt{65},\qquad
 y_0=-16+3\sqrt{65},
\]
and recall from (T1) that $A_0=(47,\sqrt{65}-21)$ and
$B_0=(7,\sqrt{65}-3)$.

\begin{lemma}[The base field]\label{lem:base-field}
\begin{enumerate}
\item[\textup{(a)}] $\OO_{F_0}=\Z[\omega]$, where $\omega^2-\omega-16=0$,
and the discriminant of $F_0$ is $65$.  The prime $2$ splits as
$2\OO_{F_0}=\mathfrak d_0\mathfrak e_0$ with
\[
 \mathfrak d_0=(2,\omega-1),\qquad \mathfrak e_0=(2,\omega),
\]
the prime $3$ is inert, and $A_0$ and $B_0$ are primes of residue degree
one above $47$ and $7$.  In particular
$(F_0)_{\mathfrak d_0}=(F_0)_{\mathfrak e_0}=\Q_2$,
$(F_0)_{A_0}=\Q_{47}$ and $(F_0)_{B_0}=\Q_7$.
\item[\textup{(b)}] $(\omega+4)=\mathfrak e_0^2$, and
$\Cl(F_0)\simeq\Z/2\Z$, generated by each of the classes $[\mathfrak e_0]$,
$[A_0]$ and $[B_0]$.
\item[\textup{(c)}] $(y_0)=A_0B_0$.
\item[\textup{(d)}] $\epsilon_0$ is a unit of norm $-1$, and
$\epsilon_0^2-1=16\epsilon_0$.  Its residue at $A_0$ is $29$, of order
$46$ in $\F_{47}^\times$, and its residue at $B_0$ is $4$, of order $3$ in
$\F_7^\times$.
\item[\textup{(e)}] $\epsilon_0\equiv1\pmod8$ at $\mathfrak d_0$ and
$\epsilon_0\equiv-1\pmod8$ at $\mathfrak e_0$.  Hence $\epsilon_0^{23}$ is a
square in $(F_0)_{\mathfrak d_0}$, lies in the square class of $-1$ in
$(F_0)_{\mathfrak e_0}$, and has nonsquare residue at $A_0$.
\end{enumerate}
\end{lemma}
\begin{proof}
We verify the statements in order; all computations are with explicit
integers.

\emph{Step 1: part (a).}  Since $65$ is squarefree and
$65\equiv1\pmod4$, the ring of integers is $\Z[\omega]$ and the
discriminant is $65$ \cite[Example~2.41]{MilneANT}.  One checks
$\omega^2-\omega-16=0$ directly.  By the Dedekind--Kummer theorem
\cite[Theorem~3.41]{MilneANT}, the factorization of a rational prime $p$
is read off from the factorization of $X^2-X-16$ modulo $p$.  Modulo $2$
this polynomial is $X(X-1)$, so $2\OO_{F_0}=(2,\omega)(2,\omega-1)$ with
both factors of residue degree one.  Modulo $3$ it is $X^2-X-1$, which
takes the values $2,2,1$ at $0,1,2$; so $3$ is inert.  Since
$21^2=65+8\cdot47$ and $3^2=65-8\cdot7$, the number $65$ is a square
modulo $47$ and modulo $7$, so $47$ and $7$ split
\cite[Example~3.44]{MilneANT}.  Since
$X^2-X-16\equiv(X-11)(X-37)\pmod{47}$ and $X^2-X-16\equiv(X-2)(X-6)\pmod7$,
Theorem~3.41 of \cite{MilneANT} shows that $A_0=(47,\omega-11)$ and
$B_0=(7,\omega-2)$ are primes of residue degree one; note that
$\sqrt{65}-21=2(\omega-11)$ and $\sqrt{65}-3=2(\omega-2)$.
A prime of residue degree one over $p$ in a quadratic field that splits
$p$ has completion $\Q_p$, which gives the last assertion.

\emph{Step 2: part (b).}  The Minkowski bound of $F_0$ is
$\frac{2!}{2^2}\sqrt{65}=\sqrt{65}/2<5$ \cite[Theorem~4.3]{MilneANT}, so
every ideal class contains an integral ideal of norm at most $4$.  Such an
ideal is a product of primes above $2$ and $3$.  The ideal $(3)$ is prime
and principal, and $\mathfrak d_0\mathfrak e_0=(2)$, so $\Cl(F_0)$ is
generated by $[\mathfrak e_0]=-[\mathfrak d_0]$.  The conjugate of $\omega$
is $1-\omega$, so $N(\omega+4)=(\omega+4)(5-\omega)=-\omega^2+\omega+20=4$.
Moreover $\omega+4\in(2,\omega)=\mathfrak e_0$, and
$\omega+4\equiv\omega\equiv1\pmod{\mathfrak d_0}$.  Hence
$(\omega+4)=\mathfrak e_0^2$, and $2[\mathfrak e_0]=0$.  It remains to show
that $\mathfrak e_0$, $A_0$ and $B_0$ are not principal.  A generator of an
ideal of norm $p$ has the form $(a+b\sqrt{65})/2$ with $a,b\in\Z$ and norm
$(a^2-65b^2)/4=\pm p$, so $a^2-65b^2=\pm4p$.  For $p=2$, $47$ and $7$ one
has $4p\equiv3\pmod5$, so $a^2\equiv\pm3\pmod5$.  The squares modulo $5$
are $0,1,4$, which excludes this.  Thus $\Cl(F_0)\simeq\Z/2\Z$, and each of
the three classes generates it.

\emph{Step 3: part (c).}  Modulo $A_0$ one has $\sqrt{65}\equiv21$, so
$y_0\equiv-16+63=47\equiv0$.  Modulo $B_0$ one has $\sqrt{65}\equiv3$, so
$y_0\equiv-16+9=-7\equiv0$.  Thus $A_0B_0$ divides $(y_0)$.  Since
$N(y_0)=256-9\cdot65=-329=-47\cdot7$, the ideal $(y_0)$ has norm $329$,
and $(y_0)=A_0B_0$.

\emph{Step 4: part (d).}  $N(\epsilon_0)=64-65=-1$ and
$\epsilon_0^2=129+16\sqrt{65}=1+16\epsilon_0$.  Modulo $A_0$ the residue
of $\epsilon_0$ is $8+21=29$.  Here $29^2=841\equiv42\not\equiv1\pmod{47}$,
and $29^{23}\equiv-1\pmod{47}$ by Euler's criterion, because $29$ is not a
square modulo $47$: by quadratic reciprocity
$\bigl(\tfrac{29}{47}\bigr)=\bigl(\tfrac{47}{29}\bigr)
=\bigl(\tfrac{18}{29}\bigr)=\bigl(\tfrac{2}{29}\bigr)=-1$, since
$29\equiv5\pmod8$.  So the order of $29$ divides $46$ but not $2$ or $23$,
and it equals $46$.  Modulo $B_0$ the residue of $\epsilon_0$ is
$8+3=11\equiv4$, and $4^3=64\equiv1\pmod7$ with $4\ne1$, so the order is
$3$.

\emph{Step 5: part (e).}  Note that $\epsilon_0=2\omega+7$ and that
$\omega(\omega-1)=16$.  At $\mathfrak d_0=(2,\omega-1)$ the element
$\omega$ is a unit, so $v_2(\omega-1)=4$, $\omega\equiv1\pmod{16}$ and
$\epsilon_0\equiv9\equiv1\pmod8$.  At $\mathfrak e_0=(2,\omega)$ the element
$\omega-1$ is a unit, so $v_2(\omega)=4$ and
$\epsilon_0\equiv7\equiv-1\pmod8$.  An element $x\in\Z_2$ with
$x\equiv1\pmod8$ is a square, by Newton's lemma
\cite[Theorem~7.32]{MilneANT} applied to $X^2-x$ at $X=1$: there
$|1-x|_2\le2^{-3}<2^{-2}=|2|_2^2$.  Since $23$ is odd, $\epsilon_0^{23}$ is
congruent to $1$ modulo $8$ at $\mathfrak d_0$, hence a square there, and
$-\epsilon_0^{23}$ is congruent to $1$ modulo $8$ at $\mathfrak e_0$, so
$\epsilon_0^{23}$ lies in the square class of $-1$ there.  Finally the
residue of $\epsilon_0^{23}$ at $A_0$ is $29^{23}\equiv-1$, which is not a
square modulo $47$ because $47\equiv3\pmod4$.
\end{proof}

In particular, $F=F_0$ with $A=A_0$, $B=B_0$ and $y=y_0$ satisfies the
standing assumptions of Section~\ref{sec:scaled-family} with $n=1$.

\subsection{The scaled family}\label{sec:scaled-family}

This subsection fixes the setting of one step and defines the family of
quadratic extensions from which each step is chosen.

\emph{Standing assumptions.}  Let $n\ge1$, and write $F=F_{n-1}$,
$A=A_{n-1}$, $B=B_{n-1}$ and $y=y_{n-1}$.  We assume:
\begin{itemize}
\item $F$ is totally real;
\item $F$ has exactly two dyadic primes $\mathfrak d$ and $\mathfrak e$,
above $\mathfrak d_0$ and $\mathfrak e_0$ respectively, and both are
absolutely unramified;
\item the completions of $F$ at $A$ and $B$ are $F_A=\Q_{47}$ and
$F_B=\Q_7$;
\item $A$ lies over $A_0$ and $B$ lies over $B_0$;
\item $(y)=AB$, and the class $[A]$ has order $2^n$ in $\Cl(F)$.
\end{itemize}
Fix an integer $H\ge5$ and set
\begin{equation}\label{eq:scaled-family}
\begin{gathered}
 h=2^H,\qquad \epsilon=\epsilon_H
   =\epsilon_0^{46\cdot2^{H-4}}=(\epsilon_0^{23})^{2^{H-3}},\qquad u=\epsilon^k,\\
 a=c(u-1)/h,\qquad \delta_-=c^2-h^2y,\qquad
 \delta_+=c^2(u-1)^2-h^2y(u+1)^2,\\
 t=\frac{y(1-u^2)-a^2}{au}
   =-\frac{h^2y+c^2}{ch}+\frac{\delta_-}{chu},\qquad
 \delta=t^2-4y=\frac{\delta_-\delta_+}{(chu)^2},\\
 L=F(\alpha),\quad\alpha^2-t\alpha+y=0,\qquad
 \beta=a+u\alpha.
\end{gathered}
\end{equation}
Here $c\in F^\times$ and the positive integer $k$ are parameters.  The
field $L$ is the candidate for $F_n$, with $y_n=\alpha$ and
$\beta_n=\beta$.  We will choose $(c)=AR$, with $R$ an \emph{auxiliary
prime} of absolute residue degree one selected in the proof of
Proposition~\ref{prop:scaled-step}, $c$ totally positive, and $k$ in an
infinite arithmetic progression of positive integers.

The formulas in \eqref{eq:scaled-family} are explained as follows.  If
$\alpha'$ is the other root of $X^2-tX+y$, then
$(a+u\alpha)(a+u\alpha')=a^2+aut+u^2y$.  So the first expression for $t$
is exactly the condition that $\beta$ have norm $y$.  Substituting
$a=c(u-1)/h$ into it gives the second expression for $t$, and expanding
$t^2-4y$ gives the factorization of $\delta$.  Thus $L=F(\sqrt\delta)$,
and $\delta$ has the same square class as $\delta_-\delta_+$.

\emph{The unit $u$ at the dyadic primes, at $A$ and at $B$.}  We use the
normalization $v_2(2)=1$ (Section~\ref{sec:notation}).  At both dyadic
primes,
\begin{equation}\label{eq:dyadic-depth}
 v_2(\epsilon^k-1)=H,\qquad (\epsilon^k-1)/2^H\equiv1\pmod2
 \quad(k\text{ odd}).
\end{equation}
This follows by repeated squaring of $\epsilon_0^{46}$, starting from the
identity $\epsilon_0^2-1=16\epsilon_0$ of
Lemma~\ref{lem:base-field}(d).  That identity gives
$v_2(\epsilon_0^2-1)=4$ at $\mathfrak d_0$ and $\mathfrak e_0$.  Since
$\epsilon_0^{46}-1=(\epsilon_0^2-1)(1+\epsilon_0^2+\cdots+\epsilon_0^{44})$
and the second factor is a sum of $23$ terms congruent to $1$ modulo $2$,
it is a unit, and
$v_2(\epsilon_0^{46}-1)=4$.  If $v_2(x-1)=s\ge2$, then $v_2(x+1)=1$ and
$v_2(x^2-1)=s+1$.  Applying this $H-4$ times gives $v_2(\epsilon-1)=H$.
The same factorization argument, with $k$ odd in place of $23$, gives
$v_2(\epsilon^k-1)=H$.  Then $(\epsilon^k-1)/2^H$ is a unit of
$(F_0)_{\mathfrak d_0}=\Q_2$, respectively $(F_0)_{\mathfrak e_0}=\Q_2$,
so it is $1$ modulo $2$.  These formulas persist in unramified
extensions, in particular at the dyadic primes of $F$.

Since $F_A=\Q_{47}$ and $F_B=\Q_7$, the residue fields at $A$ and $B$
are $\F_{47}$ and $\F_7$.  Hence $\epsilon\equiv1\pmod A$, because $46$
divides the exponent of $\epsilon$.  By Lemma~\ref{lem:base-field}(d),
the residue of $\epsilon_0$ at $B$ is a root of $X^2-16X-1$ modulo $7$,
that is, $4$ or $5$, of order three or six in $\F_7^\times$.  The
exponent of $\epsilon$ is even and prime to $3$, so the residue of
$\epsilon$ at $B$ has order three.  These facts use only the residue
fields at $A$ and $B$.  Since $A$ and $B$ lie over $A_0$ and $B_0$, the
residues of $\epsilon_0$ at $A$ and $B$ are $29$ and $4$, as at $A_0$
and $B_0$; Step~5 of the proof of Proposition~\ref{prop:scaled-step}
uses that $29$ is not a square modulo $47$.  Thus
$k\equiv1\pmod6$ ensures that $a$ is a unit at the dyadic primes and at
$B$, while $v_A(a)\ge2$.  Indeed, $k$ odd gives $v_2(u-1)=H=v_2(h)$;
$k\equiv1\pmod3$ gives $u\equiv\epsilon\not\equiv1\pmod B$; and
$v_A(u-1)\ge1$.  Here $c$ is a unit at the dyadic primes and at $B$, and
$v_A(c)=1$, as arranged in Proposition~\ref{prop:scaled-step}.

\subsection{Inert prescriptions}\label{sec:inert-prescriptions}

Hypothesis (T1) requires every branch other than $(A_n)$ and $(B_n)$ to
be eventually unramified and inert.  At
each step the construction therefore imposes the following condition at
finitely many primes.  An \emph{inert prescription} at a prime $V$ of $F$
is the requirement that $V$ be unramified and inert in $L$; the finite set
of primes carrying an inert prescription at a step is its \emph{inert
set}.  At an odd prime $V$, an inert prescription holds exactly when
$F_V(\sqrt\delta)$ is the unramified quadratic extension of $F_V$, that is,
when $\delta$ is a nonsquare unit times an even power of a uniformizer
(Section~\ref{sec:nt-background}, quadratic Kummer theory).  This is a
condition on the square class of $\delta_-\delta_+$ in $F_V$.

\emph{The Weil bound.}  To meet this condition we must show that a
polynomial in the residue of $c$ takes a nonsquare value at some point of
the residue field $\F_q$.  The residue prescriptions of
Section~\ref{sec:residue-prescriptions} need more.  There the residue of
$\alpha$ must have a prescribed two-primary component, and the residues
with this component form one coset of the subgroup of odd order in
$\F_q^\times$.  We need many elements of this coset at which a certain
discriminant is a nonzero square, so that one of them survives a few
exclusions.  In each case we count points on a double cover: of the
projective line, $Z^2=\Psi(X)$, for Lemma~\ref{lem:scaled-holding} and
Lemma~\ref{lem:scaled-residues}(1), and of a conic for
Lemma~\ref{lem:scaled-residues}(2).  For a smooth projective
geometrically connected curve $C$ of genus $g$ over $\F_q$, the Weil bound
of Section~\ref{sec:nt-background} states
\[
 \bigl|\#C(\F_q)-(q+1)\bigr|\le2g\sqrt q .
\]
We apply it to curves of the following form.  Let $q$
be odd, let $\chi$ be the quadratic character of $\F_q^\times$, extended by
$\chi(0)=0$ (this $\chi$ takes values in $\{0,\pm1\}$, unlike the residue
characters $\chi_v$ of Section~\ref{sec:nt-background}), and let $\Psi\in\F_q[X]$ be squarefree of even degree $2g+2$.
The affine curve $Z^2=\Psi(X)$ has a smooth projective geometrically connected
model $C$ of genus $g$.  Indeed, $C$ is glued from $Z^2=\Psi(X)$ and
$Z'^2=X'^{2g+2}\Psi(1/X')$ along $X'=1/X$, $Z'=Z/X^{g+1}$.  Both charts are
smooth because $\Psi$ and its reversal $X'^{2g+2}\Psi(1/X')$ are
squarefree, and $C$ is geometrically irreducible because $\Psi$ is not a
square in $\overline{\F}_q[X]$.  The double cover $X\colon C\to\mathbb P^1$
is branched exactly at the $2g+2$ roots of $\Psi$, so the
Riemann--Hurwitz formula gives $2g_C-2=-4+(2g+2)$ for the genus $g_C$ of
$C$, that is, $g_C=g$.  The points at infinity are $X'=0$,
$Z'=\pm\sqrt{\operatorname{lc}(\Psi)}$, where $\operatorname{lc}(\Psi)$ is
the leading coefficient.  Above each $x\in\F_q$ it has exactly
$1+\chi(\Psi(x))$ rational points, and it has two rational points at infinity
when the leading coefficient of $\Psi$ is a square in $\F_q^\times$.  In that
case
\[
 \#C(\F_q)=q+2+\sum_{x\in\F_q}\chi(\Psi(x)),\qquad\text{so}\qquad
 \sum_{x\in\F_q}\chi(\Psi(x))\ge-1-2g\sqrt q .
\]
Thus, if $\sqrt q$ is large compared with $g$, the sum of the values
$\chi(\Psi(x))$ is small compared with $q$, and $\Psi$ takes nonzero square
values at many points of $\F_q$.

\begin{lemma}[Inert prescription at an odd prime]\label{lem:scaled-holding}
Let $\mathcal K$ be a local field of odd residue characteristic and let
$y,h,u_0$ be units of $\mathcal K$.  Some unit $c\in\mathcal K$ makes
$\mathcal K(\sqrt\delta)$ the unramified quadratic extension of
$\mathcal K$, where $\delta$ is the discriminant in
\eqref{eq:scaled-family} formed with $u=u_0$.  This property
persists in sufficiently small neighbourhoods of $c,u_0$.
\end{lemma}
\begin{proof}
Since $\delta=\delta_-\delta_+/(chu_0)^2$, it suffices to choose $c$ so
that $\delta_-\delta_+$ is a nonsquare unit times an even power of a
uniformizer $\pi$.  We write $\F_q$ for the residue field and a bar for
residues.

\emph{Case 1: $y$ is a square in $\mathcal K$.}  Take $c$ near $h\sqrt y$
so that $v(\delta_-)=2$.  For instance, $c=h\sqrt y+\pi^2x$ with $x$ a unit
gives $\delta_-=\pi^2x(2h\sqrt y+\pi^2x)$, whose leading unit has residue
$2\bar h\sqrt{\bar y}\,\bar x$.  Since the derivative $2c$ of $\delta_-$
with respect to $c$ is a unit, this leading unit can be chosen freely by
the choice of $\bar x$.  At $c\equiv h\sqrt y$ the other factor $\delta_+$
has residue
$\bar h^2\bar y(\bar u_0-1)^2-\bar h^2\bar y(\bar u_0+1)^2
=-4\bar h^2\bar y\bar u_0\ne0$.  Choosing $\bar x$ so that the leading unit
of $\delta_-\delta_+$ is a nonsquare makes $\delta$ have nonsquare unit
part and even valuation.

\emph{Case 2: $\bar y$ is a nonsquare in $\F_q$.}  Write $c=h\xi$.  Then
$\delta_-\delta_+=h^4(\xi^2-y)\bigl(\xi^2(u_0-1)^2-y(u_0+1)^2\bigr)$, so
after rescaling $c$ by $h$ the relevant square class is that of the value
at $\bar\xi$ of
\[
 \Psi(X)=(X^2-\bar y)
       \bigl(X^2(\bar u_0-1)^2-\bar y(\bar u_0+1)^2\bigr).
\]
If $\Psi(\bar\xi)$ is a nonsquare in $\F_q^\times$ for some
$\bar\xi\in\F_q^\times$, then any lift $\xi$ of $\bar\xi$ makes
$\delta_-\delta_+$ a nonsquare unit.  We show that such a $\bar\xi$ exists.

For $\bar u_0\ne\pm1$, $\Psi$ is a squarefree quartic with no roots in
$\F_q$: its roots satisfy $X^2=\bar y$ or
$X^2=\bar y\bigl((\bar u_0+1)/(\bar u_0-1)\bigr)^2$, and $\bar y$ is a
nonsquare; the two values of $X^2$ differ because $\bar u_0\ne0$.  The
curve $Z^2=\Psi(X)$ has genus one.  It has two rational points at infinity,
since the leading coefficient $(\bar u_0-1)^2$ is a square, and two above
$0$, since $\Psi(0)=\bar y^2(\bar u_0+1)^2$ is a nonzero square.  If every
nonzero argument gave a square, the curve would have $2q+2$ rational
points.  This contradicts the genus-one bound $q+1+2\sqrt q<2q+2$.

If $\bar u_0=1$, then $\Psi(X)=-4\bar y(X^2-\bar y)$.  The identity
$\sum_{x\in\F_q}\chi(x^2-\bar y)=-1$ gives
$\sum_{x\in\F_q}\chi(\Psi(x))=-\chi(-\bar y)$, whereas $\Psi(0)=4\bar y^2$
is a nonzero square.  So the sum over the $q-1$ nonzero arguments, none of
which is a root, is $-\chi(-\bar y)-1\le0<q-1$.

If $\bar u_0=-1$, then $\Psi(X)=4X^2(X^2-\bar y)$.  At nonzero arguments
remove the square factor $4X^2$; the sum of $\chi(x^2-\bar y)$ over
$x\in\F_q^\times$ is $-1-\chi(-\bar y)\le0<q-1$.

In each case some nonzero argument has a nonsquare value.  Lift it to
$c/h$.  Finally, the square class of a nonzero element of $\mathcal K$ is
locally constant, and $\delta$ depends continuously on $c$ and $u_0$; this
proves persistence.
\end{proof}

\subsection{Residue prescriptions}\label{sec:residue-prescriptions}

Hypothesis (T6) concerns the two-primary components of the residues of the
elements $y_n$.  Let $V$ be an odd prime of $F$ other than $A$ and $B$, so
that $y$ is a unit at $V$, and suppose that $c$, $h$ and $u$ are units at
$V$.  The residue of $\alpha$ at a prime of $L$ above $V$ is a root of the
reduction of $X^2-tX+y$.  So prescribing it requires control of the trace
$t$, and this trace is not arbitrary.  Write $\F_q$ for the residue field
of $V$, and write $u$, $y$ and $\xi$ also for the residues of $u$, $y$
and $c/h$.  For fixed $u,y\in\F_q^\times$, the possible trace at
$\xi\in\F_q^\times$ is
\begin{equation}\label{eq:trace-map}
 t_u(\xi)=-\frac{u-1}{u}\xi-\frac{y(u+1)}{u\xi},
\end{equation}
by the second expression for $t$ in \eqref{eq:scaled-family}.  As in
Section~\ref{sec:notation}, $x_{(2)}$ denotes the two-primary component of
an element $x$ of a finite multiplicative group.

\begin{definition}\label{def:residue-prescription}
Let $V$ be an odd prime of $F$ other than $A$ and $B$, with residue field
$k(V)$.  A \emph{split
residue prescription} at $V$ is the requirement that $V$ split in $L$ and
that the residue of $\alpha$ at one of the two primes above $V$ have a
prescribed two-primary component $b\in k(V)^\times_{(2)}$.  An \emph{inert
residue prescription} at $V$ is the requirement that $V$ be inert in $L$
and that the residue of $\alpha$ in the quadratic extension of $k(V)$ have
a prescribed two-primary component $1$ or $-1$.  Both are called
\emph{residue prescriptions}.
\end{definition}

The next lemma realizes residue prescriptions over large residue fields.
In part (1) the root $x$ must lie in the coset of the odd-order subgroup
of $\F_q^\times$ determined by $b$.  This coset has $(q-1)/(q-1)_2$
elements, and the size condition makes it large compared with the genus of
the curve that the Weil bound is applied to.

\begin{lemma}[Residue prescriptions over large residue fields]\label{lem:scaled-residues}
Let $\F_q$ have odd characteristic, and let $u,y\in\F_q^\times$.
\begin{enumerate}
\item Put $(q-1)_2=2^{v_2(q-1)}$.  If $(q-1)/(q-1)_2>4\sqrt q+12$, then for
every $b\in\F_q^\times{}_{(2)}$ some $\xi\in\F_q^\times$ makes
$X^2-t_u(\xi)X+y$ have two distinct roots $x,y/x\in\F_q^\times$, nonzero
trace, and $x_{(2)}=b$.
\item If $q\equiv1\pmod4$, $q\ge81$, and $y_{(2)}=1$, then, for either
sign, some $\xi\in\F_q^\times$ makes that polynomial
irreducible, with nonzero trace and root $\theta\in\F_{q^2}$ satisfying
$\theta_{(2)}=\pm1$.
\end{enumerate}
\end{lemma}
\begin{proof}
In both parts we first choose the root ($x$ in part (1), $\theta$ in part (2))
with the prescribed two-primary component and then solve for $\xi$; the
Weil bound shows that a suitable root exists.

\emph{Step 1: which traces occur.}  Multiplying $t=t_u(\xi)$ by $u\xi$
shows that a given trace $t$ is attained at $\xi\in\F_q^\times$ exactly
when
\begin{equation}\label{eq:trace-inverse}
 (u-1)\xi^2+ut\xi+y(u+1)=0.
\end{equation}
For $u\ne\pm1$, this is a quadratic equation in $\xi$, and its
discriminant in the variable $\xi$ is
\[
 \operatorname{disc}_\xi=u^2t^2-4(u-1)(u+1)y=u^2(t^2-4y)+4y.
\]
A nonzero square value of $\operatorname{disc}_\xi$ suffices: then
\eqref{eq:trace-inverse} has two solutions in $\F_q$, and both are
nonzero because their product $y(u+1)/(u-1)$ is nonzero.  For $u=1$ or
$-1$, \eqref{eq:trace-map} reads $t_u(\xi)=-2y/\xi$ or $t_u(\xi)=-2\xi$,
so every nonzero trace is attained directly.

\emph{Step 2: part (1).}  Take $x=bY^{(q-1)_2}$ with
$Y\in\F_q^\times$ and $t=x+y/x$.  Since $\F_q^\times$ is cyclic of order
$q-1$, the map $Y\mapsto Y^{(q-1)_2}$ maps $\F_q^\times$ onto the subgroup of
odd order, so this parametrizes exactly the elements with two-primary
component $b$.  The polynomial $X^2-tX+y$ has the roots $x$ and $y/x$;
they are distinct when $x^2\ne y$, and the trace is nonzero when
$x^2\ne-y$.  Since $t^2-4y=(x^2-y)^2/x^2$, one has
\[
 x^2\operatorname{disc}_\xi=\Psi(x),\qquad
 \Psi(X)=u^2(X^2-y)^2+4yX^2.
\]
Suppose $u\ne\pm1$.  As a quadratic polynomial in $X^2$, $\Psi$ has
discriminant $16y^2(1-u^2)\ne0$ and constant term $u^2y^2\ne0$.  So $\Psi$
is squarefree, has nonzero constant term and has four distinct nonzero
roots.  Since $(q-1)_2$ is prime to the characteristic, the polynomial
$\Psi_b(Y)=\Psi(bY^{(q-1)_2})$ is squarefree of degree $4(q-1)_2$, with leading coefficient
the square $u^2b^4$.  Its double cover $Z^2=\Psi_b(Y)$ has genus $2(q-1)_2-1$ and two
rational points at infinity.  The curve bound gives
\[
 \sum_{Y\in\F_q}\chi\bigl(\Psi_b(Y)\bigr)\ge-1-(4(q-1)_2-2)\sqrt q .
\]
The number of $Y\in\F_q$ at which $\Psi_b(Y)$ is a nonzero square is therefore
at least $\bigl(q-4(q-1)_2-1-(4(q-1)_2-2)\sqrt q\bigr)/2$, because $\Psi_b$ has at most $4(q-1)_2$
zeros.  Discard $Y=0$, and the at most $4(q-1)_2$ arguments for which $x^2=y$ or
$x^2=-y$ (each equation has at most two solutions $x$, each with at most
$(q-1)_2$ preimages $Y$).  A lower bound for the remaining arguments with
nonzero square value is
\[
 \frac{q-3-12(q-1)_2-(4(q-1)_2-2)\sqrt q}{2}>0,
\]
where positivity follows from $q-1>(q-1)_2(4\sqrt q+12)$ and $q>1$.  At such an argument,
$\operatorname{disc}_\xi=\Psi(x)/x^2$ is a nonzero square, so Step 1 gives
$\xi$.  This proves (1).  If $u=\pm1$, the coset has more than four
elements, so excluding $x^2=\pm y$ suffices.

\emph{Step 3: part (2).}  The odd-order subgroup of $\F_q^\times$ contains
$y$ and a unique $b$ with $b^2=y$, because squaring is bijective on a group
of odd order.  Let $\mathcal C$ be the norm-one conic for
$\F_{q^2}/\F_q$; its rational points are the $\zeta\in\F_{q^2}$ with
$\zeta^{q+1}=1$, and there are $q+1$ of them.  Put $\theta=\pm b\zeta^2$.
Because $v_2(q+1)=1$, the element $\zeta^2$ has odd order, so these
elements have norm $b^2=y$ and two-primary component $\pm1$.  Since
$\zeta^q=\zeta^{-1}$, the conjugate of $\theta$ is $\pm b\zeta^{-2}$.  Their
trace is $t=\pm b(\zeta^2+\zeta^{-2})$, so
$t^2-4y=y(\zeta^2-\zeta^{-2})^2$ and
\[
 \operatorname{disc}_\xi=u^2y\bigl((\zeta^2-\zeta^{-2})^2+4/u^2\bigr).
\]
The prefactor $u^2y=(ub)^2$ is a square in $\F_q$.  Let $u\ne\pm1$.  Over
$\overline{\F}_q$, the conic $\mathcal C$ is a projective line with
coordinate $\zeta$.  Its two points at infinity, $\zeta=0$ and
$\zeta=\infty$, are defined over $\F_{q^2}$ and interchanged by the
Frobenius of $\F_{q^2}/\F_q$, so neither is $\F_q$-rational.  The function in
parentheses on $\mathcal C$ has eight simple geometric zeros: their
equation is $\zeta^8+(4/u^2-2)\zeta^4+1=0$, which has distinct roots
because, as a quadratic in $\zeta^4$, its discriminant is
$16(1-u^2)/u^4\ne0$ and its constant term is $1$.  Its two geometric poles
are $\zeta=0$ and $\zeta=\infty$, of order four, and neither is
$\F_q$-rational.  The double cover of $\mathcal C$ defined by the square
root of this function is branched exactly at the eight zeros.  Multiplying
the function by the square $\zeta^4$ shows that this double cover is the
smooth projective model of $Z^2=\zeta^8+(4/u^2-2)\zeta^4+1$, a curve of the
type above with $2g+2=8$, so it has genus three.  Above a rational point of
$\mathcal C$ at which the function takes the value $x$ there are exactly
$1+\chi(x)$ rational points of the double cover, and no rational point
lies above a pole.  Since $\#\mathcal C(\F_q)=q+1$, the
Weil bound shows that the number of conic points giving a nonzero square
is at least
\[
 \frac{q+1-6\sqrt q-8}{2}.
\]
Exclude the points with $\zeta^4=1$ or $\zeta^4=-1$.  Such a point
satisfies $\zeta^8=1$, and since $\zeta^{q+1}=1$ and $\gcd(8,q+1)=2$, it
satisfies $\zeta^2=1$; so there are only two, namely $\zeta=\pm1$.  This
leaves at least $(q-11-6\sqrt q)/2>0$ points, using $q\ge81$.  These
exclusions say exactly that $\theta\notin\F_q$ and $t\ne0$.  At a remaining
point, Step 1 gives $\xi$, and $X^2-t_u(\xi)X+y$ has the roots $\theta$ and its
conjugate, so it is irreducible.

For $u=\pm1$, every nonzero trace is attained.  The elements of
$\F_{q^2}^\times$ of norm $y$ and two-primary component $\pm1$ are the
$(q+1)/2$ elements $\pm b\zeta^2$, and at most one of them, namely
$\pm b$, fails the two exclusions.  This proves (2).
\end{proof}

\emph{Growth of residue fields.}  Every finite list of odd residue fields
reaches both size conditions after sufficiently many inert quadratic
extensions.  Indeed, for $q_j=q_0^{2^j}$ and $j\ge1$,
\[
 2^{v_2(q_j-1)}=2^{v_2(q_0^2-1)+j-1},\qquad
 \frac{q_j-1}{2^{v_2(q_j-1)}\sqrt{q_j}}\longrightarrow\infty .
\]
After at least one inert step also $q\equiv1\pmod4$.  Once both
conditions and $q\equiv1\pmod4$ hold, they persist together under
further inert or split steps.

\emph{Combining the local conditions.}  At distinct primes these
prescriptions are independent finite congruences on $c$, once the classes
of $u$ modulo fixed powers of these primes are fixed.  Indeed, by
Lemma~\ref{lem:scaled-holding} an inert prescription holds for all $c$ in
a neighbourhood of a solution, and a residue prescription depends only on
the residues of $c/h$, $u$ and $y$ at the prime; the Chinese remainder
theorem combines conditions at distinct primes.  In
Proposition~\ref{prop:scaled-step} the classes of $u$ are fixed by
$u_0=\epsilon$ and retained by the choice of $k$.  The local extensions
are then the intended unramified ones.  Under a residue prescription, the
reduction of $X^2-tX+y$ at $V$ is separable, with two roots in the residue
field for a split prescription and none for an inert one; so $V$ is
unramified in $L$ and splits or is inert as prescribed.  At a prime of
the inert set, the discriminant square class of
Lemma~\ref{lem:scaled-holding} makes $V$ unramified and inert in $L$.

\subsection{Virtual units and detecting primes}\label{sec:detecting}
Hypotheses (T3) and (T4) require that a step $L/F$ make no nontrivial ideal
class of $F$ principal and that the norms of units of $L$ be squares of
units of $F$.  Both follow from one property: every virtual unit of $F$
that is a norm from $L$ is a square in $F$ (Step~9 of the proof of
Proposition~\ref{prop:scaled-step}).  A virtual unit that is a norm from
$L$ has a square residue, after an even power of a uniformizer is
removed, at every prime of odd residue characteristic that ramifies in
$L$.  It therefore
suffices that $L$ ramify at primes whose quadratic residue characters
separate the square classes of virtual units.  This subsection gives such
primes a name.

Recall from Section~\ref{sec:nt-background} that $x\in F^\times$ is a
virtual unit if $v_{\mathfrak p}(x)$ is even at every prime $\mathfrak p$
of $F$, and that the square classes of virtual units form the group
$\mathcal V(F)$.  The exact sequence there, from units modulo squares to
$\mathcal V(F)$ to $\Cl(F)[2]$, shows that $\mathcal V(F)$ is a
finite-dimensional $\F_2$-vector space.

Let $G$ be a prime of $F$ of odd residue characteristic.  The residue
character $\chi_G$ of Section~\ref{sec:nt-background} is defined on the
elements $x\in F^\times$ with $v_G(x)$ even, is a homomorphism, and
vanishes on squares.  Since every virtual unit has even valuation at $G$,
$\chi_G$ induces a homomorphism $\mathcal V(F)\to\F_2$.

\begin{definition}\label{def:detecting}
Let $L/F$ be a quadratic extension of number fields, let $y\in F^\times$
with $[y]\notin\mathcal V(F)$, and let $r=\dim_{\F_2}\mathcal V(F)$.
Primes $G_1,\dots,G_r$ of $F$ are \emph{detecting primes} for
$(F,y)$ in $L$ if:
\begin{enumerate}
\item[(a)] each $G_i$ has absolute residue degree one and is an isolated
ramified prime of $L/F$ (Definition~\ref{def:isolated}); the $G_i$ have
distinct odd rational characteristics; and $y$ is a unit and a square
modulo each $G_i$;
\item[(b)] there is a basis $[x_1],\dots,[x_r]$ of $\mathcal V(F)$ with
each $x_j$ a unit at every $G_i$ and $\chi_{G_i}(x_j)=1$ if $i=j$,
$\chi_{G_i}(x_j)=0$ if $i\ne j$.
\end{enumerate}
\end{definition}
Condition (b) says that the characters $\chi_{G_i}$, restricted to
$\mathcal V(F)$, form the dual basis of $[x_1],\dots,[x_r]$.  Later
arguments use only its consequence that the map $\mathcal V(F)\to\F_2^r$,
$[x]\mapsto(\chi_{G_i}(x))_i$, is injective.

\subsection{The auxiliary prime: ray classes and unit orders}\label{sec:auxiliary-prime}
In a scaled step the element $c$ generates $AR$, where $R$ is the
auxiliary prime, and $c$ must be totally positive and satisfy finitely
many congruences modulo powers of fixed primes.  These requirements
become one condition on $R$.  Suppose that $c_{\rm ref}\in\OO_F$ is
totally positive, satisfies the congruences, and has $v_A(c_{\rm
ref})=1$, and write $(c_{\rm ref})=AJ$.  Let $\mathfrak f$ be a modulus
whose finite part is a high power of the product of the fixed primes and
which contains every real place.  If a prime $R$ lies in the class of
$J$ in the ray class group $\Cl_{\mathfrak f}(F)$, then $R=J\cdot(w)$
with $w$ totally positive and $w\equiv1$ modulo the finite part
of $\mathfrak f$.  The element $c=c_{\rm ref}w$ then generates $AR$,
is totally positive and satisfies the same congruences.  Thus we need a
prime in a prescribed ray class.  Class field theory identifies
$\Cl_{\mathfrak f}(F)$ with the Galois group of the ray class field of
$\mathfrak f$, and a prime lies in a given class exactly when its
Frobenius element in that field is the corresponding Galois element
(Section~\ref{sec:nt-background}; \cite[Chapter~V, Theorems~3.5
and~3.6]{MilneCFT}).  The Chebotarev density theorem \cite[Chapter~VIII,
Theorem~7.4]{MilneCFT} then supplies infinitely many such primes.

The auxiliary prime must satisfy one further condition.  The exponent
$k$ in $u=\epsilon^k$ will satisfy $k\equiv1\pmod{\mathcal M}$ for a fixed
integer $\mathcal M$, in order to keep the congruence classes of $u$, and
also $u\equiv-1\pmod R$, which makes the trace $t$ integral at $R$.  Both
are possible when the order of $\epsilon$ modulo $R$ is $2s$ with $s$ odd
and prime to $\mathcal M$.  The two-adic valuation of this order is fixed
by a Hilbert symbol computation (Step~5 of the proof of
Proposition~\ref{prop:scaled-step}).  The next lemma controls the odd
part.  For an odd prime $\ell$, the condition that the order of $\epsilon$
modulo $R$ be prime to $\ell$ is a condition on the Frobenius element of
$R$ in a field obtained by adjoining roots of unity and a root of
$\epsilon$.  Kummer theory \cite[Chapter~VII, Appendix~A,
Theorem~A.3]{MilneCFT} computes the Galois group of this field, and the
proof shows that it is independent enough of the ray class field for
Chebotarev's theorem to impose both conditions at once.

\begin{lemma}\label{lem:order-avoidance}
Let $F$ be a totally real number field, let $\epsilon\in\OO_F^\times$ have
infinite order, let $\mathcal R/F$ be a finite abelian extension, and let
$\gamma\in\Gal(\mathcal R/F)$.  Let $\mathcal P$ be a finite set of odd
rational primes.  Then infinitely many primes $R$ of $F$ of absolute
residue degree one are unramified in $\mathcal R$, have Frobenius element
$\gamma$ in $\mathcal R/F$, and have the property that the order of the
residue of $\epsilon$ modulo $R$ is prime to every $\ell\in\mathcal P$.
These primes $R$ can be chosen outside any given finite set of primes.
\end{lemma}
\begin{proof}
We fix the $\ell$-adic valuation of $NR-1$ by a Frobenius condition in a
cyclotomic extension, require $\epsilon$ to be an $\ell^{s_\ell}$th
power modulo $R$ by a Frobenius condition in a Kummer extension, and apply
Chebotarev's theorem to the compositum of these extensions with
$\mathcal R$.  Excluding a finite set of rational characteristics
excludes only finitely many primes of $F$.

\emph{Step 1: the cyclotomic condition.}  Choose a prime $R_0$ of $F$ of
absolute residue degree one, unramified in $\mathcal R$, with Frobenius
$\gamma$, prime to every $\ell\in\mathcal P$ and outside the excluded
set, and put $s_\ell=v_\ell(NR_0-1)$.  Such primes exist by Chebotarev's
theorem, since primes of absolute residue degree at least two have
Dirichlet density zero \cite[Chapter~VI, Proposition~4.5 and Corollary~4.6]{MilneCFT}.  Let
$F_\mu=F(\mu_{\ell^{s_\ell+1}}:\ell\in\mathcal P)$ and let $\gamma'$ be the
Frobenius element of $R_0$ in $F_\mu/F$.  It agrees with $\gamma$ on
$F_\mu\cap\mathcal R$, since both are Frobenius elements of $R_0$.  For a
prime $R$ of absolute residue degree one not dividing $\ell$, the
Frobenius element of $R$ acts on $\mu_{\ell^{s_\ell+1}}$ by
$\zeta\mapsto\zeta^{NR}$.  Hence $\gamma'$ fixes $\mu_{\ell^{s_\ell}}$,
and every such $R$ with Frobenius $\gamma'$ in $F_\mu/F$ satisfies
$NR\equiv NR_0\pmod{\ell^{s_\ell+1}}$, so that $v_\ell(NR-1)=s_\ell$.  This
applies to the primes eventually selected.

\emph{Step 2: roots of $\epsilon$.}  Let $\ell\in\mathcal P$ with
$s_\ell>0$.  Since $\OO_F^\times$ is finitely generated and $\epsilon$ has
infinite order, we can write $\epsilon=\eta_\ell^{\ell^{a_\ell}}$ with
$\eta_\ell\in\OO_F^\times$ and $a_\ell$ maximal.  Then $\eta_\ell$ is not
an $\ell$th power in $F$, because an $\ell$th root of a unit is a unit.
Fix an embedding of an algebraic closure of $F$ into $\mathbb C$, so
that complex conjugation acts on it.  An element $x$ of $F$ that becomes
an $\ell$th power in $F_\mu$ was already an $\ell$th power in $F$.
Indeed, multiplying an $\ell$th root of $x$ in $F_\mu$ by a root of unity
in $\mu_\ell\subseteq F_\mu$ gives the real $\ell$th root of $x$ in
$F_\mu$.  This real root lies in the real subfield of the abelian
extension $F_\mu/F$, the fixed field of complex conjugation, which is
Galois and totally real over $F$.  Every conjugate of the root over $F$ is
a real root of the same binomial $X^\ell-x$ of odd degree, so uniqueness
of the real root puts it in $F$.  Consequently $\eta_\ell$ is not an
$\ell$th power in $F_\mu$.  Let $\alpha_\ell$ be the real root of
$X^{\ell^{s_\ell}}-\eta_\ell$.  By Kummer theory, $[F_\mu(\alpha_\ell):F_\mu]$
is the order of the class of $\eta_\ell$ in
$F_\mu^\times/F_\mu^{\times\ell^{s_\ell}}$ \cite[Chapter~VII,
Appendix~A, Theorem~A.3]{MilneCFT}.  This order is $\ell^{s_\ell}$: if
$\eta_\ell^{\ell^{s_\ell-1}}=w^{\ell^{s_\ell}}$ with $w\in F_\mu$, then
$\eta_\ell=\zeta w^\ell$ with $\zeta\in\mu_{\ell^{s_\ell-1}}$, and $\zeta$
is an $\ell$th power in $\mu_{\ell^{s_\ell}}\subseteq F_\mu$, so
$\eta_\ell$ would be an $\ell$th power in $F_\mu$.

\emph{Step 3: independence from $\mathcal R$.}  The extensions
$F_\mu(\alpha_\ell)/F_\mu$ have coprime degrees.  Their compositum
$F_\mu'$ has Galois group over $F_\mu$ the product of the corresponding
cyclic groups of odd order.  The conjugates of $\alpha_\ell$ over $F$ are
the elements $\zeta\alpha_\ell$ with $\zeta\in\mu_{\ell^{s_\ell}}$, so
$F_\mu'/F$ is Galois, and complex conjugation restricts to an element
$\tilde c$ of $\Gal(F_\mu'/F)$ that fixes every $\alpha_\ell$.  If $g\in
\Gal(F_\mu'/F_\mu)$ sends $\alpha_\ell$ to $\zeta_g\alpha_\ell$, then
$\tilde cg\tilde c^{-1}$ sends $\alpha_\ell$ to
$\zeta_g^{-1}\alpha_\ell$.  So conjugation by $\tilde c$ inverts
$\Gal(F_\mu'/F_\mu)$.  Consequently
\[
 F_\mu'\cap F_\mu\mathcal R=F_\mu.
\]
Indeed, the Galois group of this intersection over $F_\mu$ is a quotient
of $\Gal(F_\mu'/F_\mu)$ on which conjugation by $\tilde c$ is inversion.
Since $F_\mu\mathcal R/F$ is abelian, this conjugation is also trivial on
it, and a group of odd order on which inversion is trivial is trivial.

\emph{Step 4: the Frobenius class.}  By the Kummer pairing,
$g\mapsto(g\alpha_\ell/\alpha_\ell)_\ell$ is an isomorphism from
$\Gal(F_\mu'/F_\mu)$ onto $\prod_\ell\mu_{\ell^{s_\ell}}$.  Hence exactly
one extension of $\gamma'$ to $F_\mu'$ fixes every $\alpha_\ell$.  By
Step~3 and Step~1, this lift and $\gamma$ on $\mathcal R$ come from a
common element $\theta$ of $\Gal(F_\mu'\mathcal R/F)$.  Every conjugate of
$\theta$ has the same restriction to the abelian extension $F_\mu\mathcal
R$, and it still fixes each $\alpha_\ell$: if $\rho$ is a Galois element,
then $\rho^{-1}\alpha_\ell=\zeta\alpha_\ell$ with
$\zeta\in\mu_{\ell^{s_\ell}}$, and $\theta$ fixes $\zeta$ because
$\gamma'$ fixes $\mu_{\ell^{s_\ell}}$.  Chebotarev's theorem now gives a
set of positive density of primes of $F$, unramified in $F_\mu'\mathcal
R$, whose Frobenius class is that of $\theta$.  Removing the primes of
absolute residue degree at least two, the finitely many primes above
$\mathcal P$ and the excluded primes leaves
infinitely many primes $R$ of absolute residue degree one.

\emph{Step 5: the order of $\epsilon$.}  Let $R$ be such a prime.  Its
Frobenius element in $\mathcal R/F$ is $\gamma$, and
$v_\ell(NR-1)=s_\ell$ by Step~1.  Let $\ell$ satisfy $s_\ell>0$.  The
element $\alpha_\ell$ is integral, since $\eta_\ell$ is a unit, and its
Frobenius image is $\alpha_\ell$ itself, so its residue $\bar\alpha_\ell$
at a prime above $R$ satisfies $\bar\alpha_\ell^{NR}=\bar\alpha_\ell$.
Hence $\bar\alpha_\ell$ lies in the residue field of $R$, and $\eta_\ell$,
and thus $\epsilon$, is an $\ell^{s_\ell}$th power in that residue field.
The multiplicative group of the residue field is cyclic of order $NR-1$,
whose $\ell$-part is $\ell^{s_\ell}$, so the order of $\epsilon$ is prime
to $\ell$.  When $s_\ell=0$, the order $NR-1$ of this group is prime to
$\ell$ and the conclusion is immediate.
\end{proof}

\subsection{One quadratic step}\label{sec:one-step}
We now carry out a single step of the tower.  Starting from a field
$F=F_{n-1}$ with the data fixed before \eqref{eq:scaled-family}
(Section~\ref{sec:scaled-family}),
Proposition~\ref{prop:scaled-step} produces $L=F_n$ and the elements
$y_n=\alpha$ and $\beta_n=\beta$ with the properties that (T2)--(T5)
require at level $n$.  It also respects any finite list of local
prescriptions.  These prescriptions are what the construction in
Sections~\ref{sec:offsets} and~\ref{sec:schedule} uses to keep primes
unramified, to split them, and to fix their residue offsets.

\begin{proposition}[One quadratic step]\label{prop:scaled-step}
Let $F$, $A$, $B$, $y$ and $H\ge5$ satisfy the assumptions stated before
\eqref{eq:scaled-family}, and fix the following data.
\begin{enumerate}
\item[(a)] A finite inert set of primes of $F$ that contains both dyadic
primes and contains neither $A$ nor $B$.
\item[(b)] A finite set of \emph{forbidden characteristics}: rational
primes that newly selected primes must avoid.
\item[(c)] Finitely many split or inert residue prescriptions
(Definition~\ref{def:residue-prescription}) of the kinds provided by
Lemma~\ref{lem:scaled-residues}, at odd primes of $F$ other than $A$ and
$B$, with at most one prescription at each prime, and with the
hypotheses of that lemma at each of them: if $V$ carries a split residue
prescription, its residue field $\F_q$ satisfies
$(q-1)/(q-1)_2>4\sqrt q+12$ as in part~(1); if $V$ carries an inert
residue prescription, then $q\equiv1\pmod4$, $q\ge81$, and the residue
of $y$ at $V$ has trivial two-primary component, as in part~(2).
Primes with a split
residue prescription lie outside the inert set; an inert residue
prescription refines the inert prescription at that prime.
\end{enumerate}
Let $\mathcal X$ be the finite set of rational primes consisting of $47$
and $7$, the rational characteristics of $A$ and $B$; the rational
characteristics of the primes of the inert set, among them $2$; the
rational characteristics of the primes carrying residue prescriptions;
and the forbidden characteristics.  Then there are $c$, $R$ and an
infinite arithmetic progression of positive integers $k$ such that, for
every such $k$, \eqref{eq:scaled-family} gives $\delta_-,\delta_+$
totally positive, $a,t\in\OO_F$, and
\begin{gather*}
 N_{L/F}\alpha=N_{L/F}\beta=y,\qquad
 A\OO_L=A'P,\quad B\OO_L=B'Q,\\
 (\alpha)=A'B',\qquad (\beta)=A'Q,\qquad
 \ord[A']=2^{n+1},\\
 \Cl(F)\hookrightarrow\Cl(L),\qquad
 N_{L/F}(\OO_L^\times)=\OO_F^{\times2}.
\end{gather*}
All prescribed inert and residue conditions hold.  The choices can be
made so that the residues of $\alpha$ at $P$ and at $Q$ are both
nonsquares.  The extension $L/F$ has an isolated ramified prime
(Definition~\ref{def:isolated}) of absolute residue degree one whose
rational characteristic is not in $\mathcal X$.
\end{proposition}

The characteristics $47$ and $7$ are placed in $\mathcal X$ so that the
detecting primes and the auxiliary prime constructed in the proof lie
over neither of them.  In Step~2 of the proof, every other prime above
the rational characteristic of a detecting prime is added to the inert
set.  Since that characteristic is outside $\mathcal X$, these primes are
neither $A$ nor $B$, which must split, and carry no earlier prescription.

\begin{proof}
We first fix the behaviour of $c$ at finitely many primes (Steps~1--3),
then realize these conditions by a generator $c$ of $AR$, with the
auxiliary prime $R$ in a prescribed ray class and with a prescribed
order of $\epsilon$ modulo $R$ (Steps~4--6).  Then we choose the exponent
$k$ (Step~7) and verify the conclusions (Steps~8 and~9).  We use the
following form of the trace from \eqref{eq:scaled-family}:
\begin{equation}\label{eq:trace-cancelled}
 t=\frac{c^2(1-u)-h^2y(1+u)}{chu}
  =\frac{c(1-u)}{hu}-\frac{hy(1+u)}{cu}.
\end{equation}

\emph{Step 1: local conditions at the prescribed primes.}  Choose all
local models with $u_0=\epsilon$.  At odd primes of the inert set without
a residue prescription use Lemma~\ref{lem:scaled-holding}.  At a prime
with a split residue prescription use Lemma~\ref{lem:scaled-residues}(1),
and for an inert residue prescription with component $1$ or $-1$ use
Lemma~\ref{lem:scaled-residues}(2); their hypotheses hold by~(c).  For each of the two dyadic primes
$v$, choose $\bar c\in k_v^\times$ with
\[
 \operatorname{Tr}_{k_v/\F_2}(\bar y/\bar c^2)=1.
\]
Then \eqref{eq:dyadic-depth} gives $\bar a=\bar c$ and $\bar t=-\bar c$,
so the reduction of $X^2-tX+y$ at $v$ is separable and irreducible, and
$v$ is unramified and inert in $L$.  At $A$ prescribe $v_A(c)=1$, and at
$B$ prescribe a unit value.  For $k\equiv1\pmod6$ the residues of the
trace at $A$ and $B$ are then
\begin{equation}\label{eq:active-traces}
 \bar t_A=-2\bar h\,\overline{y/c},\qquad
 \bar t_B=-\frac{\bar c(\bar\epsilon-1)}{\bar h\bar\epsilon}.
\end{equation}
Indeed, in \eqref{eq:trace-cancelled} at $A$ the first summand has
valuation at least two, because $u\equiv1\pmod A$, and the second is
$-2h\,y/c$ modulo $A$.  At $B$ the second summand vanishes modulo $B$,
and $\bar u=\bar\epsilon$ because $\bar\epsilon$ has order three.  The
leading coefficient of $c$ at $A$ determines $\overline{y/c}$, and the
residue of $c$ at $B$ determines $\bar t_B$, since $\bar\epsilon\ne1$.
Choosing them independently makes these two nonzero residues
nonsquares.

\emph{Step 2: detecting primes.}  We select detecting primes
(Definition~\ref{def:detecting}).  The class $[y]$ is independent of
$\mathcal V(F)$, because $v_A(y)=1$ is odd.  Choose representatives
$x_1,\dots,x_r$ of a basis of $\mathcal V(F)$.  Extend the dual basis of
$\mathcal V(F)$ to $\mathcal V(F)\oplus\langle[y]\rangle$ by zero on
$[y]$: let $\varphi_i$ be the functional with $\varphi_i(x_j)=1$ if $i=j$,
$\varphi_i(x_j)=0$ if $i\ne j$, and $\varphi_i(y)=0$.  Let
$F_{\mathcal V}=F(\sqrt{x_1},\dots,\sqrt{x_r},\sqrt y)$.  Sending $g\in\Gal(F_{\mathcal V}/F)$
to the functional $[x]\mapsto g(\sqrt x)/\sqrt x\in\{\pm1\}\cong\F_2$
is injective, and both groups have order
$|\mathcal V(F)\oplus\langle[y]\rangle|$ by Kummer theory
\cite[Chapter~VII, Appendix~A, Theorem~A.3]{MilneCFT}.  So every
functional, in particular each $\varphi_i$, comes from a unique element
of $\Gal(F_{\mathcal V}/F)$.  At a prime $\mathfrak p$ of odd
residue characteristic at which all $x_j$ and $y$ are units, the
Frobenius element sends $\sqrt x$ to $(-1)^{\chi_{\mathfrak p}(x)}\sqrt
x$, by Euler's criterion.  Chebotarev's theorem therefore gives, for each
$i$, infinitely many odd primes of absolute residue degree one whose
residue characters on $x_1,\dots,x_r,y$ are $\varphi_i$.  Choose such
primes $G_i$ in distinct rational characteristics unramified in $F$ and
not in $\mathcal X$.  Also avoid the characteristics of all primes
prescribed earlier, including $G_1,\dots,G_{i-1}$, and the supports of $\epsilon\pm1$ and of $x_1,\dots,x_r,y$.  In
particular $y$ is a unit and a square modulo every $G_i$, and
Definition~\ref{def:detecting}(b) holds for the basis $[x_j]$.

At $G_i$ prescribe
\[
 c\equiv h\sqrt y\,\frac{\epsilon+1}{\epsilon-1}\pmod{G_i},
\]
where $\sqrt y$ is a square root of $y$ modulo $G_i$.  This is a root
modulo $G_i$ of $\delta_+(\epsilon)=c^2(\epsilon-1)^2-h^2y(\epsilon+1)^2$,
viewed as a polynomial in $c$, and it is a unit at $G_i$.  The
derivative $2c(\epsilon-1)^2$ with respect to $c$ is a unit, so a
congruence for $c$ modulo $G_i^2$ makes $v_{G_i}(\delta_+(\epsilon))=1$.
Moreover $\delta_-\equiv4h^2y\epsilon/(\epsilon-1)^2\ne0$ modulo $G_i$.
Once $u\equiv\epsilon\pmod{G_i^2}$ (Step~3), it follows that
$v_{G_i}(\delta)=1$, so $G_i$ ramifies in $L$
(Section~\ref{sec:nt-background}).  Add every other prime over each
selected rational characteristic to the inert set, using
Lemma~\ref{lem:scaled-holding}.  This makes each $G_i$ an isolated
ramified prime of $L/F$: its residue characteristic is unramified in $F$
and $F_{G_i}$ is the field of $\ell$-adic numbers for that characteristic
$\ell$, since $G_i$ has absolute residue degree one.  There is at least
one such prime, since $[-1]\ne0$ in $\mathcal V(F)$, so that $r\ge1$.

\emph{Step 3: the exponent modulus and the dyadic norms.}  Choose
congruence conditions on $c$ modulo fixed powers of the primes prescribed
so far that retain all the local models of Steps~1 and~2.  Let $\mathcal
M$ be divisible by $6$ and by the orders of $\epsilon$ in the finite unit
groups of $\OO_F$ modulo these prime powers.  Then
$k\equiv1\pmod{\mathcal M}$ gives $u\equiv\epsilon$ modulo each of them,
and so retains all the specified congruence classes of $u$.  In addition
to the chosen dyadic residues of $c$, prescribe
\begin{equation}\label{eq:ray-norms}
 N_{F_{\mathfrak e}/\Q_2}c\equiv1\pmod{2^{H-1}},\qquad
 N_{F_{\mathfrak d}/\Q_2}c\equiv47(1+2^{H-2})\pmod{2^{H-1}}.
\end{equation}
These conditions are compatible with the residues of $c$, because the
norms of principal units in an unramified extension of $\Q_2$ fill
$1+2\Z_2$.  To see this, note that
\[
  N(1+2^jx)\equiv1+2^j\operatorname{Tr}(x)\pmod{2^{j+1}}\qquad(j\ge1),
\]
and that each residue trace is surjective.  So for every $j\ge1$ the norms
of elements of $1+2^j\OO$ represent both classes of $1+2^j\Z_2$ modulo
$2^{j+1}$, and successive approximation gives every element of
$1+2\Z_2$ as a norm.

\emph{Step 4: realizing the congruences by a ray class.}  Choose an
integral $c_{\rm ref}$ satisfying all these congruences by the Chinese
remainder theorem \cite[Theorem~1.14]{MilneANT}.  Adding a large positive
rational integer lying in all the moduli of these congruences makes it
totally positive.  Take a narrow ray modulus $\mathfrak f$, containing all
real places, whose finite part is supported at the prescribed primes with
sufficiently large exponents.  The ideal $J=(c_{\rm ref})A^{-1}$ is prime
to this support, since $v_A(c_{\rm ref})=1$ and $c_{\rm ref}$ is a unit
at all other prescribed primes.  As explained in
Section~\ref{sec:auxiliary-prime}, a prime $R$ in the ray class of $J$
gives a generator $c=c_{\rm ref}w$ of $AR$ that is totally positive
and satisfies all these congruences.  At $A$ the ray condition concerns
the ratio $c/c_{\rm ref}$; increasing the exponent of $A$ in $\mathfrak f$
retains the leading coefficient of the nonunit $c$.  Ray class fields
exist by \cite[Chapter~V, Theorems~3.5--3.6]{MilneCFT}, and Chebotarev's
theorem gives such primes $R$ (Section~\ref{sec:nt-background}).

\emph{Step 5: the residue field of $R$ and the order of $\epsilon$.}  Let
$R$ and $c$ be as in Step~4.  Since $c$ is totally positive and $A$ has
residue field $\F_{47}$, we have $N_{F/\Q}c=47NR$.  The global norm is the
product of the local norms at $\mathfrak d$ and $\mathfrak e$, the only
primes above $2$, so \eqref{eq:ray-norms} implies
\[
 q=NR\equiv1+2^{H-2}\pmod{2^{H-1}}.
\]
For a prime $\mathfrak p$ of odd residue characteristic, write
$\chi_{\mathfrak p}$ for the residue character of
Section~\ref{sec:nt-background}, and write Hilbert symbols additively, as
$[x,b]_v\in\F_2$ (Section~\ref{sec:notation}).  By
Section~\ref{sec:nt-background}, $[x,b]_{\mathfrak p}=v_{\mathfrak p}(b)\chi_{\mathfrak p}(x)$ when
$x$ is a unit at $\mathfrak p$, and the symbol is symmetric
\cite[Chapter~III, Theorem~4.4(b)]{MilneCFT}.  Hilbert reciprocity for
$(c,\epsilon_0^{23})$, that is, the product formula \cite[Chapter~V,
Example~5.4]{MilneCFT}, gives
\[
 \chi_R(\epsilon_0^{23})=\chi_A(\epsilon_0^{23})+[c,-1]_{\mathfrak e}=1
 \quad\text{in }\F_2.
\]
We check the contributions.  The real symbols vanish because $c$ is
totally positive.  The symbol at $\mathfrak d$ vanishes because
$\epsilon_0^{23}$ is a square in $F_{\mathfrak d}$.  At $\mathfrak e$ the
element $\epsilon_0^{23}$ has square class $-1$, so the symbol there is
$[c,-1]_{\mathfrak e}$, and we show that it vanishes.  Put
$E=F_{\mathfrak e}$, an unramified extension of $\Q_2$; then
$\sqrt{-1}\notin E$, and $E(\sqrt{-1})/E$ and $\Q_2(\sqrt{-1})/\Q_2$ are
ramified quadratic extensions.  By \cite[Chapter~III,
Theorem~4.4(b),(d)]{MilneCFT}, $[c,-1]_{\mathfrak e}=0$ exactly when $c$
is a norm from $E(\sqrt{-1})$.  By \cite[Chapter~III,
Theorem~3.4(b)]{MilneCFT}, the norm groups of the two extensions have
index two in $E^\times$ and in $\Q_2^\times$.  Both contain
$2=(1+\sqrt{-1})(1-\sqrt{-1})$, a uniformizer, and in a ramified
quadratic extension an element whose norm is a unit is itself a unit.  So
the norms of units form subgroups $\mathcal N_E\subset\OO_E^\times$ and
$\mathcal N_{\Q_2}\subset\Z_2^\times$ of index two.  Transitivity of
norms gives $N_{E/\Q_2}(\mathcal N_E)\subseteq\mathcal N_{\Q_2}$.  Since
$N_{E/\Q_2}\colon\OO_E^\times\to\Z_2^\times$ is surjective
\cite[Chapter~III, Proposition~1.2]{MilneCFT}, the preimage of
$\mathcal N_{\Q_2}$ has index two in $\OO_E^\times$; it contains
$\mathcal N_E$, so it equals $\mathcal N_E$.  By Step~1, $c$ is a unit
of $E$.  The norm $N_{E/\Q_2}c$ is congruent to $1$ modulo $2^{H-1}$ by
\eqref{eq:ray-norms}, hence modulo $8$, so it is a square in $\Q_2$
\cite[Theorem~7.32]{MilneANT}, and in particular a norm from
$\Q_2(\sqrt{-1})$.  Hence the unit $N_{E/\Q_2}c$ lies in
$\mathcal N_{\Q_2}$, so $c\in\mathcal N_E$, and $c$ is a norm from
$E(\sqrt{-1})$.  Thus $[c,-1]_{\mathfrak e}=0$.  At odd primes only $A$ and $R$ can
contribute, since $\epsilon_0^{23}$ is a unit and $(c)=AR$; there the
symbols are $\chi_A(\epsilon_0^{23})$ and $\chi_R(\epsilon_0^{23})$.
Finally $\chi_A(\epsilon_0^{23})=1$, because $\epsilon_0^{23}$ has
nonsquare residue at $A_0$ and $A$ lies above $A_0$ with the same residue
field $\F_{47}$.  So $\epsilon_0^{23}$ is a nonsquare in the cyclic group
$(\OO_F/R)^\times$ of order $q-1$, and its order has two-adic valuation
$v_2(q-1)=H-2$.  Since $\epsilon=(\epsilon_0^{23})^{2^{H-3}}$,
\[
 v_2(\ord_R\epsilon_0^{23})=H-2,\qquad v_2(\ord_R\epsilon)=1.
\]

\emph{Step 6: choice of $R$ and of $c$.}  Apply
Lemma~\ref{lem:order-avoidance} to the ray class field of $\mathfrak f$,
with $\gamma$ the image of the class of $J$ under the Artin map,
and to the odd prime divisors of $\mathcal M$, excluding the primes above
$\mathcal X$ and above the characteristics of the $G_i$.  By Step~5,
infinitely many eligible primes $R$ satisfy
\begin{equation}\label{eq:exponent-compatibility}
 \ord_R\epsilon=2s,\qquad s\text{ odd},\quad (s,\mathcal M)=1.
\end{equation}
Before $c$ is fixed, we arrange total positivity of the two discriminant
factors.  The totally positive units congruent to $1$ modulo $\mathfrak
f$ have logarithms forming a full lattice in the hyperplane of vectors
with coordinate sum zero.  Indeed, they form a subgroup of finite index
in the unit group, since this subgroup contains the squares of the units
congruent to $1$ modulo $\mathfrak f$, and the logarithms of all units
form a full lattice there \cite[Theorem~5.9]{MilneANT}.  Multiplying the
generator $c$ of $AR$ from Step~4 by one of these units keeps $(c)=AR$,
total positivity and all congruences, and makes
\[
 e^{-\Omega}(47NR)^{1/[F:\Q]}\le\tau(c)
 \le e^\Omega(47NR)^{1/[F:\Q]}
\]
at every real embedding $\tau$, with $\Omega$ depending only on the fixed
modulus.  The eligible primes have unbounded norm.  Write
$\epsilon_\tau=\tau(\epsilon)$; this is positive, since $\epsilon$ is a
square.  Take $R$ with norm large enough that
\begin{equation}\label{eq:positive-factors}
 \tau(c)^2>
 h^2\max(\tau(y),0)
 \left(\frac{\epsilon_\tau+1}{\epsilon_\tau-1}\right)^2
 \quad\text{for every }\tau.
\end{equation}
No conjugate of $\epsilon$ equals $1$.  The absolute value of the fraction
is greater than one, and it decreases when $\epsilon_\tau$ is replaced by
$\epsilon_\tau^k$ with $k\ge1$.  Hence $\delta_-=c^2-h^2y$ and every
$\delta_+(\epsilon^k)=c^2(\epsilon^k-1)^2-h^2y(\epsilon^k+1)^2$ are totally
positive.  This balancing preserves every congruence condition and the
prime $R$, and it is done before $c$ and $\delta_-$ are fixed.

\emph{Step 7: choice of $k$, integrality, and the field $L$.}  Now take
\[
 k\equiv1\pmod{\mathcal M},\qquad k\equiv s\pmod{2s}.
\]
These congruences are compatible by \eqref{eq:exponent-compatibility}:
$\gcd(\mathcal M,2s)=2$ divides $s-1$.  They define an infinite
arithmetic progression of positive integers.  Since $\epsilon^s$ has order
two in the cyclic group $(\OO_F/R)^\times$, they give $u=\epsilon^k\equiv
\epsilon^s\equiv-1\pmod R$.  Since $(c)=AR$, $v_A(y)=1$ and $R$ divides
$u+1$, the element $c$ divides $h^2y(u+1)$.  The expression
\eqref{eq:trace-cancelled} for $t$ is therefore integral at all odd
primes.  The original expression $t=(y(1-u^2)-a^2)/(au)$ proves
integrality at the dyadic primes, where $a$ is a unit.  The same holds for
$a=c(u-1)/h$, by \eqref{eq:dyadic-depth} at the dyadic primes.  The
primes $G_i$ ramify, so $\delta$ is not a square and $L/F$ is quadratic.
By \eqref{eq:positive-factors}, $\delta=\delta_-\delta_+/(chu)^2$ is
totally positive, so $L=F(\sqrt\delta)$ is totally real.

\emph{Step 8: norms and ideals.}  Direct multiplication gives
$N\alpha=N\beta=y$: the constant term of $X^2-tX+y$ is $y$, and
$N(a+u\alpha)=a^2+aut+u^2y=y$ by the first expression for $t$ in
\eqref{eq:scaled-family}.  At $A$ and $B$ the polynomial $X^2-tX+y$
reduces to $X(X-\bar t)$ with $\bar t\ne0$, so both primes split.  At
$A$, $v_A(a)\ge2$ puts $\alpha$ and $\beta=a+u\alpha$ at the same prime
of $L$ above $A$.  At $B$, $a$ is a unit, so $\alpha$ and $\beta$ vanish
at the two different primes of $L$ above $B$.  Their norm $y$ has
valuation one at $A$ and at $B$, and no other support.  This proves the
displayed ideal identities, with $A'$ and $B'$ the primes at which
$\alpha$ vanishes.  At the other primes $P$ and $Q$ above $A$ and $B$,
the residue of $\alpha$ is the other root $\bar t$, and
\eqref{eq:active-traces} shows that these residues are nonsquares.

\emph{Step 9: units and class groups.}  Let $x\in F^\times$ be a virtual
unit that is a norm from $L$, say $x=N_{L/F}w$.  Its unit representative
at each ramified $G_i$ has square residue.  Indeed, let $\mathfrak G$ be
the prime of $L$ above $G_i$.  The extension $L_{\mathfrak
G}/F_{G_i}$ is quadratic and ramified, so the residue fields agree and
the nontrivial automorphism acts trivially on the residue field.  Since
$v_{G_i}(\delta)=1$, the element $\Pi=\sqrt\delta$ is a uniformizer of
$L_{\mathfrak G}$, and $\delta=\Pi^2$ is a uniformizer at $G_i$ with
$N\Pi=-\delta$.  Writing $w=\Pi^mw_0$ with $w_0$ a unit, we get
$x=(-\delta)^mN(w_0)$, so $m=v_{G_i}(x)$, which is even because $x$ is a
virtual unit, and $x\delta^{-m}=N(w_0)$
has the residue of $w_0$ squared as its residue.  Thus all the characters
$\chi_{G_i}$ vanish on $[x]$, and by Definition~\ref{def:detecting}(b)
this virtual unit is a square in $F$.  In particular the norm of a unit
of $L$ is a unit of $F$ that is a square in $F$, hence the square of a
unit.  So the unit norm image is precisely $\OO_F^{\times2}$, since units
of $F$ supply the reverse inclusion.  If a fractional ideal $I$ of $F$
capitulates, say $I\OO_L=(x)$, then $I^2=(Nx)$.  The assertion on virtual
units writes $Nx=b^2$ with $b\in F$, and ideal factorization gives
$I=(b)$.  Thus the whole ideal class group of $F$ injects into that of
$L$.  Finally the ideals $(\alpha),(\beta)$ give $[B']=-[A']$ and
$[Q]=-[A']$, hence $[B\OO_L]=-2[A']$; since $(y)=AB$, this gives
$2[A']=[A\OO_L]$ and $[P]=[A']$.  The class $[A\OO_L]$ has order $2^n$ by
injectivity, and $n\ge1$, so $\ord[A']=2^{n+1}$.

The prime $G_1$ satisfies the last assertion of the proposition, and
Steps~1--3 give all prescribed inert and residue conditions for every
$k$ in the progression.  This completes the finite step.
\end{proof}

\subsection{Scaled steps}\label{sec:scaled-steps}
The later sections use only the following properties of the fields and
elements produced by Proposition~\ref{prop:scaled-step}.  We take them
as the definition.

\begin{definition}[Scaled step]\label{def:scaled-step}
Let $n\ge1$, and let $F=F_{n-1}$, $A=A_{n-1}$, $B=B_{n-1}$ and
$y=y_{n-1}$ satisfy the assumptions stated before
\eqref{eq:scaled-family}.  Let $r=\dim_{\F_2}\mathcal V(F)$.  A
quadratic extension $F_n/F_{n-1}$ with elements
$y_n,\beta_n\in\OO_{F_n}$ is a \emph{scaled step} if there are the
following data: an integer $H\ge5$; an inert set, forbidden
characteristics and residue prescriptions satisfying conditions
(a)--(c) of Proposition~\ref{prop:scaled-step}; an integer $k$; an
element $c\in\OO_F$; a prime $R$ of $F$; and primes $G_1,\dots,G_r$ of
$F$.  These data must have the following properties.  Form $h$,
$\epsilon$, $u$, $a$, $t$, $\delta_-$ and $\delta_+$ as in
\eqref{eq:scaled-family}.  Then $F_n=L=F(\alpha)$ for a root $\alpha$
of $X^2-tX+y$, and, with $\beta=a+u\alpha$ and the primes $A'$, $P$,
$B'$, $Q$ of $L$ named in condition (v),
\[
 y_n=\alpha,\quad \beta_n=\beta,\quad
 A_n=A',\quad P_n=P,\quad B_n=B',\quad Q_n=Q.
\]
Moreover the following conditions hold, with $\mathcal X$ the set of
Proposition~\ref{prop:scaled-step} for these data.
\begin{enumerate}
\item[(i)] The exponent $k$ is positive and $k\equiv1\pmod6$.  Hence
\eqref{eq:dyadic-depth} holds, so $u\equiv1\pmod8$ at both dyadic
primes.  Moreover $u\equiv1\pmod A$, $u-1$ is a unit at $B$, and $u$
is a square in $F$.
\item[(ii)] The element $c$ is totally positive, $(c)=AR$, and $c$ is
a unit at both dyadic primes and at $B$.  The prime $R$ has absolute
residue degree one, and its rational characteristic is neither in
$\mathcal X$ nor the characteristic of any $G_i$.  Moreover
\[
 N_{F_{\mathfrak e}/\Q_2}c\equiv1\pmod{2^{H-1}},\qquad
 N_{F_{\mathfrak d}/\Q_2}c\equiv47(1+2^{H-2})\pmod{2^{H-1}},
\]
the order of $\epsilon$ modulo $R$ is twice an odd integer, and
$u\equiv-1\pmod R$.
\item[(iii)] One has $a,t\in\OO_F$; $a$ is a unit at both dyadic
primes and at $B$, and $v_A(a)\ge2$.
\item[(iv)] The elements $\delta_-$ and $\delta_+$ are totally positive,
and $L$ is totally real.
\item[(v)] All displayed conclusions of
Proposition~\ref{prop:scaled-step} hold for $L$, $\alpha$ and $\beta$;
in them $A'$ and $P$ are the primes of $L$ above $A$, and $B'$ and $Q$
those above $B$.  Every prescribed inert and residue condition holds;
in particular both dyadic primes of $F$ are unramified and inert in
$L$.  The residues of $\alpha$ at $P$ and at $Q$ are nonsquares.
\item[(vi)] The primes $G_1,\dots,G_r$ are detecting primes for
$(F,y)$ in $L$ (Definition~\ref{def:detecting}, which applies because
$v_A(y)=1$ is odd, so that $[y]\notin\mathcal V(F)$).  Here $r\ge1$,
because $F$ is totally real and so $[-1]\ne0$ in $\mathcal V(F)$.
Each $G_i$ has absolute residue degree one and is ramified in $L$; the
$G_i$ have distinct rational characteristics, unramified in $F$ and
outside the forbidden characteristics, and every other prime of $F$
over such a characteristic is unramified in $L$.
\end{enumerate}
We then call $F_n/F_{n-1}$ a scaled step \emph{for} $H$ and for this
inert set, these forbidden characteristics and these residue
prescriptions.
\end{definition}
Conditions (v) and (vi) give \textup{(T2)}--\textup{(T5)} at level
$n$.  The congruences in (i)--(iii) are those used in
Sections~\ref{sec:dyadic} and~\ref{sec:odd} and in
Appendix~\ref{app:comparison}.

\begin{remark}\label{rem:scaled-step-outputs}
Every output of the proof of Proposition~\ref{prop:scaled-step} is a
scaled step for its input data.  Precisely, let $H$, the inert set, the
forbidden characteristics and the residue prescriptions be as in that
proposition, let $c$, $R$, the primes $G_1,\dots,G_r$ and $k$ be chosen
as in its proof, and put $F_n=L$, $y_n=\alpha$, $\beta_n=\beta$,
$A_n=A'$, $P_n=P$, $B_n=B'$ and $Q_n=Q$.  Then conditions (i)--(vi) of
Definition~\ref{def:scaled-step} hold.  For (i), Step~7 takes $k$
positive with $k\equiv1\pmod{\mathcal M}$, where $6$ divides $\mathcal
M$ by Step~3; the congruences for $u$ at $A$ and $B$ were shown after
\eqref{eq:dyadic-depth}; and $u=\epsilon^k$ is a square because
$\epsilon$ is.  For (ii), Steps~4 and~6 give a totally positive
generator $c$ of $AR$ satisfying the congruences of Steps~1--3, among
them the unit conditions at the dyadic primes and at $B$ and
\eqref{eq:ray-norms}.  In Step~6, Lemma~\ref{lem:order-avoidance} gives
$R$ of absolute residue degree one outside the primes above $\mathcal X$
and above the characteristics of the $G_i$, and
\eqref{eq:exponent-compatibility} gives the order of $\epsilon$ modulo
$R$; Step~7 gives $u\equiv-1\pmod R$.  For (iii), Step~7 gives
$a,t\in\OO_F$, and the valuations of $a$ follow, as shown after
\eqref{eq:dyadic-depth}, from $k\equiv1\pmod6$, $v_A(c)=1$ and the
unit conditions on $c$.  For (iv), see \eqref{eq:positive-factors} in
Step~6 and the end of Step~7.  For (v), Step~7 shows that $L/F$ is
quadratic, Steps~8 and~9 prove the displayed conclusions and the
nonsquare residues of $\alpha$ at $P$ and $Q$, and Steps~1--3 give
every prescribed condition, including the inertness of the dyadic primes
(Step~1).  For (vi), Step~2 chooses $r$ primes $G_i$ of absolute residue
degree one in distinct odd rational characteristics, unramified in $F$
and outside $\mathcal X$, which contains the forbidden characteristics;
$y$ is a unit and a square modulo each $G_i$, and
Definition~\ref{def:detecting}(b) holds for the basis $[x_j]$.  Each
$G_i$ ramifies in $L$ once $u\equiv\epsilon\pmod{G_i^2}$, which
$k\equiv1\pmod{\mathcal M}$ ensures, and the other primes over its
characteristic are added to the inert set in Step~2, so they are inert
in $L$.  Thus each $G_i$ is an isolated ramified prime of $L/F$, and the
$G_i$ are detecting primes for $(F,y)$ in $L$.
\end{remark}

For later use, the element $z=\alpha\beta/y$ of a scaled step has a
direct description at the dyadic primes:
\begin{equation}\label{eq:scaled-target}
 z=-u-\frac{h(u+1)}c\alpha,\qquad
 v_2(-z-1)=H,\qquad \bar z=1
 \quad\text{at both dyadic primes}.
\end{equation}
The first identity follows from $\alpha^2=t\alpha-y$ and
$a+ut=y(1-u^2)/a=-hy(u+1)/c$.  For the second, the summand
$h(u+1)\alpha/c$ has valuation $H+1$, whereas $v_2(u-1)=H$.  Also
$\delta_-/c^2\in1+2^{2H}\OO_{F_v}$ is a square at both dyadic places $v$,
and $v_A(\delta_-)=1$.

\subsection{Separating residue offsets}\label{sec:offsets}
Hypothesis (T6) asks that every maximal ideal $V\ne\EA,\EB$ of $D$ of odd
residue characteristic have a residue offset, and that distinct such $V$
over the same rational prime have distinct residue offsets.  This subsection describes
a finite procedure, built from the residue prescriptions of
Proposition~\ref{prop:scaled-step}, that achieves this for all maximal
ideals of $D$ over one rational prime.  It ends with the explicit residue offsets
over $47$ and $7$.  For a rational prime $\ell$, the \emph{fibre} over
$\ell$ at a level is the set of primes of that level above $\ell$.

\emph{Offsets along inert branches.}  Suppose that a branch $(V_n)_{n\ge
m}$ carries an inert prescription at every step after level $m$, that the
residue field of $V_m$ has size $q_m\equiv1\pmod4$, and that the
two-primary component of the residue of $y_m$ at $V_m$ has order
$2^d>1$.  The identities $N_{F_n/F_{n-1}}y_n=y_{n-1}$ imply that at level
$n\ge m$ this order is $2^{d+n-m}$.  Indeed, the residue field norm of an
inert step raises to the power $q_{n-1}+1$, whose two-adic valuation is one,
and an element of order $2^{e}$ with $e\ge1$ has a $(q_{n-1}+1)$th power of
order $2^{e-1}$.  The residue offset of the branch is therefore $d-m$.
If its component is kept trivial until level $n^*-1$, and is set to $-1$
at level $n^*$, its residue offset is $1-n^*$.

\emph{Splitting a prime with nontrivial component.}  Consider a finite
collection of odd primes other than $A,B$ that are permitted to split,
namely the primes of one fibre, and let $\mathcal F$ be the set of
residue offsets already fixed in this fibre; it is initially empty and
stays finite.  Prescribe the primes
inert until their residue fields satisfy $q\equiv1\pmod4$, $q\ge81$
and $(q-1)/(q-1)_2>4\sqrt q+12$, the residue-field hypotheses of
Lemma~\ref{lem:scaled-residues}.  After this growth phase, classify the
primes of the current level by their current two-primary components.
Throughout the procedure below, whenever a prime with trivial
component is kept inert, prescribe its new component to be $1$ using
part~(2) of that lemma.  This applies both to the initially trivial
components and to those created by splitting.
Process the primes whose component is nontrivial one at a time.  Split
the current prime, with two-primary
components $1$ and the previous component at the two primes above it.
Keep the component at the first of these primes trivial as prescribed above.  At the
other prime the order of the component is unchanged, whereas its level
increases by one.  So its \emph{provisional offset}, the value $d-m$
computed at the current level, decreases by one.  Repeated splitting
makes that offset avoid $\mathcal F$ in at most $|\mathcal F|+1$ splits.
Then prescribe this prime inert at every later step and add its offset to
$\mathcal F$.  Meanwhile the other primes with nontrivial component carry
plain inert prescriptions, which leave their provisional offsets
unchanged, since the exponent of the order of the component and the level
both increase by one.

\emph{Primes with trivial component.}  When every prime with nontrivial
component has been processed, the fibre contains finitely many primes
with trivial component: those whose component was trivial after the
growth phase, and those created with component $1$ by splitting.  They
need no splitting.  Assign them distinct, sufficiently late levels
$n^*$, with $1-n^*\notin\mathcal F$; the resulting offsets $1-n^*$ are
pairwise distinct because the levels are.
Keep each component trivial up to level $n^*-1$, impose in the step to
level $n^*$ an inert residue prescription with component $-1$, and
thereafter prescribe
the branch inert at every step.  Every residue prescription of this
procedure is imposed after the growth phase, so the residue-field
hypotheses of Lemma~\ref{lem:scaled-residues} hold at its prime, since
they persist.  Every inert residue prescription, with component $1$ or
$-1$, is imposed at a prime whose current component is trivial, that is,
at which the residue of $y$ has trivial two-primary component.  So these
prescriptions satisfy condition (c) of Proposition~\ref{prop:scaled-step},
which realizes every finite round of them simultaneously.  This
procedure terminates and gives distinct residue offsets, even for primes
ramified in an earlier step, because every step at a treated prime is
split or unramified inert.

\emph{The fibres over $47$ and $7$.}  The branches over $47$ and $7$
formed by the primes $P_m,Q_m$ ($m\ge1$) and by the conjugates of
$A_0,B_0$ need no such splitting.  At the level $m\ge1$ at which they
appear, the primes $P_m,Q_m$ have nonsquare residues of $y_m$ in
$\F_{47}$ and $\F_7$, by Definition~\ref{def:scaled-step}(v); their
two-primary components have order two.  At the first inert step the norm
exponents have two-adic valuations $v_2(48)=4$ and $v_2(8)=3$,
respectively.  The new orders are thus $2^5$ and $2^4$, and subsequent
inert steps double them.  Their residue offsets are
\[
 o_{P_m}=4-m,\qquad o_{Q_m}=3-m.
\]
At the conjugates of $A_0,B_0$ the residues of $y_0$ are $15$ modulo
$47$ and $3$ modulo $7$, both nonsquares.  Prescribing these primes inert
from the start gives residue offsets $4$ and $3$.  The residue offsets
within each of these two rational fibres are therefore pairwise distinct.

\subsection{The construction schedule}\label{sec:schedule}
It remains to choose the prescriptions at every step so that all six
hypotheses hold in the limit.  A construction schedule makes this
choice as a function of the tower built so far.

\begin{definition}\label{def:schedule}
\emph{Prescription data} for a step $F_n/F_{n-1}$ consist of an integer
$H\ge5$, an inert set, a finite set of forbidden rational
characteristics, and finitely many residue prescriptions
(Definition~\ref{def:residue-prescription}), as in
Proposition~\ref{prop:scaled-step}.  A \emph{construction schedule}
assigns prescription data to every finite tower
$F_0\subset\cdots\subset F_{n-1}$ of scaled steps.
\end{definition}

\begin{theorem}[Existence of the tower]\label{thm:construction}
There exists an infinite tower whose every step is a scaled step in
the sense of Definition~\ref{def:scaled-step}, with $H_n\ge5$ and
$H_n\to\infty$, satisfying \textup{(T1)--(T6)}.  Thus hypotheses
\textup{(T1)--(T6)} can be realized simultaneously.
\end{theorem}

\begin{proof}
We describe a construction schedule
(Definition~\ref{def:schedule}), apply
Proposition~\ref{prop:scaled-step} at every step, and then check
(T1)--(T6) for the resulting tower.

\emph{Step 1: fixed prescriptions.}  Start with $F_0$ and the data of
Section~\ref{sec:base-field}, and choose, for example, $H_n=n+5$.
Every step below is produced by Proposition~\ref{prop:scaled-step} and
is therefore a scaled step for the prescription data of the schedule
(Remark~\ref{rem:scaled-step-outputs}).  Prescribe inert at every step both dyadic
branches, the conjugates of $A_0,B_0$, and every $P_m,Q_m$ from level
$m+1$ on.  Keep every prime over $3$ unramified from the start.

\emph{Step 2: the list of characteristics.}  Begin a list with $3,5,13$;
the characteristics $5$ and $13$ are ramified in $F_0$.  Subsequently add odd rational characteristics other than $47,7$.  A
characteristic remains listed while it is being treated and after its
treatment is complete.  A listed characteristic is never allowed further
ramification: all its primes carry inert prescriptions except during the
split or inert operations of the procedure of
Section~\ref{sec:offsets}.  After each constructed step, add every newly
ramified characteristic not already listed, and also the least odd
characteristic other than $47,7$ not yet listed, the characteristics added after one step
being listed in increasing order.  Treat the listed characteristics in
order, one whole finite fibre at a time.  A new entry has only finitely
many entries before it, each requiring finitely many steps.  The least
odd characteristic other than $47,7$ that is not listed increases strictly
at every step, so every such characteristic $p$ is listed after at most
$p$ steps.  Thus every
entry is eventually treated, and all branches of its fibre then carry an
inert prescription at every later step.  At any step only finitely many
fibres are listed, so Proposition~\ref{prop:scaled-step} applies.
Selected ramified primes avoid all listed characteristics: the forbidden
characteristics passed to Proposition~\ref{prop:scaled-step} at each
step are the listed ones.  The residue prescriptions passed to it are
those of the procedure of Section~\ref{sec:offsets}, imposed only after
the growth phase of their fibre; as shown there, they satisfy its
condition (c).  In each step, the proof of
Proposition~\ref{prop:scaled-step} adds all other primes of the
characteristic of a selected detecting prime to the inert set, so they
are unramified in that step, and that characteristic is listed after
the step.
A listed fibre whose treatment has not begun carries only inert
prescriptions.  During its treatment, classify the primes by their current
two-primary components after the initial residue-field growth phase.  The
treatment of a fibre is complete only after the last level $n^*$
assigned in it; from then on all its branches carry plain inert
prescriptions.  All remaining choices, namely the integers $H_n$, the
order in which the primes of a fibre are treated, and the late levels
$n^*$, are made by a fixed rule depending only on the tower built so far:
for example $H_n=n+5$, a fixed enumeration of the primes of each fibre,
and the least admissible levels $n^*$.  So the prescription data at each
step are a function of the tower built so far, as
Definition~\ref{def:schedule} requires.  They are fixed before the next
step is constructed, and every scaled step for these data satisfies them
(Definition~\ref{def:scaled-step}(v)); in particular every output of
Proposition~\ref{prop:scaled-step} does
(Remark~\ref{rem:scaled-step-outputs}).
Hence the procedure of Section~\ref{sec:offsets} has the prescribed
effect on every fibre whichever output is chosen at each step.  The
classification of a fibre after its growth phase uses the components in
the tower actually constructed, and the resulting treatment is then
carried out to completion.

\emph{Step 3: verification.}  New ramification occurs only at previously
unramified characteristics.  It has relative index two, and no later step
ramifies there.  The absolute ramification indices at odd primes are
therefore at most two.  The primes $3,7,47$ stay absolutely unramified,
and the completions at $A_n$ and $B_n$ remain the prescribed rational
local fields.  This proves (T1).  Conditions (v) and (vi) of
Definition~\ref{def:scaled-step} give (T2)--(T4) by induction,
including the order $\ord[A_n]=2^{n+1}$ needed at the next step; for
(T3), $[A_0]$ has order two by Lemma~\ref{lem:base-field}(b).  Conditions (v)
and (vi) also give (T5) at each step.  Conditions (iv) and (v) also give the assumptions
stated before \eqref{eq:scaled-family} at level $n$, so the next step is
defined.  For (T6), let $\ell\ne47,7$ be an odd characteristic, and let
$n_\ell$ be a level at which the treatment of its fibre is complete,
that is, a level at least every level $n^*$ assigned in the fibre and
after which only plain inert prescriptions are imposed there.
The fibre over $\ell$ at level $n_\ell$ consists of all primes of
$F_{n_\ell}$ above $\ell$.  By Section~\ref{sec:offsets}, each of them has
been assigned a residue offset, computed as in the paragraph on offsets
along inert branches, and these offsets are pairwise distinct.  From level
$n_\ell$ on, every prime above $\ell$ carries an inert prescription at
every step, so every prime of the fibre at level $n_\ell$ has exactly one
prime above it at each later level.  Hence every maximal ideal $V$ of $D$
above $\ell$ is determined by its contraction $V_{n_\ell}$, and $V$ has
the residue offset assigned to $V_{n_\ell}$.  Distinct such $V$ have
distinct contractions to $F_{n_\ell}$, so their offsets are distinct.
This proves (T6) over $\ell$; the explicit residue offsets of the primes
$P_m,Q_m$ and of the conjugates of $A_0,B_0$ prove it over $47$ and $7$.
\end{proof}

The construction permits primes other than the prescribed ones to ramify
in $F_n/F_{n-1}$ at every step, and places no bound on their number.
Condition (b) of Definition~\ref{def:detecting} does not give the
Selmer condition (S), in which all ramified primes enter, and the local
prescriptions above do not give the memberships of unit images in
$\Z_p$-modules required by (O).  Sections~\ref{sec:dyadic} and~\ref{sec:odd} study these
two conditions, which are the remaining hypotheses of
Theorem~\ref{thm:B}.

\section{Subrings: fraction fields, residue fields and completed orders}\label{sec:orders}

Throughout this section the hypotheses (T1)--(T6) of
Section~\ref{sec:factorization} hold, and $S$ denotes a subring of $D$.
Section~\ref{sec:joint} proves that $S$ is atomic when the tower satisfies
an approximation condition.  This section supplies the structural
information about $S$ that the proof needs.

The argument of Section~\ref{sec:joint} runs as follows.  Let $x\in S$
have unbounded factorization lengths.  By
Proposition~\ref{prop:onesided}, $x$ has positive valuation at both
$\EA$ and $\EB$, and in $D$ it splits into two one-sided factors
(Lemma~\ref{lem:split}; compare Step~1 of the proof of
Theorem~\ref{thm:ambient}).  These
factors are determined by their divisors only up to units of $D$, and in
general they do not lie in $S$.  To move them into $S$ we must know which
elements of $D$ belong to $S$.  Every element of $D$ lies in some field
$F_N$, and, when $S$ is not a BFD, $S\cap F_N$ is an order of $F_N$
(Section~\ref{sec:orders-completed}).  An order is determined by
its completions at the rational primes, and it differs from the maximal
order at only finitely many of them.  Membership in $S$ is therefore a
finite list of local conditions.  This section shows that, when $S$ is not
a BFD, the conditions at the odd primes are under control.  There are
none at the odd primes that are unramified in the tower.  At each odd
prime that ramifies, there is one condition at each maximal ideal above
it of ramification index two, measured by a single integer.

In factorization terms the conclusion is that a subring which is not a BFD
is forced to be large locally.  Precisely, let $S\subseteq D$ be a subring
that is not a BFD.  Then:
\begin{enumerate}
\item[(i)] the fraction field of $S$ is $K$
(Proposition~\ref{prop:subfields});
\item[(ii)] $S$ maps onto the residue field of every maximal ideal of $D$
of odd residue characteristic, and distinct maximal ideals of $D$ over
the same odd prime have distinct contractions to $S$
(Proposition~\ref{prop:residues});
\item[(iii)] at an odd prime $p$ that is unramified in every $F_n$, the
completion of $S\cap F_n$ at $p$ is the maximal order
(Proposition~\ref{prop:unramified});
\item[(iv)] at an odd prime $p$ that ramifies somewhere in the tower, for
$n\ge b_p$ with $b_p$ as in Section~\ref{sec:orders-ramified}, the
completion of $S\cap F_n$ at $p$ is determined by one integer at each
maximal ideal of $D$ above $p$ of ramification index two, its conductor
exponent, and this integer does not depend on $n\ge b_p$
(Proposition~\ref{prop:orders}).
\end{enumerate}
Statement (iv) is what Section~\ref{sec:joint} needs.  There the two
factors of $x$ are moved into $S\cap F_N$ by one unit of $D$ whose images
at the dyadic places and at the odd places of ramification index two lie
in a suitable open subgroup.  At an odd prime $p$ this subgroup is read
off at one level $n\ge b_p$, where the completed order contains the
coefficient rings, and is then transported to the level $N$.

The proof of (ii) uses the Picard group and the residue offsets of (T6).
A long factorization of $x$ in $S$ contains factors supported only on late
split-off primes $P_j$ and $Q_j$ (Lemma~\ref{lem:extraction}).  Because
their divisors are principal, the classes in
\eqref{eq:classes} force their exponents to balance
(Lemma~\ref{lem:balance}).  Such a factor can then be written in terms of
the elements $y_n$ and $z_n$ of (T2), and its residue at a maximal ideal
$V$ has two-primary order $2^{o_V+r}$ for an integer $r$ that does not
depend on $V$ and becomes large as the support moves to later primes
(Lemma~\ref{lem:profile}).  Residues of unbounded two-power order force
$S$ to map onto the residue field, and the distinct offsets $o_V$ of (T6)
force $S$ to separate the maximal ideals.  Statements (iii) and (iv) then
follow from (ii) by lifting residues to Teichm\"uller representatives.

\subsection{Subfields of \texorpdfstring{$K$}{K} and subrings with a smaller fraction field}
\label{sec:orders-subfields}

We first dispose of subrings whose fraction field is smaller than $K$.
Hypothesis (T5) gives each step an isolated ramified prime
(Definition~\ref{def:isolated}).  Its role is to control the subfields of
the tower.  The point is that an automorphism fixing a subfield $F'$ of
a step $L/F$ with $F\nsubseteq F'$ moves the isolated ramified prime $G$
to an unramified prime.  The proof
makes this precise with decomposition groups.  For
a Galois extension of number fields and a prime of the larger field, the
decomposition group of the prime is the Galois group of the corresponding
extension of completions \cite[Propositions~8.10 and~8.13(a)]{MilneANT}.

\begin{lemma}\label{lem:witness}
Let $L/F$ be a quadratic extension of number fields with an isolated
ramified prime $G$ over a rational prime $\ell$
(Definition~\ref{def:isolated}).  For every subfield $F'\subseteq L$,
either $F'\subseteq F$ or $F\subseteq F'$.
\end{lemma}
\begin{proof}
We show that if $F\nsubseteq F'$, then the prime of $F'$ below the
prime of $L$ over $G$ is unramified over $\ell$, and that this is
incompatible with $F'\nsubseteq F$.

Take a Galois closure $L^{\rm gal}$ of $L$ over $\Q$ and a prime
$\mathfrak G$ of $L^{\rm gal}$ over $G$.  Set
$\Gamma_F=\Gal(L^{\rm gal}/F)$ and $\Gamma_{F'}=\Gal(L^{\rm gal}/F')$, and
let $\Gamma_{\mathfrak G}\subseteq\Gal(L^{\rm gal}/\Q)$ be the
decomposition group of $\mathfrak G$.  The group $\Gamma_{\mathfrak G}$ is
the Galois group of the completion of $L^{\rm gal}$ at $\mathfrak G$ over
$\Q_\ell$, and the completion of $F$ at $G$ is the closure of $F$ in that
completion \cite[Propositions~8.10 and~8.13(a)]{MilneANT}.  Since
$F_G=\Q_\ell$, every element of $\Gamma_{\mathfrak G}$ fixes $F$, that
is, $\Gamma_{\mathfrak G}\subseteq\Gamma_F$.

Suppose $F\nsubseteq F'$, and choose $\gamma\in\Gamma_{F'}\setminus\Gamma_F$.
The prime $\gamma\mathfrak G$ lies over a prime of $F$ different from $G$.
Otherwise, since $\Gamma_F$ acts transitively on the primes of
$L^{\rm gal}$ over $G$ \cite[proof of Theorem~3.34]{MilneANT}, we would
have $\gamma\mathfrak G=\tau\mathfrak G$ with $\tau\in\Gamma_F$, hence
$\gamma\in\Gamma_F\Gamma_{\mathfrak G}=\Gamma_F$.  Since $\gamma$ fixes
$F'$, the primes $\gamma\mathfrak G$ and $\mathfrak G$ lie over the same
prime $w$ of $F'$.  The prime of $F$ below $\gamma\mathfrak G$ is a prime
over $\ell$ other than $G$, so it is unramified in $L$ because $G$ is
isolated.  As $\ell$ is unramified in $F$, the prime of $L$ below
$\gamma\mathfrak G$ is unramified over $\ell$.  The prime $w$ lies below
it, so $e(w/\ell)=1$.

Now suppose also that $F'\nsubseteq F$.  Since $L/F$ is quadratic,
$F'F=L$.  Inside the completion of $L^{\rm gal}$ at $\mathfrak G$, the
completion of $L$ at the prime below $\mathfrak G$ is therefore the
compositum $F'_wF_G=F'_w$, which is unramified over $\Q_\ell$.  But $G$
ramifies in $L$ and $F_G=\Q_\ell$, so this completion is totally ramified
of degree two over $\Q_\ell$.  This contradiction shows that
$F'\subseteq F$.
\end{proof}

\begin{proposition}\label{prop:subfields}
Every proper subfield of $K$ is a number field.  If $S\subseteq D$ and
$\Frac S\ne K$, then $S$ is a BFD.
\end{proposition}
\begin{proof}
\emph{Step 1: the subfields of $K$.}  Let $F'\subseteq K$ be a number
field.  It is finitely generated, so it lies in some $F_n$; choose the
least such $n$.  If $n>0$, Lemma~\ref{lem:witness} applies to
$F_n/F_{n-1}$ by (T5), and minimality of $n$ excludes
$F'\subseteq F_{n-1}$.  Hence $F_{n-1}\subseteq F'\subseteq F_n$, and
$F'=F_n$.  Thus every subfield of $K$ of finite degree is either
contained in $F_0$ or equal to one of the $F_n$.  A subfield of infinite
degree is the union of its subfields of finite degree, and these have
unbounded degree.  So it contains $F_n$ for infinitely many $n$, hence for
all $n$, and it equals $K$.

\emph{Step 2: $S$ is an order.}  Let $K'=\Frac S\ne K$.  By Step~1, $K'$ is
a number field.  Every element of $S$ is an algebraic integer, so
$S\subseteq\OO_{K'}$.  Let $b\in S\setminus\{0\}$, with monic minimal
polynomial $X^d+c_{d-1}X^{d-1}+\cdots+c_0\in\Z[X]$.  Then $c_0\ne0$ and
\[
 b^{-1}=-c_0^{-1}\bigl(b^{d-1}+c_{d-1}b^{d-2}+\cdots+c_1\bigr).
\]
So $b^{-1}$ is a polynomial in $b$ with integer coefficients divided by
the nonzero integer $c_0$.  Therefore $S\otimes_\Z\Q=K'$.  As a subgroup
of the finitely generated free abelian group $\OO_{K'}$, the ring $S$ is
free of rank $[K':\Q]$, so it is a full lattice in $K'$ and hence an order
in $K'$.

\emph{Step 3: bounded factorization lengths.}  Let $b\in S$ be a unit of
$\OO_{K'}$.  The constant term $c_0$ of its minimal polynomial is
$\pm N_{\Q(b)/\Q}(b)$, and the norm of a unit is $\pm1$.  So the display
above gives $b^{-1}\in\Z[b]\subseteq S$, and $b$ is a unit of $S$.
Consequently every nonunit of $S$ is a nonunit of $\OO_{K'}$, and its norm
is an integer of absolute value at least two.  If $x\ne0$ is a product
$x=x_1\cdots x_r$ of nonunits of $S$, then
$2^r\leq|N_{K'/\Q}(x)|$.  So the lengths $r$ are bounded, and $S$ is a BFD
(Section~\ref{sec:factorization-background}).
\end{proof}

\subsection{Orders and completed orders}\label{sec:orders-completed}

From now on we assume that $S\subseteq D$ is not a BFD.  By
Proposition~\ref{prop:subfields}, $\Frac S=K$.  Write
\[
 S_n=S\cap F_n .
\]
The argument of Step~2 of the proof of Proposition~\ref{prop:subfields}
shows that $S\otimes_\Z\Q=\Frac S=K$.  Hence each element of an integral
basis of $F_n$ has a nonzero integer multiple in $S$, and one nonzero
integer multiplies the whole basis into $S$.  So $S_n$ has full rank in
$\OO_{F_n}$: it is an order in $F_n$.

An order of a number field is determined by local data at the rational
primes.  For each rational prime $p$ put
\[
 \mathcal A_{n,p}=S_n\otimes_\Z\Z_p,\qquad
 \mathcal B_{n,p}=\OO_{F_n}\otimes_\Z\Z_p.
\]
These are an order and the maximal order, respectively, in the
$\Q_p$-algebra $F_n\otimes_\Q\Q_p$, which is the product of the
completions $(F_n)_v$ over the places $v$ of $F_n$ above $p$.  We call
$\mathcal A_{n,p}$, the $p$-adic completion of the order $S_n$, the
\emph{completed order} of $S_n$ at $p$.  The next lemma records the two
facts about completed orders that we use: they are compatible with the
inclusions $F_n\subseteq F_m$, and they determine $S_n$.

\begin{lemma}[Intersection at a fixed level]\label{lem:intersection}
Let $S\subseteq D$ be a subring that is not a BFD, and let $p$ be a
rational prime.  For $m\geq n$,
\begin{equation}\label{eq:saturation}
 \mathcal A_{m,p}\cap(F_n\otimes_\Q\Q_p)=\mathcal A_{n,p}.
\end{equation}
Furthermore
\begin{equation}\label{eq:lattice}
 S_n=F_n\cap\bigcap_p\mathcal A_{n,p}.
\end{equation}
\end{lemma}
\begin{proof}
The subgroup $S_n=S_m\cap F_n$ is saturated in the free abelian group
$S_m$: if a nonzero integer multiple of an element of $S_m$ belongs to
$F_n$, so does the element.  Hence a basis of $S_n$ extends to a basis of
$S_m$.  Tensoring with $\Z_p$ gives \eqref{eq:saturation}.  In coordinates
with respect to a basis of $S_n$, which is also a $\Q$-basis of $F_n$,
the equality \eqref{eq:lattice} becomes $\Q\cap\bigcap_p\Z_p=\Z$.
\end{proof}

Equation~\eqref{eq:lattice} reduces membership in $S_n$ to membership in
the completed orders.  Since $S_n$ has finite index in $\OO_{F_n}$, we
have $\mathcal A_{n,p}=\mathcal B_{n,p}$ for all but finitely many $p$.
The rest of this section determines $\mathcal A_{n,p}$ at the odd primes
$p$.

\subsection{Elements supported on late split-off primes}
\label{sec:orders-late}

A subring that is not a BFD contains elements with unbounded
factorization lengths.  Since one-sided elements are bounded
(Proposition~\ref{prop:onesided}), such an element has positive valuation
at both $\EA$ and $\EB$.  We now extract from its long factorizations
elements of a very special shape: their divisors involve only the primes
$P_j$ and $Q_j$ split off from the two chains at late levels $j$, with
bounded exponents.  In this subsection $(f)=fD$ denotes the principal
ideal of $D$ generated by $f$.

\begin{definition}\label{def:late-supported}
Let $\kappa\geq1$ and $j_0\geq1$ be integers.  A nonzero element $f\in D$
is a \emph{late-supported element at level $j_0$ with exponent bound
$\kappa$} if there is an integer $j_1\geq j_0$ such that
\begin{equation}\label{eq:profile}
 (f)=\prod_{j=j_0}^{j_1} P_j^{a_j}Q_j^{b_j},
 \qquad 0\leq a_j,b_j\leq \kappa,
\end{equation}
and at least one exponent $a_j$ or $b_j$ is positive.
\end{definition}

In our applications the exponent bound $\kappa$ is fixed while the level
$j_0$ increases.

\begin{lemma}\label{lem:extraction}
Let $S\subseteq D$ be a subring that is not a BFD.  There is an integer
$\kappa\geq1$ such that, for every sufficiently large $j_0$, the ring $S$
contains a late-supported element at level $j_0$ with exponent bound
$\kappa$.  More precisely, if $x\in S$ has unbounded factorization
lengths in $S$, then for every $j_0\geq1$ such that \eqref{eq:tails}
holds for $x$ and all $j\geq j_0$, such an element can be chosen as a
factor of $x$ in $S$
whose exponents satisfy $a_j\leq v_{\EA}(x)$ and $b_j\leq v_{\EB}(x)$; in
particular $\kappa=\max\{v_{\EA}(x),v_{\EB}(x)\}$ works.
\end{lemma}
\begin{proof}
A long factorization of $x$ has more factors than $x$ has prime factors
away from the late split-off primes, and the surplus factors are the
required elements.

Since $S$ is not a BFD, some $x\in S$ has unbounded factorization lengths
in $S$ (Section~\ref{sec:factorization-background}), so the first
assertion follows from the second.  Let $x\in S$ have unbounded
factorization lengths.  By Proposition~\ref{prop:onesided}, $x$ is not
one-sided, so $v_{\EA}(x)$ and $v_{\EB}(x)$ are positive; in particular
$\kappa=\max\{v_{\EA}(x),v_{\EB}(x)\}\geq1$.  Let $j_0\geq1$ be such that \eqref{eq:tails} holds
for $x$ and all $j\geq j_0$; every sufficiently large $j_0$ qualifies.  By
Lemma~\ref{lem:geometry}, the set $\mathcal H$ of maximal ideals of $D$
other than $\EA$, $\EB$ and the $P_j,Q_j$ with $j\geq j_0$ at which $x$
has positive valuation is finite.  It includes those $P_j,Q_j$ with
$j<j_0$ at which $x$ has positive valuation.

Consider a factorization $x=x_1\cdots x_r$ in $S$.  Each $x_i$ is a
nonunit of $D$ by Lemma~\ref{lem:units}, and $v_V(x)=\sum_iv_V(x_i)$ at
every maximal ideal $V$.  Hence at most
$\sum_{V\in\mathcal H\cup\{\EA,\EB\}}v_V(x)$ of the factors have positive
valuation at a member of $\mathcal H\cup\{\EA,\EB\}$.  Choose a
factorization longer than this sum.  Then some factor $f=x_i$ has
valuation zero at every member of $\mathcal H\cup\{\EA,\EB\}$.  Since $f$
divides $x$ in $D$, its support is contained in the set of primes
$P_j,Q_j$ with $j\geq j_0$.  This support is finite: its valuations at
$\EA$ and $\EB$ are zero, so by \eqref{eq:tails} for $f$ it has valuation
zero at $P_j$ and $Q_j$ for all large $j$.  It is nonempty because $f$ is
not a unit of $D$, and an element of $D$ whose valuations are all zero is
a unit (Lemma~\ref{lem:geometry}).  Finally, since $f$ divides
$x$ in $D$, \eqref{eq:tails} for $x$ gives
$a_j=v_{P_j}(f)\leq v_{P_j}(x)=v_{\EA}(x)$ and
$b_j=v_{Q_j}(f)\leq v_{Q_j}(x)=v_{\EB}(x)$ for $j\geq j_0$.
\end{proof}

The divisor of a late-supported element is principal, so its class in
$\Pic(D)$ vanishes.  Lemma~\ref{lem:classes} computes this class, and
for late support the vanishing forces the exponents $a_j$ and $b_j$ to
balance exactly.

\begin{lemma}\label{lem:balance}
Let $f$ be a late-supported element at level $j_0$ with exponent bound
$\kappa$, with divisor \eqref{eq:profile}, and suppose that
$\kappa2^{-j_0}<1$.  Then the following equality holds in $\Q$, and not
only modulo $\Z$:
\begin{equation}\label{eq:balance}
 \Sigma:=\sum_{j=j_0}^{j_1} a_j2^{-j}
   =\sum_{j=j_0}^{j_1} b_j2^{-j}>0.
\end{equation}
Moreover $-v_2(\Sigma)\geq j_0-1-\log_2\kappa$.
\end{lemma}
\begin{proof}
By \eqref{eq:classes}, the class of $(f)$ in $\Pic(D)$ is
$\iota(d+\Z)$ with $d=\sum_{j=j_0}^{j_1}(a_j-b_j)2^{-j-1}$.  This class is
zero and $\iota$ is injective, so $d$ is an integer.  But
$|d|\leq\kappa\sum_{j=j_0}^{j_1}2^{-j-1}<\kappa2^{-j_0}<1$, so $d=0$.
Since $d$ is half the difference of the two sums in \eqref{eq:balance},
they are equal.  They are positive because all exponents are nonnegative
and one of them is positive.

For the last assertion, write $\Sigma=2^{v_2(\Sigma)}w$ with $w$ a
positive rational number whose numerator and denominator are odd.  Since
the denominator of $\Sigma$ is a power of two, $w$ is an odd positive
integer, and $2^{v_2(\Sigma)}\leq\Sigma$.  On the other hand
$\Sigma\leq\kappa\sum_{j\geq j_0}2^{-j}=\kappa2^{1-j_0}$.  Hence
$v_2(\Sigma)\leq1-j_0+\log_2\kappa$.
\end{proof}

In particular, for a fixed exponent bound $\kappa$, $-v_2(\Sigma)$ tends
to infinity with $j_0$.  This is the quantity that will make residues of
late-supported elements large.

\subsection{Residues at odd maximal ideals}\label{sec:orders-residues}

We now compute the residues of late-supported elements at a maximal ideal
of odd residue characteristic.  Such an element can be written, up to a
unit of $D$, as a product of the elements $y_{j-1}/y_j$ and
$z_{j-1}/z_j$ of (T2).  Hypothesis (T6) prescribes the residues of the
$y_n$.  The residues of the units of $D$ are not known, but
Lemma~\ref{lem:unitresidues} shows that their two-primary components
form a finite group, and,
for large $n$, the two-primary components of the residues of the $z_n$
have order at most two.  The result is a formula for
the two-primary order of the residue of a late-supported element.  This
argument is what gives both the surjectivity onto residue fields and the
separation of maximal ideals in Proposition~\ref{prop:residues}; neither is
assumed about $S$.

\emph{Setting.}  Throughout this subsection, $V$ is a maximal ideal of $D$
of odd residue characteristic $p$, different from $\EA$ and $\EB$.  By
(T1), there is a level from which the prime $V_n$ is unramified and inert
in $F_{n+1}$ for every $n$.  Every residue degree in the tower is a power
of two, because $F_0$ has degree two over $\Q$ and every step has degree
two, and these degrees are unbounded along $(V_n)$ because $V_n$ is inert
in every late step.  Hence the residue field
$k(V)=D/V$ is the union of the finite fields $k(V_n)$ and equals
$\bigcup_r\F_{p^{2^r}}$.  Let $k(V)^\times_{(2)}$ be the two-primary
subgroup of $k(V)^\times$.  It is the increasing union of the cyclic
groups $k(V_n)^\times_{(2)}$, whose orders are unbounded by (T6).  So
$k(V)^\times_{(2)}$ is a quasicyclic two-group: each finite subgroup is
cyclic, and for each $a\geq0$ there is exactly one subgroup of order
$2^a$, consisting of the elements of order at most $2^a$.

Let $\overline E_V\subseteq k(V)^\times_{(2)}$ be the image of $D^\times$
under reduction at $V$ followed by projection to the two-primary
component.  It contains $-1$, the residue of the unit $-1$, which is the
unique element of order two of $k(V)^\times_{(2)}$.  For an element $g$
of $K$ that is a unit at $V$, write $\langle g\rangle_V$ for the class of the
two-primary component of its residue in the quotient
$k(V)^\times_{(2)}/\overline E_V$.  We write this quotient additively, so
that $\langle gg'\rangle_V=\langle g\rangle_V+\langle g'\rangle_V$.

Two standard facts about one inert step are used repeatedly.  Suppose
that $V_n$ is unramified and inert in $F_{n+1}$, and let
$q=|k(V_n)|$.  The local extension at $V$ is unramified of degree two,
its Galois group is generated by $\sigma_{n+1}$, and this generator acts
on the residue field $k(V_{n+1})=\F_{q^2}$ as the Frobenius $x\mapsto x^q$
\cite[Propositions~7.50(b) and~8.10]{MilneANT}.  Hence, for $g\in F_{n+1}$
a unit at $V$, the residue of $N_{F_{n+1}/F_n}(g)=g\,\sigma_{n+1}(g)$ is
$\bar g^{\,q+1}$, where $\bar g$ is the residue of $g$.  If moreover
$q\equiv1\pmod4$, then $q+1$ has two-adic valuation one.

\begin{lemma}\label{lem:unitresidues}
Let $V$ be a maximal ideal of $D$ of odd residue characteristic,
different from $\EA$ and $\EB$.  The group $\overline E_V$ is finite.
Write $|\overline E_V|=2^{a_V}$, and let $\langle y_n\rangle_V$ be the class of the
two-primary component of the residue of $y_n$ in
$k(V)^\times_{(2)}/\overline E_V$, defined when $y_n$ is a unit at $V$.
Then, for $n$ sufficiently large,
\[
 2\langle y_{n+1}\rangle_V=\langle y_n\rangle_V,\qquad \ord(\langle y_n\rangle_V)=2^{n+o_V-a_V}.
\]
\end{lemma}
\begin{proof}
Both assertions come from comparing residues across one inert step, where
the relative norm becomes the map $x\mapsto x^{q+1}$.  In this proof a bar
denotes the two-primary component of a residue at $V$.

\emph{Step 1: the starting level.}  By (T1), the prime $V_n$ is
unramified and inert in $F_{n+1}$ for every $n$ from some level on.  One
inert step later the residue cardinality is a square of an odd number,
hence $1$ modulo $4$, and it remains so.  We fix a level $n_0$ from which
both properties hold, and consider $n\geq n_0$ with $q=|k(V_n)|$.

\emph{Step 2: $\overline E_V$ is finite.}  Let $\eta\in E_{n+1}$.
Hypothesis (T4) gives $N_{F_{n+1}/F_n}\eta=\xi^2$ with $\xi\in E_n$.  On
two-primary residues this reads $\bar\eta^{\,q+1}=\bar\xi^{\,2}$.  Since
$q+1$ has two-adic valuation one, $(q+1)/2$ is odd, and
$(\bar\eta^{\,(q+1)/2})^2=\bar\xi^{\,2}$.  In the cyclic group
$k(V_{n+1})^\times_{(2)}$ two elements with the same square differ by an
element of order at most two, so $\bar\eta^{\,(q+1)/2}=\pm\bar\xi$.  As
$\bar\eta$ has two-power order and $(q+1)/2$ is odd, $\bar\eta$ is a power
of $\bar\eta^{\,(q+1)/2}$.  So $\bar\eta$ lies in the subgroup generated by
the residues of the units of $E_n$ and $-1$, which is the image of $E_n$
because $-1\in E_n$.  Thus the image of $E_{n+1}$ is contained in the
image of $E_n$ for every $n\geq n_0$.  Since $D^\times=\bigcup_nE_n$, the
group $\overline E_V$ equals the image of $E_{n_0}$, a subgroup of the
finite group $k(V_{n_0})^\times$.

\emph{Step 3: the relation $2\langle y_{n+1}\rangle_V=\langle y_n\rangle_V$.}  Take $n\geq n_0$ so
large that $y_n$ and $y_{n+1}$ are units at $V$.  The norm identity
$N_{F_{n+1}/F_n}y_{n+1}=y_n$ of (T2) gives
$\bar y_{n+1}^{\,q+1}=\bar y_n$.  The element $\bar y_{n+1}^{\,q-1}$
satisfies $(\bar y_{n+1}^{\,q-1})^{q+1}=\bar y_{n+1}^{\,q^2-1}=1$, so its
order divides $q+1$.  Being a two-primary element, it has order at most
two, so it lies in $\{\pm1\}\subseteq\overline E_V$.  Hence
$\bar y_{n+1}^{\,q+1}=\bar y_{n+1}^{\,2}\,\bar y_{n+1}^{\,q-1}$ has the
same class as $\bar y_{n+1}^{\,2}$, and the norm identity becomes
$2\langle y_{n+1}\rangle_V=\langle y_n\rangle_V$.

\emph{Step 4: the orders.}  By (T6), $\bar y_n$ has order $2^{n+o_V}$ for
large $n$.  The subgroup $\overline E_V$ is the unique subgroup of order
$2^{a_V}$ of the quasicyclic group $k(V)^\times_{(2)}$, that is, the set
of elements of order at most $2^{a_V}$.  For an element $g$ of order
$2^e$, the multiple $2^kg$ lies in this subgroup exactly when
$e-k\leq a_V$.  So the class of $g$ has order $2^{\max\{e-a_V,0\}}$.
With $e=n+o_V$, which is at least $a_V$ for large $n$, this gives
$\ord(\langle y_n\rangle_V)=2^{n+o_V-a_V}$.
\end{proof}

The next lemma computes the residue of a late-supported element.  The
level $j_0$ at which its support begins must be large compared with
several quantities attached to $V$ and $\kappa$; we explain the threshold
after the statement.

\begin{lemma}\label{lem:profile}
Fix an integer $\kappa\geq1$ and a finite set of maximal ideals $V$ of $D$
of odd residue characteristic, all different from $\EA$ and $\EB$.  For
each such $V$, let $n_V$ be a level such that, for every $n\geq n_V$,
the prime $V_n$ is unramified and inert in $F_{n+1}$ with residue
cardinality congruent to $1$ modulo $4$, and the conclusions of
Lemma~\ref{lem:unitresidues} hold for $n$.  Let $j_V=j$ if $V$ is $P_j$
or $Q_j$, and $j_V=0$ otherwise.  Put
\[
 j_*(V,\kappa)=\max\{n_V+1,\ j_V+1,\
   a_V-o_V+2+\lceil\log_2\kappa\rceil,\ \lfloor\log_2\kappa\rfloor+1\},
\]
which depends only on $V$ and $\kappa$.  Let $f$ be a late-supported
element at level $j_0$ with exponent bound $\kappa$, with divisor
\eqref{eq:profile}, and suppose $j_0\geq j_*(V,\kappa)$ for all these
$V$.  Then \eqref{eq:balance} holds, $f$ is a unit at each of these $V$,
and the two-primary component of the residue of $f$ at $V$ has order
\begin{equation}\label{eq:resorder}
 2^{o_V-v_2(\Sigma)}.
\end{equation}
In particular this order does not change when $f$ is multiplied by a
unit of $D$.
\end{lemma}

The four terms in $j_*(V,\kappa)$ have the following roles.  The term
$n_V+1$ ensures that every step $F_{j-1}\subset F_j$ with $j\geq j_0$ is
one to which Lemma~\ref{lem:unitresidues} applies.  The term $j_V+1$
ensures that $V$ is not one of the primes $P_j,Q_j$ in the support of
$f$, so that $f$ and the auxiliary elements below are units at $V$.  The
term $\lfloor\log_2\kappa\rfloor+1$ gives $\kappa<2^{j_0}$, so that
Lemma~\ref{lem:balance} applies.  The term
$a_V-o_V+2+\lceil\log_2\kappa\rceil$ makes the class of $f$ modulo
$\overline E_V$ nonzero.  Residues of units of $D$ are known only modulo
$\overline E_V$, and a nonzero class modulo $\overline E_V$ determines the
order of the residue itself.

\begin{proof}
We write $f$ as a unit times a product of the elements $y_{j-1}/y_j$ and
$z_{j-1}/z_j$, compute the class of each factor in
$k(V)^\times_{(2)}/\overline E_V$, and then lift the result from the
quotient to $k(V)^\times_{(2)}$.

\emph{Step 1: a formula for $f$.}  Since
$j_0\geq\lfloor\log_2\kappa\rfloor+1>\log_2\kappa$, we have
$\kappa2^{-j_0}<1$, and Lemma~\ref{lem:balance} gives \eqref{eq:balance}.
By (T1) and (T2), the elements $y_{j-1}/y_j$ and $z_j^{-1}$ have ideals
$P_jQ_j$ and $P_jA_j^{-1}$, and for $j>j_0$ the element $z_{j-1}/z_j$ has
ideal $P_j^2P_{j-1}^{-1}$.  For $j_0<j\leq j_1+1$ put
\[
 d_j=-\sum_{i=j_0}^{j-1}(a_i-b_i)2^{j-1-i},
\]
and put $d_{j_0}=0$.  Then $d_{j+1}=2d_j+b_j-a_j$ for
$j_0\leq j\leq j_1$, and $d_{j_1+1}=-2^{j_1}\sum_{i}(a_i-b_i)2^{-i}=0$ by
\eqref{eq:balance}.  The element
\[
 \prod_{j=j_0}^{j_1} \Bigl(\frac{y_{j-1}}{y_j}\Bigr)^{b_j}
       \prod_{j=j_0+1}^{j_1} \Bigl(\frac{z_{j-1}}{z_j}\Bigr)^{d_j}
\]
has valuation $b_j$ at $Q_j$ and $b_j+2d_j-d_{j+1}=a_j$ at $P_j$ for
$j_0\leq j\leq j_1$, and valuation zero elsewhere.  It has the same
divisor as $f$, so by Lemma~\ref{lem:geometry}
\[
 f=\eta\prod_{j=j_0}^{j_1} \Bigl(\frac{y_{j-1}}{y_j}\Bigr)^{b_j}
       \prod_{j=j_0+1}^{j_1} \Bigl(\frac{z_{j-1}}{z_j}\Bigr)^{d_j},
 \qquad \eta\in D^\times.
\]

\emph{Step 2: the factors are units at $V$.}  Let $j\geq j_0$.  As
$D$-ideals, $A_jD$ has valuation one exactly at $\EA$ and at the $P_i$
with $i>j$.  Since $j\geq j_0>j_V$ and $V\neq\EA,\EB$, the ideal $V$ is
none of $\EA$, $P_i$, $Q_i$ with $i\geq j_0$.  Hence $y_{j-1}/y_j$ and
$z_j$ are units at $V$ for $j\geq j_0$, so is $z_{j-1}/z_j$ for $j>j_0$,
and $V$ is not in the support of $f$.

\emph{Step 3: the classes of the factors.}  Let $j\geq j_0$.  Since
$j-1\geq n_V$, the prime $V_{j-1}$ is inert in $F_j$ with residue
cardinality $q\equiv1\pmod4$, and $\langle y_{j-1}\rangle_V=2\langle y_j\rangle_V$.  The element
$z_j$ has relative norm one by (T2), so its residue $\bar z$ satisfies
$\bar z^{\,q+1}=1$.  The kernel of the norm map $\F_{q^2}^\times\to
\F_q^\times$ has order $q+1$ and $v_2(q+1)=1$, so the two-primary
component of $\bar z$ has order at most two.  It therefore lies in
$\{\pm1\}\subseteq\overline E_V$, and $\langle z_j\rangle_V=0$.  Consequently
$\langle z_{j-1}/z_j\rangle_V=0$ for $j>j_0$.  The class of $y_{j-1}/y_j$ is
$\langle y_{j-1}\rangle_V-\langle y_j\rangle_V=\langle y_j\rangle_V$, and $\langle \eta\rangle_V=0$ since
$\eta\in D^\times$.

\emph{Step 4: the class of $f$.}  Iterating
$\langle y_{j-1}\rangle_V=2\langle y_j\rangle_V$ gives $\langle y_j\rangle_V=2^{j_1-j}\langle y_{j_1}\rangle_V$ for
$j_0\leq j\leq j_1$.  By Steps 1 and 3 and \eqref{eq:balance},
\[
 \langle f\rangle_V=\sum_{j=j_0}^{j_1}b_j\langle y_j\rangle_V=2^{j_1}\Sigma\,\langle y_{j_1}\rangle_V .
\]
Here $2^{j_1}\Sigma$ is an integer of two-adic valuation
$j_1+v_2(\Sigma)$, and $\langle y_{j_1}\rangle_V$ has order $2^{j_1+o_V-a_V}$ by
Lemma~\ref{lem:unitresidues}, since $j_1\geq n_V$.  Lemma~\ref{lem:balance}
and $j_0\geq a_V-o_V+2+\lceil\log_2\kappa\rceil$ give
\[
 o_V-a_V-v_2(\Sigma)\geq o_V-a_V+j_0-1-\log_2\kappa\geq1 .
\]
Hence $j_1+v_2(\Sigma)<j_1+o_V-a_V$, and $\langle f\rangle_V$ has order
$2^{r}$ with $r=o_V-a_V-v_2(\Sigma)\geq1$.

\emph{Step 5: lifting to the residue.}  In the quasicyclic two-group
$k(V)^\times_{(2)}$, let $g$ be an element whose class modulo the
subgroup $\overline E_V$ of order $2^{a_V}$ has order $2^r>1$.  Then $g$
has order exactly $2^{r+a_V}$: by Step~4 of the proof of
Lemma~\ref{lem:unitresidues}, the class of an element of order $2^N$ has
order $2^{\max\{N-a_V,0\}}$.  Applied to the two-primary component of the
residue of $f$, this gives the order $2^{o_V-v_2(\Sigma)}$ of
\eqref{eq:resorder}.  Multiplying $f$ by a unit of $D$ changes neither its
divisor nor, therefore, this computation.
\end{proof}

\begin{proposition}\label{prop:residues}
Let $S\subseteq D$ be a subring that is not a BFD.  For every maximal
ideal $V$ of $D$ of odd residue characteristic, the image of $S$ in $D/V$
is the whole residue field.  Distinct maximal ideals of $D$ over the same
odd rational prime have distinct contractions to $S$.
\end{proposition}
\begin{proof}
The late-supported elements of $S$ supplied by
Lemma~\ref{lem:extraction} have residues of unbounded two-power order by
Lemma~\ref{lem:profile}; this gives surjectivity, and the offsets of (T6)
give separation.

\emph{Step 1: the image is a field.}  Let $p$ be the residue
characteristic of $V$.  The image of $S$ in $D/V$ is a subring of an
algebraic extension of $\F_p$.  For a nonzero element $r$ of this image,
$\F_p[r]$ is a finite domain, hence a field, so $r^{-1}$ lies in the
image.  Thus the image is a field.

\emph{Step 2: surjectivity.}  At $\EA$ or $\EB$ the residue field is the
prime field $\F_{47}$ or $\F_7$, and the image of $\Z\subseteq S$ already
fills it.  Let $V$ be any other maximal ideal of odd residue
characteristic, so that $k(V)=\bigcup_r\F_{p^{2^r}}$.  The subfields of
$\F_{p^{2^r}}$ are the fields $\F_{p^{2^{r'}}}$ with $r'\leq r$.  Hence a
subfield of $k(V)$ either contains $\F_{p^{2^r}}$ for every $r$, and then
equals $k(V)$, or it is one of the finite fields $\F_{p^{2^{r'}}}$.
By Lemma~\ref{lem:extraction} there is a $\kappa$ such that $S$ contains
late-supported elements $f$ with exponent bound $\kappa$ at arbitrarily
large levels $j_0$.  By Lemma~\ref{lem:profile}, applied to the set
$\{V\}$, and by Lemma~\ref{lem:balance}, the two-primary component of the
residue of such an $f$ at $V$ has order
$2^{o_V-v_2(\Sigma)}\geq2^{o_V+j_0-1-\log_2\kappa}$ once
$j_0\geq j_*(V,\kappa)$.  These orders are unbounded, whereas the elements
of a finite field have bounded orders.  So the image of $S$ is not
finite, and it is the whole residue field.

\emph{Step 3: separation.}  Let $V\neq V'$ be maximal ideals of $D$ over
the same odd rational prime, and suppose that $V\cap S=V'\cap S$.  By
Steps 1 and~2, the residue maps induce isomorphisms from the field
$S/(V\cap S)$ onto $D/V$ and onto $D/V'$.  Composing them gives an
isomorphism $D/V\to D/V'$ that sends the residue of each $f\in S$ at $V$
to its residue at $V'$.  The ideals $\EA$ and $\EB$ have different
residue characteristics, so at most one of $V,V'$ is $\EA$ or $\EB$.  If
one of them is, its residue field is finite while the other residue field
is infinite, and no such isomorphism exists.  Suppose that neither is
$\EA$ or $\EB$.  The isomorphism preserves the orders of the two-primary
components of residues.  Take $f\in S$ late-supported with exponent
bound $\kappa$ and $j_0\geq\max\{j_*(V,\kappa),j_*(V',\kappa)\}$, as in
Step~2.  Lemma~\ref{lem:profile}, applied to the set $\{V,V'\}$, gives
orders $2^{o_V-v_2(\Sigma)}$ and $2^{o_{V'}-v_2(\Sigma)}$ for the same
$\Sigma$.  These are different because $o_V\neq o_{V'}$ by (T6).  This
contradiction shows $V\cap S\neq V'\cap S$.
\end{proof}

\subsection{Teichm\"uller limits and unramified characteristics}
\label{sec:orders-unramified}

Proposition~\ref{prop:residues} is a statement about residues.  To turn
it into a statement about the completed orders $\mathcal A_{n,p}$ we lift
residues canonically.  The Teichm\"uller representatives of
Section~\ref{sec:nt-background} \cite[Remark~7.36]{MilneANT} are such
lifts, and they are limits of iterated $q$-th powers.  A completed order
is a closed subring, so it contains these limits as soon as it contains
elements with the right residues.  In an unramified complete discrete
valuation ring every element is a $p$-adic series in Teichm\"uller
representatives, so a closed subring that contains all Teichm\"uller
representatives is the whole ring.

\begin{lemma}[Teichm\"uller limits]\label{lem:teichmuller-limit}
Let $\OO$ be a complete discrete valuation ring with maximal ideal
$\mathfrak m$ and finite residue field of characteristic $p$ with $q_0$
elements, and let $q=q_0^k$ with $k\geq1$.  For every $s\in\OO$ the
sequence $s^{q^r}$ converges as $r\to\infty$, and its limit is the
Teichm\"uller representative of the residue of $s$.
\end{lemma}
\begin{proof}
If $a\equiv b\pmod{\mathfrak m^j}$ with $j\geq1$, then
$a^p\equiv b^p\pmod{\mathfrak m^{j+1}}$.  Indeed, write $a=b+c$ with
$c\in\mathfrak m^j$; in the binomial expansion of $(b+c)^p$ the terms
other than $b^p$ and $c^p$ have coefficients divisible by $p\in\mathfrak
m$, and $c^p\in\mathfrak m^{pj}\subseteq\mathfrak m^{j+1}$.  Since $q$ is
a power of $p$, it follows that $a^q\equiv b^q\pmod{\mathfrak m^{j+1}}$.
Every element $x$ of the residue field satisfies $x^q=x$, so
$s^q\equiv s\pmod{\mathfrak m}$.  By induction on $r$,
$s^{q^{r+1}}\equiv s^{q^r}\pmod{\mathfrak m^{r+1}}$.  So the sequence is
Cauchy and converges to some $\lambda\in\OO$ with
$\lambda\equiv s\pmod{\mathfrak m}$.  The Teichm\"uller representative
$\tau$ of the residue of $s$ satisfies $\tau^{q_0}=\tau$, hence
$\tau^{q^r}=\tau$ for all $r$, and $\tau\equiv s\pmod{\mathfrak m}$.  By
the congruence above, $s^{q^r}\equiv\tau^{q^r}=\tau\pmod{\mathfrak
m^{r+1}}$ for every $r$, so $\lambda=\tau$.
\end{proof}

When $\OO$ is a finite product of complete discrete valuation rings, as
$\mathcal B_{n,p}$ is, we apply the lemma in each factor.  The
\emph{Teichm\"uller representative} of a tuple of residues is then the
tuple of the Teichm\"uller representatives of its entries.

\begin{proposition}\label{prop:unramified}
Let $S\subseteq D$ be a subring that is not a BFD.  If an odd rational
prime $p$ is unramified in every $F_n$, then
$\mathcal A_{n,p}=\mathcal B_{n,p}$ for every $n$.  In particular this
holds at $47$ and $7$.
\end{proposition}
\begin{proof}
We first show that $S$ maps onto $D/pD$, using
Proposition~\ref{prop:residues} and the idempotents of $D/pD$; then we
lift with Lemma~\ref{lem:teichmuller-limit}.

\emph{Step 1: the structure of $D/pD$.}  Since $p\OO_{F_m}\cap
\OO_{F_n}=p\OO_{F_n}$ for $m\geq n$, the ring $D/pD$ is the directed
union of the rings $\OO_{F_n}/p\OO_{F_n}$.  These are finite reduced
$\F_p$-algebras, because $p$ is unramified in $F_n$.  Hence every finitely
generated subalgebra of $D/pD$ is a finite product of finite fields.  The
same is true of the image $\bar S$ of $S$ in $D/pD$.  Such rings are von
Neumann regular.  Each element $r$ has a \emph{support idempotent}
$e_r$, equal to $1$ in the factors where $r$ is nonzero and to $0$
elsewhere, and $e_r$ is a power of $r$: if $\F_p[r]\cong\prod_i\F_{q_i}$,
take $e_r=r^{M}$ with $M$ a common multiple of the $q_i-1$.  In
particular $e_r\in\bar S$ when $r\in\bar S$.  Since $D$ has dimension one, the
points of $X=\Spec(D/pD)$ are the maximal ideals of $D$ above $p$.  For each $r$,
the set of points where $r$ does not vanish equals the set where $e_r=1$,
and it is open and closed.  The space $X$ is compact, and these sets form
a basis of its topology.  An element of $D/pD$ is zero if it vanishes at
every point, because $D/pD$ is reduced.

\emph{Step 2: $\bar S=D/pD$.}  By Proposition~\ref{prop:residues}, $\bar S$
maps onto the residue field at every point of $X$, and distinct points of
$X$ have distinct contractions to $S$.  Let $x\neq y$ be points of $X$.
Some $s\in\bar S$ vanishes at exactly one of them, and then $e_s$ or
$1-e_s$ is an idempotent of $\bar S$ equal to $1$ at $x$ and to $0$ at
$y$.

Let $Y\subseteq X$ be open and closed.  For $x\in Y$ and each $y\notin
Y$, choose an idempotent of $\bar S$ equal to $1$ at $x$ and $0$ at $y$.
The sets where these idempotents vanish cover the compact set
$X\setminus Y$, so finitely many suffice.  Their product is an
idempotent of $\bar S$ equal to $1$ at $x$ and to $0$ on $X\setminus Y$.
The sets where these idempotents equal $1$, for $x\in Y$, cover the
compact set $Y$.  If $e_1,\dots,e_k$ are finitely many of them that
cover $Y$, then $1-\prod_i(1-e_i)\in\bar S$ is the characteristic
idempotent of $Y$.

Now let $r\in D/pD$.  For each $x\in X$, choose $s_x\in\bar S$ with the
same residue as $r$ at $x$.  The set where $s_x$ and $r$ agree is the set
where $s_x-r$ vanishes, so it is open and closed, and it contains $x$.
Finitely many of these sets $U_1,\dots,U_k$, for points
$x_1,\dots,x_k$, cover $X$.  The sets
$Y_i=U_i\setminus(U_1\cup\cdots\cup U_{i-1})$ are open, closed, disjoint
and cover $X$, so their characteristic idempotents $e_{Y_i}\in\bar S$ are
orthogonal and sum to $1$.  The element $\sum_ie_{Y_i}s_{x_i}\in\bar S$
agrees with $r$ at every point of $X$, so it equals $r$.  This argument
applies also to the infinite fibres at $47$ and $7$.

\emph{Step 3: Teichm\"uller representatives in $\mathcal A_{m,p}$.}  Fix
$n$.  The ring $\mathcal B_{n,p}$ is the product of the rings of integers
$\OO_{(F_n)_w}$ over the places $w$ of $F_n$ above $p$; each is unramified
over $\Z_p$, and the residue ring of $\mathcal B_{n,p}$ is
$\OO_{F_n}/p\OO_{F_n}$.  By Step~2, for each element $\bar t$ of the
finite ring $\OO_{F_n}/p\OO_{F_n}$ we may choose $t\in\OO_{F_n}$ with
residue $\bar t$ and $s\in S$ with $s\equiv t\pmod{pD}$.  These finitely
many representatives $s$ all lie in $S_m$ for some $m\geq n$, and then
$s\equiv t\pmod{p\OO_{F_m}}$.  So, in $\mathcal B_{m,p}$, the residues of
$s$ agree with those of the corresponding elements $t$ of the embedded
ring $\mathcal B_{n,p}$.  Take $q=p^f$ with $f\geq1$ divisible by
every residue degree of $F_m$ at $p$; then $q$ is also a power of the
residue cardinality of each factor of $\mathcal B_{n,p}$.  By
Lemma~\ref{lem:teichmuller-limit} in each factor, $s^{q^r}$ and $t^{q^r}$
have the same limit, and this limit is the Teichm\"uller representative
of $\bar t$ in $\mathcal B_{n,p}$.  Since $\mathcal A_{m,p}$ is a finitely
generated $\Z_p$-module, it is closed in $F_m\otimes_\Q\Q_p$.  It contains
the powers $s^{q^r}$, so it contains the Teichm\"uller representatives of
all residues of $\mathcal B_{n,p}$.

\emph{Step 4: $p$-adic expansions.}  Since $p$ is a uniformizer in each
factor of $\mathcal B_{n,p}$, every $b\in\mathcal B_{n,p}$ is a convergent
sum $\sum_{k\geq0}\omega_kp^k$ with each $\omega_k$ a Teichm\"uller
representative: take $\omega_0$ to be the Teichm\"uller representative of
the residue of $b$, and repeat with $(b-\omega_0)/p$.  The partial sums
lie in $\mathcal A_{m,p}$, which is closed, so
$\mathcal B_{n,p}\subseteq\mathcal A_{m,p}$.  Intersecting with
$F_n\otimes_\Q\Q_p$ and using \eqref{eq:saturation} gives
$\mathcal B_{n,p}\subseteq\mathcal A_{n,p}$.  The reverse inclusion holds
because $S_n\subseteq\OO_{F_n}$.  Finally, $47$ and $7$ are unramified in
every $F_n$ by (T1).
\end{proof}

\subsection{Ramified characteristics: coefficient rings and conductor exponents}
\label{sec:orders-ramified}

At a prime $p$ that ramifies in the tower the maximal order is not
generated by Teichm\"uller representatives, and the completed orders can
be smaller than the maximal order.  The Teichm\"uller representatives
generate the coefficient ring, and at a maximal ideal of ramification
index two the local ring of integers is a free module of rank two over it.  So the only freedom left for the completed order is in
the second coordinate, and we show that it is measured by one integer.

\emph{Setting.}  Let $p$ be an odd rational prime that ramifies in some
$F_n$.  By (T1), $p$ is different from $47$, $7$ and $3$, and its
absolute ramification indices are at most two.  Also by (T1), we may fix
a level $b_p$ after which all primes over $p$ are unramified and inert in
every subsequent step.  For each $n\geq b_p$, the maximal ideals $V$ of
$D$ above $p$ then correspond bijectively to the primes of $F_n$ above
$p$, via $V\mapsto V_n$, and
\[
 F_n\otimes_\Q\Q_p=\prod_{V\mid p}(F_n)_V .
\]
The idempotents of the factors $(F_n)_V$ are the primitive idempotents of
$F_n\otimes_\Q\Q_p$.  The ramification index $e(V/p)$ of $(F_n)_V$ over
$\Q_p$ is $1$ or $2$, and it does not depend on $n\geq b_p$.

For a maximal ideal $V$ of $D$ above $p$ and $n\geq b_p$, write
$W_{n,V}$ for the ring of integers of the maximal unramified subfield of
$(F_n)_V$, the \emph{coefficient ring} of $(F_n)_V$
(Section~\ref{sec:nt-background}; see \cite[Proposition~7.50 and
Corollary~7.51]{MilneANT}).  It contains the Teichm\"uller
representatives, it is unramified over $\Z_p$, and it has the same
residue field as $(F_n)_V$.  If $e(V/p)=1$, then
$W_{n,V}=\OO_{(F_n)_V}$.

Suppose that $e(V/p)=2$.  Then $(F_n)_V$ is a ramified quadratic
extension of the unramified field $\Frac(W_{n,V})$.  For such an
extension, with $p$ odd and coefficient ring $W$, recall that a square
uniformizer is a uniformizer $\pi$ with $\pi^2\in pW^\times$
(Section~\ref{sec:nt-background}).  To
obtain one, complete the square in the Eisenstein minimal polynomial,
over $\Frac(W_{b_p,V})$, of a uniformizer of $(F_{b_p})_V$; since $p$ is
odd, this gives a uniformizer $\pi_V$ with
$\pi_V^2\in pW_{b_p,V}^{\times}$.  For $n\geq b_p$ the extension
$(F_n)_V/(F_{b_p})_V$ is unramified, so $\pi_V$ is still a uniformizer
and $\pi_V^2\in pW_{n,V}^\times$: the same $\pi_V$ is a square
uniformizer of $(F_n)_V$ for every $n\geq b_p$.  It generates the local
ring of integers over the coefficient ring
\cite[Proposition~7.55 and Remark~7.56]{MilneANT}, so
\[
 \OO_{(F_n)_V}=W_{n,V}\oplus\pi_VW_{n,V}.
\]

For an order $\Lambda$ contained in $\OO_{(F_n)_V}$, the
\emph{conductor} of $\Lambda$ is the largest ideal of $\OO_{(F_n)_V}$
contained in $\Lambda$, namely
$\{x\in\OO_{(F_n)_V}:x\OO_{(F_n)_V}\subseteq\Lambda\}$.

\begin{proposition}\label{prop:orders}
Let $S\subseteq D$ be a subring that is not a BFD, and let $p$ and $b_p$
be as above.  For $n\geq b_p$, $\mathcal A_{n,p}$ contains
$\prod_{V\mid p}W_{n,V}$, including the primitive idempotents of
$F_n\otimes_\Q\Q_p$.  Hence $\mathcal A_{n,p}$ is the product of its
factors in the $(F_n)_V$.  Its factor at a maximal ideal $V$ above $p$
with $e(V/p)=1$ is the full local ring of integers $\OO_{(F_n)_V}$.  At
$V$ with $e(V/p)=2$ it is
\begin{equation}\label{eq:localorder}
 \Lambda_{n,V}=W_{n,V}+\pi_Vp^{s_V}W_{n,V}
\end{equation}
for a nonnegative integer $s_V$ independent of $n\geq b_p$;
equivalently, the conductor
$\{x\in\OO_{(F_n)_V}:x\OO_{(F_n)_V}\subseteq\Lambda_{n,V}\}$ equals
$p^{s_V}\OO_{(F_n)_V}$.  We call $s_V$ the \emph{conductor exponent} of $S$
at $V$.
\end{proposition}
\begin{proof}
We lift residue tuples to Teichm\"uller representatives, as in the proof
of Proposition~\ref{prop:unramified}, to show that the completed order
contains the coefficient rings; what remains is an ideal of the
coefficient ring in the second coordinate.

\emph{Step 1: the coefficient rings.}  By Proposition~\ref{prop:residues},
the finitely many contractions $V\cap S$ of the maximal ideals $V$ above
$p$ are distinct maximal ideals of $S$, and each residue map
$S/(V\cap S)\to D/V$ is onto.  The Chinese remainder theorem in $S$
therefore supplies, for every tuple $\rho=(\rho_V)_V$ of residues
$\rho_V\in k(V_n)$, an element $s_\rho\in S$ with residue $\rho_V$ at
each $V$.  Place these finitely many elements $s_\rho$ in some $S_m$ with
$m\geq n$, and take $q=p^f$ with $f\geq1$ divisible by every residue
degree of $F_m$ at $p$.  Lemma~\ref{lem:teichmuller-limit}, applied in
each factor of $\mathcal B_{m,p}$, ramified or not, shows that
$s_\rho^{q^r}$ converges to the tuple of Teichm\"uller representatives of
the $\rho_V$.  The Teichm\"uller representative of $\rho_V$ in
$\OO_{(F_n)_V}$ is a root of $X^q-X$ with residue $\rho_V$, so by the last
step of the proof of Lemma~\ref{lem:teichmuller-limit} it is this limit.
Hence these representatives lie in the coefficient rings $W_{n,V}$
(Section~\ref{sec:nt-background}).  As $\mathcal A_{m,p}$ is closed, the
limit lies in $\mathcal A_{m,p}$.  Every element of $\prod_VW_{n,V}$ is a convergent
sum $\sum_k\omega_kp^k$ of such tuples $\omega_k$ (as in Step~4 of the
proof of Proposition~\ref{prop:unramified}), so
$\prod_VW_{n,V}\subseteq\mathcal A_{m,p}$.  This product contains the
primitive idempotents: the idempotent of $(F_n)_V$ is the Teichm\"uller
representative of the tuple with entry $1$ at $V$ and $0$ elsewhere.
Equation~\eqref{eq:saturation} places $\prod_VW_{n,V}$ already in
$\mathcal A_{n,p}$.  This finite-level intersection is essential:
Proposition~\ref{prop:residues} produces elements of $S$ in some $S_m$,
possibly with $m>n$, and surjectivity on residue fields alone does not
provide lifts at level $n$.

\emph{Step 2: the shape of each factor.}  The primitive idempotents lie
in $\mathcal A_{n,p}$, so $\mathcal A_{n,p}$ is the product of its
factors $\Lambda_{n,V}\subseteq\OO_{(F_n)_V}$, and each $\Lambda_{n,V}$
contains $W_{n,V}$.  If $e(V/p)=1$, then $W_{n,V}=\OO_{(F_n)_V}$ and the
factor is the full local ring of integers.  Let $e(V/p)=2$.  Every
element of $\Lambda_{n,V}$ has the form $w_0+\pi_Vw_1$ with
$w_0,w_1\in W_{n,V}$, and subtracting $w_0\in\Lambda_{n,V}$ shows that
$\pi_Vw_1\in\Lambda_{n,V}$.  Hence $\Lambda_{n,V}=W_{n,V}+\pi_VJ_n$,
where $J_n=\{w\in W_{n,V}:\pi_Vw\in\Lambda_{n,V}\}$ is an ideal of
$W_{n,V}$, because $\Lambda_{n,V}$ is a $W_{n,V}$-module.  The ideal $J_n$
is nonzero because $\Lambda_{n,V}$ is a full lattice in $(F_n)_V$.  The
ring $W_{n,V}$ is a discrete valuation ring with uniformizer $p$, so
$J_n=p^{s_{n,V}}W_{n,V}$ for some integer $s_{n,V}\geq0$.

\emph{Step 3: independence of $n$.}  Let $m\geq n\geq b_p$.  The primes
of $F_n$ and of $F_m$ above $p$ correspond, so $F_n\otimes_\Q\Q_p$ and
$F_m\otimes_\Q\Q_p$ have the same primitive idempotents, and
\eqref{eq:saturation} gives
$\Lambda_{m,V}\cap(F_n)_V=\Lambda_{n,V}$ for each $V$.  The extension
$\Frac(W_{m,V})/\Frac(W_{n,V})$ is unramified, so
$p^sW_{m,V}\cap\Frac(W_{n,V})=p^sW_{n,V}$ for every $s$.  Since
$(F_m)_V=\Frac(W_{m,V})\oplus\pi_V\Frac(W_{m,V})$ with the same
$\pi_V$ as at level $n$, it follows that
$\Lambda_{m,V}\cap(F_n)_V=W_{n,V}+\pi_Vp^{s_{m,V}}W_{n,V}$.  Comparing
with $\Lambda_{n,V}$ gives $s_{m,V}=s_{n,V}$, as required.

\emph{Step 4: the conductor.}  Write $s=s_V$, $W=W_{n,V}$ and
$\pi=\pi_V$, with $\pi^2=pu_0$ for a unit $u_0\in W^\times$.  An element
$x=a+\pi b$ of $\OO_{(F_n)_V}$, with $a,b\in W$, lies in the conductor
exactly when $x$ and $x\pi=pu_0b+\pi a$ lie in $\Lambda_{n,V}$, because
$\Lambda_{n,V}$ is a $W$-module and $\OO_{(F_n)_V}=W+\pi W$; that is,
when $b\in p^sW$ and $a\in p^sW$.  So the conductor is
$p^sW\oplus\pi p^sW=p^s\OO_{(F_n)_V}$.  Conversely, the conductor
determines $s$, and $\Lambda_{n,V}=W+p^s\OO_{(F_n)_V}$.
\end{proof}

Section~\ref{sec:joint} uses the following description of the units of
the local orders $\Lambda_{n,V}$.

\begin{lemma}[Units of the local orders]\label{lem:order-units}
In the situation of Proposition~\ref{prop:orders}, let $V$ be a maximal
ideal of $D$ above $p$ with $e(V/p)=2$, and let $n\geq b_p$.  Every unit
of $\OO_{(F_n)_V}$ can be written uniquely as $x_0(1+\pi_Vb)$ with
$x_0\in W_{n,V}^{\times}$ and $b\in W_{n,V}$.  It belongs to
$\Lambda_{n,V}^{\times}$ if and only if $b\in p^{s_V}W_{n,V}$.
\end{lemma}
\begin{proof}
Write a unit as $x=a+\pi_Va'$ with $a,a'\in W_{n,V}$.  It has the same
residue as $a$, so $a\in W_{n,V}^\times$, and $x=a(1+\pi_Va^{-1}a')$.
Conversely $x_0(1+\pi_Vb)=x_0+\pi_Vx_0b$, so the coordinates of $x$
determine $x_0$ and then $b$.  If $x\in\Lambda_{n,V}$, then
$x_0b\in p^{s_V}W_{n,V}$, so $b\in p^{s_V}W_{n,V}$.  Conversely, if
$b\in p^{s_V}W_{n,V}$, then $x\in\Lambda_{n,V}$, and its inverse
\[
 x^{-1}=x_0^{-1}(1-\pi_Vb)(1-\pi_V^2b^2)^{-1}
\]
also lies in $\Lambda_{n,V}$, because $1-\pi_Vb\in\Lambda_{n,V}$ and
$1-\pi_V^2b^2\in1+pW_{n,V}$ is a unit of the coefficient ring.
\end{proof}

We have not assumed that $S_n$ is stable under any automorphism of
$F_n$, and the conductor exponents $s_V$ may differ for different maximal
ideals $V$ above the same prime $p$.

\section{Unit approximation and hereditary atomicity}\label{sec:joint}

This section proves Theorem~\ref{thm:joint}: if the elements $z_n$ of
(T2) can be approximated by global units in the sense of condition
\eqref{eq:JT} below, then $D$ is hereditarily atomic.  Since $D$ does not
satisfy ACCP (Theorem~\ref{thm:ambient}), this completes the
factorization-theoretic part of the proof of Theorem~\ref{thm:B}.
Sections~\ref{sec:dyadic} and~\ref{sec:odd} then show that, for a tower as
in Theorem~\ref{thm:B}, the conditions (S) and (O) imply \eqref{eq:JT}.

Throughout this section hypotheses (T1)--(T6) hold.  For a subring
$S\subseteq D$ with $\Frac S=K$, we use the order $S_N=S\cap F_N$, its
completion
$\mathcal A_{N,p}=S_N\otimes_\Z\Z_p$ and the maximal order
$\mathcal B_{N,p}=\OO_{F_N}\otimes_\Z\Z_p$, as in
Section~\ref{sec:orders}.  The results of Section~\ref{sec:orders} about
these completed orders are used only when $S$ is not a BFD.

\emph{The idea.}  Let $S\subseteq D$ be a subring and let $x\in S$ be a
nonzero nonunit.  If $\Frac S\ne K$, or if $x$ has bounded factorization
lengths in $S$, then $x$ is a product of atoms of $S$.  Otherwise $x$ has
positive valuation at both $\EA$ and $\EB$
(Proposition~\ref{prop:onesided}).  In $D$, such an $x$ is the product of
two one-sided factors, as in the proof of Theorem~\ref{thm:ambient}, and a
one-sided element has bounded factorization lengths in every subring that
contains it (Proposition~\ref{prop:onesided}).  Hence $x$ is a product of
atoms of $S$ as soon as both factors can be chosen in $S$.  A factor of $x$
in $D$ is determined by its divisor only up to a unit of $D$.  The proof
has two steps, and the second uses this freedom.
\begin{enumerate}
\item[(1)] \emph{Split} (Section~\ref{sec:joint-split}).
Lemma~\ref{lem:split} gives a factor $x_A$ of $x$ in $D$ with
$v_{\EA}(x_A)=v_{\EA}(x)$ and $v_{\EB}(x_A)=0$, together with a factor $f$
of $x$ in $S$, such that the ratio $\rho=f/x_A$ is a product of powers of
elements $z_n$ with large $n$.
\item[(2)] \emph{Move} (Section~\ref{sec:joint-return}).  We look for a
unit $\eta$ of $D$ such that $f'=\eta^{-1}x_A$ and $x/f'$ both lie in
$S$.  Both lie in some field $F_N$, and by \eqref{eq:lattice} an element
of $F_N$ lies in $S_N$ exactly when it lies in the completed order
$\mathcal A_{N,p}$ for every prime $p$.  Lemma~\ref{lem:return} produces
such an $\eta$ when the image of $\rho$ lies in the closure of the image of
the unit group $E_N$ in a compact group $\mathcal U_N$ made of local unit
groups.
\end{enumerate}

\emph{Why the approximation must be simultaneous.}  The unit $\eta$ in
step (2) must be one element of $E_N$ that works at all primes $p$ at once.
Suppose that for each prime $p$ we only had a unit $\eta_p$ with
$\eta_p\rho\in\mathcal A_{N,p}^\times$.  The candidate factors
$\eta_p^{-1}x_A$ would in general be different elements of $F_N$, each
satisfying the local condition at its own prime, and none of them need lie
in $S_N$.  Approximating $\rho$ by units at each place separately
therefore does not suffice for this argument: we need approximation by one
unit in the product of all the local groups.  For this reason the closure
in \eqref{eq:JT} is taken in the whole product $\mathcal U_N$.

\emph{Which places occur.}  At $47$ and $7$, the final factors are
integral and the completed orders are maximal
(Proposition~\ref{prop:unramified}), so their membership in
$\mathcal A_{N,p}$ requires no additional approximation condition.
For the local units considered in Lemma~\ref{lem:return}, there is also
no additional condition at an odd place at which $F_N$ is unramified
over $p$.  At an odd place of ramification index two, the completed order
contains the coefficient ring (Proposition~\ref{prop:orders} at a level
$M\ge\max\{N,b_p\}$, and \eqref{eq:saturation}), so only the
class of a unit modulo the units of the coefficient ring matters.  At the
dyadic places we know only that $\mathcal A_{N,2}$ contains
$1+2^e\mathcal B_{N,2}$ for some $e\ge1$, so the whole dyadic unit group
appears.  In every case membership is an open condition, so approximation
to a finite precision suffices; this is why \eqref{eq:JT} is stated with a
closure.

\subsection{The approximation group}\label{sec:joint-group}

Fix a level $N$.  For a place $v$ of $F_N$ above an odd prime $p$ with
$e(v/p)=2$, write $\OO_{N,v}=\OO_{(F_N)_v}$ for the ring of integers
of the completion and $W_{N,v}$ for its coefficient ring
(Section~\ref{sec:nt-background}).  Define the compact abelian group
\begin{equation}\label{eq:TN}
 \mathcal U_N=(\OO_{F_N}\otimes\Z_2)^\times\times
 \prod_{\substack{v\text{ odd}\\e(v/p)=2}}
       \OO_{N,v}^\times/W_{N,v}^\times.
\end{equation}
There are finitely many factors, since only finitely many primes of $F_N$
are ramified over $\Q$.  The first factor is the product of the unit
groups $\OO_{(F_N)_v}^\times$ over the dyadic places $v$.  Each factor is
compact, because $W_{N,v}^\times$ is a closed subgroup of the compact group
$\OO_{N,v}^\times$.

Let $z\in F_N^\times$ be a unit at every dyadic prime and at every odd
prime of absolute ramification index two.  We write $\jmath_N(z)$ for its
image in $\mathcal U_N$: its image in $(\OO_{F_N}\otimes\Z_2)^\times$ in
the first factor, and its class in $\OO_{N,v}^\times/W_{N,v}^\times$ in
the factor of $v$.  On the group of such elements, $\jmath_N$ is a
homomorphism.  It is defined on all global units $E_N$, and on $z_n$ and
$z_n'$ for $n\le N$.  Indeed, the divisors $(z_n)=A_nP_n^{-1}$ and
$(z_n')=B_nQ_n^{-1}$ lie above $47$ and $7$, and by (T1) the primes $47$
and $7$ are unramified in every $F_N$.  We impose the following
\emph{unit approximation condition} on the elements $z_n$; the subscript
$A$ in its name refers to the chain $(A_n)$, since $(z_n)=A_nP_n^{-1}$:
\begin{equation}\label{eq:JT}
 \tag{$\mathrm{JT}_A$}
 \exists n_0\ \forall n>n_0\ \exists N\geq n:
 \quad\jmath_N(z_n)\in\overline{\jmath_N(E_N)}.
\end{equation}
Here $\overline{\jmath_N(E_N)}$ is the closure of $\jmath_N(E_N)$ in
$\mathcal U_N$.  Thus $\jmath_N(z_n)$ lies in it exactly when, for every
neighbourhood $\mathcal N$ of the identity in $\mathcal U_N$, there is a
unit $\eta\in E_N$ with $\jmath_N(\eta z_n)\in\mathcal N$: one global unit
approximates $z_n$ in all factors of $\mathcal U_N$ at once.

The closure in the whole product can be strictly smaller than the product
of the closures of the projections: the diagonal image of $\Z$ in
$\Z_2\times\Z_2$ has closure the diagonal, while each projection has
closure $\Z_2$.

The quantifiers in \eqref{eq:JT} are exactly those that the proof of
Theorem~\ref{thm:joint} uses.  The condition concerns only $z_n$ with
$n>n_0$, because a given element $x$ requires only finitely many $z_n$,
all with $n$ larger than a cutoff $n_1$ that may be taken as large as we
like (Lemma~\ref{lem:split}).  The level $N$ may depend on $n$, because
approximation at one level persists at every later level
(Lemma~\ref{lem:propagation} below), so finitely many $z_n$ can be moved
to a common level.  In Sections~\ref{sec:dyadic} and~\ref{sec:odd} the
condition is verified with $N=n$ (Proposition~\ref{prop:assembly}).

\begin{theorem}[Unit approximation criterion]\label{thm:joint}
Assume \textup{(T1)--(T6)} and \textup{($\mathrm{JT}_A$)}.  Then $D$ is
hereditarily atomic and does not satisfy ACCP.  The same conclusion holds
if, in \textup{($\mathrm{JT}_A$)}, the elements $z_n$ are replaced by the
elements $z_n'$.
\end{theorem}

The theorem is proved in Section~\ref{sec:joint-proof}, after the two
lemmas that carry out the steps described above.  We first compare the
groups $\mathcal U_N$ at different levels.

\begin{lemma}\label{lem:propagation}
Let $N'\geq N$.  There is a continuous homomorphism
$\theta_{N,N'}:\mathcal U_N\to\mathcal U_{N'}$ such that
$\theta_{N,N'}(\jmath_N(z))=\jmath_{N'}(z)$ for every $z\in F_N^\times$
that is a unit at every prime not lying above $47$ or $7$.  This applies
to every global unit $z\in E_N$ and to $z=z_n,z_n'$ with $n\le N$; these
$z_n,z_n'$ are the elements of relative norm one of \textup{(T2)}.  The homomorphism
$\theta_{N,N'}$ maps $\overline{\jmath_N(E_N)}$ into
$\overline{\jmath_{N'}(E_{N'})}$.
\end{lemma}
\begin{proof}
We define $\theta_{N,N'}$ factor by factor, using for each factor of
$\mathcal U_{N'}$ the factor of $\mathcal U_N$ at the prime below it.

\emph{Dyadic factor.}  The inclusion
$F_N\otimes\Q_2\subseteq F_{N'}\otimes\Q_2$ maps
$(\OO_{F_N}\otimes\Z_2)^\times$ into $(\OO_{F_{N'}}\otimes\Z_2)^\times$.

\emph{Odd factors.}  Let $w$ be an odd prime of $F_{N'}$ with $e(w/p)=2$,
and let $v$ be the prime of $F_N$ below it.  Then $e(v/p)$ divides
$e(w/p)$, so $e(v/p)$ is $1$ or $2$.
\begin{enumerate}
\item[(a)] Suppose $e(v/p)=2$, so that $v$ is already ramified.  The
completion $(F_N)_v$ embeds in $(F_{N'})_w$.  The fraction field of its
coefficient ring is unramified over $\Q_p$, so it embeds in the fraction
field of the coefficient ring of $(F_{N'})_w$, and $W_{N,v}^\times\subseteq W_{N',w}^\times$.  Hence the
inclusion $\OO_{N,v}^\times\subseteq\OO_{N',w}^\times$ induces a map of
quotients $\OO_{N,v}^\times/W_{N,v}^\times\to
\OO_{N',w}^\times/W_{N',w}^\times$.  We use this map in the factor of
every prime $w$ of $F_{N'}$ above $v$.
\item[(b)] Suppose $e(v/p)=1$, so that $w$ is newly ramified.  Then
$(F_N)_v$ is unramified over $\Q_p$, so $\OO_{N,v}=W_{N,v}\subseteq
W_{N',w}$.  Every element of $F_N$ that is a unit at $v$ therefore has
trivial class in $\OO_{N',w}^\times/W_{N',w}^\times$, and we use the
trivial homomorphism in the factor of $w$.
\end{enumerate}

These maps are continuous homomorphisms, so $\theta_{N,N'}$ is one.  Let
$z\in F_N^\times$ be a unit at every prime that does not lie above $47$
or $7$; for instance $z$ is a global unit of $F_N$, or $z_n$ or $z_n'$
with $n\le N$.  By (T1), $z$ is then a unit at every dyadic prime, at
every odd prime of absolute ramification index two, and at every prime
$v$ in case~(b).  In case~(a) both
$\theta_{N,N'}(\jmath_N(z))$ and $\jmath_{N'}(z)$ have the class of $z$ in
the factor of $w$; in case~(b) both are trivial there.  This proves the
compatibility with $\jmath$.  Since $E_N\subseteq E_{N'}$, we get
$\theta_{N,N'}(\jmath_N(E_N))\subseteq\jmath_{N'}(E_{N'})$.  A continuous
map sends the closure of a set into the closure of its image, which gives
the statement about closures.
\end{proof}

\subsection{Splitting an element into two one-sided factors}
\label{sec:joint-split}

Let $x\in S$ have unbounded factorization lengths.
Lemma~\ref{lem:extraction} provides a factor $f$ of $x$ in $S$ supported
on late primes $P_j,Q_j$.  We now construct a factor $x_A$ of $x$ in $D$
that carries the whole valuation of $x$ at $\EA$ and none at $\EB$.  We
obtain it from $f$ by multiplying by powers of the elements $z_n$, with
$(z_n)=A_nP_n^{-1}$; each of them moves valuation from a prime $P_n$
split off from the chain $(A_n)$ to $A_nD$, and hence to $\EA$.  A
weighted degree, defined in Step~3 of the proof below, is unchanged by
these moves.  Together with the prescribed valuations of $x_A$ at the
$P_j$, it determines the valuation that $x_A$ acquires at $\EA$.  Thus $f/x_A$ is exactly a
product of powers of the $z_n$.

\begin{lemma}\label{lem:split}
Let $S\subseteq D$ be a subring and let $x\in S$ have unbounded
factorization lengths in $S$, with $\kappa_A=v_{\EA}(x)>0$ and
$\kappa_B=v_{\EB}(x)>0$.  For every sufficiently large integer $n_1$ there
are an element $f\in S$ with $x/f\in S$ and an element $x_A\in D$ with
$x_A\mid_Dx$, such that
\[
 \bigl(v_{\EA}(x_A),v_{\EB}(x_A)\bigr)=(\kappa_A,0),\qquad
 \bigl(v_{\EA}(x/x_A),v_{\EB}(x/x_A)\bigr)=(0,\kappa_B),
\]
and, for some integer $n_2$,
\begin{equation}\label{eq:rho}
 \rho:=f/x_A=\prod_{n_1<n\leq n_2}z_n^{m_n},
 \qquad m_n\in\Z.
\end{equation}
Both $x_A$ and $\rho$ are units at every maximal ideal of $D$ that does
not lie above $47$ or $7$.
\end{lemma}
\begin{proof}
We take for $f$ the factor of $x$ given by Lemma~\ref{lem:extraction},
and obtain $x_A$ by multiplying $f$ by powers of the $z_n$ whose
exponents are determined by bounded binary digits.

\emph{Step 1: the factor $f$.}  Choose $n_1$ so that
$v_{P_j}(x)=\kappa_A$ and $v_{Q_j}(x)=\kappa_B$ for all $j>n_1$
(see \eqref{eq:tails}) and $2^{n_1+1}>\max\{\kappa_A,\kappa_B\}$.  Every
larger $n_1$ has the same two properties.  Lemma~\ref{lem:extraction},
applied with $j_0=n_1+1$, gives a factor $f\in S$ of $x$ with $x/f\in S$
and with finite nonempty support
\[
 (f)=\prod_{j>n_1}P_j^{a_j}Q_j^{b_j},
 \qquad 0\leq a_j\leq \kappa_A,\quad0\leq b_j\leq \kappa_B.
\]
The element $f$ is late-supported at level $n_1+1$ with exponent bound
$\max\{\kappa_A,\kappa_B\}$, and $\max\{\kappa_A,\kappa_B\}\,2^{-n_1-1}<1$.
So Lemma~\ref{lem:balance} gives the balance \eqref{eq:balance}:
\[
 \Sigma:=\sum_{j>n_1}a_j2^{-j}=\sum_{j>n_1}b_j2^{-j}>0.
\]

\emph{Step 2: bounded binary digits.}  Choose $n_2$ at least the largest
index in the support of $f$ and so large that $\Sigma'=2^{n_2}\Sigma$
satisfies $\Sigma'\geq\kappa_A$; this is possible because $\Sigma>0$
by Step~1.  Then $\Sigma'$ is an integer.  Since
$a_j\le\kappa_A$ and $a_j=0$ unless $n_1<j\le n_2$, we have
$\Sigma<\kappa_A2^{-n_1}$, that is, $\Sigma'<\kappa_A2^{n_2-n_1}$.  So
the integer $\Sigma'-\kappa_A$ lies in
$[0,\kappa_A(2^{n_2-n_1}-1)]$.  Every integer in this interval has an
expansion in base two with $n_2-n_1$ digits, each between $0$ and
$\kappa_A$:
\[
 \Sigma'-\kappa_A=\sum_{n_1<j\leq n_2}a'_j2^{n_2-j},\qquad
 0\leq a'_j\leq\kappa_A.
\]
We prove this by induction on the number $k$ of digits: for $\kappa\ge1$
and $k\ge0$, every integer $M$ with $0\le M\le\kappa(2^k-1)$ has the form
$\sum_{i=0}^{k-1}d_i2^i$ with $0\le d_i\le\kappa$.  For $k=0$ we have
$M=0$.  For the inductive step let $0\le M\le\kappa(2^{k+1}-1)$, and let
$d_0$ be whichever of $\min\{\kappa,M\}$ and $\min\{\kappa,M\}-1$ has the
parity of $M$.  If $M\le\kappa$ then $d_0=M$.  Otherwise
$d_0\ge\kappa-1$, and the integer $(M-d_0)/2$ is at most
$\kappa(2^k-1)+\tfrac12$, hence at most $\kappa(2^k-1)$.  In both cases
the induction hypothesis applies to $(M-d_0)/2$.

\emph{Step 3: the factor $x_A$.}  Define integers $e_n$ for
$n_1<n\le n_2$ recursively by
\[
 e_n=a_n-a'_n+\sum_{n_1<i<n}e_i,
\]
and put
\[
 x_A=f\prod_{n_1<n\le n_2}z_n^{e_n}.
\]
We write the divisor of an invertible ideal $I$ of $D$, the family of its
valuations $(v_V(I))_V$, additively, and identify an ideal with its
divisor.  Since $A_nD=A_{n_2}D\prod_{n<j\le n_2}P_j$ and
$(z_n)=A_nP_n^{-1}$,
\[
 \divi(z_n)=A_{n_2}D+\sum_{n<j\leq n_2}P_j-P_n\qquad(n_1<n\leq n_2).
\]
The ideal $A_{n_2}D$ has valuation one at $\EA$ and at the $P_j$ with
$j>n_2$, and valuation zero elsewhere, so its support is disjoint from
$\{P_{n_1+1},\ldots,P_{n_2}\}$ and from the $Q_j$.  Hence the ideals
$A_{n_2}D$, $P_{n_1+1},\ldots,P_{n_2}$ and the $Q_j$ are multiplicatively
independent, and the coefficients of a divisor in the group they generate
are well defined.  By Step~2 the support of $f$ lies among
the $P_j,Q_j$ with $n_1<j\le n_2$.  Hence, with $E=\sum_ne_n$,
\[
 \divi(x_A)=E\,A_{n_2}D
   +\sum_{n_1<j\le n_2}\Bigl(a_j+\sum_{n_1<n<j}e_n-e_j\Bigr)P_j
   +\sum_{j>n_1}b_jQ_j
   =E\,A_{n_2}D+\sum_{n_1<j\le n_2}a'_jP_j+\sum_{j>n_1}b_jQ_j,
\]
where the second equality is the recursion defining $e_j$.  To compute
$E$, give $A_{n_2}D$ the weight $1$ and $P_j$ the weight $2^{n_2-j}$, and
call the weighted sum of the coefficients of $A_{n_2}D$ and of the $P_j$
with $n_1<j\le n_2$ in a divisor its \emph{weighted degree}.  The weighted degree of
$\divi(z_n)$ is $1+\sum_{n<j\le n_2}2^{n_2-j}-2^{n_2-n}=0$.  So $x_A$ and
$f$ have the same weighted degree, namely
$\sum_ja_j2^{n_2-j}=\Sigma'$.  Comparing with the displayed divisor of
$x_A$, whose weighted degree is $E+(\Sigma'-\kappa_A)$ by Step~2, gives
$E=\kappa_A$.

Hence $x_A$ has valuation $\kappa_A$ at $\EA$ and at $P_j$ for $j>n_2$,
valuation $a'_j\le\kappa_A$ at $P_j$ for $n_1<j\le n_2$, valuation
$b_j\le\kappa_B$ at $Q_j$, and valuation zero elsewhere, in particular at
$\EB$.  All these valuations are nonnegative, so $x_A\in D$, and
comparing them with the valuations of $x$ chosen in Step~1 gives
$x_A\mid_Dx$ (Lemma~\ref{lem:geometry}), with the stated valuations of
$x_A$ and $x/x_A$ at $\EA$ and $\EB$.

\emph{Step 4: the ratio $\rho$.}  By construction
$f/x_A=\prod_{n_1<n\le n_2}z_n^{-e_n}$, which is \eqref{eq:rho} with
$m_n=-e_n$.  The divisors of $x_A$ and $\rho$ are supported on $\EA$, the
$P_j$ and the $Q_j$, which lie above $47$ and $7$.  This proves the last
assertion.
\end{proof}

\subsection{Moving an element into a fixed order}\label{sec:joint-return}

Let $S\subseteq D$ be a subring that is not a BFD.  The next lemma turns
approximation in $\mathcal U_N$ into membership in the completed orders of
the fixed order $S_N$.  Membership of a unit in $\mathcal A_{N,p}^\times$
is a condition at finitely many primes $p$, and at each of them it is an
open condition on the image of the unit in $\mathcal U_N$.  At the odd
primes these conditions are described by
Proposition~\ref{prop:orders}, which applies only from the level $b_p$ on.
So the proof imposes each condition at a later level $M_p$, pulls it
back to level $N$ along the map $\theta_{N,M_p}$ of
Lemma~\ref{lem:propagation}, which allows the unit to be chosen at level
$N$, and then descends the resulting membership from $M_p$ to $N$ by
\eqref{eq:saturation}.

Let $z\in F_N^\times$ be a unit at every prime that does not lie above
$47$ or $7$.  By (T1) the primes $47$ and $7$ are unramified in $F_N$, so
$z$ is a unit at every dyadic prime and at every odd prime of absolute
ramification index two, and $\jmath_N(z)$ is defined.

\begin{lemma}[A unit moving an element into a fixed order]\label{lem:return}
Let $S\subseteq D$ be a subring that is not a BFD, and let
$z\in F_N^\times$ be a unit at every prime that does not lie above $47$ or
$7$.  If $\jmath_N(z)\in\overline{\jmath_N(E_N)}$, then there exists
$\eta\in E_N$ such that
\begin{equation}\label{eq:return}
 \eta z\in\mathcal A_{N,p}^{\times}
       \qquad\text{for every prime }p\notin\{47,7\}.
\end{equation}
\end{lemma}
\begin{proof}
We describe the conditions that define $\mathcal A_{N,p}^\times$ by an
open subgroup $\mathcal U_N^S$ of $\mathcal U_N$, choose $\eta$ using the
closure hypothesis, and then verify \eqref{eq:return} prime by prime.

\emph{Step 1: the primes that impose conditions.}  Since $S_N$ has finite
index in $\OO_{F_N}$, we can choose $e\geq1$ with
$2^e\mathcal B_{N,2}\subseteq\mathcal A_{N,2}$.  Let $\mathcal P$ be the
set of odd rational primes $p$ with $\mathcal A_{N,p}\ne\mathcal B_{N,p}$.
It is finite, since $\mathcal A_{N,p}=\mathcal B_{N,p}$ whenever $p$ does
not divide the finite index $[\OO_{F_N}:S_N]$.  By
Proposition~\ref{prop:unramified}, every $p\in\mathcal P$ is ramified at
some level of the tower, so Proposition~\ref{prop:orders} applies to it.

\emph{Step 2: an open condition at a later level.}  Fix $p\in\mathcal P$,
choose a level $b_p$ as in Proposition~\ref{prop:orders}, and put
$M_p=\max\{N,b_p\}$.  By Proposition~\ref{prop:orders}, at level $M_p$ the
completed order $\mathcal A_{M_p,p}$ is the product of its factors at the
maximal ideals $V$ of $D$ above $p$.  The factor at $V$ is the full local
ring of integers if $e(V/p)=1$, and the order $\Lambda_{M_p,V}$ of
\eqref{eq:localorder} if $e(V/p)=2$.  For $V$ with $e(V/p)=2$, write $\OO_{M_p,V}=\OO_{(F_{M_p})_V}$ and
$C_V=\OO_{M_p,V}^\times/W_{M_p,V}^\times$ for the factor of
$\mathcal U_{M_p}$ at the prime of $F_{M_p}$ below $V$, and let $G_V$ be
the image of $\Lambda_{M_p,V}^\times$ in $C_V$.  Two properties of $G_V$
are used.
\begin{enumerate}
\item[(i)] A unit $x$ of $\OO_{M_p,V}$ lies in $\Lambda_{M_p,V}^\times$
exactly when its class lies in $G_V$.  Indeed, if the class of $x$ is that
of $\lambda\in\Lambda_{M_p,V}^\times$, then $x=\lambda c$ with
$c\in W_{M_p,V}^\times$, and $W_{M_p,V}^\times\subseteq
\Lambda_{M_p,V}^\times$ by \eqref{eq:localorder}.
\item[(ii)] $G_V$ is an open subgroup of $C_V$.  Indeed,
\eqref{eq:localorder} gives $p^{s_V}\OO_{M_p,V}\subseteq\Lambda_{M_p,V}$,
so the open subgroup $1+p^{s_V+1}\OO_{M_p,V}$ of $\OO_{M_p,V}^\times$ lies
in $\Lambda_{M_p,V}$.  By Lemma~\ref{lem:order-units} it lies in
$\Lambda_{M_p,V}^\times$: an element $x_0(1+\pi_Vb)$ of $\Lambda_{M_p,V}$
has $\pi_V$-coordinate $x_0b\in p^{s_V}W_{M_p,V}$, so
$b\in p^{s_V}W_{M_p,V}$.  The quotient map
$\OO_{M_p,V}^\times\to C_V$ is open, so $G_V$ contains an open subgroup
of $C_V$; being a subgroup, $G_V$ is therefore open.
\end{enumerate}
Let $G_p\subseteq\mathcal U_{M_p}$ be the subgroup of elements whose
coordinate in $C_V$ lies in $G_V$ for every $V$ above $p$ with
$e(V/p)=2$, all other coordinates being unrestricted.  By~(ii), $G_p$ is
an open subgroup of $\mathcal U_{M_p}$.

\emph{Step 3: an open subgroup at level $N$ and the choice of $\eta$.}
Let $\mathcal U_N^S\subseteq\mathcal U_N$ be the subgroup of elements whose
dyadic coordinate lies in $1+2^e\mathcal B_{N,2}$ and whose image under
$\theta_{N,M_p}$ (Lemma~\ref{lem:propagation}) lies in $G_p$ for every
$p\in\mathcal P$.  Since $e\ge1$, the set $1+2^e\mathcal B_{N,2}$ is an
open subgroup of $\mathcal B_{N,2}^\times$.  Each $\theta_{N,M_p}$ is a
continuous homomorphism, so $\theta_{N,M_p}^{-1}(G_p)$ is an open subgroup
of $\mathcal U_N$.  As $\mathcal P$ is finite, $\mathcal U_N^S$ is an open
subgroup of $\mathcal U_N$.  Because $\jmath_N(z)$ lies in the closure of
$\jmath_N(E_N)$, the neighbourhood $\jmath_N(z)\,\mathcal U_N^S$ of
$\jmath_N(z)$ contains an element $\jmath_N(\eta^{-1})$ with
$\eta\in E_N$.  Then $\jmath_N(\eta z)\in\mathcal U_N^S$.

\emph{Step 4: the dyadic primes and the primes outside $\mathcal P$.}  The
element $\eta z$ lies in the subgroup $1+2^e\mathcal B_{N,2}$, and so does
its inverse.  Both lie in $\mathcal A_{N,2}$, because
$1\in\mathcal A_{N,2}$ and $2^e\mathcal B_{N,2}\subseteq\mathcal A_{N,2}$.
At an odd prime $p\notin\mathcal P\cup\{47,7\}$, the element $\eta z$ is a
unit at every prime above $p$, so $(\eta z)^{\pm1}$ already lie in the
maximal order $\mathcal B_{N,p}=\mathcal A_{N,p}$.

\emph{Step 5: the primes in $\mathcal P$.}  Fix $p\in\mathcal P$ and put
$M=M_p$.  By Step~1, $p$ ramifies at some level, so $p\ne47,7$ by (T1).
The element $\eta z$ is a unit at every prime that does not lie above
$47$ or $7$, and $\jmath_N(\eta z)\in\mathcal U_N^S$ by Step~3, so
Lemma~\ref{lem:propagation} gives
$\jmath_M(\eta z)=\theta_{N,M}(\jmath_N(\eta z))\in G_p$.  We check the
factors of $\mathcal A_{M,p}$ described in Step~2.  At $V$ with
$e(V/p)=2$, the element $\eta z$ is a unit of $\OO_{M,V}$, since
$p\ne47,7$, and its class in $C_V$ lies in $G_V$; by Step~2(i), $\eta z$
and its inverse lie in $\Lambda_{M,V}$.  At $V$ with $e(V/p)=1$, the factor
is the full local ring of integers, which contains $(\eta z)^{\pm1}$.
Thus $(\eta z)^{\pm1}\in\mathcal A_{M,p}$.  Since these elements belong to
$F_N$, \eqref{eq:saturation} returns both to $\mathcal A_{N,p}$.  This
holds for each of the finitely many $p\in\mathcal P$ with the same $\eta$,
which proves \eqref{eq:return}.
\end{proof}

The primes $p\in\mathcal P$, which divide the index
$[\OO_{F_N}:S_N]$, depend only on the fixed order $S_N$, and they are
chosen before $\eta$.  Each condition is imposed at its own level $M_p$
and transported to level $N$ by the continuous map $\theta_{N,M_p}$.  No
comparison of the orders at different levels beyond \eqref{eq:saturation}
is required.  The levels $M_p$ serve only to describe membership in the
completed orders through Proposition~\ref{prop:orders}, and
\eqref{eq:saturation} intersects this membership back to level $N$.

\subsection{Proof of the unit approximation criterion}\label{sec:joint-proof}

\begin{proof}[Proof of Theorem~\ref{thm:joint}]
We reduce to an element $x$ with unbounded factorization lengths, split
it with Lemma~\ref{lem:split}, use \eqref{eq:JT} and
Lemma~\ref{lem:propagation} to approximate the ratio $\rho$ at one common
level, and move both factors into $S$ with the unit supplied by
Lemma~\ref{lem:return}.

\emph{Step 1: reductions.}  Let $S\subseteq D$ be a subring.  If
$\Frac S\ne K$, then $S$ is a BFD by Proposition~\ref{prop:subfields}, and
a BFD is atomic.  We may therefore assume that $S$ is not a BFD and has
fraction field $K$.  An element of $S$ with bounded factorization lengths
already has a factorization into atoms: a factorization into nonunits of
maximal length has irreducible factors.  Take a nonzero nonunit $x\in S$
with unbounded factorization lengths.  By
Proposition~\ref{prop:onesided}, $x$ is not one-sided, so
$\kappa_A=v_{\EA}(x)$ and $\kappa_B=v_{\EB}(x)$ are both positive.

\emph{Step 2: splitting.}  Let $n_0$ be as in \eqref{eq:JT}.  Apply
Lemma~\ref{lem:split} with a cutoff $n_1>n_0$ that is large enough for the
lemma.  It gives $f$, $x_A$, $n_2$ and the exponents $m_n$ in
\eqref{eq:rho}.

\emph{Step 3: approximation at one common level.}  Every $z_n$ that occurs
in \eqref{eq:rho} has $n>n_1>n_0$, so \eqref{eq:JT} applies to each of
these finitely many $z_n$ and gives a level $N_n\ge n$ with
$\jmath_{N_n}(z_n)\in\overline{\jmath_{N_n}(E_{N_n})}$.  Let $N''$ be the
largest of the $N_n$.  Lemma~\ref{lem:propagation}, applied with
$\theta_{N_n,N''}$, moves these closure assertions to the common level
$N''$.  We then enlarge the level to some $N\ge N''$ such that $F_N$
contains $x,f,x/f,x_A$ and $\rho$, and apply
Lemma~\ref{lem:propagation} once more, with $\theta_{N'',N}$.  This gives
$\jmath_N(z_n)\in\overline{\jmath_N(E_N)}$ for all $n_1<n\le n_2$.  The
closure of a subgroup is a closed subgroup, so it is closed under products
and integer powers.  Since $\jmath_N$
is a homomorphism, \eqref{eq:rho} gives
\[
 \jmath_N(\rho)=\prod_{n_1<n\le n_2}\jmath_N(z_n)^{m_n}
 \in\overline{\jmath_N(E_N)}.
\]

\emph{Step 4: moving the factors into $S$.}  By Lemma~\ref{lem:split},
$\rho\in F_N^\times$ is a unit at every prime that does not lie above $47$
or $7$.  Lemma~\ref{lem:return}, applied to $z=\rho$, supplies
$\eta\in E_N$ such that $\eta\rho\in\mathcal A_{N,p}^{\times}$ for every
$p\ne47,7$.  Set
\[
 f'=\eta^{-1}x_A=f(\eta\rho)^{-1},\qquad
 g'=x/f'=(x/f)\eta\rho.
\]
Both elements belong to $D\cap F_N=\OO_{F_N}$, because $x_A\mid_Dx$ and
multiplication by a global unit changes no valuation.  Since $f$ and
$x/f$ lie in $S_N\subseteq\mathcal A_{N,p}$, the expressions on the right
place $f'$ and $g'$ in $\mathcal A_{N,p}$ for every $p\ne47,7$.  At $47$
and $7$ the completed order is maximal by
Proposition~\ref{prop:unramified}, so integrality suffices.
Equation~\eqref{eq:lattice}, applied to the one fixed order $S_N$, gives
$f',g'\in S_N\subseteq S$.  Their valuations at $(\EA,\EB)$ are those of
$x_A$ and $x/x_A$, namely $(\kappa_A,0)$ and $(0,\kappa_B)$.

\emph{Step 5: conclusion.}  The elements $f'$ and $g'$ are one-sided
nonunits of $S$.  By Proposition~\ref{prop:onesided} each has bounded
factorization lengths in $S$, so each has a factorization into atoms of
$S$.  Their product is $x$.  Thus every subring $S$ of $D$ is atomic, and
$D$ is hereditarily atomic.  Failure of ACCP was established in
Theorem~\ref{thm:ambient}.  Interchanging the roles of $\EA$ and $\EB$,
and accordingly of $A_n,P_j,z_n$ and $B_n,Q_j,z_n'$, proves the symmetric
assertion.
\end{proof}

\section{The dyadic places: a Selmer condition}\label{sec:dyadic}

Theorem~\ref{thm:joint} (unit approximation criterion) reduces hereditary
atomicity to condition \eqref{eq:JT}: for every sufficiently large $n$,
the element $z_n$ must be approximated by global units in the compact
group $\mathcal U_N$ of \eqref{eq:TN} for some $N\ge n$.  The group
$\mathcal U_N$ has a dyadic factor $(\OO_{F_N}\otimes\Z_2)^\times$ and
factors at the odd places of ramification index two.  This section treats
the dyadic factor.  Its main result, Theorem~\ref{thm:dyadic}, shows that
at a scaled step $F_n/F_{n-1}$ the Selmer condition (S) implies that the
pro-two part of $z_n$ at the dyadic places is a limit of the pro-two
parts of global units of $F_n$.  Section~\ref{sec:odd} treats the odd places and combines the two
results with $N=n$ (Proposition~\ref{prop:assembly}).

We work at one step, in the notation of Section~\ref{sec:notation}:
$F=F_{n-1}$, $L=F_n$, $\sigma=\sigma_n$, $A=A_{n-1}$, and
$A\OO_L=A'P$ with $A'=A_n$ and $P=P_n$.  The rest of this introduction
explains the strategy; the precise statements begin in
Section~\ref{sec:dyadic-selmer}.

\emph{What must be approximated.}  The prime $2$ is unramified in $L$, so
at each dyadic place $v$ of $L$ the unit group $\OO_{L_v}^\times$ is the
product of the Teichm\"uller representatives of $k_v^\times$, a cyclic
group of odd order, and the group $1+2\OO_{L_v}$ of principal units,
which is a pro-two group.  The odd-order factor concerns residues only,
and Section~\ref{sec:odd} treats it.  What remains is the group
$U_L=\prod_{v\mid2}(1+2\OO_{L_v})$ of dyadic principal units.  We must
show that the component of $z_n$ in $U_L$ lies in the closure of the
components of the global units $E_n$.  A congruence between $z_n$ and a
global unit modulo one fixed power of $2$ is not enough: we need such a
unit for every power of $2$.

\emph{Why the obstruction is an element of a Galois group.}  For each $m$,
the ray class group $\Cl_{2^m}(L)$ of modulus $2^m\OO_L$
(Section~\ref{sec:nt-background}) contains the group of residues modulo
$2^m$ of elements prime to $2$, modulo the residues of global units.  An
element $b$ prime to $2$ is congruent to a global unit modulo $2^m$
exactly when its ideal $(b)$ is trivial in $\Cl_{2^m}(L)$.  Letting $m$
tend to infinity and keeping two-primary parts gives a profinite group
$X_L$.  Class field theory identifies $X_L$ with the Galois group of the
maximal abelian pro-two extension of $L$ unramified outside $2$ and split
at the real places, and the class of an ideal with its Artin symbol.  In
this language (Lemma~\ref{lem:ray-closure}), the element $z_n$ is
approximated by global units in $U_L$ exactly when the Artin symbol of
$(z_n)=A'P^{-1}$ vanishes.  Since $P=\sigma A'$, this symbol is the
difference between the Artin symbol of $A'$ and its image under $\sigma$.
So the question is whether $\sigma$ fixes the Artin symbol of $A'$.

\emph{Why a Selmer group of square classes appears.}  The group $X_L$ is
infinite, and we do not compute it.  Instead we use the closed subgroup
$\mathcal W$ generated by elements that are visibly fixed by $\sigma$:
the image of the corresponding group $X_F$ of $F$, and the Frobenius
elements of the primes of $L$ that ramify over $F$.  On the quotient
$X_L/\mathcal W$, the automorphism $\sigma$ acts as $-1$.  By Kummer
theory, the characters of order two of $X_L$ correspond to the quadratic
extensions $L(\sqrt b)$ unramified outside $2$ and split at the real
places, and hence to square classes $[b]$ of $L$.  Those characters that
vanish on $\mathcal W$ come exactly from the square classes $[x]$ of $F$
that are totally positive, have even valuation outside the ramified and
dyadic primes, and satisfy a local condition at each ramified prime
(Lemma~\ref{lem:ray-selmer}).  A group of global square classes cut out
by local conditions at every place is called a Selmer group.  Ours is
the group $\Delta(F,\delta)$ described in Section~\ref{sec:intro-results}
and defined in Definition~\ref{def:selmer}, and
modulo $[\delta]$, which becomes a square in $L$, it is
$\mathrm{Sel}_\delta$.  Thus $\mathrm{Sel}_\delta$ is the dual of the
largest quotient of $X_L/\mathcal W$ of exponent two.

\emph{Why Hilbert symbols at the dyadic places test it.}  Duality with
$\mathrm{Sel}_\delta$ sees only the largest quotient of exponent two.
Nothing is lost if $X_L/\mathcal W$ itself has exponent two.  The
obstruction to this is a cyclic quotient of $X_L/\mathcal W$ of order
four.  Because $\sigma$ acts on it by $-1$, and inertia at a prime of $T$
supplies an involution that lifts $\sigma$, such a quotient gives a
Galois extension of $F$ with dihedral group of order eight that contains
$L(\sqrt x)$ for a nonzero class $[x]\in\mathrm{Sel}_\delta$.  An
elementary norm computation (Lemma~\ref{lem:dihedral-norm}) shows that
then $x\delta$ is a norm from $F(\sqrt x)$, so the Hilbert symbol
$(x,-\delta)_v$ equals $1$ at every place $v$ of $F$.  For classes of
$\Delta(F,\delta)$ this symbol equals $1$ at every non-dyadic place in
any case, so only the dyadic places can detect the obstruction.  The map
$\lambda_\delta$ records the symbols $(x,-\delta)_v$ at the dyadic places,
and condition (S) asks that it be injective.  Under (S), therefore,
$X_L/\mathcal W$ has no cyclic quotient of order four and has exponent two
(Proposition~\ref{prop:ray-exponent-two}).

\emph{Why exponent two suffices.}  Once $X_L/\mathcal W$ has exponent two,
its elements are detected by the classes of $\mathrm{Sel}_\delta$.  The
Artin symbol of $A'$ lies in $\mathcal W$ exactly when $A'$ splits in
$L(\sqrt x)$ for every $[x]\in\mathrm{Sel}_\delta$, that is, when every
class of $\mathrm{Sel}_\delta$ is a square in the completion
$F_A=L_{A'}$.  Elements of $\mathcal W$ are fixed by $\sigma$, so the
Artin symbol of $A'P^{-1}$ then vanishes
(Corollary~\ref{cor:active-ray}).  The local-square condition at $A$
holds at every scaled step.  An explicit norm $g_A$ from $L$ has odd
valuation at $A$ and at no other prime outside the ramified and dyadic
primes, and its quotient by $-\delta$ is a square at the dyadic places.
Hilbert reciprocity, applied to $x$ and this quotient, proves the
condition for every class of $\Delta(F,\delta)$
(Lemma~\ref{lem:A-square}).

\emph{Organization.}  Section~\ref{sec:dyadic-selmer} defines
$\Delta(F,\delta)$, $\mathrm{Sel}_\delta$, $\lambda_\delta$ and (S),
states Theorem~\ref{thm:dyadic}, and computes the dimension of
$\mathrm{Sel}_\delta$.  Section~\ref{sec:dyadic-ray} identifies the
obstruction to dyadic approximation with an Artin symbol.
Section~\ref{sec:dyadic-kummer} introduces the subgroup $\mathcal W$ and
identifies the dual of $X_L/(\mathcal W+2X_L)$ with $\mathrm{Sel}_\delta$.
Section~\ref{sec:dyadic-dihedral} shows that (S) excludes cyclic quotients
of order four.  Section~\ref{sec:dyadic-A} proves the local-square
condition at $A$ and Theorem~\ref{thm:dyadic}.

\subsection{The Selmer group and condition (S)}\label{sec:dyadic-selmer}

At a scaled step (Definition~\ref{def:scaled-step}) write
$F=F_{n-1}$ and $L=F_n=F(\sqrt\delta)$, with $\delta=t^2-4y$ as in
\eqref{eq:scaled-family}.  The extension is totally real, $2$ is
absolutely unramified in $L$, and the two dyadic primes of $F$ are inert
in $L$.  Let $T$ be the set of \emph{all} finite primes of $F$ ramified
in $L$.  It contains the detecting primes of
Definition~\ref{def:scaled-step}(vi), so it is nonempty, and all its
primes are odd.  Let $\mathcal D$ be the set of dyadic primes of $F$;
thus $|\mathcal D|=2$.  We keep $|\mathcal D|$ in the dimension count
below to show the role of the number of dyadic places.  Apart from the
restatement of (S) after Lemma~\ref{lem:selmer-dimension}, the arguments
of this section use only that $2$ is unramified in $L$, and not the
inertness of the dyadic primes of $F$ in $L$.

\begin{definition}[Selmer group and Selmer map]\label{def:selmer}
Let
\[
 \Delta(F,\delta)=\left\{[x]\in F^\times/F^{\times2}:
 \begin{array}{l}
 x\text{ is totally positive},\\
 v_{\mathfrak p}(x)\equiv0\pmod2\quad(\mathfrak p\notin T\cup\mathcal D),\\
 {}[x]_{\mathfrak p}\in\langle[\delta]_{\mathfrak p}\rangle\quad(\mathfrak p\in T)
 \end{array}\right\},\qquad
 \mathrm{Sel}_\delta=\Delta(F,\delta)/\langle[\delta]\rangle,
\]
where $[x]_{\mathfrak p}$ is the local square class of $x$ in
$F_{\mathfrak p}^\times/F_{\mathfrak p}^{\times2}$.  The \emph{Selmer map}
is
\[
 \lambda_\delta:\mathrm{Sel}_\delta\longrightarrow(\F_2^{\mathcal D})^0,\qquad
 \lambda_\delta([x])=([x,-\delta]_v)_{v\in\mathcal D},
 \qquad
 (\F_2^{\mathcal D})^0=\{(x_v)\in\F_2^{\mathcal D}:\textstyle\sum_vx_v=0\},
\]
where $[x,-\delta]_v\in\F_2$ is the additive Hilbert symbol of
Section~\ref{sec:notation}.
\end{definition}

The local condition at $\mathfrak p\in T$ says that $x$ is a square or
$\delta$ times a square in $F_{\mathfrak p}$.  Every prime of $T$ enters
this condition, including the primes of $T$ that are not detecting
primes.

The Hilbert symbols are those of Section~\ref{sec:nt-background}: the
local norm criterion and bimultiplicativity are
\cite[Chapter~III, Theorem~4.4]{MilneCFT} at the finite places and
elementary at the real places, and the global product formula
is \cite[Chapter~V, Example~5.4]{MilneCFT}.  The map $\lambda_\delta$ is
$\F_2$-linear by bimultiplicativity, and it descends from
$\Delta(F,\delta)$ to $\mathrm{Sel}_\delta$ because
$(\delta,-\delta)_v=1$.  Its values have coordinate sum zero by the
product formula applied to $x$ and $-\delta$: at each odd prime of $T$
the symbol is trivial, since the local class of $x$ is $1$ or
$[\delta]_{\mathfrak p}$ and $(\delta,-\delta)_{\mathfrak p}=1$; at every
other odd prime both $x$ and $-\delta$ have even valuation, so the symbol
is trivial (Section~\ref{sec:nt-background}); and at the real places $x$
is positive.  Hence the product of the symbols over the dyadic places is
$1$.

The \emph{Selmer condition} is
\begin{equation}\label{eq:S}
 \tag{S} \ker\lambda_\delta=0.
\end{equation}

For a number field in which $2$ is unramified, let $\operatorname{pr}_2$
remove the odd-order Teichm\"uller factor from each dyadic unit and keep
its component in the group of principal units.  This component is a
pro-two element; it includes torsion, for instance
$\operatorname{pr}_2(-1)=-1$.  For $L$ this is the projection
$\prod_{v\mid2}\OO_{L_v}^\times\to U_L$, defined on every element of $L$
that is a unit at the dyadic primes.

\begin{theorem}[Dyadic approximation]\label{thm:dyadic}
Let $F_n/F_{n-1}$ be a scaled step in the sense of
Definition~\ref{def:scaled-step}.  If this step satisfies condition
\textup{(S)}, then
\[
 \operatorname{pr}_2(z_n)\in\overline{\operatorname{pr}_2(E_n)}
 \subseteq\prod_{v\mid2}(1+2\OO_{L_v}).
\]
\end{theorem}

The proof is given at the end of Section~\ref{sec:dyadic-A}.  We first
show that (S) is a finite test and that, when it holds, the dimension of
$\mathrm{Sel}_\delta$ is forced.  For a finite set $\mathcal S$ of finite
primes of $F$, let
\[
 \mathcal V_{\mathcal S}(F)=\{[x]\in F^\times/F^{\times2}:
 v_{\mathfrak p}(x)\equiv0\pmod2\text{ for every prime }\mathfrak p\notin\mathcal S\},
\]
so that $\mathcal V_\varnothing(F)=\mathcal V(F)$ is the group of square
classes of virtual units (Section~\ref{sec:nt-background}).  Write
$\operatorname{rk}_2\Cl(F)=\dim_{\F_2}\Cl(F)/2\Cl(F)$, which equals
$\dim_{\F_2}\Cl(F)[2]$ because $\Cl(F)$ is finite, and let
$\langle[\mathfrak p]:\mathfrak p\in\mathcal S\rangle$ be the subspace
of $\Cl(F)/2\Cl(F)$ spanned by the classes of the primes in $\mathcal S$.

\begin{lemma}[The dimension of the Selmer group]\label{lem:selmer-dimension}
For every finite set $\mathcal S$ of finite primes of $F$, the group
$\mathcal V_{\mathcal S}(F)$ is finite and
\begin{equation}\label{eq:supported-dimension}
 \dim_{\F_2}\mathcal V_{\mathcal S}(F)=[F:\Q]+\operatorname{rk}_2\Cl(F)+|\mathcal S|
   -\dim_{\F_2}\langle[\mathfrak p]:\mathfrak p\in\mathcal S\rangle.
\end{equation}
The group $\Delta(F,\delta)$ is finite, contains the nonzero class
$[\delta]$, and satisfies
\begin{equation}\label{eq:selmer-dimension}
 \dim\Delta(F,\delta)\ge|\mathcal D|,\qquad
 \dim\mathrm{Sel}_\delta\ge|\mathcal D|-1.
\end{equation}
If \textup{(S)} holds, then $\dim\mathrm{Sel}_\delta=|\mathcal D|-1$.
\end{lemma}

\begin{proof}
We compute $\mathcal V_{\mathcal S}(F)$ through the parity of valuations,
and then count the conditions that cut out $\Delta(F,\delta)$.

\emph{Step 1: the formula \eqref{eq:supported-dimension}.}  Consider the
valuation-parity map
$\mathcal V_{\mathcal S}(F)\to\F_2^{\mathcal S}$,
$[x]\mapsto(v_{\mathfrak p}(x)\bmod2)_{\mathfrak p\in\mathcal S}$.  Its
kernel consists of the classes with even valuation at every prime, that
is, $\mathcal V(F)$.  Since $F$ is totally real, its torsion units are
$\pm1$, and $\dim_{\F_2}\mathcal V(F)=[F:\Q]+\operatorname{rk}_2\Cl(F)$
(Section~\ref{sec:nt-background}); the unit contribution $[F:\Q]$
includes the torsion unit $-1$.  The image of the parity map is the
kernel of the map $\F_2^{\mathcal S}\to\Cl(F)/2\Cl(F)$ that sends
$(e_{\mathfrak p})$ to the class of $\prod_{\mathfrak p\in\mathcal
S}\mathfrak p^{e_{\mathfrak p}}$.  Indeed, a parity vector lies in that
kernel exactly when the corresponding ideal
$\prod_{\mathfrak p\in\mathcal S}\mathfrak p^{e_{\mathfrak p}}$ is a square
times a principal ideal $(x)$, and then $[x]\in\mathcal V_{\mathcal S}(F)$
has the given parity vector; conversely, if
$[x]\in\mathcal V_{\mathcal S}(F)$, then $(x)$ is the product of
$\prod_{\mathfrak p\in\mathcal S}\mathfrak p^{e_{\mathfrak p}}$, with
$e_{\mathfrak p}\equiv v_{\mathfrak p}(x)$, and a square.  The image of
$\F_2^{\mathcal S}\to\Cl(F)/2\Cl(F)$ is
$\langle[\mathfrak p]:\mathfrak p\in\mathcal S\rangle$, so the kernel has
dimension $|\mathcal S|-\dim\langle[\mathfrak p]:\mathfrak p\in\mathcal
S\rangle$.  Adding the dimensions of kernel and image proves both
\eqref{eq:supported-dimension} and the finiteness of
$\mathcal V_{\mathcal S}(F)$.

\emph{Step 2: the lower bounds \eqref{eq:selmer-dimension}.}  Take
$\mathcal S=T\cup\mathcal D$.  Then $\Delta(F,\delta)$ is the subgroup of
$\mathcal V_{\mathcal S}(F)$ cut out by positivity and by the local
conditions at the primes of $T$, so it is finite.  Positivity imposes at
most $[F:\Q]$ linear conditions, one for the sign at each real place.  At
$\mathfrak q\in T$ the group $F_{\mathfrak q}^\times/F_{\mathfrak
q}^{\times2}$ has dimension two, because $\mathfrak q$ is odd, and the
condition $[x]_{\mathfrak q}\in\langle[\delta]_{\mathfrak q}\rangle$ is at
most one linear condition.  Since
$\dim_{\F_2}\langle[\mathfrak p]:\mathfrak p\in\mathcal S\rangle\le
\operatorname{rk}_2\Cl(F)$, formula \eqref{eq:supported-dimension} gives
\[
 \dim\Delta(F,\delta)\ge
 [F:\Q]+|T|+|\mathcal D|-[F:\Q]-|T|=|\mathcal D|.
\]
The class $[\delta]$ is nonzero because $L=F(\sqrt\delta)$ is a quadratic
field.  It belongs to $\Delta(F,\delta)$: the element $\delta$ is totally
positive because $L$ is totally real, its valuation is even at every odd
prime unramified in $L$, and the condition at $T$ holds trivially.  Hence
$\dim\mathrm{Sel}_\delta=\dim\Delta(F,\delta)-1\ge|\mathcal D|-1$.

\emph{Step 3: the dimension under \textup{(S)}.}  The map $\lambda_\delta$
takes values in $(\F_2^{\mathcal D})^0$, which has dimension
$|\mathcal D|-1$.  If $\lambda_\delta$ is injective, then
$\dim\mathrm{Sel}_\delta\le|\mathcal D|-1$, and with
\eqref{eq:selmer-dimension} equality holds.
\end{proof}

In the present setting $|\mathcal D|=2$, so $\mathrm{Sel}_\delta$ is never
zero and $(\F_2^{\mathcal D})^0$ is the line spanned by $(1,1)$.
Condition (S) therefore says that $\mathrm{Sel}_\delta$ has exactly two
elements and that its nonzero class $[x]$ satisfies
$\lambda_\delta([x])=(1,1)$, that is, $(x,-\delta)_v=-1$ at both dyadic
places $v$.  This is the form of (S) described in
Section~\ref{sec:intro-results}.

\subsection{Unit approximation as an Artin symbol}\label{sec:dyadic-ray}

Let $\Cl_{2^m}(F)$ be the ray class group of modulus $2^m\OO_F$, with
\emph{no infinite part}, and define $\Cl_{2^m}(L)$ likewise
(Section~\ref{sec:nt-background}).  The prime $2$ is absolutely
unramified in $F$ and in $L$, so $1+2\OO_{F_v}$ is the group of
principal units at each dyadic place $v$ of $F$, and likewise for $L$.
Since the modulus has no infinite part, the principal ideals that are
trivial in $\Cl_{2^m}(F)$ are those generated by elements congruent to
$1$ modulo $2^m$ at every dyadic place, with no positivity requirement.
Put
\[
\begin{gathered}
 U_F=\prod_{v\mid2}(1+2\OO_{F_v}),\qquad
 U_F^{(m)}=\prod_{v\mid2}(1+2^m\OO_{F_v}),\\
 \overline E_{F,2}=\overline{\operatorname{pr}_2(\OO_F^\times)}\subseteq U_F,\qquad
 X_F=\varprojlim_m \Cl_{2^m}(F)_{(2)},
\end{gathered}
\]
and define $U_L$, $U_L^{(m)}$, $\overline E_{L,2}$ and $X_L$ in the same
way.  Recall that the subscript $(2)$ denotes the two-primary subgroup of
a finite abelian group.  The groups $U_F$ and $U_L$ are written
multiplicatively; the ray class groups and their inverse limits are
written additively.

Class field theory identifies $X_F$ with the Galois group of the maximal
abelian pro-two extension $F^{\rm ab}_2/F$ unramified outside $2$ and split
at every real place \cite[Chapter~V, Theorems~3.5--3.6, Corollary~3.7 and
Example~3.9]{MilneCFT}; likewise $X_L=\Gal(L^{\rm ab}_2/L)$.  The
splitting at the real places reflects the absence of an infinite part in
the moduli.  For a fractional ideal $I$ of $F$ prime to $2$, the classes
of $I$ in the groups $\Cl_{2^m}(F)_{(2)}$ are compatible; we write
$\operatorname{Art}_F(I)\in X_F$ for the resulting element, which is the
Artin symbol of $I$ under this identification.  We define
$\operatorname{Art}_L$ likewise.

\begin{lemma}[The ray obstruction to unit approximation]\label{lem:ray-closure}
There is an exact sequence of compact groups
\begin{equation}\label{eq:ray-exact}
 0\longrightarrow U_F/\overline E_{F,2}
 \longrightarrow X_F\longrightarrow\Cl(F)_{(2)}\longrightarrow0.
\end{equation}
If $b\in F^\times$ is a unit at every dyadic prime, then
\begin{equation}\label{eq:ray-principal}
 \operatorname{Art}_F((b))=0
 \quad\Longleftrightarrow\quad
 \operatorname{pr}_2(b)\in\overline E_{F,2}.
\end{equation}
The same assertions hold with $L$ in place of $F$.
\end{lemma}

\begin{proof}
We take two-primary parts of the ray class exact sequence for the modulus
$2^m$ and pass to the limit over $m$.  We give the proof for $F$; the
proof for $L$ is identical.

\emph{Step 1: a finite level.}  The ray class exact sequence
\cite[Chapter~V, Theorem~1.7]{MilneCFT} for the modulus $2^m\OO_F$ with no
infinite part reads
\[
 0\longrightarrow(\OO_F/2^m\OO_F)^\times/\operatorname{im}(\OO_F^\times)
 \longrightarrow\Cl_{2^m}(F)\longrightarrow\Cl(F)\longrightarrow0.
\]
By the Chinese remainder theorem,
$(\OO_F/2^m\OO_F)^\times=\prod_{v\mid2}\OO_{F_v}^\times/(1+2^m\OO_{F_v})$,
and each factor is the product of the Teichm\"uller representatives of
$k_v^\times$, of odd order, and the two-group
$(1+2\OO_{F_v})/(1+2^m\OO_{F_v})$.  Taking two-primary parts is exact on
finite abelian groups.  It removes the odd residue factors and keeps the
whole group $U_F/U_F^{(m)}$.  For a subgroup $H$ of a finite abelian group
$G$, the two-primary part of $H$ is its image under the projection of $G$
onto $G_{(2)}$, so the two-primary part of the image of $\OO_F^\times$ is
the image of $\operatorname{pr}_2(\OO_F^\times)$.  Since $U_F^{(m)}$ is
open, $\operatorname{pr}_2(\OO_F^\times)U_F^{(m)}=\overline
E_{F,2}U_F^{(m)}$.  We obtain the exact sequence
\[
 0\longrightarrow
 U_F/(\overline E_{F,2}U_F^{(m)})
 \longrightarrow \Cl_{2^m}(F)_{(2)}
 \longrightarrow\Cl(F)_{(2)}\longrightarrow0.
\]

\emph{Step 2: the limit.}  These sequences are compatible as $m$ grows.
Passing to inverse limits preserves exactness, since all the groups are
finite.  The groups $U_F^{(m)}$ form a basis of neighbourhoods of $1$ in
the compact group $U_F$, and $\overline E_{F,2}$ is closed, so
$\varprojlim_mU_F/(\overline E_{F,2}U_F^{(m)})=U_F/\overline E_{F,2}$.
This proves \eqref{eq:ray-exact}.

\emph{Step 3: Artin symbols of principal ideals.}  Let $b\in F^\times$ be a
unit at every dyadic prime.  The ideal $(b)$ is prime to $2$, and the
finite ray sequence sends the residue of $b$ modulo $2^m$ to the class of
$(b)$ in $\Cl_{2^m}(F)$, up to an inverse if the reciprocal Artin
convention is used, which does not affect vanishing.  After projection to
two-primary parts, this residue becomes the class of
$\operatorname{pr}_2(b)$ in $U_F/(\overline E_{F,2}U_F^{(m)})$.  Hence
$\operatorname{Art}_F((b))=0$ exactly when $\operatorname{pr}_2(b)$ lies
in $\overline E_{F,2}U_F^{(m)}$ for every $m$, that is, in
$\overline E_{F,2}$.  This proves \eqref{eq:ray-principal}.

Because the modulus has no infinite part, the subgroup in Step 1 is the
closure of all global units, not only of the totally positive ones.
\end{proof}

\subsection{The action of \texorpdfstring{$\sigma$}{sigma} and the Selmer group as a dual}\label{sec:dyadic-kummer}

Let $X=X_L$ and let $\sigma$ be the nontrivial automorphism of $L/F$.  The
field $L^{\rm ab}_2$ is defined in terms of $L$ alone, so it is stable
under every automorphism of an algebraic closure that fixes $F$, and
$L^{\rm ab}_2/F$ is Galois.  Choose an extension $\tilde\sigma$ of
$\sigma$ to $L^{\rm ab}_2$.  Then $\sigma$ acts on $X$ by
$g\mapsto\tilde\sigma g\tilde\sigma^{-1}$.  Since $X$ is abelian, this
action does not depend on the choice of $\tilde\sigma$.  We write it
additively as $g\mapsto\sigma(g)$.  The Artin symbols are compatible with
it: $\operatorname{Art}_L(\sigma I)=\sigma(\operatorname{Art}_L(I))$ for
every fractional ideal $I$ of $L$ prime to $2$.  For a prime ideal this is
the relation between the Frobenius elements of conjugate primes in the
finite subextensions of $L^{\rm ab}_2/L$ that are Galois over $F$ (the
proof of \cite[8.14]{MilneANT} applies); such subextensions exhaust
$L^{\rm ab}_2$, since a finite subextension $M$ lies in its compositum
with $\tilde\sigma M$, which is Galois over $F$.  The general case
follows by multiplicativity.

Extension and norm of ideals induce continuous homomorphisms
\[
 \operatorname{ext}:X_F\longrightarrow X,\qquad \nu:X\longrightarrow X_F.
\]
They are defined at the finite ray levels, since congruence to $1$ modulo
$2^m$ is preserved by extension and by the norm.  They satisfy
\begin{equation}\label{eq:ray-transfer}
 \operatorname{ext}\nu=1+\sigma,\qquad
 \sigma\operatorname{ext}=\operatorname{ext}.
\end{equation}
The second identity holds because an ideal extended from $F$ is stable
under $\sigma$.  The first follows from ideal factorization: for a prime
$\mathfrak P$ of $L$, the ideal $N_{L/F}(\mathfrak P)\OO_L$ equals
$\mathfrak P\cdot\sigma\mathfrak P$, whether the prime of $F$ below
$\mathfrak P$ splits, is inert or ramifies in $L$.

For each $\mathfrak q\in T$, let $\mathfrak Q$ be the unique prime of $L$
above $\mathfrak q$.  Define the closed subgroup
\begin{equation}\label{eq:W-ray}
 \mathcal W=\overline{\langle \operatorname{ext}(X_F),
              \operatorname{Art}_L(\mathfrak Q):\mathfrak q\in T\rangle}
 \subseteq X.
\end{equation}
Every generator of $\mathcal W$ is fixed by $\sigma$: the elements of
$\operatorname{ext}(X_F)$ by \eqref{eq:ray-transfer}, and the elements
$\operatorname{Art}_L(\mathfrak Q)$ because $\sigma\mathfrak Q=\mathfrak Q$.
For every $g\in X$, the element $g+\sigma(g)=\operatorname{ext}(\nu(g))$
lies in $\mathcal W$.  Thus
\begin{equation}\label{eq:ray-inversion}
 \mathcal W\subseteq X^\sigma,\qquad \sigma=-1\text{ on }X/\mathcal W.
\end{equation}

The next lemma identifies the characters of order two of $X/\mathcal W$.
A continuous character $\psi:X\to\F_2$ corresponds to a quadratic
extension of $L$ inside $L^{\rm ab}_2$, or to $L$ itself if $\psi=0$, and
by Kummer theory (Section~\ref{sec:nt-background}) this extension is
$L(\sqrt b)$ for a square class $[b]$ of $L$.  We write $\psi_b$ for the
character, so that $\psi_b(g)=0$ exactly when $g$ fixes $\sqrt b$.  In
particular, for a prime $\mathfrak P$ of $L$ not above $2$ at which
$L(\sqrt b)/L$ is unramified, that is, at which $b$ has even valuation
(Section~\ref{sec:nt-background}), the value
$\psi_b(\operatorname{Art}_L(\mathfrak P))$ is $0$
exactly when $\mathfrak P$ splits in $L(\sqrt b)$, that is, when $b$ is a
square in $L_{\mathfrak P}$.

\begin{lemma}[The quadratic characters of the quotient]\label{lem:ray-selmer}
The Kummer pairing identifies
\begin{equation}\label{eq:ray-selmer}
 \operatorname{Hom}_{\mathrm{cts}}(X/(\mathcal W+2X),\F_2)
 \simeq \Delta(F,\delta)/\langle[\delta]\rangle=\mathrm{Sel}_\delta.
\end{equation}
A class $[x]\in\mathrm{Sel}_\delta$ pairs with a prime of $L$ through the quadratic
extension $L(\sqrt x)/L$.
\end{lemma}

\begin{proof}
We describe all quadratic characters of $X$ by square classes of $L$, and
then translate vanishing on each generator of $\mathcal W$ into a
condition on the square class.  A character with values in $\F_2$
vanishes on $2X$ automatically, so the characters of
$X/(\mathcal W+2X)$ are the characters of $X$ that vanish on
$\mathcal W$, that is, on its generators.

\emph{Step 1: characters of $X$.}  By Kummer theory
\cite[Chapter~VII, Appendix~A, Theorem~A.3]{MilneCFT}, the quadratic
characters of $X$ correspond to the square classes $[b]$ of $L$ such that
$L(\sqrt b)$ is unramified outside $2$ and split at the real places.
These are the classes of elements $b$ that are positive at all real
places and have even valuation at all odd finite primes
(Section~\ref{sec:nt-background}).

\emph{Step 2: the transfer identity.}  For such $b$, the norm $N_{L/F}b$ is
totally positive and has even valuation at every odd prime of $F$, so it
defines a character $\psi_{N_{L/F}b}$ of $X_F$.  We show
\begin{equation}\label{eq:transfer-norm-dual}
 \psi_b\circ\operatorname{ext}=\psi_{N_{L/F}b}.
\end{equation}
Exclude the finitely many primes of $F$ that lie above $2$, ramify in
$L/F$, or lie below a prime dividing $b$.  At every other prime
$\mathfrak p$ of $F$, the value of $\psi_b\circ\operatorname{ext}$ at
$\operatorname{Art}_F(\mathfrak p)$ is the sum of the quadratic residue
characters of $b$ at the primes of $L$ above $\mathfrak p$.  If
$\mathfrak p$ splits as $\mathfrak P\cdot\sigma\mathfrak P$, both residue
fields equal that of $\mathfrak p$, and the residue of $N_{L/F}b$ at
$\mathfrak p$ is the product of the residues of $b$ at $\mathfrak P$ and
at $\sigma\mathfrak P$; multiplicativity of the quadratic character gives
the character of $N_{L/F}b$ at $\mathfrak p$.  If $\mathfrak p$ is inert,
with residue field $\F_q$ and $\mathfrak P=\mathfrak p\OO_L$ with residue
field $\F_{q^2}$, then the residue of $N_{L/F}b$ is the norm of the
residue $\bar b$, and the quadratic character of an element of
$\F_{q^2}$ equals that of its norm to $\F_q$:
\[
 (N_{\F_{q^2}/\F_q}\bar b)^{(q-1)/2}
       =\bar b^{(q^2-1)/2}.
\]
Thus the two characters in \eqref{eq:transfer-norm-dual} agree on the
Artin symbols of all but finitely many primes.  Both factor through a
common finite quotient of $X_F$, and by the Chebotarev density theorem
\cite[Chapter~VIII, Theorem~7.4]{MilneCFT} every element of a finite
quotient is the Frobenius element of infinitely many primes.  This proves
\eqref{eq:transfer-norm-dual}.

\emph{Step 3: vanishing on $\operatorname{ext}(X_F)$.}  By
\eqref{eq:transfer-norm-dual}, the character $\psi_b$ vanishes on
$\operatorname{ext}(X_F)$ exactly when $N_{L/F}b=x'^2$ for some
$x'\in F^\times$.  In that case $b/x'$ has norm one, and Hilbert's
Theorem~90 writes it as $\xi/\sigma(\xi)$ with $\xi\in L^\times$.  This
instance can be checked directly: for an element $\xi'\ne-1$ of norm
one, take $\xi=1+\xi'$, and for $\xi'=-1$ take $\xi=\sqrt\delta$.  It
follows that
\[
 b=\frac{x'}{N_{L/F}\xi}\,\xi^2,
\]
so the square class of $b$ comes from $x=x'/N_{L/F}\xi\in F^\times$.
Conversely, a class coming from $x\in F^\times$ has square norm $x^2$, so
its character vanishes on $\operatorname{ext}(X_F)$.

\emph{Step 4: the local conditions.}  Let $b=x\xi^2$ with
$x\in F^\times$.  Since $L$ is totally real, $b$ is positive at all real
places of $L$ exactly when $x$ is totally positive.  At an odd prime of
$F$ outside $T$, the ramification index in $L$ is one, so the condition
that $b$ have even valuation at the primes of $L$ above it says that the
valuation of $x$ is even.  At $\mathfrak q\in T$, the ramification index
is two, so every element of $F$ has even valuation in $L_{\mathfrak Q}$.
Hence $L(\sqrt x)$ is unramified or trivial at $\mathfrak Q$, and its
character vanishes on $\operatorname{Art}_L(\mathfrak Q)$ exactly when
$x$ is a square in $L_{\mathfrak Q}=F_{\mathfrak q}(\sqrt\delta)$.  The
kernel of
\[
 F_{\mathfrak q}^\times/F_{\mathfrak q}^{\times2}
 \longrightarrow L_{\mathfrak Q}^\times/L_{\mathfrak Q}^{\times2}
\]
is $\langle[\delta]_{\mathfrak q}\rangle$.  Thus the Frobenius conditions
at the primes $\mathfrak Q$ are exactly the local conditions at $T$ in
the definition of $\Delta(F,\delta)$.  No condition arises at the dyadic
primes, since $L^{\rm ab}_2/L$ may ramify there.

\emph{Step 5: the kernel of restriction.}  By Steps~3 and~4, the
characters of $X$ vanishing on $\mathcal W$ are the images in
$L^\times/L^{\times2}$ of the classes of $\Delta(F,\delta)$.  The kernel
of the restriction map $F^\times/F^{\times2}\to L^\times/L^{\times2}$ is
$\langle[\delta]\rangle$.  This proves \eqref{eq:ray-selmer}, and the
description of the pairing follows from the remark before the lemma.
\end{proof}

\subsection{Excluding a cyclic quotient of order four}\label{sec:dyadic-dihedral}

Lemma~\ref{lem:ray-selmer} describes only the largest quotient of
$X/\mathcal W$ of exponent two.  The next elementary norm calculation is
the step that upgrades this quadratic information to a statement about
the whole pro-two group $X/\mathcal W$.  It is a form of the classical
criterion for embedding a biquadratic field in a dihedral extension of
degree eight; we prove the direction we need.

\begin{lemma}[A dihedral norm]\label{lem:dihedral-norm}
Suppose that $\widetilde L/F$ is a Galois extension with group
\[
 \mathrm{Dih}_4=\langle r,s:r^4=s^2=1,
                          \ srs=r^{-1}\rangle.
\]
Assume that the quadratic subfield fixed by the rotation $r$ is
$F(\sqrt\delta)$ and that the biquadratic subfield, fixed by $r^2$, is
$F(\sqrt\delta,\sqrt x)$.  Then $x\delta$ is a norm from $F(\sqrt x)$, and
\[
 (x,-\delta)_v=1\quad\text{at every place }v\text{ of }F.
\]
\end{lemma}

\begin{proof}
We produce an element of $F(\sqrt x)$ whose norm is $x\delta$ times a
square in $F$.  The subgroups of order four of $\mathrm{Dih}_4$ other
than $\langle r\rangle$ are $\langle r^2,s\rangle$ and
$\langle r^2,rs\rangle$; their fixed fields are $F(\sqrt x)$ and
$F(\sqrt{x\delta})$ in some order, and replacing $s$ by $rs$ exchanges
them.  Choose the reflection $s$ so that
$\widetilde L^{\langle r^2,s\rangle}=F(\sqrt x)$.  The element $r^2$ is
central and acts nontrivially on the quadratic extension
$\widetilde L^{\langle s\rangle}/F(\sqrt x)$, so we may choose
$0\ne\xi\in\widetilde L^{\langle s\rangle}$ with $r^2\xi=-\xi$.  Then
$\xi^2$ is fixed by $r^2$ and $s$, so $\xi^2\in F(\sqrt x)$.  Put
$\xi'=\xi\,r(\xi)$.  Using $r^2\xi=-\xi$ and $sr=r^{-1}s$,
\[
 r\xi'=-\xi',\qquad
 s\xi'=\xi\, r^{-1}(\xi)=-\xi'.
\]
Thus $\xi'$ is fixed by $r^2$ and by $rs$, so it lies in
$F(\sqrt{x\delta})$, and it lies in the nontrivial eigenspace of the
automorphism of $F(\sqrt{x\delta})/F$.  Write $\xi'=b\sqrt{x\delta}$ with
$b\in F^\times$.  Since $r$ restricts to the nontrivial automorphism of
$F(\sqrt x)/F$,
\[
 N_{F(\sqrt x)/F}(\xi^2)=\xi^2r(\xi^2)=\xi'^2=b^2x\delta.
\]
Dividing $\xi^2$ by $b$ gives an element of $F(\sqrt x)$ with norm
$x\delta$.  A global norm is a local norm at every place, so the local
norm criterion gives $(x,x\delta)_v=1$: at the finite places by
\cite[Chapter~III, Theorem~4.4]{MilneCFT}, and at a real place by
Section~\ref{sec:nt-background}, since a norm from $\mathbb R(\sqrt x)$ is
positive when $x<0$.  Finally $(x,x\delta)_v=(x,-\delta)_v$, since
$x\delta=(-x)(-\delta)$ and $(x,-x)_v=1$.
\end{proof}

\begin{proposition}[Exponent two under (S)]\label{prop:ray-exponent-two}
If \textup{(S)} holds, then $2X\subseteq\mathcal W$.  Consequently the
field fixed by $\mathcal W$ is the compositum of the fields $L(\sqrt x)$
for $[x]\in\mathrm{Sel}_\delta$.
\end{proposition}

\begin{proof}
We show that a failure of $2X\subseteq\mathcal W$ produces a dihedral
extension of $F$ of degree eight, which by Lemma~\ref{lem:dihedral-norm}
yields a nonzero element of $\ker\lambda_\delta$.

\emph{Step 1: a cyclic quotient of order four.}  Suppose that
$X/\mathcal W$ is not killed by two.  Then some finite quotient of the
pro-two group $X/\mathcal W$ has exponent at least four, and hence has a
cyclic quotient of order four.  Let $Y$ be the kernel of the resulting
surjection $X\to\Z/4\Z$; it is an open subgroup containing $\mathcal W$.
By \eqref{eq:ray-inversion}, $\mathcal W\subseteq X^\sigma$, so
$\sigma(\mathcal W)=\mathcal W$, and $\sigma=-1$ on $X/\mathcal W$ maps
$Y/\mathcal W$ to $-Y/\mathcal W=Y/\mathcal W$.  Hence $Y$ is stable under
$\sigma$.  The fixed field $\widetilde L$ of $Y$ is therefore Galois over
$F$, and $\widetilde L/L$ is cyclic of degree four.

\emph{Step 2: the group is dihedral.}  Take $\mathfrak q\in T$; the set $T$
is nonempty.  The extension $\widetilde L/L$ is unramified above
$\mathfrak q$, because $\widetilde L$ is contained in $L^{\rm ab}_2$,
which is unramified outside $2$, and $\mathfrak q$ is odd, since the
dyadic primes are unramified in $L$, whereas $L/F$ has
ramification index two there.  Thus an inertia group in
$\Gal(\widetilde L/F)$ above $\mathfrak q$ has order two and maps
isomorphically onto $\Gal(L/F)$.  It supplies an involution outside
$\Gal(\widetilde L/L)$.  This involution acts on the cyclic group
$\Gal(\widetilde L/L)\simeq X/Y$ as $\sigma$ does, that is, by inversion.
Hence $\Gal(\widetilde L/F)\simeq\mathrm{Dih}_4$, with rotation subgroup
$\Gal(\widetilde L/L)$, whose fixed field is $L=F(\sqrt\delta)$.

\emph{Step 3: a contradiction with \textup{(S)}.}  The nonzero quadratic
character of the cyclic quotient $X/Y$ vanishes on $\mathcal W+2X$.  By
Lemma~\ref{lem:ray-selmer}, its field is $L(\sqrt x)$ for some nonzero
$[x]\in\mathrm{Sel}_\delta$.  This field is the unique quadratic
extension of $L$ inside $\widetilde L$, so the biquadratic subfield of
$\widetilde L/F$ is $F(\sqrt\delta,\sqrt x)$.
Lemma~\ref{lem:dihedral-norm} gives $(x,-\delta)_v=1$ at every place, so
$\lambda_\delta([x])=0$, contradicting (S).

\emph{Step 4: the fixed field.}  Hence $2X\subseteq\mathcal W$.  Now
$X/\mathcal W$ is an elementary abelian two-group, and its group of
continuous characters is $\mathrm{Sel}_\delta$ by
Lemma~\ref{lem:ray-selmer}.  The corresponding quadratic fields
$L(\sqrt x)$ therefore generate exactly the fixed field of $\mathcal W$.
\end{proof}

\begin{corollary}[Obstruction at one split prime]\label{cor:active-ray}
Suppose that \textup{(S)} holds.  Let $A$ be an odd prime of $F$ that
splits in $L$ as $A\OO_L=A'P$, and suppose that $A'P^{-1}=(z)$ for some
$z\in L^\times$.  If every class in $\mathrm{Sel}_\delta$ is a square in
$F_A$, then $\operatorname{pr}_2(z)\in\overline E_{L,2}$.
\end{corollary}

\begin{proof}
We show that $\operatorname{Art}_L(A')$ lies in $\mathcal W$, hence is
fixed by $\sigma$, so that the Artin symbol of $(z)$ vanishes.  Since $A$
splits, $L_{A'}=F_A$, and the local-square assumption makes $A'$ split in
every quadratic field $L(\sqrt x)$ with $[x]\in\mathrm{Sel}_\delta$.  By
Proposition~\ref{prop:ray-exponent-two}, $X/\mathcal W$ is an elementary
abelian two-group whose characters are these quadratic fields.  These
characters all vanish on $\operatorname{Art}_L(A')$, and the continuous
characters of a profinite elementary abelian two-group separate its
points, so $\operatorname{Art}_L(A')\in\mathcal W$.  This element is fixed by
$\sigma$, since $\mathcal W\subseteq X^\sigma$.  As $\sigma A'=P$,
\[
 \operatorname{Art}_L((z))=
 \operatorname{Art}_L(A')-\operatorname{Art}_L(P)
 =\operatorname{Art}_L(A')-\sigma(\operatorname{Art}_L(A'))=0.
\]
The element $z$ is a unit at the dyadic primes, because its ideal is
prime to $2$.  Lemma~\ref{lem:ray-closure} for $L$ now gives
$\operatorname{pr}_2(z)\in\overline E_{L,2}$.
\end{proof}

\subsection{The local-square condition at \texorpdfstring{$A$}{A}}\label{sec:dyadic-A}

It remains to verify the hypothesis of Corollary~\ref{cor:active-ray} at
the prime $A=A_{n-1}$ of a scaled step.  Return to the elements
$\alpha=y_n$ and $\beta=\beta_n$ of a scaled step.  The identities below
use only the integral formulas \eqref{eq:scaled-family} over $F$, with
$u$ a global unit of $F$ different from $\pm1$, $v_A(a)\ge2$,
$v_B(a)=0$, $u\equiv1\pmod A$, the splitting of $A$ and $B$ in $L$, and
the dyadic conditions of Definition~\ref{def:scaled-step}(i)--(iii).  Put
\[
 g_+=a^2-y(u+1)^2,\qquad g_-=a^2-y(u-1)^2,
 \qquad g_A=-u g_+.
\]

\begin{lemma}[An explicit norm with one odd valuation]\label{lem:gA}
At a scaled step, $g_+$ and $g_-$ are nonzero elements of $\OO_F$, and
\begin{equation}\label{eq:active-norm}
 a^2u^2\delta=g_+g_-,\qquad
 N_{L/F}(\alpha+\beta)=-g_+/u,\qquad
 g_A=N_{L/F}\bigl(u(\alpha+\beta)\bigr).
\end{equation}
Moreover $v_A(g_A)=1$, $v_B(g_A)=0$, and $v_{\mathfrak p}(g_A)$ is even
at every odd prime $\mathfrak p\notin T\cup\{A\}$.
\end{lemma}

\begin{proof}
Neither $g_+$ nor $g_-$ vanishes, since otherwise $y$ would be a square in
$F$, contradicting its odd valuations at $A$ and $B$.  The identities
\eqref{eq:active-norm} follow by direct calculation from
\eqref{eq:scaled-family}.

At $A$, the element $u+1$ is a unit, since $u\equiv1\pmod A$ and $A$ is
odd, and $v_A(a)\ge2$, $v_A(y)=1$; so $v_A(g_+)=1$.  At $B$, the element
$a$ is a unit and $y$ has valuation one, so $v_B(g_+)=0$.  At every other
odd prime, $g_+-g_-=-4yu$ is a unit, so at least one of the integral
factors $g_+,g_-$ is a unit.  At such a prime outside $T$, the valuation
of $\delta$ is even, and \eqref{eq:active-norm} then makes both factor
valuations even, including at primes dividing $a$.  Since $u$ is a unit,
$g_A=-ug_+$ has the same valuations as $g_+$.  Thus the only odd
valuation of $g_A$ at a prime outside $T\cup\mathcal D$ is at $A$.
\end{proof}

\begin{lemma}[The local-square condition at $A$]\label{lem:A-square}
At a scaled step, every class of $\Delta(F,\delta)$ is a square in $F_A$.
\end{lemma}

\begin{proof}
We apply Hilbert reciprocity once, to $x$ and the element
$t_A=g_A/(-\delta)$, where $g_A$ is the norm of Lemma~\ref{lem:gA}; the
element $t_A$ is a square at the dyadic places.

\emph{Step 1: the dyadic places.}  At a dyadic place $v$, the elements $a$ and $y$ are units and $u\equiv1\pmod8$.  The
identity
\begin{equation}\label{eq:active-dyadic-square}
 \frac{g_A}{-\delta}
   =\frac{u^3}{1-y((u-1)/a)^2}
   \in1+8\OO_{F_v}\subseteq F_v^{\times2}
\end{equation}
follows from \eqref{eq:active-norm} by cancelling $g_+$.  The membership
in $1+8\OO_{F_v}$ holds because $(u-1)/a\in8\OO_{F_v}$ and
$u^3\in1+8\OO_{F_v}$.  The final inclusion follows from Newton's lemma
\cite[Theorem~7.32]{MilneANT}, since $2$ is unramified in $F$.  Hence
$[x,t_A]_v=0$ for every $x\in F^\times$ and every $v\in\mathcal D$.

\emph{Step 2: reciprocity.}  Let $[x]\in\Delta(F,\delta)$.  Since $A$ is
odd and splits in $L$, it lies outside $T\cup\mathcal D$, so $v_A(x)$ is
even.  Write $x=\pi^{2k}w$ with $\pi$ a uniformizer of $F_A$ and $w$ a
unit, so that the residue character $\chi_A$ of
Section~\ref{sec:nt-background} satisfies $\chi_A(x)=0$ exactly when the
residue of $w$ is a square.  We compute the additive Hilbert symbols
$[x,t_A]_v$ at the places $v\ne A$ outside $\mathcal D$.  At a prime
$\mathfrak q$ of $T$, the local class of $x$ is $1$ or
$[\delta]_{\mathfrak q}$.  Here $g_A$ is a norm from $L$, hence a local
norm from $L_{\mathfrak Q}=F_{\mathfrak q}(\sqrt\delta)$, so
$[\delta,g_A]_{\mathfrak q}=0$; also $[\delta,-\delta]_{\mathfrak q}=0$.
Hence $[x,t_A]_{\mathfrak q}=0$.  At every other odd prime
$\mathfrak q\ne A$, the element $x$ has even valuation by
Definition~\ref{def:selmer}, $g_A$ by Lemma~\ref{lem:gA}, and $\delta$
because $\mathfrak q$ is odd and unramified in $L=F(\sqrt\delta)$; so $x$
and $t_A$ have even valuation, and the symbol is zero
(Section~\ref{sec:nt-background}).  At the real places the symbol is
zero because $x$ is positive.  At $A$, the element $\delta$ is a square
in $F_A$ because $A$ splits in $L$, so $[x,t_A]_A=[x,-g_A]_A=[w,-g_A]_A$,
and this is $\chi_A(x)$ because $v_A(-g_A)=1$
(Section~\ref{sec:nt-background}).  Together with Step~1, Hilbert
reciprocity \cite[Chapter~V, Example~5.4]{MilneCFT} gives
$\chi_A(x)=0$.

\emph{Step 3: conclusion.}  Thus the unit $w$ has a square residue, so it
is a square in $F_A$ by Newton's lemma \cite[Theorem~7.32]{MilneANT},
because $A$ is odd.  Since $v_A(x)$ is even, $x$ is a square in $F_A$.
\end{proof}

The local analysis of Step~2, applied to $x$ and $g_A$ instead of $t_A$,
shows that every symbol $[x,g_A]_v$ with $v\ne A$ outside $\mathcal D$
vanishes; and $[x,g_A]_A=[w,g_A]_A=\chi_A(x)$ because $v_A(g_A)=1$.  Reciprocity therefore gives
\begin{equation}\label{eq:active-hilbert-row}
 \chi_A(x)=\sum_{v\in\mathcal D}[x,g_A]_v
 \qquad([x]\in\Delta(F,\delta)).
\end{equation}
By \eqref{eq:active-dyadic-square}, the right side equals
$\sum_{v\in\mathcal D}[x,-\delta]_v$, which vanishes by the sum-zero
property of $\lambda_\delta$; this is a second route to
Lemma~\ref{lem:A-square}.

The assertion of Lemma~\ref{lem:A-square} is well defined on
$\mathrm{Sel}_\delta$, because $\delta$ is a square in $F_A$ at the split
prime $A$.  The proof uses every prime of $T$ in the definition of
$\Delta(F,\delta)$, but it does not use the injectivity of
$\lambda_\delta$.

\begin{proof}[Proof of Theorem~\ref{thm:dyadic}]
The element $z_n$ satisfies $(z_n)=A_nP_n^{-1}=A'P^{-1}$, and $A=A_{n-1}$
is odd and splits in $L$.  By Lemma~\ref{lem:A-square}, every class of
$\mathrm{Sel}_\delta$ is a square in $F_A$; this is the remaining
hypothesis of Corollary~\ref{cor:active-ray}.  Under (S), that corollary
gives $\operatorname{pr}_2(z_n)\in\overline E_{L,2}
=\overline{\operatorname{pr}_2(E_n)}$, in the whole product of the dyadic
principal-unit groups, as asserted.
\end{proof}

\begin{remark}\label{rem:dyadic-role-of-S}
Lemma~\ref{lem:A-square} alone shows only that the image of
$\operatorname{Art}_L(A')$ vanishes in the quotient $X/(\mathcal W+2X)$
of exponent two.  Condition (S) supplies the inclusion
$2X\subseteq\mathcal W$, and this is what allows passage to the closure
of the global units in the whole group $U_L$.

We do not prove the corresponding statement that every class of
$\Delta(F,\delta)$ is a square at $B$, and it is not needed: the divisor
$A_nP_n^{-1}$ of $z_n$ involves only the primes $A_n$ and $P_n$, as used
in Theorem~\ref{thm:joint}.  Let $\chi_B$ be the residue character at $B$
(Section~\ref{sec:nt-background}).  The element $yg_A$ is a norm from
$L$, with $v_A(yg_A)=2$ and $v_B(yg_A)=1$ since $(y)=AB$, and applying reciprocity to it
as in the derivation of \eqref{eq:active-hilbert-row} gives, on the same group
$\Delta(F,\delta)$,
\[
 \chi_B(x)=\sum_{v\in\mathcal D}[x,y]_v.
\]
\end{remark}

\section{The odd ramified places and the proof of Theorem B}\label{sec:odd}

This section completes the proof of Theorem~\ref{thm:B}.  By
Theorem~\ref{thm:joint} (unit approximation criterion), it suffices to
verify condition \eqref{eq:JT} of Section~\ref{sec:joint}, and we verify
it with $N=n$ at every sufficiently late step: the image $\jmath_n(z_n)$
of $z_n$ in the compact group $\mathcal U_n$ of \eqref{eq:TN} must lie in
the closure of the image of the unit group $E_n$.  The group $\mathcal U_n$
is a finite product of pro-$\ell$ groups for finitely many primes $\ell$,
and a closed subgroup of such a product is the product of its
$\ell$-primary parts (Step~1 of the proof of
Proposition~\ref{prop:assembly}).  The approximation problem therefore
splits into one problem for each prime $\ell$.  Theorem~\ref{thm:dyadic}
(dyadic approximation) solves the $2$-primary problem under the Selmer
condition (S).  This section formulates the odd unit condition (O), shows
that it solves the $p$-primary problem for every odd prime $p$, and
combines the two.  We first explain the ideas.

\paragraph{The odd places give a problem in $\Z_p$-modules.}  The
odd-primary part of $\mathcal U_n$ consists of the quotients
$C_v=\OO_{(F_n)_v}^\times/W_{n,v}^\times$ at the places $v$ of $F_n$ of
odd residue characteristic $p$ with ramification index $e(v/p)=2$,
together with the multiplicative groups of the residue fields at the
dyadic places, which have odd order.  The quotients $C_v$ enter because of
Proposition~\ref{prop:orders}: from the level $b_p$ on, the completed
orders that arise in Section~\ref{sec:orders} have the form
$W_{n,v}+\pi_vp^{s}W_{n,v}$ at such a place, and the proof of
Lemma~\ref{lem:return} transports the resulting conditions to earlier
levels.  These orders are $W_{n,v}$-modules, so whether a unit of
$\OO_{(F_n)_v}$ lies in one of them depends only on its class in
$C_v$.  Each $C_v$ is a compact abelian pro-$p$ group, a quotient of
$1+\pi_v\OO_{(F_n)_v}$ by a closed subgroup (see the paragraph preceding
\eqref{eq:joint-primary}), hence a $\Z_p$-module.  So approximating
$z_n$ by global units at the places above $p$ means that one vector, the
image of $z_n$ in the product $M_{n,p}$ of the groups $C_v$, lies in the
closed $\Z_p$-submodule generated by the images of units.  Since $z_n$
has relative norm one, the relative units suffice; their images lie in
the submodule $M_{n,p}^-$ on which $\sigma_n$ acts by $-1$
(Lemma~\ref{lem:odd-membership}).  Of the structure of the $C_v$, this
reduction and the odd-primary part of the proof of Theorem~\ref{thm:B}
use only that they are compact pro-$p$ groups, together with the closure
identity \eqref{eq:relative-kernel-closure}; they do not use the
logarithm.  The $p$-adic logarithm, which for
$p\ge5$ identifies $C_v$ with the additive group of $W_{n,v}$, a free
$\Z_p$-module of rank $f_v$ (Lemma~\ref{lem:odd-layer}), linearizes the
problem further; it is used for the rank computation of
Lemma~\ref{lem:odd-minus-rank} and the finite test below.

\paragraph{The dyadic residue restriction.}  At a dyadic place $v$ of
$F_n$ the unit group of the completion is the product of a cyclic group
isomorphic to the multiplicative group $k_v^\times$ of the residue field,
which has odd order, and the pro-$2$ group $1+2\OO_{(F_n)_v}$.  For an odd
prime $p$, the $p$-primary part $\mathcal T_{n,p}$ of these residue
groups lies in the same $p$-primary factor of $\mathcal U_n$ as
$M_{n,p}$.  The $p$-primary approximation problem is therefore posed in
$\mathcal T_{n,p}\times M_{n,p}$.  The residue of $z_n$ at both dyadic
primes is $1$, so its component in $\mathcal T_{n,p}$ is zero, and only units whose
$p$-primary dyadic residue vanishes can approximate it.  Condition (O)
therefore uses the span of the images of those relative units.  This
restriction is vacuous unless $p$ divides $2^{2^{n-1}}+1$
\eqref{eq:dyadic-odd-restriction}.

\paragraph{Nakayama's lemma gives a finite sufficient test.}  The module
$M^-_{n,p}$ is free of finite rank over $\Z_p$, and the span in (O) is a
closed submodule of it.  A closed submodule of a finitely generated
$\Z_p$-module is the whole module as soon as it is the whole module
modulo $p$.  Reduction modulo $p$ is computed from the first two
coefficients of a unit in a $\pi_v$-adic expansion \eqref{eq:first-layer}.
Hence finitely many relative units whose reductions span an
$\F_p$-vector space of the right dimension prove (O) at $p$
(Proposition~\ref{prop:odd-certificate}).

\paragraph{Why $p=3$ does not occur.}  The logarithm argument requires
$v_p(\pi_v)>1/(p-1)$ for a uniformizer $\pi_v$ of a ramified quadratic
extension of an unramified extension of $\Q_p$.  Since $v_p(\pi_v)=1/2$,
this means $p\ge5$.  For $p=3$ the conclusion of
Lemma~\ref{lem:odd-layer} can fail.  The field $\Q_3(\sqrt{-3})$ is a
ramified quadratic extension of $\Q_3$ with coefficient ring $\Z_3$, and
it contains a primitive cube root of unity $\zeta_3$.  Since
$\zeta_3\notin\Z_3^\times$, the class of $\zeta_3$ in
$\OO_{\Q_3(\sqrt{-3})}^\times/\Z_3^\times$ has order three, so this
quotient is not a free $\Z_3$-module.  It is still pro-$3$, so this
failure concerns the logarithm, the ranks and the finite test, and not
the group-theoretic step of Lemma~\ref{lem:odd-membership}.  The
description of the completed orders (Proposition~\ref{prop:orders}) and
the construction assume \textup{(T1)} throughout, and we do not consider
ramification at $3$.  Hypothesis \textup{(T1)} keeps $3$, $7$ and $47$ unramified in
every $F_n$.  Hence every odd prime $p$ below a place of $F_n$ of
ramification index two satisfies $p\ge5$.

\paragraph{Standing notation.}  Throughout this section we fix $n\ge1$
and a scaled step $L=F_n$ over $F=F_{n-1}$
(Definition~\ref{def:scaled-step}) in a tower satisfying
\textup{(T1)--(T6)}, and we write $\sigma=\sigma_n$.  The letter $p$
denotes an odd rational prime.  By \textup{(T1)}, absolute ramification
indices at odd primes are at most two, so the places $v$ of $L$ above $p$
with $e(v/p)=2$ are exactly those at which $L_v/\Q_p$ is ramified.  By the preceding
paragraph, $p\ge5$ whenever such a place exists.

Section~\ref{sec:odd-local} treats a single place.
Section~\ref{sec:odd-modules} collects the places above $p$ into the
module $M_{n,p}$, and Section~\ref{sec:odd-condition} defines (O).
Section~\ref{sec:odd-approximation} shows that (O) is exactly the
$p$-primary part of the approximation problem,
Section~\ref{sec:odd-certificate} gives the finite test, and
Section~\ref{sec:odd-assembly} proves Theorem~\ref{thm:B}.

\subsection{The local unit quotient at an odd ramified place}\label{sec:odd-local}
Let $v$ be a place of $L$ above $p$ with $e(v/p)=2$.  Write $W_v$
for the coefficient ring of $L_v$ (Section~\ref{sec:nt-background}),
$k_v$ for its residue field, which is also the residue field of $L_v$,
and $f_v=[k_v:\F_p]$.  Since $e(v/p)=2$, the field $L_v$ is a ramified
quadratic extension of $\Frac(W_v)$.  As in
Section~\ref{sec:nt-background}, choose a square uniformizer $\pi_v$, so
that
\[
 \OO_{L_v}=W_v\oplus\pi_vW_v,\qquad
 \pi_v^2\in pW_v^\times.
\]
Let $\tau_v$ be the nontrivial automorphism of $L_v$ over $\Frac(W_v)$.
Since $\pi_v^2\in\Frac(W_v)$ and $\pi_v\notin\Frac(W_v)$, the
automorphism $\tau_v$ negates $\pi_v$.  This local involution need not be
the restriction of $\sigma$: the automorphism $\sigma$ can exchange $v$
with another place, or fix $v$ and act on $L_v$ as the nontrivial
automorphism of an unramified quadratic extension.

The next lemma linearizes $C_v$ by the logarithm, and it describes $C_v$
modulo $p$ through the first two coefficients of a unit.  The second
description is what makes the test of Section~\ref{sec:odd-certificate}
finite.

\begin{lemma}[Local unit quotients at odd ramified places]\label{lem:odd-layer}
For $p\ge5$ and $e(v/p)=2$, the quotient
$C_v=\OO_{L_v}^\times/W_v^\times$ is a free $\Z_p$-module of
rank $f_v$.  There is an isomorphism
\begin{equation}\label{eq:odd-log}
 \mathcal L_v:C_v\longrightarrow W_v,\qquad
 \mathcal L_v(x)=\frac{\log(x/\tau_vx)}{2\pi_v}.
\end{equation}
For a unit $x$, write
$x\equiv x_0+x_1\pi_v\pmod{\pi_v^2}$ with $x_i\in W_v$.
The coefficients are unique modulo $pW_v$, and $x_0$ is a unit.
The homomorphism
\begin{equation}\label{eq:first-layer}
 \varphi_v(x)=\overline{x_1/x_0}
      =\frac{1-x^{p^{f_v}-1}}{\pi_v}\bmod\pi_v
\end{equation}
identifies $C_v/pC_v$ with the additive group of $k_v$.
Moreover $\mathcal L_v(x)\bmod p=\varphi_v(x)$, and $\tau_v$ acts
as $-1$ on $C_v$.
\end{lemma}
\begin{proof}
We remove the Teichm\"uller factor, apply the logarithm, and read off
the $\pi_v$-coefficient; the statements modulo $p$ then follow from the
first two terms of the expansion.

\emph{Step 1: reduction to principal units.}  The Teichm\"uller
representatives of $L_v$ lie in $W_v$ (Section~\ref{sec:nt-background}).
Every unit of $\OO_{L_v}$ is the product of a Teichm\"uller
representative and an element of $1+\pi_v\OO_{L_v}$, and every unit of
$W_v$ is the product of a Teichm\"uller representative and an element of
$1+pW_v$.  Hence
$C_v=(1+\pi_v\OO_{L_v})/(1+pW_v)$.

\emph{Step 2: the logarithm.}  The logarithm and exponential series are
inverse homomorphisms between $1+\pi_v\OO_{L_v}$ and $\pi_v\OO_{L_v}$:
their convergence and formal composition identities apply because
$v_p(\pi_v)=1/2>1/(p-1)$ (Section~\ref{sec:nt-background}).  In
particular the terms tend to zero, and the nonleading terms have strictly
larger valuation.  The two series restrict to an isomorphism
$1+pW_v\simeq pW_v$.  Since $\OO_{L_v}=W_v\oplus\pi_vW_v$ and
$\pi_v^2W_v=pW_v$,
\[
 \pi_v\OO_{L_v}=pW_v\oplus\pi_vW_v .
\]
Therefore taking the $\pi_v$-coefficient of the logarithm gives an
isomorphism $C_v\simeq\pi_v\OO_{L_v}/pW_v\simeq W_v$.  The ring $W_v$ is
the ring of integers of an unramified extension of $\Q_p$ of degree
$f_v$, hence a free $\Z_p$-module of rank $f_v$, and so is $C_v$.

\emph{Step 3: the formula \eqref{eq:odd-log} and the involution.}
Removing the Teichm\"uller factor defines the logarithm of an arbitrary
unit.  This factor lies in $W_v$, so it is fixed by $\tau_v$, and $\tau_v$
commutes with the logarithm.  If $\log x=a+b\pi_v$ with $a\in pW_v$ and
$b\in W_v$, then $\log(\tau_vx)=a-b\pi_v$.  Subtracting the
$\tau_v$-conjugate logarithm gives $\log(x/\tau_vx)=2b\pi_v$, so
\eqref{eq:odd-log} returns the $\pi_v$-coefficient $b$, and
$\mathcal L_v$ is the isomorphism of Step~2.  Also
$x\tau_vx\in W_v^\times$, being the norm of $x$ to $\Frac(W_v)$.  Hence
the class of $\tau_vx$ in $C_v$ is the inverse of the class of $x$, which
proves the assertion about the involution.

\emph{Step 4: reduction modulo $p$.}  Modulo
$\pi_v^2\OO_{L_v}=p\OO_{L_v}$, the direct sum
$\OO_{L_v}=W_v\oplus\pi_vW_v$ gives the asserted unique coefficients
modulo $p$.  Since $x\equiv x_0\pmod{\pi_v}$, the coefficient $x_0$ is a
unit.  From
\[
 (x_0+x_1\pi_v)(y_0+y_1\pi_v)\equiv x_0y_0+(x_0y_1+x_1y_0)\pi_v
 \pmod{\pi_v^2}
\]
we see that multiplication adds the ratios $x_1/x_0$ modulo $p$.  So
$\varphi_v$ is a homomorphism.  It vanishes on $W_v^\times$, where
$x_1=0$, so it is defined on $C_v$.  It is surjective, since
$\varphi_v(1+b\pi_v)=\bar b$ for $b\in W_v$.  It kills $pC_v$, because
$k_v$ has exponent $p$.  Both $C_v/pC_v$ and $k_v$ have $p^{f_v}$
elements by Step~2, so the kernel of $\varphi_v$ is exactly $pC_v$.

\emph{Step 5: the other two expressions.}  The logarithm is congruent to
its linear term modulo $\pi_v^2$: for $w\in\pi_v\OO_{L_v}$ and $k\ge2$,
the term $w^k/k$ has $p$-adic valuation at least $k/2-v_p(k)\ge1$, since
$p\ge5$.  Let $\omega$ be the Teichm\"uller representative of the residue
of $x$, and write $x=\omega(1+w)$ with $w\in\pi_v\OO_{L_v}$.  Since
$\omega\equiv x_0\pmod{pW_v}$, we have
$w\equiv(x_1/x_0)\pi_v\pmod{\pi_v^2}$.  Hence the $\pi_v$-coefficient of
$\log x\equiv w$ is $x_1/x_0$ modulo $p$, which proves
$\mathcal L_v(x)\bmod p=\varphi_v(x)$.  Finally, put $q=p^{f_v}$.  Since
$\omega^{q-1}=1$ and every binomial term of degree at least two lies in
$\pi_v^2\OO_{L_v}$,
\[
 x^{q-1}=(1+w)^{q-1}\equiv1+(q-1)w\equiv1-(x_1/x_0)\pi_v
 \pmod{\pi_v^2}.
\]
This proves the second expression in \eqref{eq:first-layer}.
\end{proof}

\subsection{The modules \texorpdfstring{$M_{n,p}$}{M(n,p)} and their minus parts}\label{sec:odd-modules}
We now collect the quotients $C_v$ at all places $v$ above $p$ with
$e(v/p)=2$ into one $\Z_p$-module with an action of $\sigma$.  Condition
(O) concerns the part on which $\sigma$ acts by $-1$, because the images
of $z_n$ and of the relative units lie there.  For each odd rational
prime $p$ put
\begin{equation}\label{eq:odd-module}
 M_{n,p}=\prod_{\substack{v\mid p\\e(v/p)=2}}C_v,
 \qquad M_{n,p}^-=\ker(1+\sigma:M_{n,p}\to M_{n,p}).
\end{equation}
We write these groups additively; an empty product is zero.  All
places above $p$ with $e(v/p)=2$ occur, including places above primes already
ramified in $F_{n-1}$ and pairs of places exchanged by $\sigma$.  The
automorphism $\sigma$ acts on $M_{n,p}$: it extends to an isomorphism of
completions $L_v\to L_{\sigma v}$, which maps $\OO_{L_v}^\times$ onto
$\OO_{L_{\sigma v}}^\times$ and $W_v$ onto $W_{\sigma v}$, and hence
induces an isomorphism $C_v\to C_{\sigma v}$.

Form the reduction map
\[
 \Phi_p=\prod_v\varphi_v:M_{n,p}\longrightarrow k_{n,p},\qquad
 k_{n,p}=\bigoplus_vk_v .
\]
We transport the involution to the residue space by this map.  The
elements $\sigma(\pi_v)$ and $\pi_{\sigma v}$ are uniformizers of
$L_{\sigma v}$ whose squares lie in $pW_{\sigma v}^\times$.  So
$\sigma(\pi_v)=\psi_v\pi_{\sigma v}$ with $\psi_v\in W_{\sigma v}^\times$.
Indeed, $\psi_v$ is a unit whose square lies in $\Frac(W_{\sigma v})$.
Write $\psi_v=a+b\pi_{\sigma v}$ with $a,b\in\Frac(W_{\sigma v})$.  Then
$\psi_v^2=a^2+b^2\pi_{\sigma v}^2+2ab\pi_{\sigma v}$ forces $ab=0$.  If
$a=0$, then $\psi_v=b\pi_{\sigma v}$ has odd valuation and is not a unit.
So $b=0$ and $\psi_v=a\in W_{\sigma v}^\times$.  If $x\equiv x_0+x_1\pi_v\pmod{\pi_v^2}$, then
$\sigma x\equiv\sigma(x_0)+\sigma(x_1)\psi_v\pi_{\sigma v}
\pmod{\pi_{\sigma v}^2}$ with $\sigma(x_0),\sigma(x_1)\in W_{\sigma v}$.
Hence $\Phi_p\circ\sigma=\sigma_\Phi\circ\Phi_p$ for the $\F_p$-linear
involution $\sigma_\Phi$ of $k_{n,p}$ given by
\begin{equation}\label{eq:transported-sigma}
 (\sigma_\Phi a)_{\sigma v}
       =\bar\psi_v\,\bar\sigma(a_v),
\end{equation}
where $\bar\sigma:k_v\to k_{\sigma v}$ is the induced isomorphism of
residue fields.  This formula includes both the residue-field isomorphism
and the change of uniformizer.  Let $k_{n,p}^-$ be the minus part of
$k_{n,p}$ for this action.

\begin{lemma}\label{lem:odd-minus-rank}
The group $M_{n,p}^-$ is free over $\Z_p$ of rank
\begin{equation}\label{eq:odd-rank}
 r_{n,p}=\sum_{\substack{v\mid p\text{ in }L\\e(v/p)=2}}f_v
          -\sum_{\substack{w\mid p\text{ in }F\\e(w/p)=2}}f_w,
\end{equation}
and $\Phi_p$ induces an isomorphism
$M_{n,p}^-/pM_{n,p}^-\simeq k_{n,p}^-$.
If $p$ is unramified in $F$, then $M_{n,p}^-=M_{n,p}$ and
$k_{n,p}^-=k_{n,p}$.
\end{lemma}
\begin{proof}
We split $M_{n,p}$ into its $\pm1$ parts for $\sigma$ and identify the
plus part with the corresponding module for $F$.

\emph{Step 1: reduction modulo $p$.}  Lemma~\ref{lem:odd-layer}
identifies $M_{n,p}/pM_{n,p}$ with the residue space $k_{n,p}$ through
$\Phi_p$, equivariantly for \eqref{eq:transported-sigma}.  Since two is
invertible in $\Z_p$, the projectors $(1\pm\sigma)/2$ decompose
$M_{n,p}=M_{n,p}^+\oplus M_{n,p}^-$ and commute with reduction modulo
$p$.  Hence $\Phi_p$ induces $M_{n,p}^-/pM_{n,p}^-\simeq k_{n,p}^-$.  As a
direct summand of the free module $M_{n,p}$, the module $M_{n,p}^-$ is
free.

\emph{Step 2: the rank.}  Let $M_{F,p}$ be the analogous product over the
places $w$ of $F$ above $p$ with $e(w/p)=2$; Lemma~\ref{lem:odd-layer}
is a statement about a single local field and applies to these places
too.  Inclusion of completions gives
$\operatorname{ext}:M_{F,p}\to M_{n,p}$, which duplicates coordinates at
split places.  Products of local norms give $\nu:M_{n,p}\to M_{F,p}$.
These maps are defined for the following reasons.  At a place $w$ of $F$
ramified over $\Q$, the extension of completions $L_v/F_w$ is split or
unramified, because absolute ramification indices are at most two, and
norms of coefficient units are coefficient units.  At a newly ramified
place, $F_w$ is unramified over $\Q_p$, so the quotient $C_w$ for $F$ is
zero, and the norm of any unit is a coefficient unit.  Consequently
\[
 \nu\operatorname{ext}=2,\qquad \operatorname{ext}\nu=1+\sigma.
\]
It follows that $\operatorname{ext}$ identifies $M_{F,p}$ with
$M_{n,p}^+$.  It is injective because $\nu\operatorname{ext}=2$ and
$M_{F,p}$ is torsion-free, its image is fixed by $\sigma$, and for an
invariant element $x$ one has $x=\operatorname{ext}(\nu(x)/2)$.  By
Lemma~\ref{lem:odd-layer} the rank of $M_{F,p}$ is the second sum in
\eqref{eq:odd-rank}.  Subtracting it from the rank of $M_{n,p}$, the first
sum, proves \eqref{eq:odd-rank}.

\emph{Step 3: $p$ unramified in $F$.}  Let $v$ be a place
of $L$ above $p$ with $e(v/p)=2$, and let $w$ be the place of $F$ below
it.  Then $F_w$ is unramified over $\Q_p$, so $F_w=\Frac(W_v)$, and
$L_v$ is the ramified quadratic extension of $\Frac(W_v)$.
Hence $\sigma$ fixes $v$ and acts on $L_v$ as $\tau_v$, which acts as
$-1$ on $C_v$ by Lemma~\ref{lem:odd-layer}.  Hence $M_{n,p}^-=M_{n,p}$, and Step~1
gives $k_{n,p}^-=k_{n,p}$.
\end{proof}

\subsection{Relative units and the odd unit condition}\label{sec:odd-condition}
\paragraph{The primary decomposition of $\mathcal U_n$.}  Let
$\mathcal T_{n,p}$ be the $p$-primary part of $\prod_{v\mid2}k_v^\times$,
and put
$\mathcal T_{n,p}^-=\ker(1+\sigma:\mathcal T_{n,p}\to\mathcal T_{n,p})$.
Two is unramified in $F_0$, whose discriminant $65$ is odd, and by
\textup{(T1)} the dyadic primes are unramified in every step; so two is
unramified in $L$.  The factor $(\OO_L\otimes\Z_2)^\times$ of
\eqref{eq:TN} is $\prod_{v\mid2}\OO_{L_v}^\times$, and the unit group of
$\OO_{L_v}$ at a dyadic place $v$ is the product of the Teichm\"uller
representatives, a cyclic group isomorphic to $k_v^\times$ of odd order,
and the pro-$2$ group $1+2\OO_{L_v}$.  The factors $C_v$ at the places $v$
above $p$ with $e(v/p)=2$ are pro-$p$.  Indeed, $\OO_{L_v}^\times$ is the
product of the nonzero Teichm\"uller representatives and the pro-$p$ group
$1+\pi_v\OO_{L_v}$, where $\pi_v$ is a uniformizer, and the Teichm\"uller
representatives lie in $W_v^\times$; so $C_v$ is a quotient of
$1+\pi_v\OO_{L_v}$ by a closed subgroup.  This uses neither the logarithm
nor $p\ge5$.  We obtain the primary decomposition
\begin{equation}\label{eq:joint-primary}
 \mathcal U_n=U_L\times\prod_{p\text{ odd}}(\mathcal T_{n,p}\times M_{n,p}),
 \qquad U_L=\prod_{v\mid2}(1+2\OO_{L_v}).
\end{equation}
Only finitely many factors are nonzero at this fixed level.

\paragraph{The maps $\mu_{n,p}$ and $\operatorname{res}_{n,p}$.}  For
$z\in F_n^\times$ that is a unit at every place $v$ above $p$ with
$e(v/p)=2$, let $\mu_{n,p}(z)\in M_{n,p}$ be the tuple of the classes of $z$ in
the quotients $C_v$.  This is a homomorphism on the group of such
elements, and $\mu_{n,p}(\sigma z)=\sigma\mu_{n,p}(z)$.  On $z_n$ and on
global units it agrees with the projection of $\jmath_n$ to the factor
$M_{n,p}$ of \eqref{eq:joint-primary}.  Similarly, for $z$ that is a unit
at both dyadic places, let $\operatorname{res}_{n,p}(z)\in\mathcal T_{n,p}$
be the $p$-primary component of the tuple of its residues there.

Put $\mathscr E_{n,p}=E_n^1\otimes_\Z\Z_p$, where
$E_n^1=\ker(N_{L/F}:E_n\to E_{n-1})$ is the group of relative units.
Finite generation of units \cite[Theorem 5.1]{MilneANT} makes this a
finitely generated, hence compact, $\Z_p$-module.  Since $M_{n,p}$ and
$\mathcal T_{n,p}$ are $\Z_p$-modules, the restrictions of $\mu_{n,p}$
and $\operatorname{res}_{n,p}$ to $E_n^1$ extend uniquely to
$\Z_p$-linear maps on $\mathscr E_{n,p}$, which are continuous and which
we denote by the same letters:
\[
 \mu_{n,p}:\mathscr E_{n,p}\to M_{n,p}^-,\qquad
 \operatorname{res}_{n,p}:\mathscr E_{n,p}\to\mathcal T_{n,p}^-.
\]
Their images lie in the minus parts because $\sigma\eta=\eta^{-1}$ for a
relative unit $\eta$.  Define
\begin{equation}\label{eq:odd-images}
 \begin{gathered}
 U_{n,p}^{\rm all}=\mu_{n,p}(\mathscr E_{n,p}),\qquad
 U_{n,p}^0=\mu_{n,p}(\ker\operatorname{res}_{n,p}),\\
 E_{n,p}^{1,0}=\{\eta\in E_n^1:\operatorname{res}_{n,p}(\eta)=0\},\qquad
 z_{n,p}=\mu_{n,p}(z_n)\in M_{n,p}^-.
 \end{gathered}
\end{equation}
The element $z_n$ is a unit at every ramified place, because
$(z_n)=A_nP_n^{-1}$ and $47$ is unramified in $F_n$ by \textup{(T1)}.  Its
relative norm is one, since $y_n$ and $\beta_n$ both have relative norm
$y_{n-1}$ by \textup{(T2)}; so $z_{n,p}$ lies in $M_{n,p}^-$.

\paragraph{The span $U^0_{n,p}$ as a closure.}  The groups
$U_{n,p}^{\rm all}$ and $U_{n,p}^0$ are compact, hence closed.  The image
of $E_n^1$ is dense in $\mathscr E_{n,p}$, since $\Z$ is dense in $\Z_p$.
Since the group $\mathcal T^-_{n,p}$ is finite and discrete, the kernel of
$\operatorname{res}_{n,p}$ is open, and the intersection of the image of
$E_n^1$ with this kernel, which is the image of $E_{n,p}^{1,0}$, is dense
in that kernel.  By continuity and compactness, it follows that
\begin{equation}\label{eq:relative-kernel-closure}
 U_{n,p}^0=\overline{\mu_{n,p}(E_{n,p}^{1,0})}.
\end{equation}
Thus $U_{n,p}^0$ is the closed $\Z_p$-submodule of $M_{n,p}^-$
generated by the images of the relative units whose dyadic residues have
trivial $p$-primary component.

\paragraph{The odd unit condition.}  The \emph{odd unit condition} is
\begin{equation}\label{eq:O}
 \tag{O} z_{n,p}\in U_{n,p}^0
       \quad\text{for every odd rational }p\text{ ramified in }F_n.
\end{equation}
Section~\ref{sec:odd-approximation} shows that \eqref{eq:O} at $p$ is
equivalent to the $p$-primary part of the approximation problem.

\paragraph{When the dyadic restriction is vacuous.}  The two dyadic
primes of $F_0=\Q(\sqrt{65})$ have residue field $\F_2$, since
$65\equiv1\pmod 8$, and by \textup{(T1)} they are inert in every step.
So the two dyadic residue extensions at this step are $\F_{q^2}/\F_q$
with $q=2^{2^{n-1}}$, and $\sigma$ acts on them as the Frobenius
$x\mapsto x^q$.  The kernel of $1+\sigma$ on $\F_{q^2}^\times$ is
therefore the group of elements of norm one, which is cyclic of order
$q+1$.  Thus
\begin{equation}\label{eq:dyadic-odd-restriction}
 p\nmid 2^{2^{n-1}}+1
 \quad\Longrightarrow\quad
 \mathcal T_{n,p}^-=0,\quad E_{n,p}^{1,0}=E_n^1,
 \quad U_{n,p}^0=U_{n,p}^{\rm all}.
\end{equation}
Even in this case, the $M_{n,p}$-components at places above primes
already ramified in $F_{n-1}$ remain part of (O).

\begin{remark}[The residue obstruction]\label{rem:odd-residue-obstruction}
The difference between $U^0_{n,p}$ and $U^{\rm all}_{n,p}$ is measured
by a finite group.  Suppressing $n$ and $p$, there is an exact sequence
\begin{equation}\label{eq:odd-residue-obstruction}
 0\longrightarrow U^0\longrightarrow U^{\rm all}
 \longrightarrow \operatorname{res}(\mathscr E)/\operatorname{res}(\ker\mu)
 \longrightarrow0,
 \qquad \mu(e)\longmapsto[\operatorname{res}(e)].
\end{equation}
The map is well defined, because two lifts $e$ of the same element of
$U^{\rm all}$ differ by an element of $\ker\mu$.  It is onto, since the class of
$\operatorname{res}(e)$ is the image of $\mu(e)$ for every
$e\in\mathscr E$, and it vanishes on $U^0$.  Conversely, if $\operatorname{res}(e)=
\operatorname{res}(e')$ with $e'\in\ker\mu$, then $e-e'$ lies in
$\ker\operatorname{res}$ and has the same image $\mu(e)$; so the kernel
consists precisely of the images under $\mu$ whose residue can be
removed without changing the image, that is, of $U^0$.  The target is
finite because $\mathcal T^-_{n,p}$ is.  Thus (O) first requires
$z_{n,p}\in U^{\rm all}$, and then the vanishing of this finite
obstruction for a lift of $z_{n,p}$.
\end{remark}

\subsection{The odd unit condition as an approximation condition}\label{sec:odd-approximation}
The next lemma shows that (O) at $p$ says exactly that the $p$-primary
component $(0,z_{n,p})$ of $\jmath_n(z_n)$ can be approximated by global
units, simultaneously at the dyadic residues and at all places $v$ above
$p$ with $e(v/p)=2$.  It uses arbitrary units, not only relative
units, and shows that relative units with trivial $p$-primary dyadic
residue suffice.

\begin{lemma}[Membership in the closure of the unit image]\label{lem:odd-membership}
Let $\Gamma_{n,p}$ be the closure of the image of all global units $E_n$
in $\mathcal T_{n,p}\times M_{n,p}$.  For $g\in M_{n,p}^-$,
\[
 (0,g)\in\Gamma_{n,p}\quad\Longleftrightarrow\quad g\in U_{n,p}^0.
\]
\end{lemma}
\begin{proof}
We reduce to units with trivial residue component and then use
$\eta\mapsto\eta/\sigma\eta$ to pass to relative units.

\emph{Step 1: units with trivial residue.}  Let
$E'=\ker(E_n\to\mathcal T_{n,p})$, and let $\Gamma'$ be the closure of its
image in $M_{n,p}$.  The image of $E'$ lies in $\{0\}\times M_{n,p}$.
Conversely, suppose that $(0,g)$ is a limit of images of units $e_i$.
Since the residue factor $\mathcal T_{n,p}$ is finite and discrete, the
residue components of the $e_i$ are eventually $0$, so eventually
$e_i\in E'$, and $g\in\Gamma'$.  Hence
$\Gamma_{n,p}\cap(\{0\}\times M_{n,p})=\{0\}\times\Gamma'$.

\emph{Step 2: passage to relative units.}  The group $\Gamma'$ is stable
under $\sigma$.  For $\eta\in E'$, the unit $\eta/\sigma\eta$ has
relative norm one and trivial residue component, so it belongs to
$E_{n,p}^{1,0}$; its image is $(1-\sigma)\mu_{n,p}(\eta)$.  Taking
closures gives $(1-\sigma)\Gamma'\subseteq U_{n,p}^0$.  Now let
$g\in\Gamma'\cap M_{n,p}^-$.  Then $(1-\sigma)g=2g\in U_{n,p}^0$.  A
closed subgroup of a compact abelian pro-$p$ group is a
$\Z_p$-submodule, so it is closed under division by two.  Hence $g\in
U_{n,p}^0$, and $\Gamma'\cap M_{n,p}^-\subseteq U_{n,p}^0$.  With Step~1
this proves the implication from left to right.

\emph{Step 3: the converse.}  The units in $E_{n,p}^{1,0}$ lie in $E'$,
and their images in $\mathcal T_{n,p}\times M_{n,p}$ are
$(0,\mu_{n,p}(\eta))$.  So the reverse inclusion
$U_{n,p}^0\subseteq\Gamma'$ follows directly from
\eqref{eq:relative-kernel-closure}.
\end{proof}

As explained in Section~\ref{sec:joint}, approximation at each place
separately would not suffice; this is why (O) is stated in the whole
product $\mathcal T_{n,p}\times M_{n,p}$.

\subsection{A finite sufficient test}\label{sec:odd-certificate}
Condition (O) asks for membership of one vector in a closed submodule.
The following proposition gives a condition under which every vector of
$M_{n,p}^-$ is a member, and that condition can be checked modulo $p$ with
finitely many units.

\begin{proposition}[A finite sufficient test for (O)]\label{prop:odd-certificate}
For an odd ramified rational prime $p$,
\[
 U_{n,p}^0=M_{n,p}^-
 \quad\Longleftrightarrow\quad
 \operatorname{span}_{\F_p}
 \Phi_p\mu_{n,p}(E_{n,p}^{1,0})=k_{n,p}^-.
\]
In particular, any finite collection
$\eta_1,\ldots,\eta_m\in E_{n,p}^{1,0}$ whose images under
$\Phi_p\mu_{n,p}$ have rank $r_{n,p}$ proves (O) at $p$, and also the
analogue of (O) with $z_{n,p}$ replaced by any element of $M_{n,p}^-$.
\end{proposition}
\begin{proof}
We translate the right side into a statement modulo $p$ and apply
Nakayama's lemma in its form for compact $\Z_p$-modules.  Write
$M=M_{n,p}^-$ and $U=U_{n,p}^0$.

\emph{Step 1: the condition modulo $p$.}  By
Lemma~\ref{lem:odd-minus-rank}, $\Phi_p$ induces an isomorphism
$M/pM\simeq k_{n,p}^-$.  By \eqref{eq:relative-kernel-closure}, which
rests on density and the finiteness of $\mathcal T^-_{n,p}$, $U$ is the
closure of $\mu_{n,p}(E_{n,p}^{1,0})$.  Since $pM$ is open in $M$, the
sets $U$ and $\mu_{n,p}(E_{n,p}^{1,0})$ have the same image in $M/pM$.
This image is a subgroup, hence equal to its $\F_p$-span, and it lies in
$k_{n,p}^-$ after applying $\Phi_p$.  So the right side is the assertion
$U+pM=M$.

\emph{Step 2: Nakayama's lemma.}  Suppose $M=U+pM$.  Given $g\in M$,
write $g=g_0+ph_1$ with $g_0\in U$ and $h_1\in M$, then
$h_1=g_1+ph_2$, and so on.  Repeated substitution gives
\[
 g=g_0+pg_1+\cdots+p^{a-1}g_{a-1}+p^ah_a,
 \qquad g_i\in U,\quad h_a\in M,
\]
for every $a\ge1$.  The residual term $p^ah_a$ tends to zero, and $U$ is
closed, so $g\in U$.  The converse is immediate.

\emph{Step 3: the finite collection.}  By Lemma~\ref{lem:odd-minus-rank},
$k_{n,p}^-\simeq M/pM$ has dimension $r_{n,p}$ over $\F_p$.  If the
images of $\eta_1,\ldots,\eta_m$, which lie in $k_{n,p}^-$, have rank
$r_{n,p}$, they span $k_{n,p}^-$, so $U=M$ by Steps~1 and~2, and every element of $M$,
in particular $z_{n,p}$, lies in $U$.
\end{proof}

\begin{remark}[Computing the test]\label{rem:odd-certificate-computation}
The units in this test can be produced by a finite computation.  Start
with any global units $\gamma_1,\ldots,\gamma_m$ and put
$\eta_i=\gamma_i/\sigma\gamma_i\in E_n^1$.  Compute their dyadic
$p$-residues and the integer kernel lattice
\[
 \mathcal N_p=\left\{(a_i)\in\Z^m:
             \sum_i a_i\operatorname{res}_{n,p}(\eta_i)=0\right\}.
\]
It has finite index in $\Z^m$ because the residue group is finite.  For an
integer basis $a^{(1)},\ldots,a^{(m)}$ of this lattice, put
$\xi_j=\prod_i\eta_i^{a_i^{(j)}}$.  Then $\xi_j\in E_{n,p}^{1,0}$.
Compute $\Phi_p\mu_{n,p}(\xi_j)$ using \eqref{eq:first-layer}, which
requires only reduction modulo $\pi_v^2$ at every place $v$ above $p$
with $e(v/p)=2$.  After choosing $\F_p$-bases for the residue fields and
for the subspace $k_{n,p}^-$ with the action
\eqref{eq:transported-sigma}, form one matrix from these images at all
places above $p$ together.  By Proposition~\ref{prop:odd-certificate}, rank
$r_{n,p}$ certifies the equality $U_{n,p}^0=M_{n,p}^-$, and therefore (O)
at $p$.  When
\eqref{eq:dyadic-odd-restriction} applies, $\mathcal N_p=\Z^m$.
\end{remark}

\begin{remark}[Isolated ramified primes]\label{rem:odd-isolated}
For an isolated ramified prime (Definition~\ref{def:isolated}) the test is
scalar.  Suppose that $p$ is unramified in $F$, that exactly one prime
$\mathfrak p\mid p$ ramifies in $L/F$, that $F_{\mathfrak p}=\Q_p$, and
that $v_{\mathfrak p}(\delta)=2m+1$.  Then the place $v$ of $L$ above
$\mathfrak p$ is the only place of $L$ above $p$ with ramification index
two over $p$, it has $f_v=1$, and $F$ has no such place above $p$; so
$r_{n,p}=1$.  Moreover
$\pi=\sqrt\delta/p^m$ is a local uniformizer.  If
$\eta=b_0+b_1\sqrt\delta\in E_{n,p}^{1,0}$ with $b_0,b_1\in F$, then
$b_0$ and $p^mb_1$ are $\mathfrak p$-integral and $b_0$ is a unit.  This
follows by taking the half-trace and half-difference of $\eta$ and
$\sigma\eta$, and using integrality at the ramified place.  Consequently
\[
  \varphi_v(\eta)=p^mb_1/b_0\bmod\mathfrak p.
\]
By Proposition~\ref{prop:odd-certificate}, one such unit with nonzero
value proves (O) at $p$.
\end{remark}

\begin{remark}[What a deficient rank shows]\label{rem:odd-deficient}
A deficient matrix from selected units gives no negative conclusion about
the full unit image.  Even if those units generate $E_{n,p}^{1,0}$ modulo
$p$-th powers, deficient rank proves only $U_{n,p}^0\ne M_{n,p}^-$, and
the element $z_{n,p}$ may still belong to that proper submodule.  A unit
image of finite index does not suffice either: $p\Z_p\subset\Z_p$ has
finite index but does not contain $1$.
\end{remark}

\subsection{Proof of Theorem~B}\label{sec:odd-assembly}
We now combine the two halves of the approximation problem.  The
$2$-primary half is Theorem~\ref{thm:dyadic}, which uses (S); the
odd-primary half is condition (O), which Lemma~\ref{lem:odd-membership}
identifies with approximation of the odd-primary components.  The
following proposition puts them together at one step.

\begin{proposition}\label{prop:assembly}
Let $F_n/F_{n-1}$ be a scaled step in the sense of
Definition~\ref{def:scaled-step}, in a tower satisfying
\textup{(T1)--(T6)}.  If the step $F_n/F_{n-1}$ satisfies both
\textup{(S)} and \textup{(O)}, then
$\jmath_n(z_n)\in\overline{\jmath_n(E_n)}\subseteq\mathcal U_n$.
\end{proposition}
\begin{proof}
We check the membership one primary component at a time, using the
decomposition \eqref{eq:joint-primary}.

\emph{Step 1: primary components of a closed subgroup.}  A closed
subgroup of a profinite abelian group is stable under all primary
projections: in each finite quotient, a primary projection is
multiplication by an integer chosen by the Chinese remainder theorem, and
passage to the limit uses closedness.  Thus $\overline{\jmath_n(E_n)}$ is
the product of its primary subgroups, and each projection of
$\overline{\jmath_n(E_n)}$ is its intersection with the corresponding
primary factor.  Since $E_n$ is mapped into a compact group, the
projection of $\overline{\jmath_n(E_n)}$ to a factor is the closure of the
projection of $\jmath_n(E_n)$.  It therefore suffices to show that each
primary component of $\jmath_n(z_n)$ lies in the corresponding projection
of $\overline{\jmath_n(E_n)}$.

\emph{Step 2: the $2$-primary component.}  Theorem~\ref{thm:dyadic}
(dyadic approximation) uses (S) and supplies the $U_L$ component of
$\jmath_n(z_n)$ in \eqref{eq:joint-primary}, including its two-primary
torsion.

\emph{Step 3: the dyadic residues of $z_n$.}  The element $z_n$ has
residue $1$ at both dyadic primes.  By Definition~\ref{def:scaled-step}(i)
and~(iii), $u\equiv1\pmod2$ and $a$ is a dyadic unit, and $y$ is a dyadic
unit because $(y)=AB$.  Reduce modulo the dyadic primes of $L$.  In
\eqref{eq:scaled-family}, $t=(y(1-u^2)-a^2)/(au)$ with $1-u^2\equiv0$,
so $\bar t=-\bar a=\bar a$.  Since $\alpha+\sigma\alpha=t$,
$\alpha\sigma\alpha=y$ and $\beta=a+u\alpha$, we obtain
\[
 \bar t=\bar a,\qquad
 \bar\beta=\bar a+\bar\alpha=\overline{\sigma\alpha},
 \qquad \bar z_n=\overline{\alpha\beta/y}=1.
\]
Hence the $\mathcal T_{n,p}$-component of $\jmath_n(z_n)$ is zero for
every odd $p$.

\emph{Step 4: the odd components.}  For every odd ramified $p$, the
$p$-primary component of $\jmath_n(z_n)$ is therefore $(0,z_{n,p})$.
Condition (O) and Lemma~\ref{lem:odd-membership} place it in
$\Gamma_{n,p}$, which is the projection of $\overline{\jmath_n(E_n)}$ to
that primary factor.  If $p$ is unramified in $L$, the product $M_{n,p}$ is
zero, and the $p$-primary component of $\jmath_n(z_n)$ is already zero.

\emph{Step 5: conclusion.}  By Step~1, the primary components combine to
give the asserted membership in $\overline{\jmath_n(E_n)}$.
\end{proof}

\begin{proof}[Proof of Theorem~\ref{thm:B}]
Both (S) and (O) refer to the same step $F_n/F_{n-1}$.  By hypothesis
there is $n_0$ such that for every $n>n_0$ the step $F_n/F_{n-1}$ is a
scaled step satisfying (S) and (O).  For each such $n$,
Proposition~\ref{prop:assembly} gives
$\jmath_n(z_n)\in\overline{\jmath_n(E_n)}$, which is condition
\textup{($\mathrm{JT}_A$)} with $N=n$.  Theorem~\ref{thm:joint} then
proves that $D$ is hereditarily atomic.  The strict principal-ideal chain
of Theorem~\ref{thm:ambient} already proves the failure of ACCP under the
hypotheses \textup{(T1)--(T6)}.
\end{proof}

\section{The remaining existence problem}\label{sec:selection}
Theorem~\ref{thm:B} applies to one infinite tower of scaled steps that
satisfies \textup{(T1)--(T6)} and in which every sufficiently late step
satisfies both the Selmer condition (S) and the odd unit condition (O).
Theorem~\ref{thm:construction} supplies infinite towers of scaled steps
satisfying \textup{(T1)--(T6)}, but its construction imposes neither (S)
nor (O).  Each of the two conditions is a statement about a single step;
Sections~\ref{sec:dyadic} and~\ref{sec:odd} define them and give
sufficient tests.  What
remains is to make infinitely many such choices compatibly.  This section
states that problem precisely and explains what it involves.

Section~\ref{sec:selection-hypothesis} formulates a one-step extension
statement, Hypothesis~\ref{hyp:selection} (which this paper does not
prove): every finite extension of a
fixed initial tower that is built by the construction schedule and whose
added steps satisfy (S) and (O) has a further scaled step, for the
prescription data of the schedule, that satisfies both.
Proposition~\ref{prop:selection} shows that this hypothesis yields a
hereditarily atomic domain without ACCP.
Section~\ref{sec:selection-field} explains why (S) and (O) are conditions
on the field $F_n$ itself, which the local prescriptions of
Section~\ref{sec:construction} do not determine.
Section~\ref{sec:selection-precision} compares (O) with the coefficient
precision that the construction can prescribe
(Appendix~\ref{app:precision}).  Lemma~\ref{lem:precision-index} shows
that prescribed precision gives (O) at $p$ once the unit index exponent
$e_{n,p}$ of $F_n$ is small enough, and an example of lattices shows that precision without
such an index bound does not suffice.

\subsection{Extending finite towers one step at a time}
\label{sec:selection-hypothesis}
Fix an initial tower $F_0\subset\cdots\subset F_{n_0}$ obtained by the
construction of Section~\ref{sec:construction}.  Together with the fields
we retain the state of that construction at level $n_0$: its list of
characteristics, the fibres whose treatment is complete, and the current
state of the finite procedure for separating residue offsets.  Fix also
the construction schedule (Definition~\ref{def:schedule}) used in the
proof of Theorem~\ref{thm:construction}, in which the integers $H_n$
(for example $H_n=n+5$) and all other choices of that proof (the order of
treatment within a fibre and the late levels $n^*$) are made by one fixed
rule depending only on the tower built so far.  After each step this schedule
adds to its list all newly ramified characteristics and the least odd
characteristic other than $47,7$ not yet listed, and it specifies the prescription data for
the next step.  The schedule is thus a fixed rule, but the prescription
data it assigns at a later level depend on the fields already chosen.

\begin{definition}\label{def:successful-continuation}
A \emph{successful finite continuation} of the initial tower is a finite
sequence of scaled steps $F_{n_0}\subset F_{n_0+1}\subset\cdots\subset
F_m$ (Definition~\ref{def:scaled-step}) such that each step
$F_n/F_{n-1}$ with $n_0<n\le m$ is a scaled step for the prescription
data that the construction schedule assigns at level $n$, and satisfies both
(S) and (O).  The initial tower itself is a successful continuation of
length zero; no condition (S) or (O) is imposed at its levels.  A
\emph{successful infinite sequence} of scaled steps is an infinite
sequence $F_{n_0}\subset F_{n_0+1}\subset\cdots$ of scaled steps each of
whose finite initial segments is a successful finite continuation.
\end{definition}

In this definition all local conditions, sets of ramified primes and
unit groups are those of the fields $F_n$ of the sequence.

\begin{hypothesis}[Simultaneous extension]\label{hyp:selection}
Let $F_{n_0}\subset\cdots\subset F_m$, with $m\ge n_0$, be a successful
finite continuation of the initial tower.  Then there is a next step
$F_{m+1}/F_m$ that is a scaled step in the sense of
Definition~\ref{def:scaled-step} for $H=H_{m+1}$ and the next
prescription data of the construction schedule (its inert set, forbidden
characteristics and residue prescriptions), and that satisfies both (S)
and (O).  The auxiliary
prime $R$, the element $c$ and the positive exponent $k$ of this step may
depend on the whole continuation and on the next prescription data.
\end{hypothesis}

The hypothesis asks for one scaled step that satisfies (S) and (O)
together.  It does not ask that every preliminary choice of $c$ admit a
suitable exponent $k$.  The case $m=n_0$ is included: the hypothesis
asserts that the initial tower itself has a successful extension of
length one.  So the hypothesis is not satisfied vacuously by the absence
of successful continuations of positive length.

\begin{proposition}[The extension hypothesis suffices]\label{prop:selection}
If some initial tower and the construction schedule above satisfy
Hypothesis~\ref{hyp:selection}, then there exists a hereditarily
atomic domain that does not satisfy ACCP.
\end{proposition}
\begin{proof}
We build a successful infinite sequence of scaled steps by repeated use
of the hypothesis, and apply Theorem~\ref{thm:B} to the resulting tower.

\emph{Step 1: an infinite tower.}  The initial tower is a successful
continuation of length zero.  If $F_{n_0}\subset\cdots\subset F_m$ is a
successful continuation, Hypothesis~\ref{hyp:selection} gives a scaled
step $F_{m+1}/F_m$, a scaled step for the next prescription data of the
schedule, that satisfies (S) and (O); then
$F_{n_0}\subset\cdots\subset F_{m+1}$ is again a successful
continuation.  Choosing such a step at every level, by countable
dependent choice, gives an infinite tower
$F_0\subset F_1\subset\cdots$ every step of which is a scaled step; the
steps of the initial tower are scaled steps by
Remark~\ref{rem:scaled-step-outputs}, since the construction of
Theorem~\ref{thm:construction} produces them by
Proposition~\ref{prop:scaled-step}.

\emph{Step 2: hypotheses \textup{(T1)--(T6)}.}  Every step after
level $n_0$ is a scaled step for the prescription data of the schedule,
so every local prescription of the schedule is retained.  The argument of
Section~\ref{sec:offsets} that these residue prescriptions satisfy
condition (c) of Proposition~\ref{prop:scaled-step} uses only that the
prescriptions of the earlier steps hold, so it applies to this tower.  The
termination argument and the argument that each listed characteristic is
treated after finitely many steps, in the proof of
Theorem~\ref{thm:construction}, therefore apply to this tower: each
listed fibre has only finitely many predecessors in the list, and each
fibre procedure terminates.  Hence
Definition~\ref{def:scaled-step}(iv)--(vi) and the proof of
Theorem~\ref{thm:construction} give \textup{(T1)--(T6)} for the tower;
that proof uses of each step only conditions (iv)--(vi) and the
schedule: the field is totally real, every prescription holds, and newly
ramified characteristics are listed after each step.

\emph{Step 3: conclusion.}  Conditions (S) and (O) hold at every level
above $n_0$ by construction.  The tower therefore satisfies the hypotheses
of Theorem~\ref{thm:B}, and the ring of integers of $\bigcup_nF_n$ is
hereditarily atomic and does not satisfy ACCP.
\end{proof}

Hypothesis~\ref{hyp:selection} is stronger than what Theorem~\ref{thm:B}
needs.  One successful infinite sequence of scaled steps suffices, even
if some successful finite continuations have no successful extension.
Two weaker statements do not suffice by themselves.  First, the existence
of one step satisfying (S) and of another step satisfying (O) does not
give a single step satisfying both.  Second, the existence of successful
finite continuations of every finite length does not give a successful
infinite sequence.  The successful continuations form a tree, ordered by
extension, and this tree need not be finitely branching.  For a
finitely branching tree, K\"onig's lemma would turn nodes of every finite
length into an infinite path; without that property the conclusion can
fail.  For example, the tree of finite strictly decreasing sequences of
positive integers, ordered by extension, has nodes of every finite length
but no infinite path.  The example is purely combinatorial.

\subsection{Conditions (S) and (O) concern the field \texorpdfstring{$F_n$}{F\_n} itself}
\label{sec:selection-field}
The construction of Section~\ref{sec:construction} controls finitely many
local prescriptions at each step.  Conditions (S) and (O) depend in
addition on all the primes that ramify in the step and on the unit group
of the new field.  We describe both dependencies.

Consider a proposed step $L/F$.  Condition (S) is the injectivity of the
map $\lambda_\delta$ of Section~\ref{sec:dyadic}.  Its domain
$\mathrm{Sel}_\delta$ is defined using every odd prime ramified in
$L/F$, including the primes ramified in $L/F$ other than the prescribed
ones, and all the global square classes permitted by those primes.  The
construction allows such additional ramified primes at every step and
places no bound on their number.  The detecting primes for $\mathcal
V(F)$ (Definition~\ref{def:detecting}) give the conclusions of
Proposition~\ref{prop:scaled-step} about ideal classes and norms.  They do
not replace the computation of the Selmer condition, which involves the
whole set $T$ of primes ramified in $L/F$.

At the same step, (O) at $p$ asks that $z_{n,p}$ lie in the closed
$\Z_p$-submodule $U_{n,p}^0$ of $M_{n,p}^-$ (Section~\ref{sec:odd}), and
it is imposed at every odd rational prime $p$ ramified in $L$.  These
primes include the primes already ramified in $F$ as well as the newly
ramified ones, and for each $p$ the module $M_{n,p}$ involves all the
ramified places of $L$ above $p$.  The module $U^0_{n,p}$ is the closure
of the image of the relative units whose dyadic residues have trivial
$p$-primary component \eqref{eq:relative-kernel-closure}, so the dyadic
residue restriction is part of (O).  The finite
test of Proposition~\ref{prop:odd-certificate} proves the stronger
statement $U_{n,p}^0=M_{n,p}^-$, whereas (O) at $p$ asks
only that the single element $z_{n,p}$ lie in $U^0_{n,p}$.  So a step at
which the test fails may still satisfy (O).

Thus (S) and (O) are conditions on the chosen field $F_n$ and its units,
at one level at a time.  Sections~\ref{sec:dyadic} and~\ref{sec:odd} give
their definitions and sufficient tests.  Verifying (S) and (O) at finitely
many levels proves neither Hypothesis~\ref{hyp:selection}
nor the existence of a successful infinite sequence of scaled steps.

The quantifiers over the odd primes also matter.  Theorem~\ref{thm:B}
needs one level after which every step satisfies (O) at every ramified
odd prime.  Suppose instead that for each fixed $p$ the condition (O) at
$p$ holds at all steps after a level depending on $p$.  This does not
supply a common level: a new odd prime at which (O) fails at each level
is compatible with eventual success at each fixed prime.

\subsection{Coefficient precision and the unit index}
\label{sec:selection-precision}
Appendix~\ref{app:precision} shows that the construction can prescribe,
at finitely many odd primes in each step, how closely $z_n$ is
approximated by elements of the local coefficient rings
(Proposition~\ref{prop:finite-precision}).  This subsection measures what
such a prescription gives towards (O).  Throughout, the tower consists of
scaled steps and satisfies \textup{(T1)--(T6)}.  By (T1), the places
$v\mid p$ of $F_n$ that are ramified over $\Q$, for an odd prime $p$, are
exactly the places with $e(v/p)=2$, which index the product $M_{n,p}$.

\begin{definition}\label{def:coefficient-precision}
Let $v$ be a finite place of a number field $\mathcal E$, with completion
$\mathcal E_v$, coefficient ring $W_v$ (the ring of integers of the
maximal unramified subfield of $\mathcal E_v$;
Section~\ref{sec:nt-background}) and maximal ideal $\mathfrak m_v$ of
$\OO_{\mathcal E_v}$.  A unit $z\in\OO_{\mathcal E_v}^\times$ has
\emph{coefficient precision at least $d$} at $v$ if
$z\in W_v^\times(1+\mathfrak m_v^d)$.
\end{definition}

Since the Teichm\"uller representatives lie in $W_v$
(Section~\ref{sec:nt-background}), every unit has
coefficient precision at least $1$; larger $d$ means that $z$ is closer
to an element of $W_v^\times$.  The second ingredient measures how far the
unit image $U^0_{n,p}$ is from all of $M^-_{n,p}$.  For a prime $p\ge5$
ramified in $F_n$, let
\begin{equation}\label{eq:unit-index-exponent}
 e_{n,p}=\min\{e\ge0:p^eM_{n,p}^-\subseteq U_{n,p}^0\},
\end{equation}
with value $\infty$ if the set is empty.  We call $e_{n,p}$ the
\emph{unit index exponent} at $p$.  When it is finite, $p^{e_{n,p}}$ is
the exponent of the finite quotient $M_{n,p}^-/U_{n,p}^0$, not its order.
It is finite exactly when $U_{n,p}^0$ has full $\Z_p$-rank in
$M_{n,p}^-$, that is, when $U_{n,p}^0\otimes\Q_p=M_{n,p}^-\otimes\Q_p$.
Indeed $M_{n,p}^-$ is free of finite rank
(Lemma~\ref{lem:odd-minus-rank}), so a submodule of full rank has finite
index and the quotient is killed by a power of $p$; conversely a
submodule containing $p^eM_{n,p}^-$ has the same span over $\Q_p$.  The
definition uses $U^0_{n,p}$, the image of the kernel of the dyadic
residue map $\operatorname{res}_{n,p}$, and not the larger module
$U^{\rm all}_{n,p}$.

For $z\in F_n^\times$ that is a unit at the places $v\mid p$ with
$e(v/p)=2$, we write $\mu_{n,p}(z)\in M_{n,p}$ for its image, as in
Section~\ref{sec:odd}.  At such a place $v$ we use the square uniformizer
$\pi_v$ of Section~\ref{sec:odd}, so that
$\OO_{(F_n)_v}=W_v\oplus\pi_vW_v$ and $\pi_v^2\in pW_v^\times$.

\begin{lemma}[Coefficient precision and the unit index]\label{lem:precision-index}
Let $d\ge1$ and let $z\in F_n^\times$ have relative norm one
and be a unit at all places $v\mid p$ with $e(v/p)=2$, where $p\ge5$.
Write $\hat z=\mu_{n,p}(z)$ and
$\mathfrak m_v=\pi_v\OO_{(F_n)_v}$.  Then
\begin{equation}\label{eq:precision-native}
 z\in W_v^\times(1+\mathfrak m_v^d) \text{ for every }v\mid p\text{ with }e(v/p)=2
 \quad\Longleftrightarrow\quad
 \hat z\in p^{\lfloor d/2\rfloor}M_{n,p}^-.
\end{equation}
If these equivalent conditions hold and moreover
$e_{n,p}\le\lfloor d/2\rfloor$, then $\hat z\in U_{n,p}^0$.
\end{lemma}
\begin{proof}
We compute the image of $W_v^\times(1+\mathfrak m_v^d)$ in each quotient
$C_v$ by the logarithm, and then pass to the product $M_{n,p}$ and its
minus part.  Put $a=\lfloor d/2\rfloor$.

\emph{Step 1: one place.}  Let $v\mid p$ with $e(v/p)=2$.  The
logarithm isomorphism of Lemma~\ref{lem:odd-layer} between
$1+\pi_v\OO_{(F_n)_v}$ and $\pi_v\OO_{(F_n)_v}$ restricts to an
isomorphism $1+\mathfrak m_v^d\simeq\mathfrak m_v^d$, because both the
logarithm and the exponential preserve the valuation of the leading term.
For $d=2a$ and $d=2a+1$, respectively, the additive ideals are
\[
 \mathfrak m_v^{2a}=p^aW_v\oplus p^a\pi_vW_v,
 \qquad
 \mathfrak m_v^{2a+1}=p^{a+1}W_v\oplus p^a\pi_vW_v,
\]
since $\pi_v^2\in pW_v^\times$.  In both cases the $\pi_v$-coefficients
are exactly the elements of $p^aW_v$.  The isomorphism
$\mathcal L_v:C_v\to W_v$ of \eqref{eq:odd-log} takes the class of an
element of $1+\pi_v\OO_{(F_n)_v}$ to the $\pi_v$-coefficient of its
logarithm.  Hence the image of $W_v^\times(1+\mathfrak m_v^d)$ in $C_v$ is
$\mathcal L_v^{-1}(p^aW_v)=p^aC_v$.  Since $W_v^\times$ is the kernel of
$\OO_{(F_n)_v}^\times\to C_v$, a unit at $v$ lies in
$W_v^\times(1+\mathfrak m_v^d)$ if and only if its class lies in
$p^aC_v$.

\emph{Step 2: the product.}  By Step~1, the left-hand side of
\eqref{eq:precision-native} holds if and only if every component of $\hat
z$ lies in $p^aC_v$, that is, if and only if $\hat z\in p^aM_{n,p}$.
Since $z$ has relative norm one, $z\,\sigma_n(z)=1$, so
$\sigma_n\hat z=-\hat z$ and $\hat z\in M_{n,p}^-$.  It remains to show
$M_{n,p}^-\cap p^aM_{n,p}=p^aM_{n,p}^-$.  The inclusion $\supseteq$ is
clear.  Since $p$ is odd, $(1-\sigma_n)/2$ is a projector of
$M_{n,p}$ onto $M_{n,p}^-$.  If $x=p^ay\in M_{n,p}^-$ with $y\in
M_{n,p}$, then $x=\frac{1-\sigma_n}2x=p^a\cdot\frac{1-\sigma_n}2y\in
p^aM_{n,p}^-$.  This proves \eqref{eq:precision-native}.

\emph{Step 3: membership.}  If $e_{n,p}\le a$, then
$p^aM_{n,p}^-\subseteq p^{e_{n,p}}M_{n,p}^-\subseteq U_{n,p}^0$, so the
right-hand side of \eqref{eq:precision-native} gives $\hat z\in
U_{n,p}^0$.
\end{proof}

The index bound in Lemma~\ref{lem:precision-index} is also necessary, in
the following uniform sense.  Every vector of $M_{n,p}^-$ satisfying the
right-hand condition of \eqref{eq:precision-native} lies in $U^0_{n,p}$
if and only if $e_{n,p}\le\lfloor d/2\rfloor$.  This is the definition of
$e_{n,p}$, together with the inclusion $p^{a}M_{n,p}^-\subseteq
p^{e}M_{n,p}^-$ for integers $a\ge e\ge0$.  It is not a converse for the particular
element $z_n$, with divisor $A_nP_n^{-1}$.  A particular element can
belong to $U_{n,p}^0$ even when the whole subgroup
$p^{\lfloor d/2\rfloor}M_{n,p}^-$ does not.

The lemma gives one sufficient way to verify (O) at a chosen step $n$.
Choose integers $d_{n,p}\ge1$.  At every odd characteristic $p$ already
ramified in $F_{n-1}$, verify the left-hand precision condition of
\eqref{eq:precision-native} for $z_n$ with $d=d_{n,p}$, together with the
index bound $e_{n,p}\le\lfloor d_{n,p}/2\rfloor$; at every newly ramified
characteristic, verify the membership $z_{n,p}\in U^0_{n,p}$ directly.
All these conditions must hold at the same step.  Precision prescribed at
a finite list of previously treated fibres does not cover a fibre outside
that list, in particular a ramified fibre whose residue offsets are being
separated in the current step.  Proposition~\ref{prop:finite-precision}
prescribes precision only at primes that carry no residue prescription.

Increasing precision alone does not control the unit index, because the
field changes from step to step.  At the level of abstract lattices, take
\[
 M_n'=\Z_p,\qquad U_n'=p^{n+1}\Z_p,\qquad g_n'=p^n.
\]
Every $U_n'$ has finite index in $M_n'$, and $g_n'$ satisfies the
right-hand condition of \eqref{eq:precision-native} with $d=2n$, that is,
$g_n'\in p^nM_n'$; yet $g_n'\notin U_n'$ for every $n$.  This abstract
example shows that the two properties together do not imply membership.  The exponent $e_{n,p}$ depends on
the unit image of the field $F_n$ itself, which is not determined while
the parameters of the step are being chosen.  Convergence of the
successive elements $z_n$ at fixed places and finite index of the unit
images do not by themselves bound these exponents in terms of the
prescribed precisions.

To complete the construction of a counterexample it therefore suffices to
prove a simultaneous extension statement such as
Hypothesis~\ref{hyp:selection}, or to construct one successful infinite
sequence of scaled steps by another method.

\appendix
\section{Finite local prescriptions and precision}\label{app:precision}

This appendix shows how much the construction of
Section~\ref{sec:construction} can prescribe about the element
$z=\alpha\beta/y$ of a single step at finitely many odd primes.  Its
results are not used in the proofs of Theorems~\ref{thm:A}
and~\ref{thm:B}.  They are background for Section~\ref{sec:selection}.
There, Lemma~\ref{lem:precision-index} shows that the membership required
by
condition (O) at an odd prime $p$ ramified in $F_n$ follows from two ingredients: the
element $z_n$ has coefficient precision at least $d$
(Definition~\ref{def:coefficient-precision}) at every place $v$ above $p$
with $e(v/p)=2$, and the unit index exponent $e_{n,p}$ of
\eqref{eq:unit-index-exponent} is at most $\lfloor d/2\rfloor$.  This appendix shows that the first
ingredient can be prescribed by the construction, at finitely many odd
primes of the inert set in each step
(Proposition~\ref{prop:finite-precision}), and along every fixed odd
rational prime other than $47$ and $7$ from some level on (Corollary~\ref{cor:fixed-precision}).
Section~\ref{sec:selection} explains why this precision does not by
itself give (O).

The argument has three parts.  Section~\ref{sec:app-local-models} works
in one completion at an odd prime.  There we choose the parameter $c$ so
that the prime is inert in the step and $z$ lies in the coefficient ring
of the completion of the next field, and we show that small changes of
$c$ and $u$ keep $z$ in the coefficient ring up to any prescribed power
of the maximal ideal.  The tools are Teichm\"uller representatives and
Newton's lemma, which turns an approximate root of a polynomial into a
root \cite[Theorem~7.32]{MilneANT}; see
Section~\ref{sec:nt-background}.  Section~\ref{sec:app-finite-refinement}
inserts these local conditions into the proof of
Proposition~\ref{prop:scaled-step}.  Section~\ref{sec:app-dyadic}
records how closely $z$ is approximated by elements of $F$ at the dyadic
places.  The residue prescriptions used to separate residue offsets are
treated in Lemma~\ref{lem:scaled-residues}; the primes considered here lie
in the inert set and carry no residue prescription.

\subsection{Local models with \texorpdfstring{$z$}{z} in the coefficient ring at odd places}\label{sec:app-local-models}

Throughout this subsection let $\mathcal K/\Q_p$ be a finite extension
with $p$ odd, with maximal ideal $\mathfrak m_{\mathcal K}$ and residue
field $k_{\mathcal K}$.  Let $W$ be its coefficient ring, the ring of
integers of its maximal unramified subfield $\Frac W$
(Section~\ref{sec:nt-background}).  Let $W_2$ be the ring of integers of
the unramified quadratic extension of $\Frac W$, and set
$\mathcal K_2=\mathcal K\cdot\Frac W_2$.  Since $\mathcal K/\Frac W$ is
totally ramified, $\mathcal K\cap\Frac W_2=\Frac W$, and
$\mathcal K_2/\mathcal K$ is unramified quadratic.  Write $\sigma$ for the
nontrivial automorphism of $\mathcal K_2/\mathcal K$ and $N$ for the norm
from $\mathcal K_2$ to $\mathcal K$.  The automorphism $\sigma$ restricts
to the nontrivial automorphism of $\Frac W_2/\Frac W$, so it preserves
$W_2$.  The field $\Frac W_2$ is the maximal unramified subfield of
$\mathcal K_2$, so $W_2$ is the coefficient ring of $\mathcal K_2$.

The first lemma produces one choice of $c$ for which the polynomial of
\eqref{eq:scaled-family} defines the unramified quadratic extension and
$z$ lies in $W_2$.

\begin{lemma}\label{lem:coefficient-model}
Let $y,h,u\in\OO_{\mathcal K}^\times$ with $h\in W$ and $u\ne\pm1$.
Suppose that $u\in W$ or $u\equiv1\pmod{\mathfrak m_{\mathcal K}}$.
Then there is $c\in\OO_{\mathcal K}^\times$ with the following
properties.  Put $a=c(u-1)/h$ and $t=(y(1-u^2)-a^2)/(au)$.  The
polynomial $X^2-tX+y$ defines $\mathcal K_2/\mathcal K$; its roots
$\alpha,\sigma\alpha$ and the elements $\beta=a+u\alpha$ and $\sigma\beta$
are units; and $z=\alpha\beta/y$ lies in $W_2^\times$.
\end{lemma}

\begin{proof}
We first give a formula that recovers $c$, $\alpha$ and $\beta$ from a
prescribed value of $z$.  In Cases~2 and~3 we choose $z$ so that this
formula has a unit solution.  In Case~1 we instead construct $c$ over
$\Frac W$ and rescale.

\emph{Step 1 (an inverse formula).}  Let $z\in W_2\setminus W$ be a
unit with $N(z)=1$, so that $\sigma z=z^{-1}$.  Put
\[
 \mathcal Q(z)=1+u^2+u(z+z^{-1})\in\mathcal K.
\]
Suppose that the right-hand side of the first equation below is the
square of a unit of $\mathcal K$, and choose $c$, $\alpha$ and $\beta$
by
\begin{equation}\label{eq:coefficient-inverse}
 c^2=\frac{yh^2(u+1)^2}{\mathcal Q(z)},\qquad
 \alpha=-\frac{c(u+z)}{h(u+1)},\qquad
 \beta=-\frac{c(1+uz)}{h(u+1)}.
\end{equation}
Multiplication gives $N\alpha=N\beta=y$,
$\beta=c(u-1)/h+u\alpha$, and $\alpha\beta=yz$.  Indeed
$(u+z)(u+z^{-1})=(1+uz)(1+uz^{-1})=\mathcal Q(z)$ and
$(u+z)(1+uz)=z\mathcal Q(z)$.  Put $a=c(u-1)/h$; then $a\ne0$ because
$u\ne1$.  Expanding $N(a+u\alpha)=y$ gives
$a^2+au\operatorname{Tr}\alpha+u^2y=y$, so the trace of $\alpha$ is the
element $t$ of the statement, and $\alpha,\sigma\alpha$ are the roots of
$X^2-tX+y$.  The expression for $\alpha$ is linear in $z$ with nonzero
coefficient, and $z\notin\mathcal K$ because
$\mathcal K\cap\Frac W_2=\Frac W$.  Hence
$\mathcal K(\alpha)=\mathcal K(z)=\mathcal K_2$.

\emph{Case 1: $u\in W$.}  The rings $W$ and $\OO_{\mathcal K}$ have the
same residue field, so we may choose $y^\circ\in W^\times$ with the same
residue as $y$.  Since $p$ is odd and $y/y^\circ\equiv1\pmod{\mathfrak
m_{\mathcal K}}$, Newton's lemma \cite[Theorem~7.32]{MilneANT} gives a
unit $\xi\in\mathcal K$ with $\xi^2=y/y^\circ$.  Apply
Lemma~\ref{lem:scaled-holding} over the field $\Frac W$ to the units
$y^\circ,h,u$.  It supplies a unit $c^\circ\in W$ for which the
polynomial $X^2-t^\circ X+y^\circ$ of \eqref{eq:scaled-family}, formed
with $c^\circ,y^\circ$ in place of $c,y$, defines $\Frac W_2/\Frac W$.
Its trace $t^\circ$ is integral by the second expression for $t$ in
\eqref{eq:scaled-family}, whose denominators $c^\circ h$ and
$c^\circ hu$ are units.  Its roots $\alpha^\circ,\sigma\alpha^\circ$ are
therefore integral, with norm the unit $y^\circ$, and hence are units.
The element $a^\circ=c^\circ(u-1)/h$ is integral.  The corresponding
elements $\beta^\circ=a^\circ+u\alpha^\circ$ and $\sigma\beta^\circ$ are
integral, and their norm is
$(a^\circ)^2+a^\circ ut^\circ+u^2y^\circ=y^\circ$ by the definition of
$t^\circ$; so they are units.  Set
\[
 c=\xi c^\circ,\quad a=\xi a^\circ,\quad t=\xi t^\circ,\quad
 \alpha=\xi\alpha^\circ,\quad\beta=\xi\beta^\circ.
\]
These are the formulas \eqref{eq:scaled-family} for $y$: the polynomial
$X^2-tX+y$ has the root $\alpha=\xi\alpha^\circ$, and
$\beta=a+u\alpha$.  Since $\alpha^\circ\in\Frac W_2\setminus\Frac W$ and
$\mathcal K\cap\Frac W_2=\Frac W$, the element $\alpha$ generates
$\mathcal K_2$ over $\mathcal K$.  The elements $c,\alpha,\sigma\alpha,
\beta,\sigma\beta$ are units, and
\[
 z=\frac{\alpha\beta}{y}=\frac{\alpha^\circ\beta^\circ}{y^\circ}\in
 W_2^\times.
\]
Moreover $z\notin W$.  Indeed, from $\beta=a+u\alpha$ and
$\alpha^2=t\alpha-y$ one gets $\alpha\beta=(a+ut)\alpha-uy$ with
$a+ut=-yh(u+1)/c$, so that $z=-u-h(u+1)\alpha/c$ (this is the first
identity in \eqref{eq:scaled-target}).  Hence
$z=-u-h(u+1)\alpha^\circ/c^\circ\notin W$.

\emph{Case 2: $u\equiv1\pmod{\mathfrak m_{\mathcal K}}$ and $y$ is a
square.}  Here $u$ need not lie in $W$.  We choose a norm-one unit
$z\in W_2\setminus W$ arbitrarily close to $1$.  For example, let
$\theta\in W_2$ be nonzero with $\sigma\theta=-\theta$, and put
\[
 z=\frac{1+p^r\theta}{1-p^r\theta},\qquad r\gg0.
\]
Then $z$ is a unit, $\sigma z=z^{-1}$, and $z$ lies outside $W$.
Indeed, a norm-one element of $W$ satisfies $z=\sigma z=z^{-1}$, so
$z=\pm1$; but $z\equiv1\pmod p$ excludes $-1$, and $\theta\ne0$
excludes $1$.  Since $u+1\equiv2$ is a unit,
\[
 \frac{\mathcal Q(z)}{(u+1)^2}
   =1+\frac{u(z+z^{-1}-2)}{(u+1)^2},
\]
and $z+z^{-1}-2=(z-1)^2/z$ tends to $0$ as $r\to\infty$.  For
sufficiently large $r$ this quotient is therefore congruent to $1$
modulo a high power of $\mathfrak m_{\mathcal K}$, hence the square of a
unit by Newton's lemma.  Since $y$ and $h^2$ are squares of units, the
right-hand side of \eqref{eq:coefficient-inverse} is the square of a
unit, and \eqref{eq:coefficient-inverse} has a unit solution $c$.  The
numerators $u+z$ and $1+uz$ have residue $2\ne0$, and $c$ and $h(u+1)$
are units, so $\alpha$ and $\beta$, and hence $\sigma\alpha$ and
$\sigma\beta$, are units.  By Step~1, $c$ has all the required
properties.

\emph{Case 3: $u\equiv1\pmod{\mathfrak m_{\mathcal K}}$ and $y$ is not a
square.}  Since $p$ is odd and $y$ is a unit, its residue $\bar y$ is
not a square.  Put $q=|k_{\mathcal K}|$, and let $\chi$ be the quadratic
character of $\F_q^\times$.  The norm-one subgroup of $\F_{q^2}^\times$
is cyclic of order $q+1$.  Choose an element $\bar z\ne-1$ of it that is
not a square in this subgroup, that is, $\bar z^{(q+1)/2}=-1$.  The
subgroup has $(q+1)/2\ge2$ nonsquares, at most one of which is $-1$, so
this choice is possible for every odd $q$, including $q=3$.  The only norm-one elements of $\F_q$ are $\pm1$, and
$1$ is a square, so $\bar z\notin\F_q$.  The norm of $1+\bar z$ to
$\F_q$ satisfies
\[
 \chi\bigl(N_{\F_{q^2}/\F_q}(1+\bar z)\bigr)
    =\bar z^{-(q+1)/2}=-1.
\]
To check the identity, use
$N_{\F_{q^2}/\F_q}(1+\bar z)=(1+\bar z)(1+\bar z^{-1})=(1+\bar z)^2/\bar z$
and $(1+\bar z)^{q-1}=\bar z^{-1}$.  Let $z\in W_2$ be the
Teichm\"uller representative of $\bar z$.  It has norm one, since the
Teichm\"uller representatives are multiplicative and $\sigma$ acts on
them through $\bar z\mapsto\bar z^q$; and $z\notin W$ because
$\bar z\notin\F_q$.  Since $\bar u=1$, the residue of $\mathcal Q(z)$ is
$2+\bar z+\bar z^{-1}=N_{\F_{q^2}/\F_q}(1+\bar z)$, a nonsquare unit.
Hence the right-hand side of \eqref{eq:coefficient-inverse} has a
nonzero square residue, and by Newton's lemma it is the square of a unit
of $\mathcal K$.  The elements $\alpha,\beta$ are units because their
numerators have residue $1+\bar z\ne0$, and so are $\sigma\alpha$ and
$\sigma\beta$.  By Step~1, $c$ has all the required properties.

The three cases cover all $u$ allowed by the hypotheses.
\end{proof}

The second lemma shows that the conclusion of
Lemma~\ref{lem:coefficient-model} survives small changes of $c$ and $u$,
up to a prescribed precision.  In the global construction $c$ and $u$
can be controlled only modulo powers of a prime, so this is the form in
which the local model is used.

\begin{lemma}[Stability to finite precision]\label{lem:coefficient-stability}
Fix $y,h,u,c$ as in Lemma~\ref{lem:coefficient-model} and an integer
$d\ge1$.  For all $c,u$ in sufficiently small neighbourhoods of these
values, with $y,h$ unchanged, the polynomial $X^2-tX+y$ of
\eqref{eq:scaled-family} still defines $\mathcal K_2/\mathcal K$, the
elements $\alpha,\sigma\alpha,\beta,\sigma\beta$ are units, and, for
either choice of the root $\alpha$,
\[
 z\in W_2^\times(1+\mathfrak m_{\mathcal K_2}^d).
\]
\end{lemma}

\begin{proof}
We follow one root by Newton's lemma and use continuity of the formulas
for $\beta$ and $z$.  Write $t^\circ$, $\alpha^\circ$ and $z^\circ$ for
the values of $t$, $\alpha$ and $z$ given by
Lemma~\ref{lem:coefficient-model}, so that $z^\circ\in W_2^\times$.  The
derivative of $X^2-t^\circ X+y$ at $\alpha^\circ$ is
$2\alpha^\circ-t^\circ=\alpha^\circ-\sigma\alpha^\circ$, which is
nonzero, though possibly not a unit.  Let $e$ be its valuation in
$\mathcal K_2$.  The coefficient $t$ depends continuously on $c,u$ near
the fixed values.  Congruence conditions on $c,u$ modulo sufficiently
high powers of $\mathfrak m_{\mathcal K}$ therefore keep the valuation
of the derivative at $\alpha^\circ$ equal to $e$, and make the value of
$X^2-tX+y$ at $\alpha^\circ$ have valuation arbitrarily larger than
$2e$.  Newton's lemma \cite[Theorem~7.32]{MilneANT} then gives a root
$\alpha\in\mathcal K_2$ of $X^2-tX+y$ as close to $\alpha^\circ$ as
desired.  The field $\mathcal K$ is closed in $\mathcal K_2$ and does
not contain $\alpha^\circ$, so a small enough neighbourhood of
$\alpha^\circ$ avoids $\mathcal K$.  Hence $\alpha$ still generates
$\mathcal K_2$, and the polynomial defines $\mathcal K_2/\mathcal K$.
Being a unit is an open condition, so $\alpha,\sigma\alpha,\beta,
\sigma\beta$ remain units.  The rational expressions
$\beta=a+u\alpha$ and $z=\alpha\beta/y$ are continuous.  Thus
$z/z^\circ\in1+\mathfrak m_{\mathcal K_2}^d$, so
$z\in W_2^\times(1+\mathfrak m_{\mathcal K_2}^d)$.  The other choice of
root replaces $\alpha,\beta$ by $\sigma\alpha,\sigma\beta$, and hence $z$
by $\sigma z=z^{-1}$.  Conjugation preserves $W_2$ and
$1+\mathfrak m_{\mathcal K_2}^d$, so the assertion holds for either root.
\end{proof}

Since $W_2$ is the coefficient ring of $\mathcal K_2$, the conclusion of
Lemma~\ref{lem:coefficient-stability} says that $z$ has coefficient
precision at least $d$ in the sense of
Definition~\ref{def:coefficient-precision}.

\subsection{A finite refinement of the global construction}\label{sec:app-finite-refinement}

We now insert the local conditions of
Section~\ref{sec:app-local-models} into the proof of
Proposition~\ref{prop:scaled-step}.  Since both lemmas there involve
only congruence conditions on $c$ and $u$ at finitely many primes, they
fit into the ray class and Chebotarev arguments of that proof.

\begin{proposition}\label{prop:finite-precision}
Let the data be as in Proposition~\ref{prop:scaled-step}.  In addition,
let $\mathcal J$ be a finite set of primes of $F$ of odd residue
characteristic, disjoint from $\{A,B\}$ and from every prime carrying a
residue prescription, and assign an integer $d_v\ge1$ to each
$v\in\mathcal J$.  Add the primes of $\mathcal J$ to the inert set.  Then
the step can be chosen so that all conclusions of
Proposition~\ref{prop:scaled-step} hold, the step is a scaled step in the
sense of Definition~\ref{def:scaled-step} for the enlarged inert set, and,
at the unique prime $\widetilde v$ of $L$ above each $v\in\mathcal J$,
\begin{equation}\label{eq:finite-coefficient-precision}
 z,-z\in W_{\widetilde v}^\times
              (1+\mathfrak m_{\widetilde v}^{d_v}),
 \qquad z=\alpha\beta/y,
\end{equation}
where $W_{\widetilde v}$ is the coefficient ring of $L_{\widetilde v}$,
the ring of integers of its maximal unramified subfield, and
$\mathfrak m_{\widetilde v}$ is the maximal ideal of
$\OO_{L_{\widetilde v}}$.
\end{proposition}

\begin{proof}
We obtain congruence conditions on $c$ and $u$ at the primes of
$\mathcal J$ from Lemmas~\ref{lem:coefficient-model}
and~\ref{lem:coefficient-stability}, and check that the proof of
Proposition~\ref{prop:scaled-step} can respect them.

\emph{Step 1 (local models at the primes of $\mathcal J$).}  Use the
unit $u_0=\epsilon_H$ as the local starting value of $u$.  Let
$v\in\mathcal J$ and $\mathcal K=F_v$.  If the residue characteristic of
$v$ is neither $5$ nor $13$, then the prime of $F_0$ below $v$ is
unramified over $\Q$, since the discriminant of $F_0$ is $65$.  So the
completion of $F_0$ at that prime is unramified over $\Q_p$ and lies in
$\Frac W$, and $u_0\in W$.  At $5$ and $13$, the residue of
$\epsilon_0=8+\sqrt{65}$ at the prime of $F_0$ above $p$ is $8$, which
has order four modulo $5$ and modulo $13$.  Since $H\ge5$, the exponent
$46\cdot2^{H-4}$ in \eqref{eq:scaled-family} is divisible by four, and
$u_0\equiv1\pmod{\mathfrak m_{\mathcal K}}$.  Also $u_0\ne\pm1$, since
$\epsilon_0$ has infinite order.  The scale $h=2^H$ lies in $W$ and is a
unit because $v$ is odd, and $y$ is a unit at $v$ because $(y)=AB$ and
$v\notin\{A,B\}$.  Thus Lemma~\ref{lem:coefficient-model} applies at
every prime of $\mathcal J$.  In its model the polynomial $X^2-tX+y$
defines the unramified quadratic extension, so $v$ is inert, as the
inert prescription requires.  Lemma~\ref{lem:coefficient-stability},
with $d=d_v$, supplies congruence conditions on $c$ and $u$ modulo finite
powers of $v$ that ensure \eqref{eq:finite-coefficient-precision} for
$z$.  Here $\mathcal K_2=L_{\widetilde v}$, and $W_2=W_{\widetilde v}$
is its coefficient ring.

\emph{Step 2 (compatibility with the proof of
Proposition~\ref{prop:scaled-step}).}  Insert these congruence conditions
into the proof of Proposition~\ref{prop:scaled-step}.  They are
compatible with the other local models of that proof: those models are
imposed at other primes, except for the inert prescriptions at the
primes of $\mathcal J$, which the new conditions refine.  Enlarge the
ray modulus $\mathfrak f$ so that it retains the congruence conditions on
$c$.  Replace the integer $\mathcal M$ by a multiple that is also divisible
by the order of $\epsilon_H$ in each finite group
$(\OO_F/v^{m_v})^\times$, $v\in\mathcal J$, where $v^{m_v}$ is the
modulus of the congruence condition on $u$ at $v$.  Then
$k\equiv1\pmod{\mathcal M}$ gives $u=\epsilon_H^k\equiv u_0
\pmod{v^{m_v}}$, so it retains these conditions.
Lemma~\ref{lem:order-avoidance} applies to the enlarged finite modulus
and integer, so the auxiliary prime $R$ and the condition
$u\equiv-1\pmod R$ can still be selected.  The argument that makes
$\delta_-$ and $\delta_+$ totally positive multiplies $c$ by units
congruent to $1$ modulo the enlarged modulus, and so also preserves
these congruence conditions.  The rest of the proof of
Proposition~\ref{prop:scaled-step} is unchanged.  The new congruence
conditions are local conditions at primes of $\mathcal J$ only, and the
enlarged $\mathcal M$ is still divisible by $6$.  Hence conditions
(i)--(vi) of Definition~\ref{def:scaled-step} still hold, and the step
is a scaled step for the enlarged inert set.

\emph{Step 3 (the element $-z$).}  Finally, $-1$ is a unit of the
coefficient ring at every place of odd residue characteristic, so the
approximation \eqref{eq:finite-coefficient-precision} for $z$ gives the
same approximation for $-z$.
\end{proof}

Applying the proposition at every step, with growing $d_v$, gives
precision along each fixed odd rational prime other than $47$ and $7$.  We use the construction
schedule of the proof of Theorem~\ref{thm:construction}, in which the
odd rational characteristics other than $47$ and $7$ are listed and
treated one fibre at a time (fibres as in Section~\ref{sec:offsets}).

\begin{corollary}[Coefficient precision along a fixed rational prime]\label{cor:fixed-precision}
The tower of Theorem~\ref{thm:construction} can be chosen so that the
following holds.  For every odd rational prime $p\ne47,7$ there is a
level $n_p$ such that, for all $n>n_p$ and every place $v\mid p$ of
$F_n$,
\[
 z_n\in W_{n,v}^\times(1+\mathfrak m_{n,v}^{n}),
\]
where $W_{n,v}$ is the coefficient ring of $(F_n)_v$, the ring of
integers of its maximal unramified subfield, and $\mathfrak m_{n,v}$ is
the maximal ideal of $\OO_{(F_n)_v}$.
\end{corollary}

\begin{proof}
We run the construction schedule of Theorem~\ref{thm:construction} and
add, at each step, the precision conditions of
Proposition~\ref{prop:finite-precision} at all fibres whose treatment is
complete.  At the step that constructs $F_n$, apply
Proposition~\ref{prop:finite-precision} to all primes of $F_{n-1}$ in
the finitely many odd fibres whose treatment is already complete, with
$d_v=n$.  Such primes carry an inert prescription at every later step
and have no further residue prescriptions, so they are admissible as
$\mathcal J$.  The fibre currently being treated is excluded until its
treatment is complete.  The construction schedule still completes the
treatment of each fixed fibre in finitely many steps; the level at which
the treatment of the fibre over $p$ is complete gives $n_p$.  Each step
is a scaled step for the prescription data of the schedule, by
Proposition~\ref{prop:finite-precision}, since the primes of $\mathcal
J$ already carry inert prescriptions.  So the verification in the proof
of Theorem~\ref{thm:construction} applies.
\end{proof}

\subsection{Dyadic approximation of \texorpdfstring{$z$}{z} by elements of \texorpdfstring{$F$}{F}}\label{sec:app-dyadic}

At the dyadic places the situation is different: $z$ is close to an
element of $F$ only up to a bounded depth.  For a scaled step,
\eqref{eq:scaled-target} gives
\begin{equation}\label{eq:dyadic-scalar-depth}
 v_2(-z-1)=H,\qquad
 v_2\bigl(z/(-u)-1\bigr)=H+1
\end{equation}
at both dyadic places.  The second identity follows from
$z+u=-h(u+1)\alpha/c$, since $v_2(h)=H$, $v_2(u+1)=1$, and $\alpha,c,u$
are units at the dyadic places.

The valuation $H+1$ is the largest possible among elements of the dyadic
completion $F_v$.  The residue polynomial of $X^2-tX+y$ at $v$ is
separable, so $\alpha-\sigma\alpha$ is a unit.  Applying $\sigma$ to the
expression for $z+u$ gives $z-\sigma z=-h(u+1)(\alpha-\sigma\alpha)/c$,
and hence
\[
 v_2(z-\sigma z)=H+1.
\]
For an element $b$ of $F_v$, the elements $z-b$ and $\sigma z-b$ are
conjugate and so have the same valuation.  If this valuation were
greater than $H+1$, so would be the valuation of their difference
$z-\sigma z$, a contradiction.  The element $b=-u$ attains $H+1$.

\begin{remark}\label{rem:precision-scope}
The global unit $-1$ gives the congruence $z\equiv-1\pmod{2^H}$ at the
dyadic places.  By Proposition~\ref{prop:finite-precision} it also
preserves every odd approximation in
\eqref{eq:finite-coefficient-precision}, since $-z$ satisfies the same
condition as $z$.

The integers $d_v$ in Proposition~\ref{prop:finite-precision} are finite
inputs to the construction of one step, and changing them can change $c$
and the field $L$.  Likewise, Corollary~\ref{cor:fixed-precision}
concerns the varying fields $F_n$ at each fixed rational prime.  Neither
result gives the membership required by (O) for a fixed step, or a bound
on the unit index exponent $e_{n,p}$ of \eqref{eq:unit-index-exponent}.  Neither
result gives coefficient precision at a fibre while its residues are
being prescribed, whether by a split residue prescription or by an
inert residue prescription (with component $1$ or $-1$).  Section~\ref{sec:selection} discusses
what more is needed.
\end{remark}

\section{Comparison fields and valuation parities}
\label{app:comparison}

Condition (S) of Section~\ref{sec:dyadic} asks that the map
$\lambda_\delta$ on the Selmer group $\mathrm{Sel}_\delta$ of a scaled step
$L=F(\sqrt\delta)$ be injective.  Among the conditions defining
$\mathrm{Sel}_\delta$ is a local condition, stated in terms of $\delta$,
at every prime of $F$ that ramifies in $L$.  The discriminant
factors as $\delta=\delta_-\delta_+/(chu)^2$, and the ramified primes fall
into two disjoint sets: the set $T_-$ of those at which $\delta_-$ has odd
valuation and the set $T_+$ of those at which $\delta_+$ has odd valuation
(Lemma~\ref{lem:comparison-supports}).  The factor $\delta_-=c^2-h^2y$
does not involve the unit $u$, so $T_-$ is determined once $c$ is chosen.

This appendix expresses (S) through the auxiliary quadratic extension
$L'=F(\sqrt{\delta_+})$, which we call the \emph{comparison field}, and
through the parities of valuations at the primes of $T_-$.  The odd primes
ramified in $L'$ are those of $T_+$ together with $A$.  In the
reformulation, the classes of $\mathrm{Sel}_\delta$ are sorted by the
parities of their valuations at the primes of $T_-$, and those with even
valuation at all these primes are described through a Selmer group
attached to $L'$, in which the primes of $T_-$ enter only through
quadratic residue symbols.  The argument has three
parts.  First, the classes of $\Delta(F,\delta)$ with even valuation at
every prime of $T_-$ form a subgroup $\Delta_0$, which injects into
$\mathrm{Sel}_\delta$ with quotient the group of realized valuation-parity
vectors on $T_-$ modulo the all-one vector
(Lemma~\ref{lem:comparison-parity}).  Second, $\Delta_0$ is identified
with the subgroup of the Selmer group of $L'$ on which the residue symbols
at $T_-$ vanish, compatibly with the dyadic Hilbert symbols
(Lemma~\ref{lem:comparison-zero}).  Third, the two pieces combine into a
short exact sequence of kernels (Proposition~\ref{prop:comparison}).
Because this sequence is exact, the reformulation is an equivalence: (S)
holds for the step exactly when two explicit maps, one on the Selmer group
of $L'$ and one on the parity vectors, are injective.  When $|T_-|=1$, the
Selmer data of $L$ and of $L'$ coincide.  The equivalence concerns a given
step; it does not by itself produce a step that satisfies (S)
(Remark~\ref{rem:comparison-scope}).

\emph{Setting.}  Throughout this appendix we fix a scaled step $L/F$
(Definition~\ref{def:scaled-step}) of a tower satisfying
\textup{(T1)--(T6)}, with the notation of \eqref{eq:scaled-family} and of
Section~\ref{sec:notation}.  In particular $(y)=AB$, $(c)=AR$, the unit $u$
is a square in $F$, the elements $\delta_-$ and $\delta_+$ are totally
positive, and
\[
 \delta_-=c^2-h^2y,\qquad \delta_+=c^2(u-1)^2-h^2y(u+1)^2,
 \qquad \delta=\frac{\delta_-\delta_+}{(chu)^2}.
\]
The subscripts $L$ and $L'$ distinguish the Selmer group of the
constructed step $L$ from the analogous group of the comparison field
$L'$.  The odd unit condition (O) is not affected: it always concerns $L$.

\subsection{Ramification sets of \texorpdfstring{$L$ and $L'$}{L and L'}}\label{sec:comparison-ramification}

Define finite sets of odd primes of $F$ by
\[
 T_- =\{\mathfrak p\ne A,B:v_{\mathfrak p}(\delta_-)\equiv1\pmod2\},\qquad
 T_+ =\{\mathfrak p\ne A,B:v_{\mathfrak p}(\delta_+)\equiv1\pmod2\}.
\]
These sets contain every prime with the indicated valuation, whether it
was prescribed in the construction or arose from the choice of the step.

\begin{lemma}\label{lem:comparison-supports}
The sets $T_-$ and $T_+$ are disjoint.  The sets of all odd primes of $F$
ramified in $L/F$ and in $L'/F$ are, respectively,
\[
 T_L=T_-\sqcup T_+,\qquad T_{L'}=\{A\}\sqcup T_+.
\]
The element $\delta_-$ is a square in $F_v$ at every dyadic place $v$ and
in $F_{\mathfrak p}$ at every prime $\mathfrak p\in T_+$.  At $A$ one has
$v_A(\delta_-)=v_A(\delta_+)=1$ and $[\delta_-]_A=[\delta_+]_A$.  Moreover
$T_-$ is nonempty.
\end{lemma}

\begin{proof}
We first compute valuations prime by prime and read off the ramification
sets, then show that $\delta_-$ is a local square at the dyadic places and
at $T_+$ using that it is a norm from $L$, and finally deduce
$T_-\ne\varnothing$ from the class of $A$ in $\Cl(F)/2\Cl(F)$.

\emph{Step 1: odd primes other than $A$ and $B$.}  Expanding the
definitions gives the identity
\begin{equation}\label{eq:comparison-difference}
 \delta_+=(u-1)^2\delta_--4h^2yu.
\end{equation}
At an odd prime $\mathfrak p\ne A,B$, the last term $4h^2yu$ is a unit,
because $h$ is a power of $2$, $(y)=AB$ and $u$ is a unit.  The elements
$\delta_-$ and $\delta_+$ are integral, so they cannot both lie in
$\mathfrak p$: if one of them lies in $\mathfrak p$, the identity shows
that the other is congruent to a unit modulo $\mathfrak p$.  An element
with odd valuation at $\mathfrak p$ lies in $\mathfrak p$, so $T_-\cap T_+=\varnothing$.

\emph{Step 2: the primes $A$ and $B$.}  At $A$ we have
$v_A(c)=v_A(y)=1$ and $u\equiv1\pmod A$.  In $\delta_-=c^2-h^2y$ the first
term has valuation $2$ and the second valuation $1$, so $v_A(\delta_-)=1$.
In $\delta_+$ the first term has valuation at least $2$ and the second
valuation $1$, since $u+1\equiv2\pmod A$ is a unit; so $v_A(\delta_+)=1$.
The prime $A$ splits in $L$, so $\delta$ is a square in $F_A$, and
$[\delta]=[\delta_-\delta_+]$ gives $[\delta_-]_A=[\delta_+]_A$.  At $B$ the
elements $c$, $h$ and $u-1$ are units and $v_B(y)=1$, so $\delta_-$ and
$\delta_+$ are units at $B$.

\emph{Step 3: the ramification sets.}  At an odd prime, a quadratic
extension $F(\sqrt x)$ is ramified exactly when $v_{\mathfrak p}(x)$ is odd
(Section~\ref{sec:nt-background}, quadratic Kummer theory).  The
denominator $(chu)^2$ changes no valuation parity, so
$v_{\mathfrak p}(\delta)\equiv v_{\mathfrak p}(\delta_-)+v_{\mathfrak
p}(\delta_+)\pmod 2$ at every prime; this includes the auxiliary prime
$R$ dividing $c$, which belongs to $T_+$ exactly when $v_R(\delta_+)$ is
odd.  By Steps~1 and~2, $v_{\mathfrak p}(\delta)$ is odd exactly at the
primes of $T_-\sqcup T_+$, and $v_{\mathfrak p}(\delta_+)$ is odd exactly
at $A$ and at the primes of $T_+$.  This proves the two ramification
formulas.

\emph{Step 4: the dyadic places.}  At a dyadic place $v$ we have
$\delta_-/c^2\in1+2^{2H}\OO_{F_v}\subseteq F_v^{\times2}$, as noted after
Definition~\ref{def:scaled-step}.  Hence $\delta_-$ is a square in $F_v$,
and $[\delta_+]_v=[\delta]_v$.  Thus $L'$ has the same unramified quadratic
completion at $v$ as $L$.

\emph{Step 5: the primes of $T_+$.}  A direct calculation from the
identities for $\alpha$ and $\beta$ in \eqref{eq:scaled-family}, using
$N_{L/F}(\alpha)=y$, $\operatorname{Tr}_{L/F}(\alpha)=t$ and
$a=c(u-1)/h$, gives
\begin{equation}\label{eq:comparison-norm}
 N_{L/F}(\alpha-\beta)
   =\frac{a^2-y(u-1)^2}{u}
   =\frac{(u-1)^2\delta_-}{h^2u}.
\end{equation}
The factor $(u-1)^2/(h^2u)$ is a square in $F^\times$, because $u$ is a
square, and a square $s^2$ of $F^\times$ is the norm $N_{L/F}(s)$.
Consequently $\delta_-$ is a norm from $L$.  Let $\mathfrak p\in T_+$.  By
Step~1, $\delta_-$ is a unit at $\mathfrak p$, and $L_{\mathfrak
p}/F_{\mathfrak p}$ is a ramified quadratic extension of odd residue
characteristic.  Since $\mathfrak p$ has a single prime of $L$ above it,
$\delta_-$ is also a norm from the completion $L_{\mathfrak p}$; write
$\delta_-=N_{L_{\mathfrak p}/F_{\mathfrak p}}(w)$ with $w\in L_{\mathfrak p}$.  Then $w$ has valuation zero.  Since
$L_{\mathfrak p}/F_{\mathfrak p}$ is totally ramified, the two fields have
the same residue field and the nontrivial automorphism acts trivially on
it, so $w$ and its conjugate have the same residue.  Hence the residue of
$\delta_-$ is the square of the residue of $w$.  A unit with square residue
at an odd prime is a square, by Newton's lemma \cite[Theorem~7.32]{MilneANT}.
So $\delta_-$ is a square in $F_{\mathfrak p}$.  The argument does not
divide by $c$, so it applies also when $\mathfrak p$ is the auxiliary
prime $R$.

\emph{Step 6: $T_-$ is nonempty.}  Lemma~\ref{lem:base-field}(b) shows
that $\Cl(F_0)\simeq\Z/2\Z$ with $[A_0]\ne0$, so $[A_0]\notin2\Cl(F_0)$.
For $j\ge1$ the prime $A_{j-1}$ splits in $F_j$ and $A_j$ has residue
degree one over it, so the norm map on ideal classes satisfies
\[
 N_{F_j/F_{j-1}}[A_j]=[A_{j-1}].
\]
If $[A_j]=2x$ with $x\in\Cl(F_j)$, then
$[A_{j-1}]=2N_{F_j/F_{j-1}}(x)\in2\Cl(F_{j-1})$.  By induction on $j$,
therefore, $[A_j]\notin2\Cl(F_j)$ for every $j$; in particular
$[A]\notin2\Cl(F)$.  By Step~2 and the definition of $T_-$, the only
primes at which $\delta_-$ has odd valuation are $A$ and the primes of
$T_-$.  Reducing the principal ideal $(\delta_-)$ modulo squares of ideals
and taking classes in $\Cl(F)/2\Cl(F)$ gives
\begin{equation}\label{eq:minus-class-relation}
 [A]+\sum_{\mathfrak p\in T_-}[\mathfrak p]=0
       \quad\text{in }\Cl(F)/2\Cl(F).
\end{equation}
Since $[A]\ne0$ in $\Cl(F)/2\Cl(F)$, the set $T_-$ is nonempty.
\end{proof}

We keep the set $\mathcal D$ of the two dyadic primes of $F$.  By
Lemma~\ref{lem:comparison-supports}, $T_L$ is the set $T$ of all finite
primes ramified in $L$ used in Section~\ref{sec:dyadic}, since no dyadic
prime ramifies in $L$.  Both $L$ and $L'$ are totally real, because
$\delta$ and $\delta_+$ are totally positive.  The comparison field is
ramified at $A$, whereas the constructed step splits at $A$, so $L'$
cannot replace $L$ as a step of the tower; it serves only to reformulate
(S).

\subsection{The Selmer group of the comparison field}\label{sec:comparison-selmer}

We now set up the objects on the side of $L'$: its Selmer group with the
dyadic Hilbert-symbol map, a canonical representative of each Selmer class,
and the residue symbols at $T_-$.  For the constructed step put
\[
 \Delta_L=\Delta(F,\delta),\qquad
 \mathrm{Sel}_L=\Delta_L/\langle[\delta]\rangle=\mathrm{Sel}_\delta,
 \qquad \lambda_L=\lambda_\delta.
\]
Define $\Delta_{L'}$ by the conditions of Section~\ref{sec:dyadic}, with
$\delta_+$ in place of $\delta$ and the set $T_{L'}$ of all odd primes
ramified in $L'/F$ in place of $T$.  Thus $\Delta_{L'}$ consists of the
square classes $[x]\in F^\times/F^{\times2}$ of totally positive elements
$x$ with even valuation at every prime outside $\mathcal D\cup T_{L'}$ and
with $[x]_{\mathfrak p}\in\langle[\delta_+]_{\mathfrak p}\rangle$ at every
$\mathfrak p\in T_{L'}$.  Let
\[
 \mathrm{Sel}_{L'}=\Delta_{L'}/\langle[\delta_+]\rangle,
 \qquad \lambda_{L'}:\mathrm{Sel}_{L'}\longrightarrow (\F_2^{\mathcal D})^0,
 \qquad \lambda_{L'}([x])=([x,-\delta_+]_v)_{v\in \mathcal D}.
\]
Here and below, a Hilbert-symbol formula on a quotient is evaluated on any
representative.  The map $\lambda_{L'}$ is well defined because
$(\delta_+,-\delta_+)_v=1$.  Its values lie in $(\F_2^{\mathcal D})^0$ by
the argument used for $\lambda_\delta$: for $[x]\in\Delta_{L'}$ the symbol
$(x,-\delta_+)_v$ is trivial at every $\mathfrak p\in T_{L'}$ because
$[x]_{\mathfrak p}$ is $1$ or $[\delta_+]_{\mathfrak p}$; it is trivial at
every other odd prime because $x$ and $\delta_+$ both have even valuation
there; and it is trivial at the real places because $x$ is totally
positive.  The product formula for Hilbert symbols \cite[Chapter~V,
Example~5.4]{MilneCFT} then gives $\sum_{v\in\mathcal
D}[x,-\delta_+]_v=0$.

\emph{The $A$-normalized representative.}  At $A$ the two local classes
allowed in $\Delta_{L'}$ are $1$ and $[\delta_+]_A$, and the second has odd
valuation by Lemma~\ref{lem:comparison-supports}.  A class $r$ of
$\mathrm{Sel}_{L'}$ has exactly two representatives in $\Delta_{L'}$,
namely $[x]$ and $[x\delta_+]$, and exactly one of them is a square in
$F_A$: the one with even valuation at $A$.  We call it the
\emph{$A$-normalized representative} $\hat r$ of $r$.  Explicitly,
\begin{equation}\label{eq:A-normalization}
 \mathrm{Sel}_{L'}\longrightarrow\Delta_{L'},
 \qquad r=[x]\longmapsto\hat r=[x\delta_+^{\,v_A(x)\bmod2}].
\end{equation}
Replacing $x$ by $x\delta_+$ changes the parity of $v_A(x)$ and leaves the
right side unchanged.  The classes of $\Delta_{L'}$ that are squares in
$F_A$ form a subgroup that meets $\langle[\delta_+]\rangle$ trivially and
maps onto $\mathrm{Sel}_{L'}$.  So \eqref{eq:A-normalization} is the
inverse of the resulting isomorphism, and it is a linear section of the
quotient map $\Delta_{L'}\to\mathrm{Sel}_{L'}$.  The
uniqueness here concerns square classes, not elements of $F^\times$.

\emph{Residue symbols at $T_-$.}  Let $\mathfrak p$ be an odd prime and
let $[x]$ be a square class with $v_{\mathfrak p}(x)$ even.  The residue
character $\chi_{\mathfrak p}$ of Section~\ref{sec:nt-background} is
defined on such classes, and $[x]$ is a square in $F_{\mathfrak p}$
exactly when $\chi_{\mathfrak p}(x)=0$.  The primes of
$T_-$ lie outside $\mathcal D\cup T_{L'}$ by
Lemma~\ref{lem:comparison-supports}, so every class of $\Delta_{L'}$ has
even valuation at them.  Define
\begin{equation}\label{eq:Theta}
 \Theta:\mathrm{Sel}_{L'}\longrightarrow\F_2^{T_-},
 \qquad \Theta(r)=\bigl(\chi_{\mathfrak p}(\hat r)\bigr)_{\mathfrak p\in T_-}.
\end{equation}
This map is linear, because $r\mapsto\hat r$ is linear and
$\chi_{\mathfrak p}$ is a homomorphism on the classes of even valuation at
$\mathfrak p$.  We also write
\[
 (\Theta,\lambda_{L'}):\mathrm{Sel}_{L'}\longrightarrow
 \F_2^{T_-}\times(\F_2^{\mathcal D})^0,\qquad
 r\longmapsto(\Theta(r),\lambda_{L'}(r)).
\]

\begin{lemma}\label{lem:comparison-reciprocity}
The image of $\Theta$ lies in the hyperplane of vectors with coordinate sum
zero.  In particular $\operatorname{rank}\Theta\le|T_-|-1$.
\end{lemma}

\begin{proof}
We apply the product formula to the pair $(x,\delta_-)$, where $x$ is the
$A$-normalized representative of $r\in\mathrm{Sel}_{L'}$, and check that
only the symbols at $T_-$ can be nontrivial.  The symbol $(x,\delta_-)_v$
is trivial at $A$, because $x$ is a square in $F_A$.  It is trivial at the
dyadic places and at the primes of $T_+$, because $\delta_-$ is a square
there by Lemma~\ref{lem:comparison-supports}.  At an odd prime outside
$\{A\}\cup T_+\cup T_-$, both $x$ and $\delta_-$ have even valuation, so
the symbol is trivial.  The real symbols are trivial because $x$ is
totally positive.  At $\mathfrak p\in T_-$, the valuation
$v_{\mathfrak p}(\delta_-)$ is odd and $v_{\mathfrak p}(x)$ is even, so the
additive symbol $[x,\delta_-]_{\mathfrak p}$ equals $\chi_{\mathfrak p}(x)$
(Section~\ref{sec:nt-background}, Hilbert symbols).  The product formula
\cite[Chapter~V, Example~5.4]{MilneCFT} yields
$\sum_{\mathfrak p\in T_-}\chi_{\mathfrak p}(x)=0$.  The hyperplane of
vectors with coordinate sum zero in $\F_2^{T_-}$ has dimension $|T_-|-1$.
\end{proof}

\subsection{Vanishing valuation parities and the comparison}\label{sec:comparison-parities}

We now compare the two Selmer groups.  Define the
\emph{valuation-parity map} on $T_-$ and its image by
\[
 \operatorname{par}:\Delta_L\longrightarrow\F_2^{T_-},
 \qquad \operatorname{par}([x])=(v_{\mathfrak p}(x)\bmod2)_{\mathfrak p\in T_-},
 \qquad \mathrm{Par}=\im \operatorname{par},
\]
and let
\[
 \Delta_0=\ker\operatorname{par}=\{[x]\in\Delta_L:
          v_{\mathfrak p}(x)\equiv0\pmod2\text{ for all }\mathfrak p\in T_-\}.
\]
Let $\mathbf1\in\F_2^{T_-}$ denote the all-one vector, and set
$\overline{\mathrm{Par}}=\mathrm{Par}/\langle\mathbf1\rangle$.  At every
$\mathfrak p\in T_-$ the valuation $v_{\mathfrak p}(\delta)$ is odd,
because $\delta_+$ is a unit there by Lemma~\ref{lem:comparison-supports};
hence $\operatorname{par}([\delta])=\mathbf1$, and
$\mathbf1\in\mathrm{Par}$.  Since $T_-$ is nonempty, $\mathbf1\ne0$.

\begin{lemma}\label{lem:comparison-parity}
The composite $\Delta_0\to\Delta_L\to\mathrm{Sel}_L$ is injective, and the
map $\operatorname{par}$ induces an exact sequence
\begin{equation}\label{eq:comparison-pattern-sequence}
 0\longrightarrow\Delta_0\longrightarrow \mathrm{Sel}_L
       \longrightarrow\overline{\mathrm{Par}}\longrightarrow0,
 \qquad \dim\overline{\mathrm{Par}}\le|T_-|-1.
\end{equation}
\end{lemma}

\begin{proof}
The kernel of $\Delta_L\to\mathrm{Sel}_L$ is $\langle[\delta]\rangle$, and
$\operatorname{par}([\delta])=\mathbf1\ne0$, so it meets $\Delta_0$
trivially.  Since $\operatorname{par}([\delta])=\mathbf1$, the map
$\operatorname{par}$ descends to a map $\mathrm{Sel}_L\to
\overline{\mathrm{Par}}$, which is surjective by the definition of
$\mathrm{Par}$.  Let a class of $\mathrm{Sel}_L$, represented by $[x]$, lie
in its kernel.  Then $\operatorname{par}([x])$ is $0$ or $\mathbf1$.  In
the second case we replace $x$ by $x\delta$, whose parity vector is $0$.
So the class has a representative in $\Delta_0$, and this representative
is the unique one of $[x]$ and $[x\delta]$ with parity vector $0$.  This
proves exactness.  Finally $\mathrm{Par}\subseteq\F_2^{T_-}$ and
$\mathbf1\ne0$, so $\dim\overline{\mathrm{Par}}\le|T_-|-1$.
\end{proof}

Through Lemma~\ref{lem:comparison-parity} we regard $\Delta_0$ as a
subgroup of $\mathrm{Sel}_L$, and we write
$\lambda_0=\lambda_L|_{\Delta_0}$, that is,
$\lambda_0([x])=([x,-\delta]_v)_{v\in\mathcal D}$ for $[x]\in\Delta_0$.

\begin{lemma}\label{lem:comparison-zero}
Reduction modulo $\langle[\delta_+]\rangle$ defines a canonical isomorphism
\begin{equation}\label{eq:zero-pattern-comparison}
 \Delta_0\xrightarrow{\ \sim\ }\ker\Theta\subseteq \mathrm{Sel}_{L'},
 \qquad [x]\longmapsto[x]\bmod\langle[\delta_+]\rangle,
\end{equation}
whose inverse sends $r\in\ker\Theta$ to its $A$-normalized representative
$\hat r$.  Under this isomorphism $\lambda_0$ corresponds to
$\lambda_{L'}|_{\ker\Theta}$, and therefore
\begin{equation}\label{eq:comparison-zero-kernel}
 \ker\lambda_0\simeq\ker(\Theta,\lambda_{L'}).
\end{equation}
\end{lemma}

\begin{proof}
We compare the local conditions defining $\Delta_0$ and $\Delta_{L'}$
prime by prime, then check that the two maps are mutually inverse and
compare the dyadic symbols.

\emph{Step 1: local behaviour of $\Delta_0$.}  Let $[x]\in\Delta_0$.  At a
prime $\mathfrak p\in T_-$, the two local classes allowed in $\Delta_L$
are $1$ and $[\delta]_{\mathfrak p}$, and the second has odd valuation.
Since $v_{\mathfrak p}(x)$ is even, $x$ is a square in $F_{\mathfrak p}$.
Also $x$ is a square in $F_A$ by Lemma~\ref{lem:A-square}, which holds
for every class of $\Delta(F,\delta)$ without assuming (S) (its proof
uses \eqref{eq:active-dyadic-square} and Hilbert reciprocity, but not the
injectivity of $\lambda_\delta$).  At a
prime $\mathfrak p\in T_+$ the local conditions of $\Delta_L$ and of
$\Delta_{L'}$ agree: $\delta_-$ is a square in $F_{\mathfrak p}$ by
Lemma~\ref{lem:comparison-supports}, so
$[\delta]_{\mathfrak p}=[\delta_+]_{\mathfrak p}$.

\emph{Step 2: the forward map.}  Let $[x]\in\Delta_0$.  Then $x$ is
totally positive and has even valuation outside $\mathcal D\cup T_+$, since
$T_L=T_-\sqcup T_+$ and the valuations at $T_-$ are even; in particular it
has even valuation outside $\mathcal D\cup T_{L'}$.  At $A$ the class
$[x]_A=1$ is allowed, and at $T_+$ the local conditions agree by Step~1.
Hence $[x]\in\Delta_{L'}$.  Let $r$ be its class in $\mathrm{Sel}_{L'}$.
Since $x$ is a square in $F_A$, we have $\hat r=[x]$.  Since $x$ is a
square at every prime of $T_-$, $\Theta(r)=0$.  The forward map is
injective: if $[x]\in\langle[\delta_+]\rangle$ then $[x]=1$, because
$[x]\ne[\delta_+]$, as $v_A(x)$ is even while $v_A(\delta_+)=1$.

\emph{Step 3: the inverse map.}  Let $r\in\ker\Theta$.  Its
$A$-normalized representative $\hat r$ is totally positive, has even
valuation outside $\mathcal D\cup\{A\}\cup T_+$, and is a square in $F_A$;
so it has even valuation outside $\mathcal D\cup T_+$, and in particular
outside $\mathcal D\cup T_L$.  At $T_+$ its local class is allowed in
$\Delta_L$ by Step~1.  At $\mathfrak p\in T_-$ it has even valuation and
$\chi_{\mathfrak p}(\hat r)=0$, so it is a square in $F_{\mathfrak p}$,
which is allowed in $\Delta_L$.  Hence $\hat r\in\Delta_L$, and
$\hat r\in\Delta_0$.  The two maps are mutually inverse: the class of
$\hat r$ modulo $\langle[\delta_+]\rangle$ is $r$, and by Step~2 the
$A$-normalized representative of the image of $[x]\in\Delta_0$ is $[x]$.

\emph{Step 4: the dyadic symbols.}  At a dyadic place $v$, $\delta_-$ is a
square in $F_v$ by Lemma~\ref{lem:comparison-supports}, so
$[\delta]_v=[\delta_+]_v$ and $[x,-\delta]_v=[x,-\delta_+]_v$.  Hence
$\lambda_0([x])=\lambda_{L'}(r)$ when $[x]\in\Delta_0$ corresponds to
$r\in\ker\Theta$.  Thus $\ker\lambda_0$ corresponds to the set of
$r\in\ker\Theta$ with $\lambda_{L'}(r)=0$, which is
$\ker(\Theta,\lambda_{L'})$.
\end{proof}

The map $\lambda_L$ induces a map
\[
 \bar\lambda:\overline{\mathrm{Par}}\longrightarrow
           (\F_2^{\mathcal D})^0/\lambda_0(\Delta_0)
\]
by sending a class of $\overline{\mathrm{Par}}$ to the image of
$\lambda_L(s)$, where $s\in\mathrm{Sel}_L$ is any lift under
\eqref{eq:comparison-pattern-sequence}.  Two lifts differ by an element of
$\Delta_0$, so their values under $\lambda_L$ differ by an element of
$\lambda_0(\Delta_0)$, and $\bar\lambda$ is well defined.

\begin{proposition}[Comparison formulation of (S)]
\label{prop:comparison}
There is an exact sequence
\begin{equation}\label{eq:comparison-kernel-sequence}
 0\longrightarrow\ker(\Theta,\lambda_{L'})
  \longrightarrow\ker\lambda_L
  \longrightarrow\ker\bar\lambda\longrightarrow0.
\end{equation}
Consequently (S) is equivalent to the conjunction of
$\ker(\Theta,\lambda_{L'})=0$ and $\ker\bar\lambda=0$.
If $|T_-|=1$, there is a canonical isomorphism
\[
 (\mathrm{Sel}_L,\lambda_L)\simeq(\mathrm{Sel}_{L'},\lambda_{L'}).
\]
\end{proposition}

\begin{proof}
We restrict the exact sequence \eqref{eq:comparison-pattern-sequence} to
the kernels of $\lambda_0$, $\lambda_L$ and $\bar\lambda$.

\emph{Step 1: exactness.}  The intersection of $\ker\lambda_L$ with the
subgroup $\Delta_0$ is $\ker\lambda_0$, which is isomorphic to
$\ker(\Theta,\lambda_{L'})$ by \eqref{eq:comparison-zero-kernel}; this
gives the injective map on the left.  If $s\in\ker\lambda_L$, then the
image of $s$ in $\overline{\mathrm{Par}}$ lies in $\ker\bar\lambda$, by the
definition of $\bar\lambda$; this gives the map on the right.  An element
$s\in\ker\lambda_L$ maps to zero in $\overline{\mathrm{Par}}$ exactly when
$s\in\Delta_0$, by \eqref{eq:comparison-pattern-sequence}, that is, exactly
when $s\in\ker\lambda_0$.  This is exactness in the middle.  For
surjectivity on the right, let a class of $\overline{\mathrm{Par}}$ lie in
$\ker\bar\lambda$, and choose a lift $s\in\mathrm{Sel}_L$.  Then
$\lambda_L(s)\in\lambda_0(\Delta_0)$, say $\lambda_L(s)=\lambda_0(d)$ with
$d\in\Delta_0$.  The element $s-d$ lies in $\ker\lambda_L$ and has the same
image in $\overline{\mathrm{Par}}$.

\emph{Step 2: the equivalence.}  Condition (S) states that
$\ker\lambda_L=0$.  By the exact sequence \eqref{eq:comparison-kernel-sequence},
this holds exactly when both outer terms vanish.

\emph{Step 3: the case $|T_-|=1$.}  The hyperplane of vectors with
coordinate sum zero in $\F_2^{T_-}=\F_2$ is zero, so
Lemma~\ref{lem:comparison-reciprocity} forces $\Theta=0$ and
$\ker\Theta=\mathrm{Sel}_{L'}$.  Also $\mathrm{Par}=\langle\mathbf1\rangle$,
because $\mathrm{Par}$ contains the nonzero vector $\mathbf1$ of the
one-dimensional space $\F_2^{T_-}$; hence $\overline{\mathrm{Par}}=0$.  By
Lemmas~\ref{lem:comparison-parity} and~\ref{lem:comparison-zero},
\[
 \mathrm{Sel}_L\simeq\Delta_0\simeq \mathrm{Sel}_{L'},
\]
and these isomorphisms carry $\lambda_L$ to $\lambda_0$ and $\lambda_0$ to
$\lambda_{L'}$.
\end{proof}

In this tower $|\mathcal D|=2$, so $(\F_2^{\mathcal D})^0$ is one
dimensional, and Proposition~\ref{prop:comparison} takes the following
form.  If $\lambda_0\ne0$, its image is the whole codomain, so the target
of $\bar\lambda$ is zero and $\bar\lambda=0$.  Condition (S) is then
equivalent to
\[
 \dim\Delta_0=1,\qquad
 \mathrm{Par}=\langle\mathbf1\rangle.
\]
Indeed $\ker\lambda_0=0$ means that $\Delta_0$ embeds in the
one-dimensional space $(\F_2^{\mathcal D})^0$, while $\lambda_0\ne0$ forces
$\Delta_0\ne0$; and $\ker\bar\lambda=\overline{\mathrm{Par}}$ vanishes
exactly when $\mathrm{Par}=\langle\mathbf1\rangle$.  If $\lambda_0=0$, then
$\ker\lambda_0=\Delta_0$, and (S) is equivalent to $\Delta_0=0$ together
with the injectivity of $\bar\lambda$.

\begin{remark}\label{rem:comparison-scope}
The conditions above are equivalent forms of (S) for the chosen step; they
do not produce a step with $|T_-|=1$.  Although $\delta_-$ is fixed once $c$ has been
chosen, changing $c$ to meet another finite prescription can change $T_-$,
and the realized image $\mathrm{Par}$ also depends on $\delta_+$.  The
reformulation therefore leaves the simultaneous extension hypothesis of
Section~\ref{sec:selection} unchanged.
\end{remark}

\section{Proof dependencies}
\label{app:dependencies}

This appendix lists, for each main result, the results and facts that its
proof uses; Section~\ref{sec:intro-strategy} describes how the arguments
fit together.  Theorem~\ref{thm:A} rests on the construction of
Theorem~\ref{thm:construction} and on the factorization results of
Section~\ref{sec:factorization}.  The proof of Theorem~\ref{thm:B} has two
halves.  Theorem~\ref{thm:joint} deduces hereditary atomicity from
\textup{(T1)--(T6)} and the unit approximation condition
\textup{($\mathrm{JT}_A$)}.  Proposition~\ref{prop:assembly} shows that
(S) and (O) at one scaled step (Definition~\ref{def:scaled-step}) of a
tower satisfying \textup{(T1)--(T6)} give the unit approximation of $z_n$
in $\mathcal U_n$ at that step; if this holds at every sufficiently late
step, it gives \textup{($\mathrm{JT}_A$)} with approximation level $N=n$.  All results in
the table are proved in this paper.  The standard facts from algebraic
number theory that the proofs use most often are collected, with
references, in Section~\ref{sec:nt-background}; the others, such as
Minkowski's bound and the Chinese remainder theorem, are cited where they
are used.  A few elementary facts, such as the criterion that
$\mathcal E(\sqrt x)$ is split at a real place of $\mathcal E$ exactly
when $x$ is positive there, are used without citation.

\begin{longtable}{@{}>{\raggedright\arraybackslash}p{0.24\textwidth}>{\raggedright\arraybackslash}p{0.44\textwidth}>{\raggedright\arraybackslash}p{0.25\textwidth}@{}}
\toprule
Result & Inputs used & Role\\
\midrule
\endhead
Theorem~\ref{thm:construction} &
The base field $F_0$, Proposition~\ref{prop:scaled-step} (whose
outputs are scaled steps by Remark~\ref{rem:scaled-step-outputs}), and
the construction schedule (Definition~\ref{def:schedule}) of
Section~\ref{sec:construction}, and the procedure of
Section~\ref{sec:offsets}. &
Realizes \textup{(T1)--(T6)}, without (S) or (O).\\[4pt]
Lemma~\ref{lem:classes}, Proposition~\ref{prop:onesided}, and
Theorem~\ref{thm:ambient} &
\textup{(T1)--(T3)}, descent of principality to a finite level, the
injection $\iota$ of Lemma~\ref{lem:classes}, integer valuations,
Lemma~\ref{lem:geometry}, and Lemma~\ref{lem:units}. &
One-sided boundedness; an atomic almost Dedekind domain without ACCP.\\[4pt]
Theorem~\ref{thm:A} &
Theorem~\ref{thm:construction}, Lemma~\ref{lem:geometry}, and
Theorem~\ref{thm:ambient}. &
An explicit atomic almost Dedekind domain without ACCP.\\[4pt]
Proposition~\ref{prop:subfields} &
The isolated ramified primes of \textup{(T5)} with Lemma~\ref{lem:witness},
integrality, and Lemma~\ref{lem:units} ($S^\times=S\cap D^\times$). &
Reduces a subring that is not a BFD to fraction field $K$.\\[4pt]
Proposition~\ref{prop:residues} &
A subring that is not a BFD, Lemmas~\ref{lem:extraction},
\ref{lem:balance}, \ref{lem:unitresidues} and~\ref{lem:profile}, the classes
of Lemma~\ref{lem:classes}, \textup{(T4)}, the identities $N_{F_n/F_{n-1}}y_n=y_{n-1}$, and the
distinct residue offsets of \textup{(T6)}. &
Recovers all odd residue fields and separates the maximal ideals over
each odd prime.\\[4pt]
Propositions~\ref{prop:unramified} and~\ref{prop:orders} &
\textup{(T1)}, Proposition~\ref{prop:residues},
Lemma~\ref{lem:teichmuller-limit}, compactness of the
infinite fibres, saturated intersection, and eventual unramified
extension. &
Describes the completed orders and fixes their conductor exponents.\\[4pt]
Lemma~\ref{lem:return} and Theorem~\ref{thm:joint} &
Theorem~\ref{thm:ambient} (failure of ACCP),
Proposition~\ref{prop:subfields}, Propositions~\ref{prop:unramified}
and~\ref{prop:orders} with Lemma~\ref{lem:order-units} (the completed
orders at one level $\ge b_p$), Lemma~\ref{lem:intersection}, closure in the whole product $\mathcal U_N$ and its propagation
(Lemma~\ref{lem:propagation}), Lemma~\ref{lem:split}, and one-sided
boundedness (Proposition~\ref{prop:onesided}). &
Deduces hereditary atomicity from \textup{($\mathrm{JT}_A$)}.\\[4pt]
Theorem~\ref{thm:dyadic} &
The Selmer condition (S), the ray class and Kummer descriptions of
Lemmas~\ref{lem:ray-closure} and~\ref{lem:ray-selmer}, the dihedral norm
of Lemma~\ref{lem:dihedral-norm}, the presence of odd ramified primes in
the step, Proposition~\ref{prop:ray-exponent-two} and
Corollary~\ref{cor:active-ray}, Lemma~\ref{lem:gA}, and
Lemma~\ref{lem:A-square}. &
Gives the dyadic principal-unit component of the approximation.\\[4pt]
Lemma~\ref{lem:odd-membership} and
Proposition~\ref{prop:odd-certificate} &
Absolute odd ramification index at most two, the pro-$p$ structure of
the quotients $C_v$, and
the image of the relative units in the whole product $M_{n,p}$, with its
dyadic residue restriction; for Proposition~\ref{prop:odd-certificate}
also $3$ unramified and Lemmas~\ref{lem:odd-layer}
and~\ref{lem:odd-minus-rank}. &
Identifies (O); supplies a finite sufficient test.\\[4pt]
Proposition~\ref{prop:assembly} and Theorem~\ref{thm:B} &
(S) and (O) on the same scaled step, Theorem~\ref{thm:dyadic},
Lemma~\ref{lem:odd-membership}, residue $1$ of $z_n$ at the dyadic
primes, primary decomposition, Theorems~\ref{thm:joint}
and~\ref{thm:ambient}. &
Proves Theorem~\ref{thm:B}.\\[4pt]
Proposition~\ref{prop:selection} &
Hypothesis~\ref{hyp:selection}, dependent choice, the construction
schedule, Definition~\ref{def:scaled-step}(iv)--(vi) and the
verification of \textup{(T1)--(T6)} in the proof of
Theorem~\ref{thm:construction}, and
Theorem~\ref{thm:B}. &
Reduces the hypotheses of Theorem~\ref{thm:B} to a one-step extension
statement.\\[4pt]
Lemma~\ref{lem:precision-index} and
Appendix~\ref{app:precision} &
The logarithm on the quotients $C_v$ (Lemma~\ref{lem:odd-layer}), the
unit index exponent $e_{n,p}$, local models with $z$ in the coefficient
ring, congruence conditions modulo prime powers, and the construction
schedule. &
Relates prescribed coefficient precision to membership in the unit
image.\\[4pt]
Proposition~\ref{prop:comparison} &
The totally positive elements $\delta_-,\delta_+$, the ramification
sets of $L$ and $L'$, Lemma~\ref{lem:A-square}, and Hilbert reciprocity. &
Reformulates (S) by the comparison field and valuation parities.\\
\bottomrule
\end{longtable}

\begin{thebibliography}{99}
\bibitem{AAZ}
D.~D.~Anderson, D.~F.~Anderson, and M.~Zafrullah,
\emph{Factorization in integral domains},
J. Pure Appl. Algebra \textbf{69} (1990), 1--19.
\href{https://doi.org/10.1016/0022-4049(90)90074-R}{doi:10.1016/0022-4049(90)90074-R}.

\bibitem{CGH}
J.~Coykendall, F.~Gotti, and R.~E.~Hasenauer,
\emph{Hereditary atomicity in integral domains},
J. Pure Appl. Algebra \textbf{227} (2023), 107249.
\href{https://doi.org/10.1016/j.jpaa.2022.107249}{doi:10.1016/j.jpaa.2022.107249}.
\href{https://arxiv.org/abs/2112.00264v1}{arXiv:2112.00264v1}.
The result numbers cited here refer to this arXiv version.

\bibitem{FHL}
M.~Fontana, E.~Houston, and T.~Lucas,
\emph{Factoring Ideals in Integral Domains},
Lecture Notes of the Unione Matematica Italiana \textbf{14},
Springer, Heidelberg, 2013.
\href{https://doi.org/10.1007/978-3-642-31712-5}{doi:10.1007/978-3-642-31712-5}.

\bibitem{GHK}
A.~Geroldinger and F.~Halter-Koch,
\emph{Non-Unique Factorizations: Algebraic, Combinatorial and Analytic
Theory}, Pure and Applied Mathematics \textbf{278},
Chapman \& Hall/CRC, 2006.
\href{https://doi.org/10.1201/9781420003208}{doi:10.1201/9781420003208}.

\bibitem{Gilmer}
R.~Gilmer, \emph{Multiplicative Ideal Theory}, corrected reprint,
Queen's Papers in Pure and Applied Mathematics \textbf{90},
Queen's University, 1992.

\bibitem{GL}
F.~Gotti and B.~Li, \emph{Arithmetic properties encoded in undermonoids},
Semigroup Forum \textbf{112} (2026), 427--449.
\href{https://doi.org/10.1007/s00233-025-10578-3}{doi:10.1007/s00233-025-10578-3}.
\href{https://arxiv.org/abs/2412.11199v1}{arXiv:2412.11199v1}.
The result numbers cited here refer to this arXiv version.

\bibitem{Grams}
A.~Grams, \emph{Atomic rings and the ascending chain condition for
principal ideals}, Proc. Cambridge Philos. Soc. \textbf{75} (1974), 321--329.
\href{https://doi.org/10.1017/S0305004100048532}{doi:10.1017/S0305004100048532}.

\bibitem{MilneANT}
J.~S.~Milne, \emph{Algebraic Number Theory}, version 3.08 (2020).
\url{https://www.jmilne.org/math/CourseNotes/ANT.pdf}.

\bibitem{MilneCFT}
J.~S.~Milne, \emph{Class Field Theory}, version 4.03 (2020).
\url{https://www.jmilne.org/math/CourseNotes/CFT.pdf}.

\bibitem{MilneWeil}
J.~S.~Milne, \emph{The Riemann Hypothesis over Finite Fields: From Weil
to the Present Day}, in \emph{The Legacy of Bernhard Riemann after One
Hundred and Fifty Years}, Vol.~II, Adv. Lect. Math. (ALM) \textbf{35.2},
International Press, 2016, 487--565.
Reprinted in ICCM Notices \textbf{4} (2016), no.~2, 14--52.
\href{https://arxiv.org/abs/1509.00797v1}{arXiv:1509.00797v1} (2015).
The Section 1 pinpoint refers to the arXiv version.
\end{thebibliography}
\end{document}
